\chapter{\texorpdfstring{Mathematics}{6. Mathematics}}

\section{\texorpdfstring{Math Utilities}{6.1 Math Utilities}}
\setcounter{section}{1}
\setcounter{subsection}{0}
\subsection{Floating-Point Comparison}
\label{sec:6.1.1}
Floating-point results carry rounding error, so two values that should be equal in exact arithmetic usually differ in their last bits. Comparing them with \inlinecode{==} is therefore a bug in almost every numeric setting, and the helpers below compare within a tolerance instead. This section also collects the mathematical constants and the sign inspections that the rest of the anthology uses.

\begin{compactitem}
  \item \inlinecode{M\_PI}, \inlinecode{M\_E}, \inlinecode{M\_PHI}, \inlinecode{M\_INF}, and \inlinecode{M\_NAN} are the usual constants, computed rather than assumed so that no platform-specific macro is required.
  \item \inlinecode{EQ()}, \inlinecode{NE()}, \inlinecode{LT()}, \inlinecode{GT()}, \inlinecode{LE()}, and \inlinecode{GE()} relationally compare two values $x$ and $y$. Arguments may be of different types; the common type governs behavior. If the common type is a floating-point type, exactly equal values (including same-signed infinities) compare equal; otherwise absolute-error epsilon comparison is used. Values within \inlinecode{EPS} of each other are considered equal, and \inlinecode{LT}/\inlinecode{GT}/\inlinecode{LE}/\inlinecode{GE} shift the boundary by \inlinecode{EPS} accordingly. Otherwise exact comparison is used (\inlinecode{==}, \inlinecode{\textless{}}, etc.). The branch is selected with \inlinecode{if constexpr}, so exact types without floating-point arithmetic (e.g. integers, \inlinecode{Modular}, \inlinecode{Rational}) compose as well.
  \item \inlinecode{rEQ(ref, val)} returns whether \inlinecode{val} equals reference \inlinecode{ref} within relative error \inlinecode{EPS}. The tolerance scales with $|{\inlinecodemath{ref}}|$, so \inlinecode{rEQ(ref, val)} is NOT the same as \inlinecode{rEQ(val, ref)}. Use this when one argument is a known exact value and the other is a computed approximation. Degenerates to exact comparison when \inlinecode{ref} is $0$, since the tolerance collapses to $0$; use \inlinecode{EQ} near zero.
  \item \inlinecode{rEQ\_sym(a, b)} is the symmetric (commutative) variant: tolerance scales with $\max(|{\inlinecodemath{a}}|, |{\inlinecodemath{b}}|)$, so the result is the same regardless of argument order. Still degenerates near zero when both arguments are close to $0$. For both relative comparisons, exactly equal infinities compare equal, while unequal infinities and comparisons between finite and non-finite values compare unequal. Non-floating common types use exact equality.
  \item \inlinecode{sgn(x)} returns $-1$ (if $x < 0$), $0$ (if $x = 0$), or $1$ (if $x > 0$). Unlike \inlinecode{std::signbit()} or \inlinecode{std::copysign()}, this does not handle the sign of \inlinecode{NaN}.
  \item \inlinecode{sign\_bit(x)} is analogous to \inlinecode{std::signbit()}, returning whether the sign bit of the floating point number is set to true. If so, then \inlinecode{x} is considered ``negative.'' Note that this works as expected on \inlinecode{+0.0}, \inlinecode{-0.0}, \inlinecode{Inf}, \inlinecode{-Inf}, \inlinecode{NaN}, as well as \inlinecode{-NaN}. Warning: This assumes that the sign bit is the leading (most significant) bit in the internal IEEE representation and that bytes are stored in little-endian order.
  \item \inlinecode{copy\_sign(x, y)} is analogous to \inlinecode{std::copysign()}, returning a number with the magnitude of \inlinecode{x} but the sign of \inlinecode{y}.
\end{compactitem}

\textbf{Time Complexity}: $O(1)$ per call to all operations.

\textbf{Space Complexity}: $O(1)$ auxiliary for all operations.

\begin{lstlisting}[style=cpp]
#include <algorithm>
#include <climits>
#include <cmath>
#include <limits>
#include <type_traits>

#ifndef M_PI
const double M_PI = std::acos(-1.0);  // Or std::numbers::pi in C++20 and later.
#endif
#ifndef M_E
const double M_E = std::exp(1.0);  // or std::numbers::e and std::numbers::e_v<> in C++20 and later
#endif
const double M_PHI = (1.0 + std::sqrt(5.0)) / 2.0;
const double M_INF = std::numeric_limits<double>::infinity();
const double M_NAN = std::numeric_limits<double>::quiet_NaN();

const double EPS = 1e-9;

template<typename T, typename U, typename C = std::common_type_t<T, U>>
bool EQ(T a, U b) {
  if constexpr (std::is_floating_point_v<C>) return C(a) == C(b) || std::fabs(C(a) - C(b)) <= EPS;
  return C(a) == C(b);
}

template<typename T, typename U, typename C = std::common_type_t<T, U>>
bool LT(T a, U b) {
  if constexpr (std::is_floating_point_v<C>) return C(a) < C(b) - EPS;
  return C(a) < C(b);
}

template<typename T, typename U> bool NE(T a, U b) { return !EQ(a, b); }
template<typename T, typename U> bool GT(T a, U b) { return LT(b, a); }
template<typename T, typename U> bool LE(T a, U b) { return !LT(b, a); }
template<typename T, typename U> bool GE(T a, U b) { return !LT(a, b); }

template<typename T, typename U, typename C = std::common_type_t<T, U>>
bool rEQ(T ref, U val) {
  C x = C(ref), y = C(val);
  if (x == y) return true;
  if constexpr (!std::is_floating_point_v<C>) {
    return false;
  }
  if (!std::isfinite(x) || !std::isfinite(y)) return false;
  return std::fabs(x - y) <= EPS * std::fabs(x);
}

template<typename T, typename U, typename C = std::common_type_t<T, U>>
bool rEQ_sym(T a, U b) {
  C x = C(a), y = C(b);
  if (x == y) return true;
  if constexpr (!std::is_floating_point_v<C>) {
    return false;
  }
  if (!std::isfinite(x) || !std::isfinite(y)) return false;
  return std::fabs(x - y) <= EPS * std::max(std::fabs(x), std::fabs(y));
}

template<typename T>
int sgn(const T &x) {
  return (T{0} < x) - (x < T{0});
}

template<typename Dbl>
bool sign_bit(Dbl x) {
  return (((unsigned char *)&x)[sizeof(x) - 1] >> (CHAR_BIT - 1)) & 1;
}

template<typename Dbl>
Dbl copy_sign(Dbl x, Dbl y) {
  return sign_bit(y) ? -std::fabs(x) : std::fabs(x);
}
\end{lstlisting}
\Needspace*{4\baselineskip}
\textbf{Example Usage}\par
\begin{lstlisting}[style=cpp]
#include <cassert>
using namespace std;

int main() {
  assert(EQ(M_PI, 3.14159265359));
  assert(EQ(M_E, 2.718281828459) && EQ(M_PHI, 1.61803398875));
  assert(EQ(M_INF, M_INF) && rEQ(M_INF, M_INF) && rEQ_sym(M_INF, M_INF));
  assert(!rEQ(M_INF, 0.0) && !rEQ_sym(M_INF, -M_INF));
  assert(!rEQ(1000000000000LL, 1000000000001LL));

  double x = -12345.6789;
  assert((-M_INF < x) && (x < M_INF));
  assert((M_INF + x == M_INF) && (M_INF - x == M_INF));
  assert((M_NAN != x) && (M_NAN != M_INF) && (M_NAN != M_NAN));
  assert(!(M_NAN < x) && !(M_NAN > x) && !(M_NAN <= x) && !(M_NAN >= x));
  assert(isnan(0.0 * M_INF) && isnan(M_INF - M_INF));

  assert(sgn(x) == -1 && sgn(0.0) == 0 && sgn(5678) == 1);
  assert(sign_bit(x) && !sign_bit(0.0) && sign_bit(-0.0));
  assert(!sign_bit(M_INF) && sign_bit(-M_INF));
  assert(!sign_bit(M_NAN) && sign_bit(-M_NAN));
  assert(copy_sign(1.0, +2.0) == +1.0 && copy_sign(M_INF, -2.0) == -M_INF);
  assert(copy_sign(1.0, -2.0) == -1.0 && sign_bit(copy_sign(M_NAN, -2.0)));
  return 0;
}
\end{lstlisting}
\subsection{Rounding}
\label{sec:6.1.2}
Rounding is only ambiguous at a tie, when a value sits exactly halfway between two integers, and the rule chosen for that case is what distinguishes the functions below. The choice matters more than it appears: always rounding ties away from zero biases a sum of many rounded values upward, which is why financial and statistical work usually rounds ties to even instead, spreading them evenly between the two neighbors.

\begin{compactitem}
  \item \inlinecode{floor0(x)} returns \inlinecode{x} rounded down, symmetrically towards zero. This function is analogous to \inlinecode{std::trunc()}.
  \item \inlinecode{ceil0(x)} returns \inlinecode{x} rounded up, symmetrically away from zero. It mirrors \inlinecode{floor0()} and has no \inlinecode{\textless{}cmath\textgreater{}} equivalent, since \inlinecode{std::ceil()} rounds towards positive infinity rather than outward.
  \item \inlinecode{round\_half\_up(x)} returns \inlinecode{x} rounded half up, towards positive infinity.
  \item \inlinecode{round\_half\_down(x)} returns \inlinecode{x} rounded half down, towards negative infinity.
  \item \inlinecode{round\_half\_to0(x)} returns \inlinecode{x} rounded half down, symmetrically towards zero.
  \item \inlinecode{round\_half\_from0(x)} returns \inlinecode{x} rounded half up, symmetrically away from zero. This function is analogous to \inlinecode{std::round()}.
  \item \inlinecode{round\_half\_even(x, eps = 1e-9)} returns \inlinecode{x} rounded half to even, using \inlinecode{eps} to detect ties for banker's rounding.
  \item \inlinecode{round\_half\_alternate(x)} returns \inlinecode{x} rounded, where ties are broken by alternating rounds towards positive and negative infinity.
  \item \inlinecode{round\_half\_alternate0(x)} returns \inlinecode{x} rounded, where ties are broken by alternating symmetric rounds towards and away from zero.
  \item \inlinecode{round\_half\_random(x)} returns \inlinecode{x} rounded, where ties are broken randomly.
  \item \inlinecode{round\_n\_places(x, n, round)} returns \inlinecode{x} rounded to \inlinecode{n} digits after the decimal, using the specified rounding function \inlinecode{round(x)}.
\end{compactitem}

Every function propagates \inlinecode{+0.0}, \inlinecode{-0.0}, \inlinecode{Inf}, \inlinecode{-Inf}, \inlinecode{NaN}, and \inlinecode{-NaN} exactly with the same sign, as \inlinecode{std::round()} does. The four built on adding or subtracting $0.5$ before flooring are otherwise exact except at $0.5 - 2^{-54}$, which rounds up to $1$ because the sum is already $1$ before flooring; prefer \inlinecode{std::round()} or \inlinecode{round\_half\_even()} where that last bit matters.

\textbf{Time Complexity}: $O(1)$ per call to all operations.

\textbf{Space Complexity}: $O(1)$ auxiliary for all operations.

\begin{lstlisting}[style=cpp]
#include <chrono>
#include <cmath>
#include <random>

template<typename Dbl>
Dbl floor0(const Dbl &x) {
  Dbl res = std::floor(std::fabs(x));
  return std::signbit(x) ? -res : res;
}

template<typename Dbl>
Dbl ceil0(const Dbl &x) {
  Dbl res = std::ceil(std::fabs(x));
  return std::signbit(x) ? -res : res;
}

template<typename Dbl>
Dbl round_half_up(const Dbl &x) {
  return std::copysign(std::floor(x + 0.5), x);  // copysign() only ever fixes a zero result.
}

template<typename Dbl>
Dbl round_half_down(const Dbl &x) {
  return std::copysign(std::ceil(x - 0.5), x);  // copysign() only ever fixes a zero result.
}

template<typename Dbl>
Dbl round_half_to0(const Dbl &x) {
  Dbl res = round_half_down(std::fabs(x));
  return std::signbit(x) ? -res : res;
}

template<typename Dbl>
Dbl round_half_from0(const Dbl &x) {
  Dbl res = round_half_up(std::fabs(x));
  return std::signbit(x) ? -res : res;
}

template<typename Dbl>
Dbl round_half_even(const Dbl &x, const Dbl &eps = 1e-9) {
  if (std::signbit(x)) {
    return -round_half_even(-x, eps);
  }
  Dbl ipart;
  std::modf(x, &ipart);
  if (std::fabs(x - (ipart + 0.5)) < eps) {  // exactly halfway: break the tie towards even
    return (std::fmod(ipart, 2.0) < eps) ? ipart : ceil0(ipart + 0.5);
  }
  return round_half_from0(x);
}

template<typename Dbl>
Dbl round_half_alternate(const Dbl &x) {
  Dbl up = round_half_up(x), down = round_half_down(x);
  if (std::isnan(x) || up == down) {  // A NaN must not consume a toggle: NaN == NaN is false.
    return up;
  }
  static bool round_up = false;
  return (round_up = !round_up) ? up : down;
}

template<typename Dbl>
Dbl round_half_alternate0(const Dbl &x) {
  Dbl away = round_half_from0(x), toward = round_half_to0(x);
  if (std::isnan(x) || away == toward) {  // A NaN must not consume a toggle.
    return away;
  }
  static bool round_away = false;
  return (round_away = !round_away) ? away : toward;
}

template<typename Dbl>
Dbl round_half_random(const Dbl &x) {
  Dbl away = round_half_from0(x), toward = round_half_to0(x);
  if (std::isnan(x) || away == toward) {
    return away;
  }
  static std::mt19937 rng(std::chrono::steady_clock::now().time_since_epoch().count());
  return (rng() % 2 == 0) ? away : toward;
}

template<typename Dbl, typename RoundFn>
Dbl round_n_places(const Dbl &x, unsigned int n, RoundFn round) {
  Dbl scale = std::pow(Dbl{10}, n);
  return round(x * scale) / scale;
}
\end{lstlisting}
\Needspace*{4\baselineskip}
\textbf{Example Usage}\par
\begin{lstlisting}[style=cpp]
#include <cassert>
#include <limits>
using namespace std;

bool EQ(double a, double b) {
  return fabs(a - b) < 1e-9;
}

int main() {
  assert(EQ(floor0(1.5), 1.0) && EQ(ceil0(1.5), 2.0));
  assert(EQ(floor0(-1.5), -1.0) && EQ(ceil0(-1.5), -2.0));
  // Both round outward at any fraction, unlike std::round() which rounds to the nearest.
  assert(EQ(floor0(1.8), 1.0) && EQ(ceil0(1.2), 2.0) && EQ(ceil0(-1.2), -2.0));

  // NaN, the infinities, and the sign of zero are all preserved, as std::round() does.
  const double M_INF = numeric_limits<double>::infinity();
  const double M_NAN = numeric_limits<double>::quiet_NaN();
  assert(isnan(round_half_even(M_NAN)) && isinf(ceil0(-M_INF)));
  assert(signbit(round_half_from0(-0.0)) && !signbit(round_half_down(0.0)));
  assert(EQ(round_half_up(+1.5), +2) && EQ(round_half_down(+1.5), +1));
  assert(EQ(round_half_up(-1.5), -1) && EQ(round_half_down(-1.5), -2));
  assert(EQ(round_half_to0(+1.5), +1) && EQ(round_half_from0(+1.5), +2));
  assert(EQ(round_half_to0(-1.5), -1) && EQ(round_half_from0(-1.5), -2));
  assert(EQ(round_half_even(+1.5), +2) && EQ(round_half_even(-1.5), -2));
  assert(EQ(round_half_even(3.1), 3) && EQ(round_half_even(3.4), 3));  // Non-ties round normally.

  double alt1 = round_half_alternate(+1.5);
  assert(EQ(round_half_alternate(+1.2), +1));  // Non-ties do not consume an alternating turn.
  double alt2 = round_half_alternate(+1.5);
  assert(!EQ(alt1, alt2));
  double alt01 = round_half_alternate0(-1.5);
  assert(EQ(round_half_alternate0(-1.2), -1));
  double alt02 = round_half_alternate0(-1.5);
  assert(!EQ(alt01, alt02));

  assert(EQ(round_n_places(-1.23456, 3, round_half_to0<double>), -1.235));
  return 0;
}
\end{lstlisting}
\subsection{Base Conversion}
\label{sec:6.1.3}
A positional numeral is a sequence of digits weighted by powers of its base, so converting between bases is repeated division: divide by the target base, record the remainder as the next digit, and continue with the quotient. Doing that division on a digit vector rather than on a machine integer is what lets \inlinecode{convert\_base()} handle values far larger than any built-in type. Roman numerals are the counterexample that shows what positional notation buys, being additive with a subtractive special case and having no digit for zero.

\begin{compactitem}
  \item \inlinecode{to\_base(x, b = 10)} returns the digits of the unsigned integer \inlinecode{x} in base \inlinecode{b}, where index $0$ of the result stores the least significant digit.
  \item \inlinecode{to\_roman(x)} returns the Roman numeral representation of the unsigned integer \inlinecode{x} as a string.
  \item \inlinecode{from\_roman(s)} returns the value of the Roman numeral \inlinecode{s}, which must consist of Roman digits. Subtractive pairs such as \inlinecode{"IX"} are handled, but the numeral is not otherwise validated, so a malformed one such as \inlinecode{"IC"} simply returns the value its symbols imply.
  \item \inlinecode{convert\_base(d, a, b)} converts an integer in base \inlinecode{a} as a vector \inlinecode{d} of digits (where \inlinecode{d[0]} is the least significant digit) to base \inlinecode{b} as a vector of digits (again with index $0$ holding the least significant digit). This uses repeated long division, so the value itself does not need to fit in a machine integer.
\end{compactitem}

Overflow warning: Each intermediate \inlinecode{rem*a + digit} must fit in \inlinecode{uint64\_t}.

\Needspace*{3\baselineskip}
\textbf{Time Complexity}:
\begin{compactitem}
  \item $O(\log_{b}\left(x + 1\right) + 1)$ per call to \inlinecode{to\_base()}.
  \item $O(x / 1000 + 1)$ per call to \inlinecode{to\_roman()} due to the repeated \inlinecode{M} prefix.
  \item $O(n)$ per call to \inlinecode{from\_roman()}, where $n$ is the length of the numeral.
  \item $O(d \cdot e)$ per call to \inlinecode{convert\_base(d, a, b)}, where $d$ is the number of input digits and $e$ is the number of output digits.
\end{compactitem}

\Needspace*{3\baselineskip}
\textbf{Space Complexity}:
\begin{compactitem}
  \item $O(e)$ auxiliary for \inlinecode{to\_base(x, b)}, where $e$ is the number of output digits.
  \item $O(x / 1000 + 1)$ auxiliary for \inlinecode{to\_roman(x)}.
  \item $O(1)$ auxiliary for \inlinecode{from\_roman()}.
  \item $O(d + e)$ auxiliary for \inlinecode{convert\_base(d, a, b)}.
\end{compactitem}

\begin{lstlisting}[style=cpp]
#include <algorithm>
#include <cassert>
#include <cstdint>
#include <map>
#include <string>
#include <utility>
#include <vector>

std::vector<int> to_base(unsigned int x, int b = 10) {
  assert(b >= 2);
  std::vector<int> res;
  do {  // do-while so that a value of 0 yields the single digit {0}
    res.push_back(x % b);
    x /= b;
  } while (x != 0);
  return res;
}

std::string to_roman(unsigned int x) {
  static const std::string h[] = {"", "C", "CC", "CCC", "CD", "D", "DC", "DCC", "DCCC", "CM"};
  static const std::string t[] = {"", "X", "XX", "XXX", "XL", "L", "LX", "LXX", "LXXX", "XC"};
  static const std::string o[] = {"", "I", "II", "III", "IV", "V", "VI", "VII", "VIII", "IX"};
  std::string prefix(x / 1000, 'M');
  x %= 1000;
  return prefix + h[x / 100] + t[x / 10 % 10] + o[x % 10];
}

unsigned int from_roman(const std::string &s) {
  static const std::map<char, unsigned int> VALUES = {{'I', 1},   {'V', 5},   {'X', 10},  {'L', 50},
                                                      {'C', 100}, {'D', 500}, {'M', 1000}};
  unsigned int res = 0;
  for (int i = 0; i < static_cast<int>(s.size()); i++) {
    auto it = VALUES.find(s[i]);
    assert(it != VALUES.end());
    auto next = i + 1 < static_cast<int>(s.size()) ? VALUES.find(s[i + 1]) : VALUES.end();
    // A smaller symbol before a larger one is subtracted, as the IV in XIV.
    res += next != VALUES.end() && it->second < next->second ? -it->second : it->second;
  }
  return res;
}

std::vector<int> convert_base(const std::vector<int> &d, int a, int b) {
  assert(a >= 2 && b >= 2);
  assert(std::all_of(d.begin(), d.end(), [a](int x) { return 0 <= x && x < a; }));
  std::vector<int> cur = d, res;
  auto trim = [](std::vector<int> &v) {
    while (v.size() > 1 && v.back() == 0) {
      v.pop_back();
    }
  };
  trim(cur);
  if (cur.empty() || (cur.size() == 1 && cur[0] == 0)) {
    return {0};
  }
  while (!(cur.size() == 1 && cur[0] == 0)) {
    std::vector<int> q(cur.size());
    uint64_t rem = 0;
    for (int i = static_cast<int>(cur.size()) - 1; i >= 0; i--) {
      uint64_t x = rem * static_cast<uint64_t>(a) + static_cast<uint64_t>(cur[i]);
      q[i] = static_cast<int>(x / static_cast<uint64_t>(b));
      rem = x % static_cast<uint64_t>(b);
    }
    res.push_back(static_cast<int>(rem));
    trim(q);
    cur = std::move(q);
  }
  return res;
}
\end{lstlisting}
\Needspace*{4\baselineskip}
\textbf{Example Usage}\par
\begin{lstlisting}[style=cpp]
#include <cassert>
using namespace std;

int main() {
  assert(to_roman(1234) == "MCCXXXIV");
  assert(to_roman(5678) == "MMMMMDCLXXVIII");
  assert(from_roman("MCCXXXIV") == 1234);
  assert(from_roman("IX") == 9 && from_roman("MMXXVI") == 2026);

  vector<int> digits{6, 5, 4, 3, 2, 1};
  assert(convert_base(to_base(123456, 20), 20, 10) == digits);
  assert(convert_base(vector<int>{0, 0, 0}, 10, 2) == vector<int>{0});

  vector<int> big_decimal(30, 9);  // 10^30 - 1, larger than uint64_t.
  assert(convert_base(convert_base(big_decimal, 10, 2), 2, 10) == big_decimal);
  return 0;
}
\end{lstlisting}
\subsection{Date and Time}
\label{sec:6.1.4}
Date questions become ordinary arithmetic once dates are converted to a day count. These functions use the proleptic Gregorian calendar, which applies the modern leap rule to every year, including those before the calendar was adopted, and they count days from the Unix epoch of 1970-01-01.

The conversion follows Howard Hinnant's closed-form algorithm, which shifts the year to start in March so that the leap day falls at the end of the year and never interrupts the month lengths. The remaining months then have lengths that repeat in a pattern whose running total is exactly $\lfloor (153m + 2)/5 \rfloor$ days, and whole $400$-year eras of $146097$ days handle the leap rule without a single division-by-$100$ special case. Both directions are branch-light formulas rather than loops over years, so a date thousands of years away costs the same as tomorrow.

\begin{compactitem}
  \item \inlinecode{is\_leap\_year(y)} returns whether year \inlinecode{y} has a February 29.
  \item \inlinecode{days\_in\_month(y, m)} returns the length of month \inlinecode{m} of year \inlinecode{y}, where \inlinecode{m} is in $[1, 12]$.
  \item \inlinecode{days\_from\_civil(y, m, d)} returns the number of days from 1970-01-01 to the date \inlinecode{y}-\inlinecode{m}-\inlinecode{d}, which is negative for earlier dates. The date must be valid.
  \item \inlinecode{civil\_from\_days(days)} returns the tuple (\inlinecode{year}, \inlinecode{month}, \inlinecode{day}) of the date that many days from 1970-01-01, inverting \inlinecode{days\_from\_civil()}.
  \item \inlinecode{day\_of\_week(y, m, d)} returns the weekday of the date, from $0$ for Sunday to $6$ for Saturday.
  \item \inlinecode{days\_between(y1, m1, d1, y2, m2, d2)} returns the number of days from the first date to the second, which is negative when the second date is earlier.
\end{compactitem}

Unix time counts seconds from the same epoch, and is exactly $86400$ times the day count plus the seconds elapsed within the day. That decomposition is exact only because Unix time ignores leap seconds: it is not a count of elapsed SI seconds but a label that repeats one value during a positive leap second, which is what keeps every day $86400$ seconds long.

\begin{compactitem}
  \item \inlinecode{seconds\_from\_hms(h, m, s)} returns the seconds elapsed since midnight, and \inlinecode{hms\_from\_seconds(seconds)} returns the tuple (\inlinecode{hour}, \inlinecode{minute}, \inlinecode{second}) inverting it.
  \item \inlinecode{timestamp\_from\_civil(y, m, d, hh = 0, mm = 0, ss = 0)} returns the Unix timestamp of that instant.
  \item \inlinecode{civil\_from\_timestamp(t)} returns the \inlinecode{DateTime} of a Unix timestamp, which may be negative for instants before the epoch. Its day count must fit in \inlinecode{int}.
\end{compactitem}

ISO 8601 numbers weeks from the one containing the first Thursday of the year, and assigns a date to the year of its own week's Thursday. Early January therefore often belongs to the previous ISO year and late December to the next, which is the mistake this replaces: 2021-01-01 is ISO week $53$ of $2020$.

\begin{compactitem}
  \item \inlinecode{iso\_week\_date(y, m, d)} returns the tuple (\inlinecode{iso\_year}, \inlinecode{week}, \inlinecode{weekday}) of the date, where \inlinecode{week} counts from $1$ and \inlinecode{weekday} runs from $1$ for Monday to $7$ for Sunday.
\end{compactitem}

Easter follows the anonymous Gregorian computus, which locates the paschal full moon from the year's place in the $19$-year Metonic cycle, corrects for the century's leap and lunar drift, then advances to the following Sunday.

\begin{compactitem}
  \item \inlinecode{easter(y)} returns the (\inlinecode{month}, \inlinecode{day}) of Easter Sunday in the Gregorian year \inlinecode{y}, which must be at least $1583$.
\end{compactitem}

Dates are ordinary \inlinecode{int} values with no range checking, and the era arithmetic stays within \inlinecode{int} for any year within a few million of the epoch. The proleptic calendar disagrees with history before a country's adoption of the Gregorian one in or after 1582, where problems usually specify Julian rules instead: the simpler leap rule of every fourth year.

\textbf{Time Complexity}: $O(1)$ per call to all operations.

\textbf{Space Complexity}: $O(1)$ auxiliary for all operations.

\begin{lstlisting}[style=cpp]
#include <cassert>
#include <climits>
#include <cstdint>
#include <tuple>
#include <utility>

struct DateTime {
  int year, month, day, hour, minute, second;
};

bool is_leap_year(int y) {
  return y % 4 == 0 && (y % 100 != 0 || y % 400 == 0);
}

int days_in_month(int y, int m) {
  assert(m >= 1 && m <= 12);
  static const int lengths[] = {31, 28, 31, 30, 31, 30, 31, 31, 30, 31, 30, 31};
  return m == 2 && is_leap_year(y) ? 29 : lengths[m - 1];
}

int days_from_civil(int y, int m, int d) {
  assert(1 <= m && m <= 12 && 1 <= d && d <= days_in_month(y, m));
  y -= m <= 2;  // Start the year in March, so that February 29 is the final day of a year.
  int era = (y >= 0 ? y : y - 399) / 400;
  int yoe = y - era * 400;  // Year of era, in [0, 399].
  int doy =
      (153 * (m + (m > 2 ? -3 : 9)) + 2) / 5 + d - 1;  // Day of March-based year, in [0, 365].
  int doe = yoe * 365 + yoe / 4 - yoe / 100 + doy;     // Day of era, in [0, 146096].
  return era * 146097 + doe - 719468;                  // 719468 is the day of era of 1970-01-01.
}

std::tuple<int, int, int> civil_from_days(int days) {
  days += 719468;
  int era = (days >= 0 ? days : days - 146096) / 146097;
  int doe = days - era * 146097;
  int yoe = (doe - doe / 1460 + doe / 36524 - doe / 146096) / 365;
  int y = yoe + era * 400;
  int doy = doe - (365 * yoe + yoe / 4 - yoe / 100);
  int mp = (5 * doy + 2) / 153;  // Month index of the March-based year, in [0, 11].
  int d = doy - (153 * mp + 2) / 5 + 1;
  int m = mp + (mp < 10 ? 3 : -9);
  return {y + (m <= 2), m, d};
}

int day_of_week(int y, int m, int d) {
  // 1970-01-01 was a Thursday, which is weekday 4.
  return ((days_from_civil(y, m, d) + 4) % 7 + 7) % 7;
}

int days_between(int y1, int m1, int d1, int y2, int m2, int d2) {
  return days_from_civil(y2, m2, d2) - days_from_civil(y1, m1, d1);
}

int seconds_from_hms(int h, int m, int s) {
  assert(h >= 0 && h < 24 && m >= 0 && m < 60 && s >= 0 && s < 60);
  return (h * 60 + m) * 60 + s;
}

std::tuple<int, int, int> hms_from_seconds(int seconds) {
  assert(seconds >= 0 && seconds < 86400);
  return {seconds / 3600, seconds / 60 % 60, seconds % 60};
}

int64_t timestamp_from_civil(int y, int m, int d, int hh = 0, int mm = 0, int ss = 0) {
  return 86400LL * days_from_civil(y, m, d) + seconds_from_hms(hh, mm, ss);
}

DateTime civil_from_timestamp(int64_t t) {
  int64_t days = t / 86400, rest = t % 86400;
  if (rest < 0) {  // Truncating division rounds toward zero, but the day count must floor.
    rest += 86400;
    days--;
  }
  assert(INT_MIN <= days && days <= INT_MAX);
  auto [y, m, d] = civil_from_days(static_cast<int>(days));
  auto [hh, mm, ss] = hms_from_seconds(static_cast<int>(rest));
  return {y, m, d, hh, mm, ss};
}

std::tuple<int, int, int> iso_week_date(int y, int m, int d) {
  int days = days_from_civil(y, m, d);
  int weekday = ((days + 3) % 7 + 7) % 7 + 1;  // The epoch was a Thursday, ISO weekday 4.
  int thursday = days - weekday + 4;           // This week's Thursday fixes the ISO year.
  int iso_year = std::get<0>(civil_from_days(thursday));
  return {iso_year, (thursday - days_from_civil(iso_year, 1, 1)) / 7 + 1, weekday};
}

std::pair<int, int> easter(int y) {
  assert(y >= 1583);
  int a = y % 19;                // Position in the Metonic cycle.
  int b = y / 100, c = y % 100;  // Century and year within the century.
  int d = b / 4, e = b % 4;      // Leap corrections for the century.
  int f = (b + 8) / 25, g = (b - f + 1) / 3;
  int h = (19 * a + b - d - g + 15) % 30;  // Days from March 21 to the paschal full moon.
  int i = c / 4, k = c % 4;
  int l = (32 + 2 * e + 2 * i - h - k) % 7;  // Days from the full moon to the next Sunday.
  int mm = (a + 11 * h + 22 * l) / 451;      // Correction for the two latest possible dates.
  return {(h + l - 7 * mm + 114) / 31, (h + l - 7 * mm + 114) % 31 + 1};
}
\end{lstlisting}
\Needspace*{4\baselineskip}
\textbf{Example Usage}\par
\begin{lstlisting}[style=cpp]
#include <cassert>
using namespace std;

int main() {
  assert(is_leap_year(2024) && is_leap_year(2000));
  assert(!is_leap_year(2023) && !is_leap_year(1900));
  assert(days_in_month(2024, 2) == 29 && days_in_month(2023, 2) == 28);
  assert(days_in_month(2023, 12) == 31 && days_in_month(2023, 4) == 30);

  assert(days_from_civil(1970, 1, 1) == 0);
  assert(days_from_civil(1969, 12, 31) == -1);
  assert(days_from_civil(2000, 1, 1) == 10957);
  assert(civil_from_days(0) == make_tuple(1970, 1, 1));
  assert(civil_from_days(-1) == make_tuple(1969, 12, 31));
  assert(civil_from_days(10957) == make_tuple(2000, 1, 1));

  assert(day_of_week(1970, 1, 1) == 4);   // A Thursday.
  assert(day_of_week(2000, 1, 1) == 6);   // A Saturday.
  assert(day_of_week(1969, 7, 20) == 0);  // A Sunday.

  // 2024 was a leap year, so February contributed 29 days.
  assert(days_between(2024, 1, 1, 2025, 1, 1) == 366);
  assert(days_between(2023, 1, 1, 2024, 1, 1) == 365);
  assert(days_between(2024, 3, 1, 2024, 2, 1) == -29);

  assert(seconds_from_hms(23, 31, 30) == 84690);
  assert(hms_from_seconds(84690) == make_tuple(23, 31, 30));
  assert(timestamp_from_civil(1970, 1, 1) == 0);
  assert(timestamp_from_civil(2009, 2, 13, 23, 31, 30) == 1234567890);

  auto [y, mo, d, hh, mm, ss] = civil_from_timestamp(1234567890);
  assert(y == 2009 && mo == 2 && d == 13 && hh == 23 && mm == 31 && ss == 30);
  auto before = civil_from_timestamp(-1);  // One second before the epoch.
  assert(before.year == 1969 && before.month == 12 && before.day == 31);
  assert(before.hour == 23 && before.minute == 59 && before.second == 59);

  assert(iso_week_date(1970, 1, 1) == make_tuple(1970, 1, 4));
  // An ISO year runs from the Monday of the week holding the first Thursday, so it straddles
  // the calendar year on both sides.
  assert(iso_week_date(2019, 12, 30) == make_tuple(2020, 1, 1));
  assert(iso_week_date(2021, 1, 1) == make_tuple(2020, 53, 5));
  assert(iso_week_date(2026, 8, 17) == make_tuple(2026, 34, 1));

  assert(easter(2000) == make_pair(4, 23));
  assert(easter(2024) == make_pair(3, 31));
  assert(easter(2025) == make_pair(4, 20));
  auto [month, day] = easter(2026);
  assert(day_of_week(2026, month, day) == 0);  // Easter always falls on a Sunday.

  // The century rule, where 2000 is a leap year but 1900 is not.
  assert(days_between(2000, 2, 28, 2000, 3, 1) == 2);
  assert(days_between(1900, 2, 28, 1900, 3, 1) == 1);
  assert(civil_from_days(days_from_civil(2000, 2, 29)) == make_tuple(2000, 2, 29));
  assert(civil_from_days(days_from_civil(1, 1, 1)) == make_tuple(1, 1, 1));
  assert(civil_from_days(days_from_civil(-44, 3, 15)) == make_tuple(-44, 3, 15));
  return 0;
}
\end{lstlisting}

\section{\texorpdfstring{Combinatorics}{6.2 Combinatorics}}
\setcounter{section}{2}
\setcounter{subsection}{0}
\subsection{Combinatorial Calculations}
\label{sec:6.2.1}
The following functions implement common operations in combinatorics. All size and count inputs must be nonnegative, and all moduli must be positive. All return values and table entries are computed as 64-bit integers modulo an input argument $m$ or $p$.

\begin{compactitem}
  \item \inlinecode{factorial(n, m = MOD)} returns $n! \bmod m$.
  \item \inlinecode{factorial\_without\_p(n, p = MOD)} returns $n!$ modulo the prime $p$ after removing every factor of $p$ from the factorial.
  \item \inlinecode{legendre(n, p)} returns the exponent of the prime $p$ in the factorization of $n!$, where $p$ must be at least $2$.
  \item \inlinecode{binomial\_table(n, m = MOD)} returns the first $n + 1$ rows of Pascal's triangle as a two-dimensional vector $t$ such that $t[i][j] = \binom{i}{j} \bmod m$.
  \item \inlinecode{permute(n, k, m = MOD)} returns $(n \mathbin{\text{permute}} k) \bmod m$.
  \item \inlinecode{choose(n, k, p = MOD)} returns $\binom{n}{k} \bmod p$, where $p$ is prime and $n < p$.
  \item \inlinecode{choose\_lucas(n, k, p = MOD)} returns $\binom{n}{k} \bmod p$ for any 64-bit \inlinecode{n} and \inlinecode{k}, where $p$ is prime. It is only practical for a small $p$, since each base-$p$ digit costs a call to \inlinecode{choose()}.
  \item \inlinecode{multichoose(n, k, p = MOD)} returns $(n \mathbin{\text{multichoose}} k) \bmod p$, where $p$ is prime and $n + k - 1 < p$ when $n > 0$.
  \item \inlinecode{catalan(n, p = MOD)} returns the $n$-th Catalan number mod $p$, where $p$ is prime and $2n < p$.
  \item \inlinecode{partitions(n, m = MOD)} returns the number of partitions of $n$, mod $m$.
  \item \inlinecode{partitions(n, k, m = MOD)} returns the number of partitions of $n$ into $k$ parts, mod $m$.
  \item \inlinecode{stirling1(n, k, m = MOD)} returns the $(n, k)$ unsigned Stirling number of the 1st kind mod $m$.
  \item \inlinecode{stirling2(n, k, m = MOD)} returns the $(n, k)$ Stirling number of the 2nd kind mod $m$.
  \item \inlinecode{eulerian1(n, k, m = MOD)} returns the $(n, k)$ Eulerian number of the 1st kind mod $m$, where $n > k$.
  \item \inlinecode{eulerian2(n, k, m = MOD)} returns the $(n, k)$ Eulerian number of the 2nd kind mod $m$, where $n > k$.
\end{compactitem}

Legendre's formula sums $\lfloor n/p^i \rfloor$ over every $i \geq 1$, counting how many of $1, \dots, n$ are divisible by each power of $p$, which totals the exponent of $p$ in $n!$. It answers questions such as how many trailing zeros $n!$ has in base $p$. Paired with Kummer's theorem, which equates the exponent of $p$ in $\binom{n}{k}$ with the number of carries when adding $k$ and $n - k$ in base $p$, it also decides when a binomial coefficient is divisible by $p$.

Lucas' theorem reduces a binomial coefficient modulo a prime to a product over base-$p$ digits: $\binom{n}{k} \equiv \prod_i \binom{n_i}{k_i} \pmod p$, where $n_i$ and $k_i$ are the $i$-th digits of $n$ and $k$ in base $p$. The product is $0$ as soon as some $k_i$ exceeds $n_i$. For a composite modulus, factor it into prime powers, apply the generalized version of the theorem to each factor, and recombine the residues with the Chinese remainder theorem of section \secref{6.3.2}.

Every routine here recomputes its answer from scratch, which is what suits a handful of values or a modulus that varies between calls. When a solution instead needs thousands of binomial coefficients against one fixed modulus, precompute factorials and inverse factorials: the combinatorics tables of section \secref{6.3.3} amortize $O(n)$ of table growth and then answer each query in $O(1)$.

Overflow warning: Modular products use ordinary \inlinecode{int64\_t} multiplication, so the square of the chosen modulus must fit in \inlinecode{int64\_t}.

\Needspace*{3\baselineskip}
\textbf{Time Complexity}:
\begin{compactitem}
  \item $O(n)$ per call to \inlinecode{factorial()}.
  \item $O(p \log_{p}\left(n\right))$ per call to \inlinecode{factorial\_without\_p()}.
  \item $O(\log_{p}\left(n\right))$ per call to \inlinecode{legendre()}.
  \item $O(n^{2})$ per call to \inlinecode{binomial\_table()}.
  \item $O(k)$ per call to \inlinecode{permute()}.
  \item $O(\min\left(k, n - k\right))$ per call to \inlinecode{choose()}.
  \item $O(p \log_{p}\left(n\right))$ per call to \inlinecode{choose\_lucas()}.
  \item $O(k)$ per call to \inlinecode{multichoose()}.
  \item $O(n)$ per call to \inlinecode{catalan()}.
  \item $O(n^{2})$ per call to \inlinecode{partitions(n, m)}.
  \item $O(n \cdot k)$ per call to \inlinecode{partitions(n, k, m)}, \inlinecode{stirling1()}, \inlinecode{stirling2()}, \inlinecode{eulerian1()}, and \inlinecode{eulerian2()}.
\end{compactitem}

\Needspace*{3\baselineskip}
\textbf{Space Complexity}:
\begin{compactitem}
  \item $O(n^{2})$ auxiliary for \inlinecode{binomial\_table(n, m)}.
  \item $O(n \cdot k)$ auxiliary for \inlinecode{partitions(n, k, m)}, \inlinecode{stirling1(n, k, m)}, \inlinecode{stirling2(n, k, m)}, \inlinecode{eulerian1(n, k, m)}, and \inlinecode{eulerian2(n, k, m)}.
  \item $O(1)$ auxiliary for all other operations.
\end{compactitem}

\begin{lstlisting}[style=cpp]
#include <cassert>
#include <cstdint>
#include <vector>

const int64_t MOD = 1000000007;

int64_t factorial(int n, int64_t m = MOD) {
  assert(n >= 0 && m > 0);
  int64_t res = 1;
  for (int i = 2; i <= n; i++) {
    res = (res * i) % m;
  }
  return res % m;
}

int64_t factorial_without_p(int64_t n, int64_t p = MOD) {
  assert(n >= 0 && p >= 2);
  int64_t res = 1;
  while (n > 1) {
    if (n / p % 2 == 1) {
      res = res * (p - 1) % p;
    }
    int64_t max = n % p;
    for (int64_t i = 2; i <= max; i++) {
      res = (res * i) % p;
    }
    n /= p;
  }
  return res % p;
}

int64_t legendre(int64_t n, int64_t p) {
  assert(n >= 0 && p >= 2);
  int64_t res = 0;
  for (int64_t q = n / p; q > 0; q /= p) {  // Dividing repeatedly avoids overflowing p^i.
    res += q;
  }
  return res;
}

std::vector<std::vector<int64_t>> binomial_table(int n, int64_t m = MOD) {
  assert(n >= 0 && m > 0);
  std::vector<std::vector<int64_t>> t(n + 1);
  for (int i = 0; i <= n; i++) {
    for (int j = 0; j <= i; j++) {
      if (i < 2 || j == 0 || i == j) {
        t[i].push_back(1);
      } else {
        t[i].push_back((t[i - 1][j - 1] + t[i - 1][j]) % m);
      }
    }
  }
  return t;
}

int64_t permute(int n, int k, int64_t m = MOD) {
  assert(n >= 0 && k >= 0 && m > 0);
  if (n < k) {
    return 0;
  }
  int64_t res = 1;
  for (int i = 0; i < k; i++) {
    res = res * (n - i) % m;
  }
  return res % m;
}

int64_t powmod(int64_t x, int64_t n, int64_t m) {
  int64_t res = 1 % m;
  for (x %= m; n > 0; n >>= 1) {
    if (n & 1) {
      res = res * x % m;
    }
    x = x * x % m;
  }
  return res;
}

int64_t choose(int64_t n, int64_t k, int64_t p = MOD) {
  assert(n >= 0 && k >= 0 && p >= 2 && n < p);
  if (n < k) {
    return 0;
  }
  if (k > n - k) {
    k = n - k;
  }
  int64_t num = 1, den = 1;
  for (int64_t i = 0; i < k; i++) {
    num = num * (n - i) % p;
  }
  for (int64_t i = 1; i <= k; i++) {
    den = den * i % p;
  }
  return num * powmod(den, p - 2, p) % p;
}

int64_t choose_lucas(int64_t n, int64_t k, int64_t p = MOD) {
  assert(n >= 0 && p >= 2);
  if (k < 0 || k > n) {
    return 0;
  }
  int64_t res = 1;
  while (n > 0) {  // Since k never exceeds n, k is also exhausted when n reaches 0.
    int64_t ni = n % p, ki = k % p;
    if (ki > ni) {
      return 0;
    }
    res = res * choose(ni, ki, p) % p;
    n /= p;
    k /= p;
  }
  return res;
}

int64_t multichoose(int n, int k, int64_t p = MOD) {
  assert(n >= 0 && k >= 0 && p >= 2);
  if (n == 0) {
    return k == 0;
  }
  return choose(static_cast<int64_t>(n) + k - 1, k, p);
}

int64_t catalan(int n, int64_t p = MOD) {
  assert(n >= 0 && p >= 2);
  return choose(2LL * n, n, p) * powmod(n + 1, p - 2, p) % p;
}

int64_t partitions(int n, int64_t m = MOD) {
  assert(n >= 0 && m > 0);
  std::vector<int64_t> t(n + 1);
  t[0] = 1;
  for (int i = 1; i <= n; i++) {
    for (int j = i; j <= n; j++) {
      t[j] = (t[j] + t[j - i]) % m;
    }
  }
  return t[n] % m;
}

int64_t partitions(int n, int k, int64_t m = MOD) {
  assert(n >= 0 && k >= 0 && m > 0);
  std::vector<std::vector<int64_t>> t(n + 1, std::vector<int64_t>(k + 1));
  t[0][0] = 1;
  for (int i = 1; i <= n; i++) {
    for (int j = 1, h = k < i ? k : i; j <= h; j++) {
      t[i][j] = (t[i - 1][j - 1] + t[i - j][j]) % m;
    }
  }
  return t[n][k] % m;
}

int64_t stirling1(int n, int k, int64_t m = MOD) {
  assert(n >= 0 && k >= 0 && m > 0);
  std::vector<std::vector<int64_t>> t(n + 1, std::vector<int64_t>(k + 1));
  t[0][0] = 1;
  for (int i = 1; i <= n; i++) {
    for (int j = 1; j <= k; j++) {
      t[i][j] = (i - 1) * t[i - 1][j] % m;
      t[i][j] = (t[i][j] + t[i - 1][j - 1]) % m;
    }
  }
  return t[n][k] % m;
}

int64_t stirling2(int n, int k, int64_t m = MOD) {
  assert(n >= 0 && k >= 0 && m > 0);
  std::vector<std::vector<int64_t>> t(n + 1, std::vector<int64_t>(k + 1));
  t[0][0] = 1;
  for (int i = 1; i <= n; i++) {
    for (int j = 1; j <= k; j++) {
      t[i][j] = j * t[i - 1][j] % m;
      t[i][j] = (t[i][j] + t[i - 1][j - 1]) % m;
    }
  }
  return t[n][k] % m;
}

int64_t eulerian1(int n, int k, int64_t m = MOD) {
  assert(n > k && k >= 0 && m > 0);
  if (k > n - 1 - k) {
    k = n - 1 - k;
  }
  std::vector<std::vector<int64_t>> t(n + 1, std::vector<int64_t>(k + 1, 1));
  for (int j = 1; j <= k; j++) {
    t[0][j] = 0;
  }
  for (int i = 1; i <= n; i++) {
    for (int j = 1; j <= k; j++) {
      t[i][j] = (i - j) * t[i - 1][j - 1] % m;
      t[i][j] = (t[i][j] + ((j + 1) * t[i - 1][j] % m)) % m;
    }
  }
  return t[n][k] % m;
}

int64_t eulerian2(int n, int k, int64_t m = MOD) {
  assert(n > k && k >= 0 && m > 0);
  std::vector<std::vector<int64_t>> t(n + 1, std::vector<int64_t>(k + 1, 1));
  for (int i = 1; i <= n; i++) {
    for (int j = 1; j <= k; j++) {
      if (i == j) {
        t[i][j] = 0;
      } else {
        t[i][j] = (j + 1) * t[i - 1][j] % m;
        t[i][j] = ((2 * i - 1 - j) * t[i - 1][j - 1] % m + t[i][j]) % m;
      }
    }
  }
  return t[n][k] % m;
}
\end{lstlisting}
\Needspace*{4\baselineskip}
\textbf{Example Usage}\par
\begin{lstlisting}[style=cpp]
#include <cassert>

int main() {
  auto t = binomial_table(10);
  for (int i = 0; i < static_cast<int>(t.size()); i++) {
    for (int j = 0; j < static_cast<int>(t[i].size()); j++) {
      assert(t[i][j] == choose(i, j));
    }
  }
  assert(factorial(10) == 3628800);
  assert(factorial_without_p(123456) == 639390503);
  assert(factorial_without_p(7, 5) == 3);  // 7! / 5 = 1008 = 3 (mod 5).
  assert(legendre(10, 2) == 8);            // 10! = 2^8 * 3^4 * 5^2 * 7.
  assert(legendre(100, 5) == 24);          // 100! ends in 24 zeros.
  assert(permute(10, 4) == 5040);
  assert(choose(20, 7) == 77520);
  assert(choose_lucas(20, 7) == 77520);
  auto mod7 = binomial_table(30, 7);
  for (int i = 0; i <= 30; i++) {
    for (int j = 0; j <= i; j++) {
      assert(choose_lucas(i, j, 7) == mod7[i][j]);
    }
  }
  // By Lucas' theorem mod 2, a binomial coefficient is odd exactly when k is a submask of n.
  assert(choose_lucas((1LL << 40) - 1, 12345, 2) == 1);
  assert(choose_lucas(1LL << 40, 1, 2) == 0);
  assert(multichoose(20, 7) == 657800);
  assert(catalan(10) == 16796);
  assert(partitions(4) == 5);
  assert(partitions(0, 0) == 1);
  assert(partitions(100, 5) == 38225);
  assert(stirling1(4, 2) == 11);
  assert(stirling2(4, 3) == 6);
  assert(eulerian1(9, 5) == 88234);
  assert(eulerian2(8, 3) == 195800);
  return 0;
}
\end{lstlisting}
\subsection{Enumerating Arrangements}
\label{sec:6.2.2}
For the purposes of this section, we define a ``size $k$ arrangement of $n$'' to be a permutation of a size $k$ subset of the integers in $[0, n)$, for $0 \leq k \leq n$. There are $n \mathbin{\text{permute}} k$ possible arrangements, but $n^k$ possible arrangements if repeated values are allowed.

For arrangements without repeats, each prefix fixes some used values and leaves a shrinking set of choices for the remaining positions. Ranking counts how many unused smaller choices could have been placed at the current position, multiplied by the number of possible suffix arrangements; unranking performs the inverse search by subtracting those block sizes. With repeats allowed, the sequence is just a base-$n$ counter.

\begin{compactitem}
  \item \inlinecode{next\_arrangement(n, a)} tries to rearrange \inlinecode{a} to the next lexicographically greater arrangement, returning true if such an arrangement exists or false if the array is already in descending order (in which case \inlinecode{a} is unchanged). The input \inlinecode{a} must consist of distinct integers in the range $[0, {\inlinecodemath{n}})$.
  \item \inlinecode{arrangement\_by\_rank(n, k, r)} returns the size $k$ arrangement of $n$ which is lexicographically ranked $r$ out of all size $k$ arrangements of $n$, where $r$ is a 0-based rank in the range $[0, n \mathbin{\text{permute}} k)$.
  \item \inlinecode{rank\_by\_arrangement(n, a)} returns an integer representing the 0-based rank of arrangement \inlinecode{a}, which must consist of distinct integers in the range $[0, {\inlinecodemath{n}})$.
  \item \inlinecode{next\_arrangement\_with\_repeats(n, a)} tries to rearrange \inlinecode{a} to the next lexicographically greater arrangement with repeats, returning true if such an arrangement exists or false if the array is already in descending order (in which case \inlinecode{a} is unchanged). The input \inlinecode{a} must consist of integers in the range $[0, {\inlinecodemath{n}})$. If \inlinecode{a} were interpreted as a $k$ digit integer in base $n$, this function could be thought of as incrementing the integer.
\end{compactitem}

Overflow warning: Exact counts and ranks used by the ranking operations must fit in \inlinecode{int64\_t}.

\Needspace*{3\baselineskip}
\textbf{Time Complexity}:
\begin{compactitem}
  \item $O(n \cdot k)$ per call to \inlinecode{next\_arrangement()}, \inlinecode{arrangement\_by\_rank()}, and \inlinecode{rank\_by\_arrangement(n, a)}, where $k$ is the size of \inlinecode{a}.
  \item $O(k)$ per call to \inlinecode{next\_arrangement\_with\_repeats()}.
\end{compactitem}

\Needspace*{3\baselineskip}
\textbf{Space Complexity}:
\begin{compactitem}
  \item $O(n)$ auxiliary for \inlinecode{next\_arrangement()}, \inlinecode{arrangement\_by\_rank()}, and \inlinecode{rank\_by\_arrangement()}.
  \item $O(1)$ auxiliary for \inlinecode{next\_arrangement\_with\_repeats()}.
\end{compactitem}

\begin{lstlisting}[style=cpp]
#include <algorithm>
#include <cassert>
#include <cstdint>
#include <numeric>
#include <vector>

bool next_arrangement(int n, std::vector<int> &a) {
  int k = static_cast<int>(a.size());
  assert(n >= 0 && k <= n);
  std::vector<char> used(n);
  for (int i = 0; i < k; i++) {
    assert(0 <= a[i] && a[i] < n && !used[a[i]]);
    used[a[i]] = true;
  }
  for (int i = k - 1; i >= 0; i--) {
    used[a[i]] = false;
    for (int j = a[i] + 1; j < n; j++) {
      if (!used[j]) {
        a[i++] = j;
        used[j] = true;
        for (int x = 0; i < k; x++) {
          if (!used[x]) {
            a[i++] = x;
          }
        }
        return true;
      }
    }
  }
  return false;
}

int64_t n_permute_k(int n, int k) {
  assert(n >= 0 && 0 <= k && k <= n);
  int64_t res = 1;
  for (int i = 0; i < k; i++) {
    res *= n - i;
  }
  return res;
}

std::vector<int> arrangement_by_rank(int n, int k, int64_t r) {
  assert(0 <= r && r < n_permute_k(n, k));
  std::vector<int> values(n), res(k);
  std::iota(values.begin(), values.end(), 0);
  for (int i = 0; i < k; i++) {
    int64_t count = n_permute_k(n - 1 - i, k - 1 - i);
    int pos = r / count;
    res[i] = values[pos];
    values.erase(values.begin() + pos);
    r %= count;
  }
  return res;
}

int64_t rank_by_arrangement(int n, const std::vector<int> &a) {
  int k = static_cast<int>(a.size());
  assert(n >= 0 && k <= n);
  assert(std::all_of(a.begin(), a.end(), [n](int x) { return 0 <= x && x < n; }));
  int64_t res = 0;
  std::vector<char> used(n);
  for (int i = 0; i < k; i++) {
    assert(!used[a[i]]);
    int count = 0;
    for (int j = 0; j < a[i]; j++) {
      if (!used[j]) {
        count++;
      }
    }
    res += count * n_permute_k(n - i - 1, k - i - 1);
    used[a[i]] = true;
  }
  return res;
}

bool next_arrangement_with_repeats(int n, std::vector<int> &a) {
  int k = static_cast<int>(a.size());
  assert(n >= 0 && std::all_of(a.begin(), a.end(), [n](int x) { return 0 <= x && x < n; }));
  for (int i = k - 1; i >= 0; i--) {
    if (a[i] < n - 1) {
      a[i]++;
      std::fill(a.begin() + i + 1, a.end(), 0);
      return true;
    }
  }
  return false;
}
\end{lstlisting}
\Needspace*{4\baselineskip}
\textbf{Example Usage}\par
\begin{lstlisting}[style=cpp]
#include <cassert>
#include <iostream>
using namespace std;

template<typename It>
void print_range(It lo, It hi) {
  cout << "{";
  for (; lo != hi; ++lo) {
    cout << *lo << (lo == hi - 1 ? "" : ",");
  }
  cout << "} ";
}

int main() {
  {
    int n = 4, k = 3;
    vector<int> a{0, 1, 2};
    cout << n << " permute " << k << " arrangements:" << endl;
    int count = 0;
    do {
      print_range(a.begin(), a.end());
      assert(a == arrangement_by_rank(n, k, count));
      assert(rank_by_arrangement(n, a) == count);
      count++;
      if (count == 10 || count == 20) {
        cout << endl;
      }
    } while (next_arrangement(n, a));
    assert(count == 24);
    cout << endl;
  }
  {
    int n = 4, k = 2;
    vector<int> a{0, 0};
    cout << endl << n << "^" << k << " arrangements with repeats:" << endl;
    int count = 0;
    do {
      print_range(a.begin(), a.end());
      count++;
      if (count == 13) {
        cout << endl;
      }
    } while (next_arrangement_with_repeats(n, a));
    assert(count == 16);
    cout << endl;
  }
  return 0;
}
\end{lstlisting}
\begin{exampleoutput}
4 permute 3 arrangements:
{0,1,2} {0,1,3} {0,2,1} {0,2,3} {0,3,1} {0,3,2} {1,0,2} {1,0,3} {1,2,0} {1,2,3}
{1,3,0} {1,3,2} {2,0,1} {2,0,3} {2,1,0} {2,1,3} {2,3,0} {2,3,1} {3,0,1} {3,0,2}
{3,1,0} {3,1,2} {3,2,0} {3,2,1}

4^2 arrangements with repeats:
{0,0} {0,1} {0,2} {0,3} {1,0} {1,1} {1,2} {1,3} {2,0} {2,1} {2,2} {2,3} {3,0}
{3,1} {3,2} {3,3}
\end{exampleoutput}
\subsection{Enumerating Permutations}
\label{sec:6.2.3}
These routines enumerate and analyze reorderings of a sequence of $n$ elements; the values need not be distinct except where noted.

The lexicographic successor is found by locating the rightmost ascent, increasing that position by the smallest possible amount, then reversing the suffix into its minimum order. Ranking and unranking permutations use the factorial number system: each position contributes the number of unused smaller values before the chosen one, scaled by the number of possible suffix permutations. The cycle decomposition views a permutation as a function on indices and follows each unvisited orbit until it returns to its start. Sorting by swaps is that same decomposition read differently: a cycle of length $k$ is resolved by $k - 1$ swaps and no fewer, so the whole array needs $n$ minus its number of cycles.

\begin{compactitem}
  \item \inlinecode{next\_permutation2(lo, hi, comp = std::less\textless{}\textgreater{}())} is analogous to \inlinecode{std::next\_permutation()}, taking two BidirectionalIterators as a range $[{\inlinecodemath{lo}}, {\inlinecodemath{hi}})$ for which the function tries to rearrange to the next lexicographically greater permutation according to \inlinecode{comp}. The function returns true if such a permutation exists, or false if the range is already in reverse comparator order, in which case it is rearranged into comparator order.
  \item \inlinecode{next\_permutation(a, comp = std::less\textless{}\textgreater{}())} is analogous to \inlinecode{next\_permutation2()}, except that it takes a vector instead of a range.
  \item \inlinecode{next\_permutation\_mask(x)} returns the next integer having the same number of 1-bits. Treating each 1-bit as whether to take its corresponding item generates combinations of a set of $n$ items. It returns $0$ if no successor fits in \inlinecode{uint64\_t}.
  \item \inlinecode{permutation\_by\_rank(n, r)} returns the permutation of the integers in the range $[0, {\inlinecodemath{n}})$ which is lexicographically ranked $r$, where $r$ is a 0-based rank in the range $[0, n!)$.
  \item \inlinecode{rank\_by\_permutation(a)} returns an integer representing the 0-based rank of permutation \inlinecode{a}, which must be a permutation of the integers $[0, n)$.
  \item \inlinecode{permutation\_cycles(a)} returns the orbits of the index mapping \inlinecode{i} $\mapsto$ \inlinecode{a[i]}. For example, $\{3, 1, 0, 2\}$ decomposes into cycles $\{0, 3, 2\}$ and $\{1\}$ because following the mapping from index $0$ gives $0 \to 3 \to 2 \to 0$, while index $1$ maps to itself.
  \item \inlinecode{min\_swaps\_to\_sort(a)} returns the fewest swaps of any two elements that sort \inlinecode{a}, whose values must be distinct. This is unrelated to the inversion count of section \secref{1.1.5}, which measures the number of swaps when only adjacent elements may be exchanged.
\end{compactitem}

The two ranking operations require $n \leq 20$, so every possible rank fits in \inlinecode{int64\_t}.

\Needspace*{3\baselineskip}
\textbf{Time Complexity}:
\begin{compactitem}
  \item $O(n)$ per call to \inlinecode{next\_permutation2()} and \inlinecode{next\_permutation()}, where $n$ is the input size.
  \item $O(n^{2})$ per call to \inlinecode{permutation\_by\_rank()} and \inlinecode{rank\_by\_permutation()}.
  \item $O(1)$ per call to \inlinecode{next\_permutation\_mask()}.
  \item $O(n)$ per call to \inlinecode{permutation\_cycles()}.
  \item $O(n \log n)$ per call to \inlinecode{min\_swaps\_to\_sort()}, from ranking the values.
\end{compactitem}

\Needspace*{3\baselineskip}
\textbf{Space Complexity}:
\begin{compactitem}
  \item $O(1)$ auxiliary for \inlinecode{next\_permutation2()} and \inlinecode{next\_permutation()}.
  \item $O(n)$ auxiliary for \inlinecode{permutation\_by\_rank()}, \inlinecode{rank\_by\_permutation()}, \inlinecode{permutation\_cycles()}, and \inlinecode{min\_swaps\_to\_sort()}.
\end{compactitem}

\begin{lstlisting}[style=cpp]
#include <algorithm>
#include <cassert>
#include <cstdint>
#include <functional>
#include <numeric>
#include <utility>
#include <vector>

template<typename It, typename Compare = std::less<>>
bool next_permutation2(It lo, It hi, Compare comp = Compare{}) {
  if (lo == hi) {
    return false;
  }
  It i = lo;
  if (++i == hi) {
    return false;
  }
  i = hi;
  --i;
  while (true) {
    It j = i;
    if (comp(*--i, *j)) {
      It k = hi;
      while (!comp(*i, *--k)) {
      }
      std::iter_swap(i, k);
      std::reverse(j, hi);
      return true;
    }
    if (i == lo) {
      std::reverse(lo, hi);
      return false;
    }
  }
}

template<typename T, typename Compare = std::less<>>
bool next_permutation(std::vector<T> &a, Compare comp = Compare{}) {
  int n = static_cast<int>(a.size());
  for (int i = n - 2; i >= 0; i--) {
    if (comp(a[i], a[i + 1])) {
      for (int j = n - 1;; j--) {
        if (comp(a[i], a[j])) {
          std::swap(a[i++], a[j]);
          for (j = n - 1; i < j; i++, j--) {
            std::swap(a[i], a[j]);
          }
          return true;
        }
      }
    }
  }
  std::reverse(a.begin(), a.end());
  return false;
}

uint64_t next_permutation_mask(uint64_t x) {
  if (x == 0) {
    return 0;
  }
  uint64_t s = x & -x, r = x + s;
  if (r == 0) {
    return 0;
  }
  return r | (((x ^ r) >> 2) / s);
}

std::vector<int> permutation_by_rank(int n, int64_t r) {
  assert(0 <= n && n <= 20 && r >= 0);
  std::vector<int64_t> factorial(n + 1, 1);
  std::vector<int> values(n), res(n);
  for (int i = 1; i <= n; i++) {
    factorial[i] = i * factorial[i - 1];
  }
  assert(r < factorial[n]);
  std::iota(values.begin(), values.end(), 0);
  for (int i = 0; i < n; i++) {
    int pos = r / factorial[n - 1 - i];
    res[i] = values[pos];
    values.erase(values.begin() + pos);
    r %= factorial[n - 1 - i];
  }
  return res;
}

int64_t rank_by_permutation(const std::vector<int> &a) {
  int n = static_cast<int>(a.size());
  assert(n <= 20);
  std::vector<char> seen(n);
  for (int x : a) {
    assert(0 <= x && x < n && !seen[x]);
    seen[x] = true;
  }
  if (n == 0) {
    return 0;
  }
  std::vector<int64_t> factorial(n);
  factorial[0] = 1;
  for (int i = 1; i < n; i++) {
    factorial[i] = i * factorial[i - 1];
  }
  int64_t res = 0;
  for (int i = 0; i < n; i++) {
    int v = a[i];
    for (int j = 0; j < i; j++) {
      if (a[j] < a[i]) {
        v--;
      }
    }
    res += v * factorial[n - 1 - i];
  }
  return res;
}

using Cycles = std::vector<std::vector<int>>;

Cycles permutation_cycles(const std::vector<int> &a) {
  int n = static_cast<int>(a.size());
  std::vector<char> visit(n);
  for (int x : a) {
    assert(0 <= x && x < n && !visit[x]);
    visit[x] = true;
  }
  std::fill(visit.begin(), visit.end(), false);
  Cycles res;
  for (int i = 0; i < n; i++) {
    if (!visit[i]) {
      int j = i;
      std::vector<int> curr;
      do {
        curr.push_back(j);
        visit[j] = true;
        j = a[j];
      } while (j != i);
      res.push_back(std::move(curr));
    }
  }
  return res;
}

int min_swaps_to_sort(const std::vector<int> &a) {
  int n = static_cast<int>(a.size());
  std::vector<int> order(n);
  std::iota(order.begin(), order.end(), 0);
  std::sort(order.begin(), order.end(), [&](int i, int j) { return a[i] < a[j]; });
  assert(std::adjacent_find(order.begin(), order.end(), [&](int i, int j) {
           return a[i] == a[j];
         }) == order.end());
  return n - static_cast<int>(permutation_cycles(order).size());
}
\end{lstlisting}
\Needspace*{4\baselineskip}
\textbf{Example Usage}\par
\begin{lstlisting}[style=cpp]
#include <bitset>
#include <cassert>
#include <iostream>
using namespace std;

template<typename It>
void print_range(It lo, It hi) {
  cout << "{";
  for (; lo != hi; ++lo) {
    cout << *lo << (lo == hi - 1 ? "" : ",");
  }
  cout << "} ";
}

int main() {
  {
    const int n = 4;
    vector<int> a{0, 1, 2, 3}, b = a, c = a;
    cout << "Permutations of [0, " << n << "):" << endl;
    int count = 0;
    do {
      print_range(a.begin(), a.end());
      assert(b == a);
      assert(c == a);
      assert(permutation_by_rank(n, count) == a);
      assert(rank_by_permutation(a) == count);
      count++;
      if (count == 8 || count == 16) {
        cout << endl;
      }
      std::next_permutation(b.begin(), b.end());
      next_permutation(c);
    } while (next_permutation(a));
    assert(count == 24);
    assert((a == vector<int>{0, 1, 2, 3}));
    cout << endl;
  }
  {  // Permutations of binary digits.
    const int n = 5;
    cout << "\nPermutations of 2 zeros and 3 ones:" << endl;
    uint64_t lo = bitset<5>(string("00111")).to_ullong();
    uint64_t hi = bitset<6>(string("100011")).to_ullong();
    int count = 0;
    do {
      cout << bitset<n>(lo).to_string() << " ";
      count++;
    } while ((lo = next_permutation_mask(lo)) != hi);
    assert(count == 10);
    assert(next_permutation_mask(1ULL << 63) == 0);
    cout << endl;
  }
  {
    vector<int> a{3, 2, 1}, b = a;
    assert(next_permutation(a, greater<int>()));
    assert(next_permutation2(b.begin(), b.end(), greater<int>()));
    assert((a == vector<int>{3, 1, 2} && b == a));
  }
  {  // Decomposition into cycles.
    vector<int> a{3, 1, 0, 2};
    cout << "\nDecomposition of {3,1,0,2} into cycles:" << endl;
    Cycles c = permutation_cycles(a);
    assert((c == Cycles{{0, 3, 2}, {1}}));
    for (const auto &cycle : c) {
      print_range(cycle.begin(), cycle.end());
    }
    cout << endl;
    // Two cycles over four elements, so two swaps sort it.
    assert(min_swaps_to_sort(a) == 2);
    assert(min_swaps_to_sort({50, 20, 40, 10}) == 1);  // Values need not be a permutation.
    assert(min_swaps_to_sort({1, 2, 3}) == 0);
    assert(min_swaps_to_sort({}) == 0);
  }
  return 0;
}
\end{lstlisting}
\begin{exampleoutput}
Permutations of [0, 4):
{0,1,2,3} {0,1,3,2} {0,2,1,3} {0,2,3,1} {0,3,1,2} {0,3,2,1} {1,0,2,3} {1,0,3,2}
{1,2,0,3} {1,2,3,0} {1,3,0,2} {1,3,2,0} {2,0,1,3} {2,0,3,1} {2,1,0,3} {2,1,3,0}
{2,3,0,1} {2,3,1,0} {3,0,1,2} {3,0,2,1} {3,1,0,2} {3,1,2,0} {3,2,0,1} {3,2,1,0}

Permutations of 2 zeros and 3 ones:
00111 01011 01101 01110 10011 10101 10110 11001 11010 11100

Decomposition of {3,1,0,2} into cycles:
{0,3,2} {1}
\end{exampleoutput}
\subsection{Enumerating Combinations}
\label{sec:6.2.4}
A combination is a subset of size $k$ chosen from a total of $n$ (not necessarily distinct) elements, where order does not matter.

The lexicographic successor advances the rightmost chosen position that can still move right, then packs every later position as far left as possible. Ranking and unranking use the combinatorial number system: at each position, count how many combinations would be skipped by choosing a smaller next value, then either add that count to the rank or subtract it while searching for the requested rank. Bitmask successors use the same order as increasing integers with a fixed popcount.

\begin{compactitem}
  \item \inlinecode{next\_combination(lo, mid, hi, comp = std::less\textless{}\textgreater{}())} takes random-access iterators \inlinecode{lo}, \inlinecode{mid}, and \inlinecode{hi} as a range $[{\inlinecodemath{lo}}, {\inlinecodemath{hi}})$ of $n$ elements for which the function will rearrange such that the $k$ elements in $[{\inlinecodemath{lo}}, {\inlinecodemath{mid}})$ become the next lexicographically greater combination. The function returns true if such a combination exists, or false if $[{\inlinecodemath{lo}}, {\inlinecodemath{mid}})$ already consists of the lexicographically greatest combination of the elements in $[{\inlinecodemath{lo}}, {\inlinecodemath{hi}})$, in which case the range is reset to its first combination. The comparator \inlinecode{comp} defines the element ordering.
  \item \inlinecode{next\_combination(n, a)} rearranges \inlinecode{a} to become the next lexicographically greater combination of distinct integers in the range $[0, {\inlinecodemath{n}})$. The vector \inlinecode{a} must be sorted and contain distinct integers in the range $[0, {\inlinecodemath{n}})$.
  \item \inlinecode{next\_combination\_mask(x)} interprets the bits of an integer \inlinecode{x} as a mask with 1-bits specifying the chosen items for a combination and returns the mask of the next lexicographically greater combination (that is, the lowest integer greater than \inlinecode{x} with the same number of 1 bits). Note that this does not generate combinations in the same order as \inlinecode{next\_combination()}, nor does it work if the corresponding $n$ items are not distinct (in that case, duplicate combinations will be generated). It returns $0$ if no successor fits in \inlinecode{uint64\_t}.
  \item \inlinecode{combination\_by\_rank(n, k, r)} returns the combination of $k$ distinct integers in the range $[0, {\inlinecodemath{n}})$ that is lexicographically ranked $r$, where $r$ is a 0-based rank in the range $[0, \binom{n}{k})$.
  \item \inlinecode{rank\_by\_combination(n, a)} returns an integer representing the 0-based rank of combination \inlinecode{a}, which must contain sorted distinct integers in $[0, {\inlinecodemath{n}})$.
  \item \inlinecode{next\_combination\_with\_repeats(n, a)} rearranges \inlinecode{a} to become the next lexicographically greater combination of not necessarily distinct integers in the range $[0, {\inlinecodemath{n}})$. The vector \inlinecode{a} must be sorted. Note that there is a total of $n \mathbin{\text{multichoose}} k$ combinations if repetition is allowed, where $n \mathbin{\text{multichoose}} k = \binom{n + k - 1}{k}$.
\end{compactitem}

Overflow warning: Exact counts and ranks used by the ranking operations must fit in \inlinecode{int64\_t}.

\Needspace*{3\baselineskip}
\textbf{Time Complexity}:
\begin{compactitem}
  \item $O(n)$ per call to \inlinecode{next\_combination(lo, mid, hi)}, where $n$ is the distance between \inlinecode{lo} and \inlinecode{hi}.
  \item $O(k)$ per call to \inlinecode{next\_combination(n, a)} and \inlinecode{next\_combination\_with\_repeats(n, a)}, where $k$ is the size of \inlinecode{a}.
  \item $O(1)$ per call to \inlinecode{next\_combination\_mask()}.
  \item $O(n \cdot k)$ per call to \inlinecode{combination\_by\_rank()} and \inlinecode{rank\_by\_combination()}.
\end{compactitem}

\Needspace*{3\baselineskip}
\textbf{Space Complexity}:
\begin{compactitem}
  \item $O(k)$ auxiliary for \inlinecode{combination\_by\_rank(n, k, r)}.
  \item $O(1)$ auxiliary for all other operations.
\end{compactitem}

\begin{lstlisting}[style=cpp]
#include <algorithm>
#include <cassert>
#include <cstdint>
#include <functional>
#include <iterator>
#include <numeric>
#include <utility>
#include <vector>

template<typename It, typename Compare = std::less<>>
bool next_combination(It lo, It mid, It hi, Compare comp = Compare{}) {
  using T = typename std::iterator_traits<It>::value_type;
  if (lo == mid || mid == hi) {
    return false;
  }
  It l = mid - 1, h = hi - 1;
  int len1 = 1, len2 = 1;
  while (l != lo && !comp(*l, *h)) {
    --l;
    ++len1;
  }
  if (l == lo && !comp(*l, *h)) {
    std::rotate(lo, mid, hi);
    return false;
  }
  for (; mid < h; ++len2) {
    if (!comp(*l, *--h)) {
      ++h;
      break;
    }
  }
  if (len1 == 1 || len2 == 1) {
    std::iter_swap(l, h);
  } else if (len1 == len2) {
    std::swap_ranges(l, mid, h);
  } else {
    std::iter_swap(l++, h++);
    int total = (--len1) + (--len2), gcd = total;
    for (int i = len1; i != 0;) {
      std::swap(gcd %= i, i);
    }
    int skip = total / gcd - 1;
    for (int i = 0; i < gcd; i++) {
      It curr = (i < len1) ? (l + i) : (h + (i - len1));
      int k = i;
      T prev = *curr;
      for (int j = 0; j < skip; j++) {
        k = (k + len1) % total;
        It next = (k < len1) ? (l + k) : (h + (k - len1));
        *curr = *next;
        curr = next;
      }
      *curr = prev;
    }
  }
  return true;
}

bool next_combination(int n, std::vector<int> &a) {
  int k = static_cast<int>(a.size());
  for (int i = k - 1; i >= 0; i--) {
    if (a[i] < n - k + i) {
      a[i]++;
      while (++i < k) {
        a[i] = a[i - 1] + 1;
      }
      return true;
    }
  }
  return false;
}

uint64_t next_combination_mask(uint64_t x) {
  if (x == 0) {
    return 0;
  }
  uint64_t s = x & -x, r = x + s;
  if (r == 0) {
    return 0;
  }
  return r | (((x ^ r) >> 2) / s);
}

int64_t n_choose_k(int64_t n, int64_t k) {
  assert(n >= 0 && 0 <= k && k <= n);
  if (k > n - k) {
    k = n - k;
  }
  int64_t res = 1;
  for (int i = 0; i < k; i++) {
    int64_t num = n - i, den = i + 1;
    int64_t g = std::gcd(num, den);
    num /= g;
    den /= g;
    res /= den;
    res *= num;  // Overflow warning.
  }
  return res;
}

std::vector<int> combination_by_rank(int n, int k, int64_t r) {
  assert(0 <= r && r < n_choose_k(n, k));
  std::vector<int> res(k);
  int count = n;
  for (int i = 0; i < k; i++) {
    int j = 1;
    for (;; j++) {
      int64_t am = n_choose_k(count - j, k - 1 - i);
      if (r < am) {
        break;
      }
      r -= am;
    }
    res[i] = (i > 0) ? (res[i - 1] + j) : (j - 1);
    count -= j;
  }
  return res;
}

int64_t rank_by_combination(int n, const std::vector<int> &a) {
  int k = static_cast<int>(a.size());
  assert(n >= 0 && k <= n);
  assert(std::all_of(a.begin(), a.end(), [n](int x) { return 0 <= x && x < n; }));
  assert(std::adjacent_find(a.begin(), a.end(), std::greater_equal<int>()) == a.end());
  int64_t res = 0;
  int prev = -1;
  for (int i = 0; i < k; i++) {
    for (int j = prev + 1; j < a[i]; j++) {
      res += n_choose_k(n - 1 - j, k - 1 - i);
    }
    prev = a[i];
  }
  return res;
}

bool next_combination_with_repeats(int n, std::vector<int> &a) {
  int k = static_cast<int>(a.size());
  for (int i = k - 1; i >= 0; i--) {
    if (a[i] < n - 1) {
      for (++a[i]; ++i < k;) {
        a[i] = a[i - 1];
      }
      return true;
    }
  }
  return false;
}
\end{lstlisting}
\Needspace*{4\baselineskip}
\textbf{Example Usage}\par
\begin{lstlisting}[style=cpp]
#include <cassert>
#include <iostream>
using namespace std;

template<typename It>
void print_range(It lo, It hi) {
  cout << "{";
  for (; lo != hi; ++lo) {
    cout << *lo << (lo == hi - 1 ? "" : ",");
  }
  cout << "} ";
}

int main() {
  {
    int k = 3;
    string s = "11234";
    cout << "\"" << s << "\" choose " << k << ":" << endl;
    int count = 0;
    do {
      cout << s.substr(0, k) << " ";
      count++;
    } while (next_combination(s.begin(), s.begin() + k, s.end()));
    assert(count == 7);  // Duplicate '1' values collapse equal output strings.
    cout << endl;
  }
  {  // Unordered combinations using masks.
    int n = 5, k = 3;
    string char_set = "abcde";  // Must be distinct.
    cout << "\n\"" << char_set << "\" choose " << k << " with masks:" << endl;
    uint64_t mask = (1ULL << k) - 1, limit = 1ULL << n;
    int count = 0;
    do {
      for (int i = 0; i < n; i++) {
        if ((mask >> i) & 1) {
          cout << char_set[i];
        }
      }
      cout << " ";
      count++;
      mask = next_combination_mask(mask);
    } while (mask < limit);
    assert(count == 10);
    assert(next_combination_mask(1ULL << 63) == 0);
    cout << endl;
  }
  {
    string s = "4321";
    assert(next_combination(s.begin(), s.begin() + 2, s.end(), greater<char>()));
    assert(s.substr(0, 2) == "42");
  }
  {  // Combinations of distinct integers in [0, n).
    int n = 5, k = 3;
    vector<int> a{0, 1, 2};
    cout << endl << n << " choose " << k << ":" << endl;
    int count = 0;
    do {
      print_range(a.begin(), a.end());
      assert(a == combination_by_rank(n, k, count));
      assert(rank_by_combination(n, a) == count);
      count++;
    } while (next_combination(n, a));
    assert(count == 10);
    assert(n_choose_k(66, 33) == 7219428434016265740LL);
    cout << endl;
  }
  {  // Combinations with repeats.
    int n = 3, k = 2;
    vector<int> a{0, 0};
    cout << endl << n << " multichoose " << k << ":" << endl;
    int count = 0;
    do {
      print_range(a.begin(), a.end());
      count++;
    } while (next_combination_with_repeats(n, a));
    assert(count == 6);
    cout << endl;
  }
  return 0;
}
\end{lstlisting}
\begin{exampleoutput}
"11234" choose 3:
112 113 114 123 124 134 234

"abcde" choose 3 with masks:
abc abd acd bcd abe ace bce ade bde cde

5 choose 3:
{0,1,2} {0,1,3} {0,1,4} {0,2,3} {0,2,4} {0,3,4} {1,2,3} {1,2,4} {1,3,4} {2,3,4}

3 multichoose 2:
{0,0} {0,1} {0,2} {1,1} {1,2} {2,2}
\end{exampleoutput}
\subsection{Enumerating Partitions}
\label{sec:6.2.5}
A partition of a natural number $n$ is a way to write $n$ as a sum of positive integers where the order of the addends does not matter.

The lexicographic successor edits the suffix of a non-increasing partition: increase the rightmost part that can grow without breaking the order, then refill the remaining sum with ones. Ranking and unranking use a dynamic-programming table that counts partitions of a remaining sum by their next part. This lets the algorithms skip whole blocks of partitions sharing a common prefix.

\begin{compactitem}
  \item \inlinecode{next\_partition(p)} takes a reference to a vector \inlinecode{p} of positive integers as a partition of $n$ for which the function will re-assign to become the next lexicographically greater partition. The function returns true if such a partition exists, or false if \inlinecode{p} already consists of the lexicographically greatest partition (i.e. the single integer $n$).
  \item \inlinecode{partition\_by\_rank(n, r)} returns the partition of $n$ that is lexicographically ranked $r$ if addends in each partition were sorted in non-increasing order, where $r$ is a 0-based rank in the range $[0, p(n))$ where $p(n)$ is the number of partitions of $n$.
  \item \inlinecode{rank\_by\_partition(p)} returns an integer representing the 0-based rank of the partition specified by vector \inlinecode{p}, which must consist of positive integers sorted in non-increasing order.
  \item \inlinecode{generate\_increasing\_partitions(n, f)} calls the function \inlinecode{f(lo, hi)} on strictly increasing partitions of $n$ in lexicographically increasing order of partition, where \inlinecode{lo} and \inlinecode{hi} are random-access iterators to a range $[{\inlinecodemath{lo}}, {\inlinecodemath{hi}})$ of integers. Note that non-strictly increasing partitions like $\{1, 1, 1, 1\}$ are skipped.
\end{compactitem}

Overflow warning: Exact partition counts and ranks must fit in \inlinecode{int64\_t}.

\Needspace*{3\baselineskip}
\textbf{Time Complexity}:
\begin{compactitem}
  \item $O(n)$ per call to \inlinecode{next\_partition()}.
  \item $O(n^{2})$ per call to \inlinecode{partition\_by\_rank()} and \inlinecode{rank\_by\_partition()}.
  \item $O(p(n))$ per call to \inlinecode{generate\_increasing\_partitions()}, where $p(n)$ is the number of partitions of $n$.
\end{compactitem}

\Needspace*{3\baselineskip}
\textbf{Space Complexity}:
\begin{compactitem}
  \item $O(1)$ auxiliary for \inlinecode{next\_partition()}.
  \item $O(n^{2})$ auxiliary heap space for \inlinecode{partition\_function()}, \inlinecode{partition\_by\_rank()}, and \inlinecode{rank\_by\_partition()}.
  \item $O(n)$ auxiliary stack space for \inlinecode{generate\_increasing\_partitions()}.
\end{compactitem}

\begin{lstlisting}[style=cpp]
#include <algorithm>
#include <cassert>
#include <cstdint>
#include <functional>
#include <numeric>
#include <vector>

bool next_partition(std::vector<int> &p) {
  int n = static_cast<int>(p.size());
  if (n <= 1) {
    return false;
  }
  int s = p[n - 1] - 1, i = n - 2;
  p.pop_back();
  for (; i > 0 && p[i] == p[i - 1]; i--) {
    s += p[i];
    p.pop_back();
  }
  for (p[i]++; s > 0; s--) {
    p.push_back(1);
  }
  return true;
}

int64_t partition_function(int a, int b) {
  assert(a >= 0 && 0 <= b && b <= a);
  static std::vector<std::vector<int64_t>> p(1, std::vector<int64_t>(1, 1));
  if (a >= static_cast<int>(p.size())) {
    int old = static_cast<int>(p.size());
    p.resize(a + 1);
    for (int i = 0; i < old; i++) {
      p[i].resize(a + 1);
    }
    for (int i = old; i <= a; i++) {
      p[i].resize(a + 1);
      for (int j = 1; j <= i; j++) {
        p[i][j] = p[i - 1][j - 1] + p[i - j][j];
      }
    }
  }
  return p[a][b];
}

std::vector<int> partition_by_rank(int n, int64_t r) {
  assert(n >= 0 && r >= 0);
  std::vector<int> res;
  for (int i = n, j; i > 0; i -= j) {
    for (j = 1;; j++) {
      int64_t count = partition_function(i, j);
      if (r < count) {
        break;
      }
      r -= count;
    }
    res.push_back(j);
  }
  return res;
}

int64_t rank_by_partition(const std::vector<int> &p) {
  assert(std::all_of(p.begin(), p.end(), [](int x) { return x > 0; }));
  assert(std::is_sorted(p.begin(), p.end(), std::greater<int>()));
  int64_t res = 0;
  int sum = std::accumulate(p.begin(), p.end(), 0);
  for (int x : p) {
    for (int j = 0; j < x; j++) {
      res += partition_function(sum, j);
    }
    sum -= x;
  }
  return res;
}

template<typename Fn>
void generate_increasing_partitions(int left, int prev, int i, std::vector<int> &p, Fn f) {
  if (left == 0) {
    f(p.begin(), p.begin() + i);
    return;
  }
  for (p[i] = prev + 1; p[i] <= left; p[i]++) {
    generate_increasing_partitions(left - p[i], p[i], i + 1, p, f);
  }
}

template<typename Fn>
void generate_increasing_partitions(int n, Fn f) {
  std::vector<int> p(n);
  generate_increasing_partitions(n, 0, 0, p, f);
}
\end{lstlisting}
\Needspace*{4\baselineskip}
\textbf{Example Usage}\par
\begin{lstlisting}[style=cpp]
#include <cassert>
#include <iostream>
using namespace std;

template<typename It>
void print_range(It lo, It hi) {
  cout << "{";
  for (; lo != hi; ++lo) {
    cout << *lo << (lo == hi - 1 ? "" : ",");
  }
  cout << "} ";
}

int main() {
  // Exercise incremental growth of the cached dynamic-programming table.
  assert(partition_function(3, 1) == 1);
  assert(partition_function(8, 3) == 5);
  {
    int n = 4;
    vector<int> a(n, 1);
    cout << "Partitions of " << n << ":" << endl;
    int count = 0;
    do {
      print_range(a.begin(), a.end());
      assert(a == partition_by_rank(n, count));
      assert(rank_by_partition(a) == count);
      count++;
    } while (next_partition(a));
    assert(count == 5);
    cout << endl;
  }
  {
    int n = 8;
    cout << "\nIncreasing partitions of " << n << ":" << endl;
    int count = 0;
    generate_increasing_partitions(n, [&](auto lo, auto hi) {
      print_range(lo, hi);
      count++;
    });
    assert(count == 6);
    cout << endl;
  }
  return 0;
}
\end{lstlisting}
\begin{exampleoutput}
Partitions of 4:
{1,1,1,1} {2,1,1} {2,2} {3,1} {4}

Increasing partitions of 8:
{1,2,5} {1,3,4} {1,7} {2,6} {3,5} {8}
\end{exampleoutput}
\subsection{Generic Combinatorial Enumeration}
\label{sec:6.2.6}
Ranks, unranks, and enumerates combinatorial sequences in lexicographic order using only a prefix counting oracle. This is useful when the objects are too numerous to precompute, but there is a simple dynamic-programming formula for the number of valid completions after a fixed prefix.

\inlinecode{Enumerator(range, length, count)} considers values in $[0, {\inlinecodemath{range}})$ at each position of a length-\inlinecode{length} sequence. The callback \inlinecode{count(prefix)} returns the number of valid full sequences beginning with \inlinecode{prefix}; it may capture a precomputed dynamic-programming table or other state. The factory functions below supply this oracle for the principal combinatorial families, while the class remains available directly for custom families. The range and length must be nonnegative; the factory arguments must satisfy $0 \leq k \leq n$.

\begin{compactitem}
  \item \inlinecode{arrangements(n, k)} returns an enumerator for the length-$k$ arrangements of $[0, n)$.
  \item \inlinecode{permutations(n)} returns an enumerator for the permutations of $[0, n)$.
  \item \inlinecode{combinations(n, k)} returns an enumerator for the size-$k$ subsets of $[0, n)$, represented as strictly increasing sequences.
  \item \inlinecode{partitions(n)} returns an enumerator for the additive partitions of $n$, represented as non-increasing length-$n$ sequences padded with zeros.
  \item \inlinecode{to\_rank(a)} returns the 0-based rank of the valid combinatorial sequence \inlinecode{a}.
  \item \inlinecode{from\_rank(r)} returns a combinatorial sequence of integers that is lexicographically ranked \inlinecode{r}, where \inlinecode{r} is a 0-based rank in the range $[0, {\inlinecodemath{total\_count()}})$.
  \item \inlinecode{total\_count()} returns the number of valid sequences.
  \item \inlinecode{enumerate(f)} calls the function \inlinecode{f(a)} on every specified combinatorial sequence \inlinecode{a} in lexicographically increasing order.
\end{compactitem}

Overflow warning: All exact counts and ranks must fit in \inlinecode{int64\_t}.

\Needspace*{3\baselineskip}
\textbf{Time Complexity}:
\begin{compactitem}
  \item $O(A \cdot L \cdot C)$ per call to \inlinecode{to\_rank()} and \inlinecode{from\_rank()}, where $A$ is \inlinecode{range}, $L$ is \inlinecode{length}, and $C$ is the cost of \inlinecode{count()}.
  \item $O(T \cdot A \cdot L \cdot C)$ per call to \inlinecode{enumerate()}, excluding the cost of the callback \inlinecode{f}, where $T$ is \inlinecode{total\_count()}.
  \item $O(A \cdot (\min\left(L, A - L\right) + 1))$ per call to \inlinecode{combinations()} and $O(L^{2})$ per call to \inlinecode{partitions()}; the other factories take $O(1)$ time per call.
\end{compactitem}

\Needspace*{3\baselineskip}
\textbf{Space Complexity}:
\begin{compactitem}
  \item $O(L)$ auxiliary for \inlinecode{to\_rank()}, \inlinecode{from\_rank()}, and \inlinecode{enumerate()}, excluding output storage.
  \item $O(S)$ object storage, where $S$ is the state captured by \inlinecode{count()}.
  \item $O(A \cdot (\min\left(L, A - L\right) + 1))$ object storage for \inlinecode{combinations()} and $O(L^{2})$ for \inlinecode{partitions()}.
\end{compactitem}

\begin{lstlisting}[style=cpp]
#include <algorithm>
#include <cassert>
#include <cstdint>
#include <numeric>
#include <utility>
#include <vector>

template<typename Count>
class Enumerator {
  int range, length;
  Count count;

 public:
  Enumerator(int range, int length, Count count)
      : range(range), length(length), count(std::move(count)) {
    assert(range >= 0 && length >= 0);
  }

  int64_t total_count() { return count({}); }

  int64_t to_rank(const std::vector<int> &a) {
    assert(static_cast<int>(a.size()) == length);
    assert(std::all_of(a.begin(), a.end(), [&](int x) { return 0 <= x && x < range; }));
    int64_t res = 0;
    std::vector<int> prefix;
    for (int i = 0; i < static_cast<int>(a.size()); i++) {
      prefix.push_back(0);
      for (; prefix.back() < a[i]; prefix.back()++) {
        res += count(prefix);
      }
    }
    return res;
  }

  std::vector<int> from_rank(int64_t r) {
    assert(0 <= r && r < total_count());
    std::vector<int> a;
    for (int i = 0; i < length; i++) {
      a.push_back(0);
      for (; a.back() < range; a.back()++) {
        int64_t curr = count(a);
        if (r < curr) {
          break;
        }
        r -= curr;
      }
    }
    return a;
  }

  // Accepts any callable f(a), including capturing lambdas and functors.
  template<typename Fn>
  void enumerate(Fn f) {
    int64_t total = total_count();
    for (int64_t i = 0; i < total; i++) {
      std::vector<int> curr = from_rank(i);
      f(curr);
    }
  }
};

auto arrangements(int n, int k) {
  assert(n >= 0 && 0 <= k && k <= n);
  auto count = [=](const std::vector<int> &prefix) -> int64_t {
    int len = static_cast<int>(prefix.size());
    for (int i = 0; i < len - 1; i++) {
      if (prefix[i] == prefix[len - 1]) {
        return 0;
      }
    }
    int64_t res = 1;
    for (int i = 0; i < k - len; i++) {
      res *= n - len - i;  // Overflow warning.
    }
    return res;
  };
  return Enumerator(n, k, count);
}

auto permutations(int n) {
  return arrangements(n, n);
}

auto combinations(int n, int k) {
  assert(n >= 0 && 0 <= k && k <= n);
  std::vector<std::vector<int64_t>> table(n + 1, std::vector<int64_t>(std::min(k, n - k) + 1));
  int max_col = static_cast<int>(table[0].size()) - 1;
  for (int i = 0; i <= n; i++) {
    for (int j = 0; j <= std::min(i, max_col); j++) {
      table[i][j] = (j == 0) ? 1 : table[i - 1][j - 1] + table[i - 1][j];  // Overflow warning.
    }
  }
  auto count = [n, k, table = std::move(table)](const std::vector<int> &pref) -> int64_t {
    int len = static_cast<int>(pref.size());
    if (len >= 2 && pref[len - 1] <= pref[len - 2]) {
      return 0;
    }
    int rem = len == 0 ? n : n - pref.back() - 1;
    int choose = k - len;
    return (choose < 0 || choose > rem) ? 0 : table[rem][std::min(choose, rem - choose)];
  };
  return Enumerator(n, k, std::move(count));
}

auto partitions(int n) {
  assert(n >= 0);
  std::vector<std::vector<int64_t>> table(n + 1, std::vector<int64_t>(n + 1));
  std::fill(table[0].begin(), table[0].end(), 1);
  for (int sum = 1; sum <= n; sum++) {
    for (int max_part = 1; max_part <= n; max_part++) {
      table[sum][max_part] = table[sum][max_part - 1];
      if (max_part <= sum) {
        table[sum][max_part] += table[sum - max_part][max_part];  // Overflow warning.
      }
    }
  }
  auto count = [n, table = std::move(table)](const std::vector<int> &pref) -> int64_t {
    int len = static_cast<int>(pref.size());
    int sum = std::accumulate(pref.begin(), pref.end(), 0);
    if (sum == n) {
      return 1;
    }
    if (sum > n || (len > 0 && pref.back() == 0) || (len >= 2 && pref[len - 1] > pref[len - 2])) {
      return 0;
    }
    return table[n - sum][len == 0 ? n : pref.back()];
  };
  return Enumerator(n + 1, n, std::move(count));
}
\end{lstlisting}
\Needspace*{4\baselineskip}
\textbf{Example Usage}\par
\begin{lstlisting}[style=cpp]
#include <iostream>
using namespace std;

void print_sequence(const vector<int> &a) {
  cout << "{";
  for (int i = 0; i < static_cast<int>(a.size()); i++) {
    cout << a[i] << (i + 1 == static_cast<int>(a.size()) ? "" : ",");
  }
  cout << "} ";
}

int main() {
  cout << "3 permute 2 arrangements:" << endl;
  auto arr = arrangements(3, 2);
  int count = 0;
  arr.enumerate([&](const auto &a) {
    print_sequence(a);
    count++;
  });
  assert(count == 6);
  assert((arr.from_rank(5) == vector<int>{2, 1}));
  assert(arr.to_rank(vector<int>{2, 1}) == 5);

  cout << "\n\nPermutations of [0, 3):" << endl;
  auto perm = permutations(3);
  count = 0;
  perm.enumerate([&](const auto &a) {
    print_sequence(a);
    count++;
  });
  assert(count == 6);

  cout << "\n\n4 choose 3 combinations:" << endl;
  auto comb = combinations(4, 3);
  count = 0;
  comb.enumerate([&](const auto &a) {
    print_sequence(a);
    count++;
  });
  assert(count == 4);
  assert((comb.from_rank(3) == vector<int>{1, 2, 3}));
  assert(comb.to_rank(vector<int>{1, 2, 3}) == 3);
  assert(combinations(67, 1).total_count() == 67);
  assert(combinations(67, 67).total_count() == 1);
  assert(combinations(66, 33).total_count() == 7219428434016265740LL);

  cout << "\n\nPartitions of 4:" << endl;
  auto part = partitions(4);
  count = 0;
  part.enumerate([&](const auto &a) {
    print_sequence(a);
    count++;
  });
  assert(count == 5);

  int length = 4;
  auto count_completions = [length](const vector<int> &pref) -> int64_t {
    if (static_cast<int>(pref.size()) > length) {
      return 0;
    }
    for (int i = 0; i < static_cast<int>(pref.size()); i++) {
      if (pref[i] < 0 || pref[i] > 1 || (i > 0 && pref[i - 1] == 1 && pref[i] == 1)) {
        return 0;
      }
    }
    int64_t after_zero = 1, after_one = 1;
    for (int i = static_cast<int>(pref.size()); i < length; i++) {
      int64_t next_zero = after_zero + after_one;  // Overflow warning.
      after_one = after_zero;
      after_zero = next_zero;
    }
    return pref.empty() || pref.back() == 0 ? after_zero : after_one;
  };
  Enumerator e(2, length, count_completions);
  cout << "\n\nLength-4 binary strings without consecutive ones:" << endl;
  count = 0;
  e.enumerate([&](const auto &a) {
    print_sequence(a);
    count++;
  });
  assert(count == 8 && e.total_count() == 8);
  assert((e.from_rank(7) == vector<int>{1, 0, 1, 0}));
  assert(e.to_rank(vector<int>{1, 0, 1, 0}) == 7);
  cout << endl;
  return 0;
}
\end{lstlisting}
\begin{exampleoutput}
3 permute 2 arrangements:
{0,1} {0,2} {1,0} {1,2} {2,0} {2,1}

Permutations of [0, 3):
{0,1,2} {0,2,1} {1,0,2} {1,2,0} {2,0,1} {2,1,0}

4 choose 3 combinations:
{0,1,2} {0,1,3} {0,2,3} {1,2,3}

Partitions of 4:
{1,1,1,1} {2,1,1,0} {2,2,0,0} {3,1,0,0} {4,0,0,0}

Length-4 binary strings without consecutive ones:
{0,0,0,0} {0,0,0,1} {0,0,1,0} {0,1,0,0} {0,1,0,1} {1,0,0,0} {1,0,0,1} {1,0,1,0}
\end{exampleoutput}

\section{\texorpdfstring{Number Theory}{6.3 Number Theory}}
\setcounter{section}{3}
\setcounter{subsection}{0}
\subsection{Mod Helpers and Binary Exponentiation}
\label{sec:6.3.1}
Modular arithmetic replaces every value by its remainder modulo \inlinecode{m}, which keeps intermediate results bounded and is what makes counting problems with astronomically large answers tractable. The operations are straightforward except for two hazards, both handled below: a sum or product can overflow the underlying integer type before the remainder is taken, and a negative operand leaves a negative remainder under C++ truncating division. Binary exponentiation squares the base while walking the bits of the exponent, so a power costs $O(\log n)$ multiplications rather than $n$; the same loop computes powers in any monoid, which is how the matrix exponentiation of section \secref{6.5.1} and the linear recurrence of section \secref{6.6.6} are evaluated.

\begin{compactitem}
  \item \inlinecode{addmod(a, b, m)} and \inlinecode{submod(a, b, m)} respectively return addition and subtraction modulo \inlinecode{m}, each result in $[0, {\inlinecodemath{m}})$. Both operands may be negative or unreduced; they are normalized before arithmetic to avoid signed overflow. The modulus \inlinecode{m} must be positive.
  \item \inlinecode{mulmod(a, b, m)} returns \inlinecode{a} multiplied by \inlinecode{b}, modulo \inlinecode{m}. This is done in a way to avoid overflow: on compilers with \inlinecode{\_\_uint128\_t} it uses one wide product, while the portable fallback uses double-and-add multiplication. The fallback is slower by a logarithmic factor, but avoids relying on nonstandard 128-bit integers. Unlike \inlinecode{addmod}/\inlinecode{submod}, this takes unsigned operands and supports a full 64-bit modulus.
  \item \inlinecode{powmod(x, n, m)} returns \inlinecode{x} raised to the power \inlinecode{n}, modulo \inlinecode{m}.
\end{compactitem}

Paste these when a solution needs a handful of modular operations; when modular arithmetic pervades one instead, the \inlinecode{Modular} value type of section \secref{6.3.3} overloads the operators so expressions read normally. Only \inlinecode{mulmod()} has no counterpart there, since a value type multiplying in \inlinecode{int64\_t} overflows once the modulus passes about three billion, which is why it underlies the Miller-Rabin test and Pollard's rho of section \secref{6.3.9}.

\Needspace*{3\baselineskip}
\textbf{Time Complexity}:
\begin{compactitem}
  \item $O(1)$ per call to \inlinecode{addmod()} and \inlinecode{submod()}.
  \item $O(1)$ per call to \inlinecode{mulmod()} with \inlinecode{\_\_uint128\_t}, or $O(\log b)$ with the portable fallback, where $b$ is the second argument.
  \item $O(\log n)$ calls to \inlinecode{mulmod()} per call to \inlinecode{powmod(x, n, m)}.
\end{compactitem}

\textbf{Space Complexity}: $O(1)$ auxiliary for all operations.

\begin{lstlisting}[style=cpp]
#include <cassert>
#include <cstdint>

int64_t addmod(int64_t a, int64_t b, int64_t m) {
  assert(m > 0);
  if ((a %= m) < 0) a += m;
  if ((b %= m) < 0) b += m;
  return a >= m - b ? a - (m - b) : a + b;
}

int64_t submod(int64_t a, int64_t b, int64_t m) {
  assert(m > 0);
  if ((a %= m) < 0) a += m;
  if ((b %= m) < 0) b += m;
  return a >= b ? a - b : m - (b - a);
}

uint64_t mulmod(uint64_t a, uint64_t b, uint64_t m) {
  assert(m > 0);
#if defined(__SIZEOF_INT128__)
  return static_cast<uint64_t>(static_cast<__uint128_t>(a) * b % m);
#else
  uint64_t res = 0, cur = a % m;
  for (; b > 0; b >>= 1) {
    if (b & 1) {
      res = res >= m - cur ? res - (m - cur) : res + cur;
    }
    cur = cur >= m - cur ? cur - (m - cur) : cur + cur;
  }
  return res;
#endif
}

uint64_t powmod(uint64_t x, uint64_t n, uint64_t m) {
  assert(m > 0);
  uint64_t res = 1 % m;
  for (; n > 0; n >>= 1) {
    if (n & 1) {
      res = mulmod(res, x, m);
    }
    x = mulmod(x, x, m);
  }
  return res;
}
\end{lstlisting}
\Needspace*{4\baselineskip}
\textbf{Example Usage}\par
\begin{lstlisting}[style=cpp]
#include <cassert>

int main() {
  assert(addmod(7, 8, 10) == 5 && submod(2, 5, 10) == 7);
  // Negative and unreduced operands are normalized into [0, m).
  assert(addmod(-3, -4, 10) == 3 && submod(-3, 4, 10) == 3);
  assert(addmod(25, -7, 10) == 8);
  assert(addmod(INT64_MAX - 1, INT64_MAX - 1, INT64_MAX) == INT64_MAX - 2);

  assert(mulmod(INT64_MAX - 1, INT64_MAX - 1, INT64_MAX) == 1);
  assert(powmod(2, 10, 1000000007) == 1024);
  assert(powmod(2, 62, 1000000) == 387904);
  assert(powmod(10001, 10001, 100000) == 10001);
  return 0;
}
\end{lstlisting}
\subsection{GCD, LCM, Mod Inverse, Chinese Remainder}
\label{sec:6.3.2}
Core integer and modular-arithmetic operations built around the Euclidean algorithm. Extended GCD produces Bezout coefficients, which solve linear Diophantine equations and modular inverses, while the Chinese remainder routines combine compatible congruences into one residue class.

\begin{compactitem}
  \item \inlinecode{gcd(a, b)} returns the greatest common divisor of \inlinecode{a} and \inlinecode{b} using the Euclidean algorithm. This is mainly for educational purposes, as \inlinecode{std::gcd(a, b)} from \inlinecode{\textless{}numeric\textgreater{}} is available as of C++17 (\inlinecode{\_\_gcd(a, b)} from \inlinecode{\textless{}algorithm\textgreater{}} in C++14 and earlier).
  \item \inlinecode{lcm(a, b)} returns the least common multiple of \inlinecode{a} and \inlinecode{b}. This implementation is mainly for educational purposes, as \inlinecode{std::lcm(a, b)} from \inlinecode{\textless{}numeric\textgreater{}} is available as of C++17.
  \item \inlinecode{extended\_gcd(a, b)} returns a pair $(x, y)$ of integers such that $\gcd(a, b) = ax + by$.
  \item \inlinecode{diophantine(a, b, c, \&g, \&x, \&y)} solves the linear Diophantine equation $ax + by = c$, returning whether a solution exists (one does if and only if $\gcd(a, b)$ divides $c$). \inlinecode{g} is always set to $\gcd(a, b)$, while (\inlinecode{x}, \inlinecode{y}) is set only on success to a particular solution bounded by $\max(|a|, |b|, |c|)$ in magnitude. Every other solution is $({\inlinecodemath{x}} + t({\inlinecodemath{b}}/{\inlinecodemath{g}}), {\inlinecodemath{y}} - t({\inlinecodemath{a}}/{\inlinecodemath{g}}))$ for an integer $t$. A 128-bit intermediate is used where available to keep the scaling step overflow-free.
  \item \inlinecode{mod(a, m)} returns the least nonnegative residue of \inlinecode{a} modulo \inlinecode{m}, that is, the unique value in $[0, {\inlinecodemath{m}})$ congruent to \inlinecode{a}, where \inlinecode{m} must be positive. Unlike the C++ remainder operator \inlinecode{\%}, whose result follows the sign of \inlinecode{a}, the result is never negative.
  \item \inlinecode{mod\_inverse(a, m)} returns an integer $x$ such that $ax \equiv 1 \pmod m$, where the arguments must satisfy $m > 0$ and $\gcd(a, m) = 1$.
  \item \inlinecode{mod\_inverse\_table(p)} returns a vector \inlinecode{v} where \inlinecode{v[i]} is the modular inverse of \inlinecode{i} modulo the prime \inlinecode{p} for every valid index \inlinecode{i}.
  \item \inlinecode{crt(r1, m1, r2, m2, \&r, \&m)} merges the two congruences $x \equiv r_1 \pmod{m_1}$ and $x \equiv r_2 \pmod{m_2}$ for arbitrary moduli (not necessarily coprime). It returns whether the system is consistent, and on success sets \inlinecode{m} to \inlinecode{lcm(m\_1, m\_2)} and \inlinecode{r} to the unique solution in $[0, {\inlinecodemath{m}})$. Both moduli must be positive, and their least common multiple must fit in \inlinecode{int64\_t}. Fold it pairwise to merge more than two congruences.
  \item \inlinecode{garner\_restore(a, p)} returns the smallest nonnegative solution whose residue modulo \inlinecode{p[i]} is \inlinecode{a[i]} at every valid index. The entries of \inlinecode{p} must be pairwise coprime (unlike \inlinecode{crt}, which allows shared factors) and positive, and \inlinecode{a} and \inlinecode{p} must have equal sizes. The exact solution is unique modulo the product of all moduli, so that product and the final answer must fit in \inlinecode{int64\_t}.
  \item \inlinecode{garner\_restore\_mod(a, p, m)} returns that same CRT solution modulo \inlinecode{m}. This is the right variant when the product of the moduli is too large to fit in \inlinecode{int64\_t}; \inlinecode{m} must be positive.
\end{compactitem}

Overflow warning: Every intermediate and result of the templated signed-integer helpers must be representable by \inlinecode{Int}; in particular, the minimum value of \inlinecode{Int} cannot be negated or divided by $-1$.

\Needspace*{3\baselineskip}
\textbf{Time Complexity}:
\begin{compactitem}
  \item $O(\log M)$ per call to \inlinecode{gcd()}, \inlinecode{lcm()}, \inlinecode{extended\_gcd()}, \inlinecode{diophantine()}, \inlinecode{mod\_inverse()}, and \inlinecode{crt()}, where $M$ is the largest relevant input magnitude.
  \item $O(1)$ per call to \inlinecode{mod()}.
  \item $O(p)$ per call to \inlinecode{mod\_inverse\_table()}.
  \item $O(n^{2})$ per call to \inlinecode{garner\_restore()} and \inlinecode{garner\_restore\_mod()}.
\end{compactitem}

\Needspace*{3\baselineskip}
\textbf{Space Complexity}:
\begin{compactitem}
  \item $O(p)$ auxiliary for \inlinecode{mod\_inverse\_table()}.
  \item $O(n)$ auxiliary for \inlinecode{garner\_restore()} and \inlinecode{garner\_restore\_mod()}.
  \item $O(1)$ auxiliary for all other operations.
\end{compactitem}

\begin{lstlisting}[style=cpp]
#include <algorithm>
#include <cassert>
#include <cstdint>
#include <type_traits>
#include <utility>
#include <vector>

template<typename Int>
Int gcd(Int a, Int b) {
  while (b != 0) {
    Int t = b;
    b = a % b;
    a = t;
  }
  return (a < 0 ? -a : a);
}

template<typename Int>
Int lcm(Int a, Int b) {
  if (a == 0 || b == 0) {
    return 0;
  }
  Int res = a / gcd(a, b) * b;
  return (res < 0 ? -res : res);
}

template<typename Int>
std::pair<Int, Int> extended_gcd(Int a, Int b) {
  Int x = 1, y = 0, x1 = 0, y1 = 1;
  while (b != 0) {
    Int q = a / b, prev_x1 = x1, prev_y1 = y1, prev_b = b;
    x1 = x - q * x1;
    y1 = y - q * y1;
    b = a - q * b;
    x = prev_x1;
    y = prev_y1;
    a = prev_b;
  }
  return (a > 0) ? std::pair<Int, Int>{x, y} : std::pair<Int, Int>{-x, -y};
}

template<typename Int>
bool diophantine(Int a, Int b, Int c, Int *g, Int *x, Int *y) {
  if (a == 0 && b == 0) {
    *g = *x = *y = 0;
    return c == 0;
  }
  if (a == 0) {
    *g = (b < 0) ? -b : b;
    if (c % b != 0) {
      return false;
    }
    *x = 0;
    *y = c / b;
    return true;
  }
  if (b == 0) {
    *g = (a < 0) ? -a : a;
    if (c % a != 0) {
      return false;
    }
    *x = c / a;
    *y = 0;
    return true;
  }
  *g = gcd(a, b);
  if (c % *g != 0) {
    return false;
  }
  // Absorb the bulk of c into a*dx + b*dy, then scale the base solution by the small remainder so
  // the reported (x, y) stay bounded by max(|a|, |b|, |c|). The scaled product is taken modulo b
  // (resp. a) through a 128-bit intermediate where available, to avoid overflow.
  std::pair<Int, Int> base = extended_gcd(a, b);
  Int dx = c / a;
  c -= dx * a;
  Int dy = c / b;
  c -= dy * b;
  Int f = c / *g;
#if defined(__SIZEOF_INT128__)
  __extension__ typedef std::conditional_t<sizeof(Int) <= 4, int64_t, __int128> wide_t;
#else
  using wide_t = int64_t;
#endif
  *x = dx + static_cast<Int>(static_cast<wide_t>(base.first) * f % b);
  *y = dy + static_cast<Int>(static_cast<wide_t>(base.second) * f % a);
  return true;
}

template<typename Int>
Int mod(Int a, Int m) {
  assert(m > 0);
  Int r = a % m;
  return (r < 0) ? (r + m) : r;
}

template<typename Int>
Int mod_inverse(Int a, Int m) {
  assert(m > 0 && gcd(a, m) == 1);
  return mod(extended_gcd(a, m).first, m);
}

std::vector<int> mod_inverse_table(int p) {
  assert(p >= 2);
  std::vector<int> res(p);
  res[1] = 1;
  for (int i = 2; i < p; i++) {
    res[i] = (p - (p / i) * res[p % i] % p) % p;
  }
  return res;
}

bool crt(int64_t r1, int64_t m1, int64_t r2, int64_t m2, int64_t *r, int64_t *m) {
  r1 = mod(r1, m1);
  r2 = mod(r2, m2);
  int64_t g, x, y;
  if (!diophantine(m1, -m2, r2 - r1, &g, &x, &y)) {
    return false;
  }
  *m = m1 / g * m2;
#if defined(__SIZEOF_INT128__)
  __extension__ typedef __int128 int128_t;
  *r = static_cast<int64_t>(
      (static_cast<int128_t>(r1) + static_cast<int128_t>(m1) * mod(x, m2 / g)) % *m
  );
#else
  *r = mod(r1 + m1 * mod(x, m2 / g), *m);  // Overflow warning.
#endif
  return true;
}

std::vector<int64_t> garner_digits(const std::vector<int> &a, const std::vector<int> &p) {
  assert(a.size() == p.size());
  assert(std::all_of(p.begin(), p.end(), [](int m) { return m > 0; }));
  int n = static_cast<int>(a.size());
  std::vector<int64_t> x(n);
  for (int i = 0; i < n; i++) {
    x[i] = mod(static_cast<int64_t>(a[i]), static_cast<int64_t>(p[i]));
  }
  for (int i = 0; i < n; i++) {
    for (int j = 0; j < i; j++) {
      // Reduce mod p[i] each step; otherwise x[i] compounds to ~p^i and overflows int64_t.
      x[i] = mod_inverse(static_cast<int64_t>(p[j]), static_cast<int64_t>(p[i])) * (x[i] - x[j]);
      x[i] = (x[i] % p[i] + p[i]) % p[i];
    }
  }
  return x;
}

int64_t garner_restore(const std::vector<int> &a, const std::vector<int> &p) {
  assert(a.size() == p.size());
  int n = static_cast<int>(a.size());
  if (n == 0) {
    return 0;
  }
  std::vector<int64_t> x = garner_digits(a, p);
  int64_t res = x[0], m = 1;
  for (int i = 1; i < n; i++) {
    m *= p[i - 1];
    res += x[i] * m;
  }
  return res;
}

int64_t garner_restore_mod(const std::vector<int> &a, const std::vector<int> &p, int64_t m) {
  assert(a.size() == p.size() && m > 0);
  int n = static_cast<int>(a.size());
  if (n == 0) {
    return 0;
  }
  std::vector<int64_t> x = garner_digits(a, p);
  int64_t res = 0;
  for (int i = n - 1; i >= 0; i--) {
#if defined(__SIZEOF_INT128__)
    __extension__ typedef __int128 int128_t;
    res = static_cast<int64_t>((static_cast<int128_t>(res) * p[i] + x[i]) % m);
#else
    res = (res * p[i] + x[i]) % m;  // Requires the intermediate product to fit without int128.
#endif
  }
  return res;
}
\end{lstlisting}
\Needspace*{4\baselineskip}
\textbf{Example Usage}\par
\begin{lstlisting}[style=cpp]
#include <cassert>
using namespace std;

int main() {
  assert(mod(-5, 3) == 1);  // Negative dividends wrap up into [0, m).
  assert(mod(5, 3) == 2);
  assert(mod(-6, 3) == 0);
  {
    for (int a = -20; a <= 20; a++) {
      for (int b = -20; b <= 20; b++) {
        int g = gcd(a, b);
        auto [x, y] = extended_gcd(a, b);
        assert(g == a * x + b * y);
        if (g == 1 && b > 1) {
          assert(mod(a * mod_inverse(a, b), b) == 1);
        }
      }
    }
  }
  {
    int p = 17;
    auto res = mod_inverse_table(p);
    for (int i = 0; i < p; i++) {
      if (i > 0) {
        assert(mod(i * res[i], p) == 1);
      }
    }
  }
  {
    vector<int> a{2, 3, 1}, m{3, 4, 5};
    int x = garner_restore(a, m);
    assert(x == garner_restore_mod(a, m, 1000000007));
    int n = static_cast<int>(a.size());
    for (int i = 0; i < n; i++) {
      assert(mod(x, m[i]) == a[i]);
    }
    assert(x == 11);
  }
  assert(garner_restore(vector<int>{-1}, vector<int>{5}) == 4);
#if defined(__SIZEOF_INT128__)
  {  // Garner modulo another number works even when the product of CRT moduli is too large.
    vector<int> p{1000000007, 1000000009, 1000000033};
    __extension__ typedef __int128 int128_t;
    int128_t x = (int128_t{1} << 80) + 123456789;
    vector<int> a;
    for (int m : p) {
      a.push_back(static_cast<int>(x % m));
    }
    int64_t m = 998244353;
    assert(garner_restore_mod(a, p, m) == static_cast<int64_t>(x % m));
  }
#endif
  {  // Diophantine: exhaustively verify solvability and that solutions satisfy a*x + b*y == c.
    for (int a = -8; a <= 8; a++) {
      for (int b = -8; b <= 8; b++) {
        for (int c = -8; c <= 8; c++) {
          int g, x, y, gg = gcd(a, b);
          bool ok = diophantine(a, b, c, &g, &x, &y);
          assert(ok == (gg == 0 ? c == 0 : c % gg == 0));
          assert(g == gg);  // g is set even when no solution exists.
          if (ok) {
            assert(a * x + b * y == c);
          }
        }
      }
    }
  }
  {  // CRT: merge two congruences, including non-coprime and inconsistent moduli.
    for (int m1 = 1; m1 <= 12; m1++) {
      for (int m2 = 1; m2 <= 12; m2++) {
        for (int r1 = 0; r1 < m1; r1++) {
          for (int r2 = 0; r2 < m2; r2++) {
            int64_t r, m;
            bool ok = crt(r1, m1, r2, m2, &r, &m);
            int limit = lcm(m1, m2), want = -1;
            for (int v = 0; v < limit && want < 0; v++) {
              if (v % m1 == r1 && v % m2 == r2) {
                want = v;
              }
            }
            assert(ok == (want >= 0));
            if (ok) {
              assert(m == limit && r == want);
            }
          }
        }
      }
    }
  }
  return 0;
}
\end{lstlisting}
\subsection{Modular Integer and Combinatorics}
\label{sec:6.3.3}
Provides a modular integer value type together with lazy factorial-based combinatorics tables. Both are common contest helpers for dynamic programming, counting, polynomial operations, and any calculation whose answers are taken modulo some number such as $10^9 + 7$. Where the helpers of section \secref{6.3.1} suit a few scattered operations, a value type is what makes an entire solution read in ordinary arithmetic notation while every operation reduces silently.

The modulus is a template parameter rather than a member, which is a deliberate speed choice: a compile-time constant lets the compiler replace each division by a multiplication and a shift, while a modulus only known at run time forces a hardware divide costing several times as much. Storing it in a \inlinecode{static} member instead makes the same class work with a runtime modulus at that price, and is worth doing when the modulus is read from the input. A modulus above roughly $3 \times 10^9$ needs the \inlinecode{mulmod()} of section \secref{6.3.1} in place of the plain products used here.

The \inlinecode{Modular\textless{}MOD\textgreater{}} value type wraps arithmetic modulo the positive compile-time constant \inlinecode{MOD}. Its signed \inlinecode{auto} argument determines the storage type: \inlinecode{Modular\textless{}1000000007\textgreater{}} uses 32-bit storage, while \inlinecode{Modular\textless{}(1LL \textless{}\textless{} 61) - 1\textgreater{}} uses 64-bit storage. The implementation follows the conventional \inlinecode{Mint} wrapper: construction normalizes values, while hidden friend operators support mixed expressions such as \inlinecode{2 + Mint(3)} through implicit conversion.

\begin{compactitem}
  \item \inlinecode{Modular\textless{}MOD\textgreater{}(x = 0)} constructs the residue class of integer \inlinecode{x} modulo \inlinecode{MOD}.
  \item \inlinecode{Modular\textless{}MOD\textgreater{}::mod()} returns the modulus \inlinecode{MOD}.
  \item \inlinecode{value()} and \inlinecode{operator()()} return the stored representative in $[0, {\inlinecodemath{MOD}})$.
  \item Explicit casts to \inlinecode{int}, \inlinecode{long long}, \inlinecode{double}, and \inlinecode{long double} convert that stored representative to the corresponding primitive.
  \item \inlinecode{pow(n)} returns this value raised to nonnegative integer exponent \inlinecode{n}.
  \item \inlinecode{inv()} returns the multiplicative inverse, asserting the value is coprime to \inlinecode{MOD}.
  \item Operators \inlinecode{+}, \inlinecode{-}, \inlinecode{*}, \inlinecode{/}, comparison, increment, decrement, and stream I/O are overloaded. Division likewise requires the divisor to be coprime to \inlinecode{MOD}.
\end{compactitem}

Overflow warning: Construction, multiplication, and inverses widen through \inlinecode{int64\_t} for 32-bit storage or \inlinecode{\_\_int128} for 64-bit storage. Thus $2p$ must fit the storage type and $(p - 1)^2$ must fit the intermediate type. On compilers without \inlinecode{\_\_int128}, only 32-bit moduli are supported. Estimating a 64-bit product with \inlinecode{long double} is not used because it silently loses the needed precision on platforms where \inlinecode{long double} is only 64 bits.

\Needspace*{3\baselineskip}
\textbf{Time Complexity}:
\begin{compactitem}
  \item $O(1)$ per addition, subtraction, multiplication, comparison, and stream output.
  \item $O(\log n)$ per call to \inlinecode{pow()}.
  \item $O(\log m)$ per call to \inlinecode{inv()} and division, where $m = {\inlinecodemath{MOD}}$.
\end{compactitem}

\textbf{Space Complexity}: $O(1)$ auxiliary for Modular arithmetic.

\begin{lstlisting}[style=cpp]
#include <cassert>
#include <cstdint>
#include <iostream>
#include <limits>
#include <type_traits>
#include <utility>
#include <vector>

template<typename T>
T inverse(T a, T m) {
  T original_m = m;
  T u = 0, v = 1;
  while (a != 0) {
    T q = m / a;
    m -= q * a;
    std::swap(a, m);
    u -= q * v;
    std::swap(u, v);
  }
  assert(m == 1);
  u %= original_m;
  return u < 0 ? u + original_m : u;
}

template<auto MOD>
class Modular {
  using T = decltype(MOD);
  static_assert(std::is_integral_v<T> && std::is_signed_v<T>);
  static_assert(MOD > 0 && MOD <= std::numeric_limits<T>::max() / 2);

  // Widen through int64_t for 32-bit storage, or __int128 for 64-bit storage where it exists
  // (__extension__ silences -pedantic). Without __int128, only 32-bit moduli are supported.
#if defined(__SIZEOF_INT128__)
  __extension__ typedef std::conditional_t<sizeof(T) <= 4, int64_t, __int128> wide_t;
#else
  using wide_t = int64_t;
  static_assert(sizeof(T) <= 4, "64-bit moduli require __int128, which is unavailable");
#endif

  T v;

  static T normalize(wide_t x) {
    if (-static_cast<wide_t>(mod()) <= x && x < static_cast<wide_t>(mod())) {
      T y = static_cast<T>(x);
      return y < 0 ? y + mod() : y;
    }
    x %= mod();
    if (x < 0) {
      x += mod();
    }
    return static_cast<T>(x);
  }

 public:
  Modular(wide_t x = 0) : v(normalize(x)) {}

  static T mod() { return MOD; }
  T value() const { return v; }
  T operator()() const { return v; }
  explicit operator int() const { return static_cast<int>(v); }
  explicit operator long long() const { return static_cast<long long>(v); }
  explicit operator double() const { return static_cast<double>(v); }
  explicit operator long double() const { return static_cast<long double>(v); }

  Modular pow(int64_t n) const {
    assert(n >= 0);
    Modular base = *this, res = 1;
    while (n > 0) {
      if (n & 1) {
        res *= base;
      }
      base *= base;
      n >>= 1;
    }
    return res;
  }

  Modular inv() const {
    assert(v != 0);
    return Modular(inverse(static_cast<wide_t>(v), static_cast<wide_t>(mod())));
  }

  Modular &operator+=(const Modular &other) {
    v += other.v;
    if (v >= mod()) {
      v -= mod();
    }
    return *this;
  }

  Modular &operator-=(const Modular &other) {
    v -= other.v;
    if (v < 0) {
      v += mod();
    }
    return *this;
  }

  Modular &operator*=(const Modular &other) {
    v = normalize(static_cast<wide_t>(v) * other.v);
    return *this;
  }

  Modular operator-() const { return Modular(-v); }
  Modular &operator++() { return *this += 1; }
  Modular &operator--() { return *this -= 1; }
  Modular &operator/=(const Modular &other) { return *this *= other.inv(); }
  Modular operator++(int) { Modular res = *this; ++*this; return res; }
  Modular operator--(int) { Modular res = *this; --*this; return res; }
  friend Modular operator+(Modular a, const Modular &b) { return a += b; }
  friend Modular operator-(Modular a, const Modular &b) { return a -= b; }
  friend Modular operator*(Modular a, const Modular &b) { return a *= b; }
  friend Modular operator/(Modular a, const Modular &b) { return a /= b; }
  friend bool operator==(const Modular &a, const Modular &b) { return a.v == b.v; }
  friend bool operator!=(const Modular &a, const Modular &b) { return !(a == b); }
  friend bool operator<(const Modular &a, const Modular &b) { return a.v < b.v; }
  friend std::ostream &operator<<(std::ostream &os, const Modular &x) { return os << x.v; }

  friend std::istream &operator>>(std::istream &is, Modular &x) {
    int64_t value;
    is >> value;
    x = Modular(value);
    return is;
  }
};
\end{lstlisting}
\inlinecode{ModCombinatorics\textless{}Mint\textgreater{}} maintains lazy factorial and inverse-factorial tables for a modular type \inlinecode{Mint}. It assumes the modulus is prime and all requested factorials are invertible. Growing the tables costs one modular inverse rather than one per entry: after extending the factorials by repeated multiplication, a single inversion of the largest factorial seeds a backward pass in which each inverse factorial follows from the next by one multiplication.

\begin{compactitem}
  \item \inlinecode{factorial(n)} returns $n!$ using a lazy factorial table.
  \item \inlinecode{choose(n, k)} returns $\binom n k$ using lazy factorial and inverse-factorial tables.
  \item \inlinecode{permute(n, k)} returns the number of ordered selections of \inlinecode{k} distinct elements from \inlinecode{n}.
  \item \inlinecode{multichoose(n, k)} returns the number of size-\inlinecode{k} multisets drawn from \inlinecode{n} types.
\end{compactitem}

Tables pay off once a solution asks for many values against one fixed modulus, which is the usual counting or dynamic programming case. For a handful of values, or for a modulus that changes between calls, the table-free routines of section \secref{6.2.1} compute the same quantities in $O(\min\left(k, n - k\right))$ per call and allocate nothing.

\textbf{Time Complexity}: $O(n)$ total table growth to answer calls with arguments up to $n$, then $O(1)$ per call.

\textbf{Space Complexity}: $O(n)$ for the factorial and inverse-factorial tables.

\begin{lstlisting}[style=cpp]
template<typename Mint>
class ModCombinatorics {
  std::vector<Mint> fact, inv_fact;

  void ensure(int n) {
    auto old = static_cast<int>(fact.size());
    if (old > n) {
      return;
    }
    fact.resize(n + 1);
    inv_fact.resize(n + 1);
    for (int i = old; i <= n; i++) {
      fact[i] = fact[i - 1] * i;
    }
    inv_fact[n] = fact[n].inv();
    for (int i = n; i > old; i--) {
      inv_fact[i - 1] = inv_fact[i] * i;
    }
  }

 public:
  ModCombinatorics() : fact(1, Mint{1}), inv_fact(1, Mint{1}) {}

  Mint factorial(int n) {
    ensure(n);
    return fact[n];
  }

  Mint choose(int n, int k) {
    if (k < 0 || k > n) {
      return 0;
    }
    ensure(n);
    return fact[n] * inv_fact[k] * inv_fact[n - k];
  }

  Mint permute(int n, int k) {
    if (k < 0 || k > n) {
      return 0;
    }
    ensure(n);
    return fact[n] * inv_fact[n - k];
  }

  Mint multichoose(int n, int k) { return n == 0 ? Mint(k == 0) : choose(n + k - 1, k); }
};

using Mint = Modular<1000000007>;
\end{lstlisting}
\Needspace*{4\baselineskip}
\textbf{Example Usage}\par
\begin{lstlisting}[style=cpp]
#include <cassert>
#include <sstream>
using namespace std;

int main() {
  Mint a = 1000000008LL;
  Mint b = -2;
  assert(a.value() == 1);
  assert(b() == Mint::mod() - 2);
  assert(static_cast<int>(a) == 1);
  assert(static_cast<long long>(a) == 1LL);
  assert(static_cast<double>(a) == 1.0);
  assert(a + b == -1);
  assert(2 + Mint(3) == 5);
  assert(Mint(2) * 3 == 6);
  assert(Mint(2).pow(10) == 1024);
  assert(Mint(5) / 5 == 1);

  Mint x = 0;
  x++;
  ++x;
  assert(x == 2);
  stringstream ss("1000000008");
  ss >> x;
  assert(x == 1);

  ModCombinatorics<Mint> comb;
  assert(comb.factorial(5) == 120);
  assert(comb.choose(5, 2) == 10);
  assert(comb.permute(5, 2) == 20);
  assert(comb.multichoose(3, 2) == 6);
  assert(comb.multichoose(0, 0) == 1);

  // Each distinct modulus is just another instantiation.
  using Mint2 = Modular<998244353>;
  assert(Mint2(2).pow(10) == 1024);
  assert(Mint2(-1) == Mint2::mod() - 1);

#if defined(__SIZEOF_INT128__)
  // A 64-bit modulus (deduced as long long) automatically widens through __int128.
  using Mint3 = Modular<(1LL << 61) - 1>;
  assert(Mint3(2).pow(62) == 2);  // 2^62 mod (2^61 - 1)
  assert(Mint3(3).inv() * 3 == 1);
#endif
  return 0;
}
\end{lstlisting}
\subsection{Primitive Roots}
\label{sec:6.3.4}
A primitive root modulo $m$ is a number $g$ whose powers generate every invertible residue modulo $m$. Equivalently, the multiplicative order of $g$ modulo $m$ is exactly $\phi(m)$. Primitive roots are useful for constructing generators of finite multiplicative groups, choosing roots for number theoretic transforms, discrete logarithm setups, and randomized modular hashing or testing schemes that need a full-period multiplicative generator.

Primitive roots do not exist for every $m$. They exist exactly for $m \in \{1, 2, 4, p^k, 2p^k\}$, where $p$ is an odd prime and $k \geq 1$. This implementation first checks that classification by factoring $m$, then factors $\phi(m)$. A candidate $g$ is primitive exactly when $g^{\phi(m)} \equiv 1 \pmod m$ and $g^{\phi(m)/q} \not\equiv 1 \pmod m$ for every prime factor $q$ of $\phi(m)$.

For a prime modulus $p$, this reduces to the common contest test: factor $p - 1$, then scan for the first $g$ such that $g^{(p - 1)/q} \not\equiv 1 \pmod p$ for every prime factor $q$ of $p - 1$.

\begin{compactitem}
  \item \inlinecode{primitive\_root(m)} returns the smallest primitive root modulo \inlinecode{m}, or $-1$ if no primitive root exists.
\end{compactitem}

Modular multiplication uses \inlinecode{\_\_uint128\_t} when available for speed, with a slower portable double-and-add fallback to avoid overflow on compilers without 128-bit integers.

\textbf{Time Complexity}: $O(\sqrt{m} + g \cdot k \cdot \log m)$ per call, where $g$ is the answer candidate tested and $k$ is the number of distinct prime factors of $\phi(m)$.

\textbf{Space Complexity}: $O(k)$ auxiliary.

\begin{lstlisting}[style=cpp]
#include <cassert>
#include <cstdint>
#include <numeric>
#include <set>
#include <utility>
#include <vector>

int64_t mulmod(int64_t a, int64_t b, int64_t m) {
  assert(a >= 0 && b >= 0 && m > 0);
#if defined(__SIZEOF_INT128__)
  return static_cast<int64_t>(static_cast<__uint128_t>(a) * b % m);
#else
  int64_t res = 0;
  for (a %= m; b > 0; b >>= 1) {
    if (b & 1) {
      res = res >= m - a ? res - (m - a) : res + a;
    }
    a = a >= m - a ? a - (m - a) : a + a;
  }
  return res;
#endif
}

int64_t powmod(int64_t x, int64_t n, int64_t m) {
  int64_t res = 1 % m;
  for (x %= m; n > 0; n >>= 1) {
    if (n & 1) {
      res = mulmod(res, x, m);
    }
    x = mulmod(x, x, m);
  }
  return res;
}

std::vector<std::pair<int64_t, int>> factorize(int64_t n) {
  std::vector<std::pair<int64_t, int>> res;
  for (int64_t p = 2; p <= n / p; p += (p == 2 ? 1 : 2)) {
    if (n % p != 0) {
      continue;
    }
    int count = 0;
    while (n % p == 0) {
      n /= p;
      count++;
    }
    res.emplace_back(p, count);
  }
  if (n > 1) {
    res.emplace_back(n, 1);
  }
  return res;
}

int64_t primitive_root(int64_t m) {
  if (m <= 0) {
    return -1;
  }
  if (m == 1 || m == 2 || m == 4) {
    return m - 1;
  }
  std::vector<std::pair<int64_t, int>> factors = factorize(m);
  bool has_root = (factors.size() == 1 && factors[0].first != 2) ||
                  (factors.size() == 2 && factors[0].first == 2 && factors[0].second == 1);
  if (!has_root) {
    return -1;
  }
  int64_t phi = m;
  std::set<int64_t> phi_factors;
  for (auto [p, e] : factors) {
    phi = phi / p * (p - 1);
    if (e > 1) {
      phi_factors.insert(p);
    }
    for (auto [q, count] : factorize(p - 1)) {
      phi_factors.insert(q);
    }
  }
  for (int64_t g = 1; g < m; g++) {
    if (std::gcd(g, m) != 1 || powmod(g, phi, m) != 1) {
      continue;
    }
    bool ok = true;
    for (int64_t q : phi_factors) {
      if (powmod(g, phi / q, m) == 1) {
        ok = false;
        break;
      }
    }
    if (ok) {
      return g;
    }
  }
  return -1;
}
\end{lstlisting}
\Needspace*{4\baselineskip}
\textbf{Example Usage}\par
\begin{lstlisting}[style=cpp]
#include <cassert>

int main() {
  assert(primitive_root(1) == 0);
  assert(primitive_root(2) == 1);
  assert(primitive_root(4) == 3);

  assert(primitive_root(17) == 3);         // 3 has order phi(17) = 16.
  assert(primitive_root(998244353) == 3);  // Common NTT prime.
  assert(primitive_root(9) == 2);          // Odd prime power.
  assert(primitive_root(18) == 5);         // Twice an odd prime power.

  assert(primitive_root(8) == -1);
  assert(primitive_root(12) == -1);
  return 0;
}
\end{lstlisting}
\subsection{Discrete Logarithm (Baby-Step Giant-Step)}
\label{sec:6.3.5}
Given integers $a$, $b$, and a modulus $m$, the discrete logarithm problem asks for the smallest nonnegative integer $x$ with $a^x$ congruent to $b$ modulo $m$. This is the modular analogue of a logarithm and underlies Diffie-Hellman key exchange, order computations in the multiplicative group, and many number-theory reductions.

The baby-step giant-step method writes $x = np - q$ for a block size $n$ near the square root of the modulus. It tabulates the ``baby steps'' $b \cdot a^q$ for $q$ in $[0, n]$, then walks the ``giant steps'' $a^{np}$ for increasing $p$ until one lands in the table, yielding $x = np - q$. This trades the linear scan of all exponents for $O(\sqrt{m})$ work and memory. A short reduction loop first strips common factors between $a$ and $m$, so the routine works for any modulus, not only when $a$ and $m$ are coprime.

\begin{compactitem}
  \item \inlinecode{discrete\_log(a, b, m)} returns the smallest nonnegative \inlinecode{x} with \inlinecode{a\textasciicircum{}x} congruent to \inlinecode{b} modulo \inlinecode{m}, or $-1$ if no such \inlinecode{x} exists. Requires \inlinecode{m} $\geq 1$; \inlinecode{a} and \inlinecode{b} are reduced modulo \inlinecode{m} internally.
\end{compactitem}

\textbf{Time Complexity}: $O(\sqrt{m})$ expected per call with constant-time modular multiplication, and $O(\sqrt{m} \cdot \log m)$ with the portable double-and-add fallback.

\textbf{Space Complexity}: $O(\sqrt{m})$ auxiliary for the table.

\begin{lstlisting}[style=cpp]
#include <cassert>
#include <cmath>
#include <cstdint>
#include <numeric>
#include <unordered_map>

int64_t mulmod(int64_t a, int64_t b, int64_t m) {
  assert(a >= 0 && b >= 0 && m > 0);
#if defined(__SIZEOF_INT128__)
  return static_cast<int64_t>(static_cast<__uint128_t>(a) * b % m);
#else
  int64_t res = 0;
  for (a %= m; b > 0; b >>= 1) {
    if (b & 1) {
      res = res >= m - a ? res - (m - a) : res + a;
    }
    a = a >= m - a ? a - (m - a) : a + a;
  }
  return res;
#endif
}

int64_t discrete_log(int64_t a, int64_t b, int64_t m) {
  assert(m >= 1);
  a %= m;
  b %= m;
  if (a < 0) {
    a += m;
  }
  if (b < 0) {
    b += m;
  }
  // Strip common factors of a and m so that a becomes invertible; add counts the exponents removed.
  int64_t k = 1 % m, add = 0, g;
  while ((g = std::gcd(a, m)) > 1) {
    if (b == k) {
      return add;
    }
    if (b % g != 0) {
      return -1;
    }
    b /= g;
    m /= g;
    add++;
    k = mulmod(k, a / g, m);
  }
  int64_t n = static_cast<int64_t>(std::sqrt(static_cast<double>(m))) + 1;
  int64_t an = 1;
  for (int64_t i = 0; i < n; i++) {
    an = mulmod(an, a, m);  // an = a^n, the giant step.
  }
  std::unordered_map<int64_t, int64_t> baby;
  int64_t cur = b;
  for (int64_t q = 0; q <= n; q++) {
    baby[cur] = q;  // Keep the largest q for each value to make the answer minimal.
    cur = mulmod(cur, a, m);
  }
  cur = k;
  for (int64_t p = 1; p <= n; p++) {
    cur = mulmod(cur, an, m);
    auto it = baby.find(cur);
    if (it != baby.end()) {
      return n * p - it->second + add;
    }
  }
  return -1;
}
\end{lstlisting}
\Needspace*{4\baselineskip}
\textbf{Example Usage}\par
\begin{lstlisting}[style=cpp]
#include <cassert>

int main() {
  // 2^x = 9 (mod 11): 2^6 = 64 = 9, and 6 is the smallest such exponent.
  assert(discrete_log(2, 9, 11) == 6);

  assert(discrete_log(3, 1, 7) == 0);    // a^0 = 1 always.
  assert(discrete_log(2, 3, 5) == 3);    // 2^3 = 8 = 3 (mod 5).
  assert(discrete_log(2, 1, 6) == 0);    // Works for composite, non-coprime moduli.
  assert(discrete_log(4, 6, 8) == -1);   // No power of 4 is 6 (mod 8).
  assert(discrete_log(5, 33, 58) == 9);  // 5^9 = 33 (mod 58).
  return 0;
}
\end{lstlisting}
\subsection{Modular Square Root (Tonelli-Shanks)}
\label{sec:6.3.6}
Given an integer $a$ and an odd prime $p$, find a modular square root: an $x$ with $x^2$ congruent to $a$ modulo $p$. Exactly half of the nonzero residues modulo a prime are squares (quadratic residues), so a solution exists only for those; when it does, the two roots are $x$ and $p - x$.

The Tonelli-Shanks algorithm handles the general case. Euler's criterion, $a^{(p-1)/2}$ modulo $p$, first tests whether $a$ is a residue at all. When $p$ is congruent to $3 \pmod{4}$, the root is simply $a^{(p+1)/4}$. Otherwise the algorithm factors $p - 1$ as $s \cdot 2^e$, picks any quadratic non-residue to generate the $2$-power part of the group, and repeatedly squares a running value to find the order of the current error term, correcting the candidate root until the error becomes $1$.

\begin{compactitem}
  \item \inlinecode{mod\_sqrt(a, p)} returns the smaller of the two square roots of \inlinecode{a} modulo the odd prime \inlinecode{p}, or $-1$ if \inlinecode{a} is not a quadratic residue. The value $0$ maps to $0$. The argument \inlinecode{a} is reduced modulo \inlinecode{p} internally. The modulus \inlinecode{p} must be an odd prime.
\end{compactitem}

\textbf{Time Complexity}: $O(\log^{2} p + z \log p)$ per call, where $z$ is the smallest integer at least $2$ that is a quadratic non-residue modulo $p$; thus the unconditional worst-case bound is $O(p \log p)$.

\textbf{Space Complexity}: $O(1)$ auxiliary.

\begin{lstlisting}[style=cpp]
#include <cassert>
#include <cstdint>

int64_t mulmod(int64_t a, int64_t b, int64_t m) {
  assert(a >= 0 && b >= 0 && m > 0);
#if defined(__SIZEOF_INT128__)
  return static_cast<int64_t>(static_cast<__uint128_t>(a) * b % m);
#else
  int64_t res = 0;
  for (a %= m; b > 0; b >>= 1) {
    if (b & 1) {
      res = res >= m - a ? res - (m - a) : res + a;
    }
    a = a >= m - a ? a - (m - a) : a + a;
  }
  return res;
#endif
}

int64_t powmod(int64_t x, int64_t n, int64_t m) {
  int64_t res = 1 % m;
  for (x %= m; n > 0; n >>= 1) {
    if (n & 1) {
      res = mulmod(res, x, m);
    }
    x = mulmod(x, x, m);
  }
  return res;
}

int64_t mod_sqrt(int64_t a, int64_t p) {
  a %= p;
  if (a < 0) {
    a += p;
  }
  if (a == 0) {
    return 0;
  }
  if (powmod(a, (p - 1) / 2, p) != 1) {
    return -1;  // a is a quadratic non-residue.
  }
  if (p % 4 == 3) {
    int64_t x = powmod(a, p / 4 + 1, p);
    return x < p - x ? x : p - x;
  }
  // Write p - 1 = s * 2^e with s odd.
  int64_t s = p - 1, e = 0;
  while (s % 2 == 0) {
    s /= 2;
    e++;
  }
  // Any quadratic non-residue generates the 2-power part of the group.
  int64_t z = 2;
  while (powmod(z, (p - 1) / 2, p) != p - 1) {
    z++;
  }
  int64_t x = powmod(a, (s + 1) / 2, p);  // Candidate root, off by a 2-power factor.
  int64_t b = powmod(a, s, p);            // Error term, a power of two in order.
  int64_t g = powmod(z, s, p);            // Generator of the relevant 2-group.
  int64_t r = e;
  while (true) {
    int64_t t = b, order = 0;
    while (order < r && t != 1) {
      t = mulmod(t, t, p);
      order++;
    }
    if (order == 0) {
      return x < p - x ? x : p - x;
    }
    int64_t gs = powmod(g, 1LL << (r - order - 1), p);
    g = mulmod(gs, gs, p);
    x = mulmod(x, gs, p);
    b = mulmod(b, g, p);
    r = order;
  }
}
\end{lstlisting}
\Needspace*{4\baselineskip}
\textbf{Example Usage}\par
\begin{lstlisting}[style=cpp]
#include <cassert>

int main() {
  int64_t r = mod_sqrt(2, 113);
  assert(r != -1 && mulmod(r, r, 113) == 2);
  assert(r == 51);  // 51^2 = 2601 = 2 (mod 113); the other root is 62.

  assert(mod_sqrt(0, 7) == 0);
  assert(mod_sqrt(5, 41) != -1);  // p = 1 (mod 4) exercises the full algorithm.
  assert(mod_sqrt(4, 7) == 2);    // p = 3 (mod 4) fast path.
  assert(mod_sqrt(3, 7) == -1);   // 3 is a non-residue modulo 7.

  // Every residue's reported root squares back to it.
  for (int64_t a = 0; a < 41; a++) {
    int64_t x = mod_sqrt(a, 41);
    if (x != -1) {
      assert(mulmod(x, x, 41) == a);
    }
  }
  return 0;
}
\end{lstlisting}
\subsection{Jacobi Symbol}
\label{sec:6.3.7}
The Jacobi symbol $(a / n)$, defined for an odd positive integer $n$, generalizes the Legendre symbol. For an odd prime $p$ the Legendre symbol $(a / p)$ is $0$ when $p$ divides $a$, $1$ when $a$ is a nonzero quadratic residue modulo $p$, and $-1$ when $a$ is a non-residue. The Jacobi symbol extends this to composite $n$ by multiplying the Legendre symbols of the prime factors of $n$, counted with multiplicity. It is computed without factoring $n$, using the law of quadratic reciprocity together with the supplementary rules for $-1$ and $2$, in the style of the binary GCD.

The symbol is fully multiplicative in $a$ and in $n$. For prime $n$ it exactly decides quadratic residuosity, so it is the fast way to test whether a modular square root exists. For composite $n$, beware that $(a / n) = 1$ does not guarantee that $a$ is a quadratic residue; this asymmetry is what the Solovay-Strassen primality test exploits.

\begin{compactitem}
  \item \inlinecode{jacobi(a, n)} returns the Jacobi symbol (\inlinecode{a} / \inlinecode{n}) as one of $-1$, $0$, or $1$. The modulus \inlinecode{n} must be a positive odd integer; \inlinecode{a} is reduced modulo \inlinecode{n} internally and may be negative. The result is $0$ exactly when \inlinecode{a} and \inlinecode{n} share a common factor.
\end{compactitem}

\textbf{Time Complexity}: $O(\log n)$ arithmetic operations per call.

\textbf{Space Complexity}: $O(1)$ auxiliary.

\begin{lstlisting}[style=cpp]
#include <cassert>
#include <cstdint>
#include <utility>

int jacobi(int64_t a, int64_t n) {
  assert(n > 0 && n % 2 == 1);
  a %= n;
  if (a < 0) {
    a += n;
  }
  int result = 1;
  while (a != 0) {
    while (a % 2 == 0) {
      a /= 2;
      // (2 / n) is -1 when n is 3 or 5 modulo 8.
      if (n % 8 == 3 || n % 8 == 5) {
        result = -result;
      }
    }
    std::swap(a, n);
    // Quadratic reciprocity flips the sign when both are 3 modulo 4.
    if (a % 4 == 3 && n % 4 == 3) {
      result = -result;
    }
    a %= n;
  }
  return n == 1 ? result : 0;
}
\end{lstlisting}
\Needspace*{4\baselineskip}
\textbf{Example Usage}\par
\begin{lstlisting}[style=cpp]
int64_t legendre_via_euler(int64_t a, int64_t p) {
  int64_t r = 1 % p;
  for (int64_t b = ((a % p) + p) % p, e = (p - 1) / 2; e > 0; e >>= 1) {
    if (e & 1) {
      r = r * b % p;
    }
    b = b * b % p;
  }
  return r == p - 1 ? -1 : r;  // r is 0, 1, or p - 1.
}

int main() {
  assert(jacobi(2, 15) == 1);  // (2/15) = (2/3)(2/5) = (-1)(-1).
  assert(jacobi(7, 15) == -1);
  assert(jacobi(3, 15) == 0);  // gcd(3, 15) = 3.
  assert(jacobi(1001, 9907) == -1);

  // For odd primes the Jacobi symbol equals the Legendre symbol from Euler's criterion.
  for (int64_t p : {3, 5, 7, 11, 13, 101, 9973}) {
    for (int64_t a = 0; a < 50; a++) {
      assert(jacobi(a, p) == legendre_via_euler(a, p));
    }
  }
  return 0;
}
\end{lstlisting}
\subsection{Prime Generation}
\label{sec:6.3.8}
Generate prime numbers using the sieve of Eratosthenes, which repeatedly marks the multiples of each prime as composite. The linear sieve technique optimizes the classic sieve by finding the least prime factor of each integer. When processing \inlinecode{i}, it marks \inlinecode{i * p} for each known prime \inlinecode{p} no greater than \inlinecode{least[i]}. Every composite \inlinecode{x} is therefore generated exactly once: writing \inlinecode{x = i*p}, where \inlinecode{p} is its least prime factor, gives \inlinecode{p} $\leq$ \inlinecode{least[i]}, while any representation using a larger prime is skipped. This single visit per composite gives the $O(n)$ running time.

The segmented sieve technique computes primes in a high but narrow window without sieving everything below it. Since every composite at most \inlinecode{hi} has a prime factor at most $\sqrt{{\inlinecodemath{hi}}}$, it first sieves the primes up to $\sqrt{{\inlinecodemath{hi}}}$, then uses only those primes to mark composites within $[{\inlinecodemath{lo}}, {\inlinecodemath{hi}}]$ (each prime $p$ starting at the first of its multiples that is at least both $p^2$ and \inlinecode{lo}). This needs only $O(w + \sqrt{{\inlinecodemath{hi}}})$ space, where $w = {\inlinecodemath{hi}} - {\inlinecodemath{lo}} + 1$, instead of the $O({\inlinecodemath{hi}})$ space that sieving all of $[2, {\inlinecodemath{hi}}]$ would require.

\begin{compactitem}
  \item \inlinecode{sieve(n)} returns a vector of all the primes less than or equal to \inlinecode{n}.
  \item \inlinecode{linear\_sieve(n, \&least\_out)} returns a vector of all primes less than or equal to \inlinecode{n} in linear time. If \inlinecode{least\_out} is not null, it is filled so that \inlinecode{least\_out[x]} is the least prime factor of \inlinecode{x} for every \inlinecode{x} $\geq 2$.
  \item \inlinecode{segmented\_sieve(lo, hi)} returns a vector of all the primes in the range $[{\inlinecodemath{lo}}, {\inlinecodemath{hi}}]$.
\end{compactitem}

\Needspace*{3\baselineskip}
\textbf{Time Complexity}:
\begin{compactitem}
  \item $O(n \log \log n)$ per call to \inlinecode{sieve(n)}.
  \item $O(n)$ per call to \inlinecode{linear\_sieve()}.
  \item $O((w + \sqrt{{\inlinecodemath{hi}}}) \cdot \log\left(\log\left({\inlinecodemath{hi}}\right)\right))$ per call to \inlinecode{segmented\_sieve(lo, hi)}, where $w = {\inlinecodemath{hi}} - {\inlinecodemath{lo}} + 1$.
\end{compactitem}

\Needspace*{3\baselineskip}
\textbf{Space Complexity}:
\begin{compactitem}
  \item $O(n)$ auxiliary for \inlinecode{sieve(n)}.
  \item $O(n)$ auxiliary for \inlinecode{linear\_sieve()}.
  \item $O(w + \sqrt{{\inlinecodemath{hi}}})$ auxiliary for \inlinecode{segmented\_sieve(lo, hi)}.
\end{compactitem}

\begin{lstlisting}[style=cpp]
#include <algorithm>
#include <cmath>
#include <cstdint>
#include <vector>

std::vector<int> sieve(int n) {
  if (n < 2) {
    return {};
  }
  std::vector<char> prime(n + 1, true);
  for (int i = 2; i <= n / i; i++) {
    if (prime[i]) {
      for (int64_t j = 1LL * i * i; j <= n; j += i) {
        prime[j] = false;
      }
    }
  }
  std::vector<int> res;
  for (int i = 2; i <= n; i++) {
    if (prime[i]) {
      res.push_back(i);
    }
  }
  return res;
}

std::vector<int> linear_sieve(int n, std::vector<int> *least_out = nullptr) {
  if (n < 2) {
    if (least_out != nullptr) {
      least_out->assign(std::max(n + 1, 0), 0);
    }
    return {};
  }
  std::vector<int> least(n + 1), primes;
  for (int i = 2; i <= n; i++) {
    if (least[i] == 0) {
      least[i] = i;
      primes.push_back(i);
    }
    for (int p : primes) {
      if (p > least[i] || i > n / p) {
        break;
      }
      least[i * p] = p;
    }
  }
  if (least_out != nullptr) {
    *least_out = least;
  }
  return primes;
}

std::vector<int> segmented_sieve(int lo, int hi) {
  if (hi < 2 || lo > hi) {
    return {};
  }
  int sqrt_hi = std::sqrt(hi);
  while (1LL * (sqrt_hi + 1) * (sqrt_hi + 1) <= hi) {
    sqrt_hi++;
  }
  int fourth_root_hi = std::sqrt(sqrt_hi);
  while (1LL * (fourth_root_hi + 1) * (fourth_root_hi + 1) <= sqrt_hi) {
    fourth_root_hi++;
  }
  std::vector<char> prime1(sqrt_hi + 1, true), prime2(hi - lo + 1, true);
  for (int i = 2; i <= fourth_root_hi; i++) {
    if (prime1[i]) {
      for (int64_t j = 1LL * i * i; j <= sqrt_hi; j += i) {
        prime1[j] = false;
      }
    }
  }
  for (int i = 2; i <= sqrt_hi; i++) {
    if (prime1[i]) {
      int64_t first_multiple = ((static_cast<int64_t>(lo) + i - 1) / i) * i;
      int64_t start = std::max(static_cast<int64_t>(i) * i, first_multiple);
      for (int64_t j = start; j <= hi; j += i) {
        prime2[j - lo] = false;
      }
    }
  }
  std::vector<int> res;
  for (int i = (lo > 1) ? lo : 2; i <= hi; i++) {
    if (prime2[i - lo]) {
      res.push_back(i);
    }
  }
  return res;
}
\end{lstlisting}
\Needspace*{4\baselineskip}
\textbf{Example Usage}\par
\begin{lstlisting}[style=cpp]
#include <cassert>
#include <ctime>
#include <iostream>
using namespace std;

int main() {
  int pmax = 10000000;
  time_t start;
  double delta;

  start = clock();
  auto p = sieve(pmax);
  delta = static_cast<double>((clock() - start)) / CLOCKS_PER_SEC;
  cout << "sieve(n=" << pmax << "): " << delta << "s" << endl;
  assert((sieve(1) == vector<int>{}));
  assert((sieve(10) == vector<int>{2, 3, 5, 7}));

  vector<int> least;
  assert(linear_sieve(pmax, &least) == p);
  assert(least[999983] == 999983);
  assert(least[999984] == 2);

  int lo = 1000000000, hi = 1005000000;
  start = clock();
  p = segmented_sieve(lo, hi);
  delta = static_cast<double>((clock() - start)) / CLOCKS_PER_SEC;
  cout << "segmented_sieve([" << lo << ", " << hi << "]): " << delta << "s" << endl;
  assert((segmented_sieve(14, 16) == vector<int>{}));
  assert((segmented_sieve(17, 19) == vector<int>{17, 19}));
  assert(!p.empty() && p.front() >= lo && p.back() <= hi);
  return 0;
}
\end{lstlisting}
\begin{exampleoutput}
sieve(n=10000000): 0.042641s
segmented_sieve([1000000000, 1005000000]): 0.01969s
\end{exampleoutput}
\subsection{Prime Tests and Factorization}
\label{sec:6.3.9}
Provides primality tests, prime factorization, and divisor generation. Alongside generic trial division, it includes deterministic Miller-Rabin and Pollard's rho for signed 64-bit integers. Prime factorizations are represented as sorted vectors of (\inlinecode{prime}, \inlinecode{exponent}) pairs. For $0$ and $1$, the prime factorization is empty.

Trial division tests primality and computes prime factorizations by trying possible divisors through the square root.

\begin{compactitem}
  \item \inlinecode{is\_prime\_slow(n)} returns whether the integer \inlinecode{n} is prime using trial division.
  \item \inlinecode{factorize\_slow(n)} returns the prime factorization of \inlinecode{n} using trial division.
\end{compactitem}

\textbf{Time Complexity}: $O(\sqrt{n})$ per call to \inlinecode{is\_prime\_slow()} and \inlinecode{factorize\_slow()}.

\textbf{Space Complexity}: $O(1)$ auxiliary for \inlinecode{is\_prime\_slow()} and $O(f)$ for the returned factorization, where $f$ is the number of compressed prime factors.

\begin{lstlisting}[style=cpp]
#include <algorithm>
#include <cassert>
#include <chrono>
#include <cstdint>
#include <numeric>
#include <random>
#include <utility>
#include <vector>

template<typename Int>
bool is_prime_slow(Int n) {
  if (n == 2 || n == 3) {
    return true;
  }
  if (n < 2 || n % 2 == 0 || n % 3 == 0) {
    return false;
  }
  for (Int i = 5, w = 4; i <= n / i; i += w) {
    if (n % i == 0) {
      return false;
    }
    w = 6 - w;
  }
  return true;
}

template<typename Int>
std::vector<std::pair<Int, int>> factorize_slow(Int n) {
  if (n <= 1) {
    return {};
  }
  std::vector<std::pair<Int, int>> res;
  for (Int i = 2; i <= n / i; i++) {
    if (n % i == 0) {
      int exponent = 0;
      do {
        n /= i;
        exponent++;
      } while (n % i == 0);
      res.emplace_back(i, exponent);
    }
  }
  if (n > 1) {
    res.emplace_back(n, 1);
  }
  return res;
}
\end{lstlisting}
The Miller-Rabin test writes $n - 1 = d \cdot 2^r$ and repeatedly squares $a^d$; a base that breaks the pattern required of a prime proves that $n$ is composite.

\begin{compactitem}
  \item \inlinecode{is\_probable\_prime(n, k = 20)} returns whether \inlinecode{n} is probably prime using \inlinecode{k} random Miller-Rabin bases. The result is guaranteed correct when \inlinecode{n} is prime. For composite \inlinecode{n}, a true result has one-sided error probability at most $(1 / 4)^k$.
  \item \inlinecode{is\_prime(n)} returns whether the signed 64-bit integer \inlinecode{n} is prime using deterministic Miller-Rabin bases sufficient up to and including $2^{63} - 1$.
\end{compactitem}

Modular multiplication uses \inlinecode{\_\_uint128\_t} when available, with a slower portable double-and-add fallback that avoids overflow on compilers without 128-bit integers.

\textbf{Time Complexity}: $O(k \cdot \log n)$ modular multiplications per call to \inlinecode{is\_probable\_prime()} and $O(\log n)$ per call to \inlinecode{is\_prime()}. The portable multiplication fallback adds another $O(\log n)$ factor.

\textbf{Space Complexity}: $O(1)$ auxiliary.

\begin{lstlisting}[style=cpp]
uint64_t mulmod(uint64_t a, uint64_t b, uint64_t m) {
  assert(m > 0);
#if defined(__SIZEOF_INT128__)
  return static_cast<uint64_t>(static_cast<__uint128_t>(a) * b % m);
#else
  uint64_t res = 0, cur = a % m;
  for (; b > 0; b >>= 1) {
    if (b & 1) {
      res = res >= m - cur ? res - (m - cur) : res + cur;
    }
    cur = cur >= m - cur ? cur - (m - cur) : cur + cur;
  }
  return res;
#endif
}

uint64_t powmod(uint64_t x, uint64_t n, uint64_t m) {
  assert(m > 0);
  uint64_t res = 1 % m;
  for (; n > 0; n >>= 1) {
    if (n & 1) {
      res = mulmod(res, x, m);
    }
    x = mulmod(x, x, m);
  }
  return res;
}

uint64_t rand64u() {
  static std::mt19937_64 rng(std::chrono::steady_clock::now().time_since_epoch().count());
  return rng();
}

bool is_probable_prime(int64_t n, int k = 20) {
  if (n == 2 || n == 3) {
    return true;
  }
  if (n < 2 || n % 2 == 0 || n % 3 == 0) {
    return false;
  }
  uint64_t s = n - 1, p = n - 1;
  while (!(s & 1)) {
    s >>= 1;
  }
  for (int i = 0; i < k; i++) {
    uint64_t x, r = powmod(rand64u() % (n - 3) + 2, s, n);  // random witness base in [2, n - 2]
    for (x = s; x != p && r != 1 && r != p; x <<= 1) {
      r = mulmod(r, r, n);
    }
    if (r != p && !(x & 1)) {
      return false;
    }
  }
  return true;
}

bool is_prime(int64_t n) {
  if (n < 2) {
    return false;
  }
  static const int small_primes[] = {2, 3, 5, 7, 11, 13, 17, 19, 23, 29};
  for (int p : small_primes) {
    if (n % p == 0) {
      return n == p;
    }
  }
  if (n < 31 * 31) {
    return true;
  }
  uint64_t t;
  int s = 0;
  for (t = n - 1; !(t & 1); t >>= 1) {
    s++;
  }
  static const uint64_t bases[] = {2, 325, 9375, 28178, 450775, 9780504, 1795265022};
  for (uint64_t a : bases) {
    if (a % n == 0) {
      continue;
    }
    uint64_t r = powmod(a, t, n);
    if (r == 1) {
      continue;
    }
    bool ok = false;
    for (int j = 0; j < s && !ok; j++) {
      ok |= (r == static_cast<uint64_t>(n) - 1);
      r = mulmod(r, r, n);
    }
    if (!ok) {
      return false;
    }
  }
  return true;
}
\end{lstlisting}
Pollard's rho iterates a pseudorandom map modulo $n$; two iterates that collide modulo a hidden prime factor reveal that factor through a GCD with their difference. The complete factorization routine first removes small prime factors with a cached linear sieve, then uses Miller-Rabin and Pollard's rho on the remaining cofactor. Factorizations remain sorted and compressed as (\inlinecode{prime}, \inlinecode{exponent}) pairs.

\begin{compactitem}
  \item \inlinecode{rho\_factor(n)} returns a factor of \inlinecode{n} that is not necessarily prime using Pollard's rho with Brent's optimization. If \inlinecode{n} is prime, then \inlinecode{n} itself is returned. The algorithm may need to be retried to find a nontrivial factor, as done in \inlinecode{factorize\_rho()}.
  \item \inlinecode{factorize\_rho(n)} returns the prime factorization of a 64-bit integer using Miller-Rabin and Pollard's rho without initial trial division.
  \item \inlinecode{factorize(n, small\_prime\_limit = 1000000)} returns the prime factorization of a 64-bit integer \inlinecode{n}, first testing primes through \inlinecode{small\_prime\_limit} and then falling back to Pollard's rho. It supports integers up to and including $2^{63} - 1$.
  \item \inlinecode{divisors\_from\_factors(factors)} returns all divisors from a compressed factorization.
  \item \inlinecode{divisors(n)} returns all divisors of a 64-bit integer using \inlinecode{factorize()} followed by \inlinecode{divisors\_from\_factors()}.
\end{compactitem}

\Needspace*{3\baselineskip}
\textbf{Time Complexity}:
\begin{compactitem}
  \item $O(\sqrt{p})$ expected modular multiplications per call to \inlinecode{rho\_factor(n)}, where $p$ is the smallest prime factor of \inlinecode{n}. For composite \inlinecode{n}, this is at most $O(n^{1/4})$.
  \item $O(n^{1/4})$ expected per call to \inlinecode{factorize\_rho()}. \inlinecode{factorize()} additionally performs $O(L/\log L)$ trial divisions, where $L$ is \inlinecode{small\_prime\_limit}, and takes $O(L)$ time if the cache is rebuilt.
  \item $O(d \log d)$ per call to \inlinecode{divisors\_from\_factors()}, where $d$ is the number of divisors generated.
  \item $O(L/\log L + n^{1/4} + d \log d)$ expected per call to \inlinecode{divisors()}, plus $O(L)$ if the cache is rebuilt.
\end{compactitem}

\Needspace*{3\baselineskip}
\textbf{Space Complexity}:
\begin{compactitem}
  \item $O(c)$ cached space after sieving through $c$.
  \item $O(f + \log n)$ auxiliary for recursive factorization, where $f$ is the number of compressed prime factors returned.
  \item $O(d)$ output space and $O(\log d)$ auxiliary stack space for \inlinecode{divisors\_from\_factors()} and \inlinecode{divisors()}.
\end{compactitem}

\begin{lstlisting}[style=cpp]
using Factors = std::vector<std::pair<int64_t, int>>;

int64_t rho_factor(int64_t n) {
  if (n % 2 == 0) {
    return 2;
  }
  uint64_t y = rand64u() % (n - 1) + 1;
  uint64_t c = rand64u() % (n - 1) + 1;
  uint64_t m = rand64u() % (n - 1) + 1;
  uint64_t g = 1, r = 1, q = 1, ys = 0, x = 0;
  for (r = 1; g == 1; r <<= 1) {
    x = y;
    for (uint64_t i = 0; i < r; i++) {
      y = (mulmod(y, y, n) + c) % n;
    }
    for (uint64_t k = 0; k < r && g == 1; k += m) {
      ys = y;
      int64_t lim = std::min(m, r - k);
      for (int j = 0; j < lim; j++) {
        y = (mulmod(y, y, n) + c) % n;
        q = mulmod(q, (x > y) ? (x - y) : (y - x), n);
      }
      g = std::gcd(q, n);
    }
  }
  if (g == static_cast<uint64_t>(n)) {
    do {
      ys = (mulmod(ys, ys, n) + c) % n;
      g = std::gcd((x > ys) ? (x - ys) : (ys - x), n);
    } while (g <= 1);
  }
  return g;
}

Factors factorize_rho(int64_t n) {
  std::vector<int64_t> factors;
  auto collect = [&](auto &&collect, int64_t value) {
    if (value <= 1) {
      return;
    }
    if (is_prime(value)) {
      factors.push_back(value);
      return;
    }
    int64_t p;
    do {
      p = rho_factor(value);
    } while (p == value);
    collect(collect, p);
    collect(collect, value / p);
  };
  collect(collect, n);
  std::sort(factors.begin(), factors.end());
  Factors res;
  for (int64_t p : factors) {
    if (res.empty() || res.back().first != p) {
      res.emplace_back(p, 1);
    } else {
      res.back().second++;
    }
  }
  return res;
}

static const auto &cached_sieve(int n) {
  struct Cache {
    int limit = 1;
    std::vector<int> least{0, 1}, primes;
  };
  static Cache cache;
  if (n <= cache.limit) {
    return cache;
  }
  cache.limit = std::max(1, n);
  cache.least.assign(cache.limit + 1, 0);
  cache.primes.clear();
  for (int i = 2; i <= cache.limit; i++) {
    if (cache.least[i] == 0) {
      cache.least[i] = i;
      cache.primes.push_back(i);
    }
    for (int p : cache.primes) {
      if (p > cache.least[i] || i > cache.limit / p) {
        break;
      }
      cache.least[i * p] = p;
    }
  }
  return cache;
}

Factors factorize(int64_t n, int small_prime_limit = 1000000) {
  assert(small_prime_limit >= 1);
  if (n <= 1) {
    return {};
  }
  const auto &sieve = cached_sieve(small_prime_limit);
  if (n <= small_prime_limit) {
    Factors res;
    while (n > 1) {
      int p = sieve.least[n], cnt = 0;
      do {
        n /= p;
        cnt++;
      } while (n > 1 && sieve.least[n] == p);
      res.emplace_back(p, cnt);
    }
    return res;
  }
  Factors res;
  const std::vector<int> &primes = sieve.primes;
  for (int p : primes) {
    if (p > small_prime_limit || static_cast<int64_t>(p) > n / p) {
      break;
    }
    if (n % p == 0) {
      int cnt = 0;
      do {
        n /= p;
        cnt++;
      } while (n % p == 0);
      res.emplace_back(p, cnt);
    }
  }
  Factors tail = factorize_rho(n);
  res.insert(res.end(), tail.begin(), tail.end());
  return res;
}

std::vector<int64_t> divisors_from_factors(const Factors &factors) {
  std::vector<int64_t> res{1};
  for (const auto &factor : factors) {
    int old_size = static_cast<int>(res.size());
    int64_t power = 1;
    for (int e = 1; e <= factor.second; e++) {
      power *= factor.first;
      for (int i = 0; i < old_size; i++) {
        res.push_back(res[i] * power);
      }
    }
  }
  std::sort(res.begin(), res.end());
  return res;
}

std::vector<int64_t> divisors(int64_t n) {
  if (n <= 1) {
    return (n < 1) ? std::vector<int64_t>() : std::vector<int64_t>(1, 1);
  }
  return divisors_from_factors(factorize(n));
}
\end{lstlisting}
\Needspace*{4\baselineskip}
\textbf{Example Usage}\par
\begin{lstlisting}[style=cpp]
#include <ctime>
#include <iomanip>
#include <iostream>
#include <set>
using namespace std;

void validate(int64_t n, const Factors &factors) {
  if (n <= 1) {
    assert(factors.empty());
    return;
  }
  if (is_prime(n)) {
    assert((factors == Factors{{n, 1}}));
    return;
  }
  int64_t prod = 1;
  for (const auto &factor : factors) {
    assert(is_prime(factor.first));
    for (int i = 0; i < factor.second; i++) {
      prod *= factor.first;
    }
  }
  assert(prod == n);
}

int main() {
  {  // Small primality tests.
    vector<pair<int, bool>> cases{
        {-10, false}, {-1, false}, {0, false}, {1, false}, // non-primes
        {2, true}, {3, true}, {4, false}, {5, true},
        {9, false},    // divisible by 3
        {25, false},   // square of small prime
        {29, true},    // last small-prime table entry
        {31, true},    // just past small-prime table
        {49, false},   // square, not divisible by 2 or 3
        {121, false},  // 11^2
        {169, false},  // 13^2
        {961, false},  // 31^2, boundary for n < 31*31 shortcut
        {997, true},   // prime near 1000
        {1001, false}, // 7*11*13
        {2047, false}, // 23*89, pseudoprime to base 2
        {341, false},  // 11*31, Fermat pseudoprime to base 2
        {561, false},  // Carmichael number
    };
    for (const auto &[n, expected] : cases) {
      assert(is_prime_slow(n) == expected);
      assert(is_probable_prime(n) == expected);
      assert(is_prime(n) == expected);
    }
  }
  {  // Primality test benchmark.
    vector<int64_t> nums{
        772023803LL, 2147483647LL, 5705234089LL, 6339503641LL, 999966000289LL,
    };
    cout << "Primality test:" << endl;
    cout << setw(15) << left << "num:";
    cout << setw(18) << left << "is_prime_slow():";
    cout << setw(20) << left << "is_probable_prime():";
    cout << setw(13) << right << "is_prime():" << endl;
    cout << fixed << setprecision(3);
    for (int64_t n : nums) {
      clock_t start = clock();
      bool p1 = is_prime_slow(n);
      double t1 = static_cast<double>(clock() - start) / CLOCKS_PER_SEC;
      start = clock();
      bool p2 = is_probable_prime(n);
      double t2 = static_cast<double>(clock() - start) / CLOCKS_PER_SEC;
      start = clock();
      bool p3 = is_prime(n);
      double t3 = static_cast<double>(clock() - start) / CLOCKS_PER_SEC;
      assert(p1 == p2 && p1 == p3);
      cout << setw(12) << left << n << right;
      cout << setw(14) << 1000 * t1 << "ms";
      cout << setw(16) << 1000 * t2 << "ms";
      cout << setw(16) << 1000 * t3 << "ms" << endl;
    }
    cout << endl;
  }
  {  // Small factorization tests.
    for (int64_t i = 1; i <= 10000; i++) {
      auto v1 = factorize_slow(i);
      auto v2 = factorize(i);
      validate(i, v1);
      assert(v1 == v2);
      auto d = divisors(i);
      set<int> s(d.begin(), d.end());
      assert(d.size() == s.size());
      for (int j = 1; j <= i; j++) {
        if (i % j == 0) {
          assert(s.count(j));
        }
      }
    }
  }
  {  // Compressed factors are convenient for building divisors.
    Factors factors{{2, 3}, {3, 2}, {5, 1}};
    vector<int64_t> divs = divisors_from_factors(factors);
    assert(
        (divs == vector<int64_t>{
                     1,  2,  3,  4,  5,  6,  8,  9,  10, 12,  15,  18,
                     20, 24, 30, 36, 40, 45, 60, 72, 90, 120, 180, 360,
                 })
    );
    assert(factorize(360) == factors);
    assert(divisors(360) == divs);
  }
  {  // Large factorization tests.
    const vector<int64_t> nums{
        (1LL << 62),                    // high power of two
        9223372036854775807LL,          // 2^63 - 1, many small factors
        999983LL * 999983,              // square of a large prime
        1900009LL * 1910009 * 2541547,  // three medium-size prime factors
        50000017LL * 50001037,          // 16-digit semiprime
    };
    cout << "Factorization test:" << endl;
    cout << setw(24) << left << "num:";
    cout << setw(20) << left << "factorize_slow():";
    cout << "factorize():" << endl;
    cout << fixed << setprecision(3);
    for (int64_t n : nums) {
      clock_t start = clock();
      Factors factors1 = factorize_slow(n);
      double t1 = static_cast<double>(clock() - start) / CLOCKS_PER_SEC;
      start = clock();
      Factors factors2 = factorize(n);
      double t2 = static_cast<double>(clock() - start) / CLOCKS_PER_SEC;
      validate(n, factors1);
      validate(n, factors2);
      assert(factors1 == factors2);
      cout << setw(26) << left << n << right << setw(8) << 1000 * t1 << "ms";
      cout << setw(16) << 1000 * t2 << "ms" << endl;
      cout << left;
    }
  }
  return 0;
}
\end{lstlisting}
\begin{exampleoutput}
Primality test:
num:           is_prime_slow():  is_probable_prime():  is_prime():
772023803            0.006ms           0.005ms           0.001ms
2147483647           0.009ms           0.006ms           0.002ms
5705234089           0.014ms           0.000ms           0.001ms
6339503641           0.013ms           0.001ms           0.001ms
999966000289         0.171ms           0.001ms           0.001ms

Factorization test:
num:                    factorize_slow():   factorize():
4611686018427387904          0.002ms           0.002ms
9223372036854775807          0.049ms           0.007ms
999966000289                 0.503ms           0.043ms
9223361212852495307          0.952ms           0.096ms
2500052700017629            24.587ms           0.194ms
\end{exampleoutput}
\subsection{Euler's Totient Function}
\label{sec:6.3.10}
Euler's totient function $\phi(n)$ returns the number of positive integers less than or equal to $n$ that are relatively prime to $n$; that is, the number of integers $k$ in the range $[1, n]$ for which $\gcd(n, k) = 1$.

The function is multiplicative, with $\phi(p^k) = p^k - p^{k-1}$ for a prime power, which yields the product formula $\phi(n) = n \prod_{p \mid n} (1 - 1/p)$ taken over the distinct primes $p$ dividing $n$. \inlinecode{phi(n)} applies this formula directly, factoring \inlinecode{n} by trial division. \inlinecode{phi\_table(n)} computes the whole range at once with a sieve of Eratosthenes: starting from \inlinecode{v[i] = i}, it sweeps each prime $p$ across its multiples \inlinecode{j} and folds in the factor $(1 - 1/p)$ by replacing \inlinecode{v[j]} with ${\inlinecodemath{v[j]}} - {\inlinecodemath{v[j]}} / p$. This is faster than the alternative that uses the divisor-sum identity $\sum_{d \mid n} \phi(d) = n$ to subtract each \inlinecode{v[i]} from its multiples, which costs $O(n \log n)$.

The overload \inlinecode{phi\_table(lo, hi)} computes the totients of a high but narrow range with a segmented sieve, analogous to the segmented prime sieve. It sieves the primes up to $\sqrt{{\inlinecodemath{hi}}}$, applies each one's $(1 - 1/p)$ factor to its multiples within $[{\inlinecodemath{lo}}, {\inlinecodemath{hi}}]$ while dividing that prime out of a parallel copy of every value, and finally folds in any single leftover prime factor greater than $\sqrt{{\inlinecodemath{hi}}}$ (a number at most \inlinecode{hi} can have at most one such factor).

\begin{compactitem}
  \item \inlinecode{phi(n)} returns Euler's totient of \inlinecode{n}.
  \item \inlinecode{phi\_table(n)} returns a vector \inlinecode{v} of length \inlinecode{n + 1} such that \inlinecode{v[i]} stores \inlinecode{phi(i)} for every \inlinecode{i} in the range $[0, {\inlinecodemath{n}}]$.
  \item \inlinecode{phi\_table(lo, hi)} returns a vector \inlinecode{v} of length ${\inlinecodemath{hi}} - {\inlinecodemath{lo}} + 1$ such that \inlinecode{v[i - lo]} stores \inlinecode{phi(i)} for every \inlinecode{i} in the range $[{\inlinecodemath{lo}}, {\inlinecodemath{hi}}]$. This overload assumes \inlinecode{lo} $\geq 0$; if \inlinecode{hi} \textless{} \inlinecode{lo}, it returns an empty vector.
\end{compactitem}

\Needspace*{3\baselineskip}
\textbf{Time Complexity}:
\begin{compactitem}
  \item $O(\sqrt{n})$ per call to \inlinecode{phi(n)}.
  \item $O(n \log \log n)$ per call to \inlinecode{phi\_table(n)}.
  \item $O((w + \sqrt{{\inlinecodemath{hi}}}) \cdot \log\left(\log\left({\inlinecodemath{hi}}\right)\right))$ per call to \inlinecode{phi\_table(lo, hi)}, where $w = {\inlinecodemath{hi}} - {\inlinecodemath{lo}} + 1$.
\end{compactitem}

\Needspace*{3\baselineskip}
\textbf{Space Complexity}:
\begin{compactitem}
  \item $O(1)$ auxiliary for \inlinecode{phi(n)}.
  \item $O(n)$ auxiliary for \inlinecode{phi\_table(n)}.
  \item $O(w + \sqrt{{\inlinecodemath{hi}}})$ auxiliary for \inlinecode{phi\_table(lo, hi)}.
\end{compactitem}

\begin{lstlisting}[style=cpp]
#include <algorithm>
#include <cmath>
#include <cstdint>
#include <numeric>
#include <vector>

int phi(int n) {
  int res = n;
  for (int i = 2; i <= n / i; i++) {
    if (n % i == 0) {
      while (n % i == 0) {
        n /= i;
      }
      res -= res / i;
    }
  }
  if (n > 1) {
    res -= res / n;
  }
  return res;
}

std::vector<int> phi_table(int n) {
  std::vector<int> res(n + 1);
  std::iota(res.begin(), res.end(), 0);
  for (int i = 2; i <= n; i++) {
    if (res[i] == i) {  // i is prime, so apply its (1 - 1/i) factor to every multiple.
      for (int j = i; j <= n; j += i) {
        res[j] -= res[j] / i;
      }
    }
  }
  return res;
}

std::vector<int> phi_table(int lo, int hi) {
  if (hi < lo) {
    return {};
  }
  int root = static_cast<int>(std::sqrt(hi));
  while (1LL * (root + 1) * (root + 1) <= hi) {
    root++;
  }
  // Sieve the primes up to sqrt(hi); any larger prime divisor is the one leftover factor.
  std::vector<char> composite(root + 1);
  std::vector<int> primes;
  for (int i = 2; i <= root; i++) {
    if (!composite[i]) {
      primes.push_back(i);
      for (int64_t j = 1LL * i * i; j <= root; j += i) {
        composite[j] = true;
      }
    }
  }
  std::vector<int> res(hi - lo + 1);
  std::vector<int64_t> remaining(hi - lo + 1);
  for (int i = lo; i <= hi; i++) {
    res[i - lo] = i;
    remaining[i - lo] = i;
  }
  for (int p : primes) {
    // Start at the first multiple of p at least lo, but never below p itself (this skips 0).
    int64_t start = std::max(1LL * p, ((1LL * lo + p - 1) / p) * p);
    for (int64_t j = start; j <= hi; j += p) {
      int idx = static_cast<int>(j - lo);
      res[idx] -= res[idx] / p;
      while (remaining[idx] % p == 0) {
        remaining[idx] /= p;
      }
    }
  }
  for (int i = lo; i <= hi; i++) {
    int idx = i - lo;
    if (remaining[idx] > 1) {  // a single prime factor greater than sqrt(hi) is left over
      res[idx] -= static_cast<int>(res[idx] / remaining[idx]);
    }
  }
  return res;
}
\end{lstlisting}
\Needspace*{4\baselineskip}
\textbf{Example Usage}\par
\begin{lstlisting}[style=cpp]
#include <cassert>
using namespace std;

int main() {
  assert(phi(1) == 1);
  assert(phi(9) == 6);
  assert(phi(1234567) == 1224720);
  const int n = 1000;
  vector<int> pt = phi_table(n);
  for (int i = 0; i <= n; i++) {
    assert(pt[i] == phi(i));
  }
  // The segmented overload agrees with phi(), including over a low range (with 0 and 1) and a
  // high, narrow window where many values have a prime factor exceeding sqrt(hi).
  vector<int> lphi = phi_table(0, 50);
  for (int i = 0; i <= 50; i++) {
    assert(lphi[i] == phi(i));
  }
  int lo = 1000000000, hi = 1000000100;
  vector<int> hphi = phi_table(lo, hi);
  for (int i = lo; i <= hi; i++) {
    assert(hphi[i - lo] == phi(i));
  }
  return 0;
}
\end{lstlisting}
\subsection{Divisor Functions}
\label{sec:6.3.11}
The divisor-counting function $d(n)$ counts the positive divisors of $n$, and the divisor-sum function $\sigma(n)$ adds them up. Both are multiplicative, so each is determined by the prime factorization $n = \prod p_i^{e_i}$: choosing an exponent in $[0, e_i]$ for every prime gives $d(n) = \prod (e_i + 1)$, and summing the resulting geometric series for each prime gives $\sigma(n) = \prod (1 + p_i + \dots + p_i^{e_i})$. A single value is therefore computed by trial division, factoring $n$ as the divisors are accumulated.

\begin{compactitem}
  \item \inlinecode{divisor\_count(n)} returns $d(n)$, the number of positive divisors of \inlinecode{n}, which must be positive.
  \item \inlinecode{divisor\_sum(n)} returns $\sigma(n)$, the sum of the positive divisors of \inlinecode{n}, which must be positive.
\end{compactitem}

For every value in a range, factoring each number separately is wasteful. Instead, sweep over each candidate divisor $d$ and add its contribution to every multiple of $d$. The total work is the harmonic sum $n/1 + n/2 + \dots + n/n$, which is $O(n \log n)$. The same sweep computes any divisor-indexed accumulation, such as the divisor sums used by the Mobius inversion of section \secref{6.3.12}.

\begin{compactitem}
  \item \inlinecode{divisor\_count\_table(n)} returns a vector \inlinecode{v} of length \inlinecode{n + 1} such that \inlinecode{v[i]} is the number of divisors of \inlinecode{i} for every \inlinecode{i} in $[1, {\inlinecodemath{n}}]$, with \inlinecode{v[0]} left as $0$.
  \item \inlinecode{divisor\_sum\_table(n)} returns a vector \inlinecode{v} of length \inlinecode{n + 1} such that \inlinecode{v[i]} is the sum of the divisors of \inlinecode{i} for every \inlinecode{i} in $[1, {\inlinecodemath{n}}]$, with \inlinecode{v[0]} left as $0$.
\end{compactitem}

The aliquot sum $\sigma(n) - n$ adds the proper divisors, marking \inlinecode{n} perfect when it equals \inlinecode{n} and pairing amicable numbers when each is the other's. For an \inlinecode{n} too large for trial division, factor it with Pollard's rho in section \secref{6.3.9} and apply the product formulas; to tabulate a high narrow window, sieve it with the primes up to $\sqrt{{\inlinecodemath{hi}}}$ as in section \secref{6.3.8}.

Overflow warning: $\sigma(n)$ can be several times larger than \inlinecode{n}, so \inlinecode{divisor\_sum()} overflows \inlinecode{int64\_t} somewhat before \inlinecode{n} itself does. Trial division binds first in practice, since $\sqrt{n}$ divisions are already too slow well below that point; factor with section \secref{6.3.9} instead.

\Needspace*{3\baselineskip}
\textbf{Time Complexity}:
\begin{compactitem}
  \item $O(\sqrt{n})$ per call to \inlinecode{divisor\_count(n)} and \inlinecode{divisor\_sum(n)}.
  \item $O(n \log n)$ per call to \inlinecode{divisor\_count\_table(n)} and \inlinecode{divisor\_sum\_table(n)}.
\end{compactitem}

\Needspace*{3\baselineskip}
\textbf{Space Complexity}:
\begin{compactitem}
  \item $O(1)$ auxiliary for \inlinecode{divisor\_count()} and \inlinecode{divisor\_sum()}.
  \item $O(n)$ for the vectors returned by \inlinecode{divisor\_count\_table(n)} and \inlinecode{divisor\_sum\_table(n)}.
\end{compactitem}

\begin{lstlisting}[style=cpp]
#include <cassert>
#include <cstdint>
#include <vector>

int64_t divisor_count(int64_t n) {
  assert(n > 0);
  int64_t res = 1;
  for (int64_t p = 2; p <= n / p; p++) {
    if (n % p == 0) {
      int e = 0;
      for (; n % p == 0; n /= p) {
        e++;
      }
      res *= e + 1;
    }
  }
  if (n > 1) {  // A prime cofactor larger than the square root is left over.
    res *= 2;
  }
  return res;
}

int64_t divisor_sum(int64_t n) {
  assert(n > 0);
  int64_t res = 1;
  for (int64_t p = 2; p <= n / p; p++) {
    if (n % p == 0) {
      int64_t power = 1, series = 1;
      for (; n % p == 0; n /= p) {
        power *= p;
        series += power;
      }
      res *= series;  // Overflow warning.
    }
  }
  if (n > 1) {
    res *= n + 1;
  }
  return res;
}

std::vector<int> divisor_count_table(int n) {
  assert(n >= 0);
  std::vector<int> res(n + 1);
  for (int d = 1; d <= n; d++) {
    for (int i = d; i <= n; i += d) {
      res[i]++;
    }
  }
  return res;
}

std::vector<int64_t> divisor_sum_table(int n) {
  assert(n >= 0);
  std::vector<int64_t> res(n + 1);
  for (int d = 1; d <= n; d++) {
    for (int i = d; i <= n; i += d) {
      res[i] += d;
    }
  }
  return res;
}
\end{lstlisting}
\Needspace*{4\baselineskip}
\textbf{Example Usage}\par
\begin{lstlisting}[style=cpp]
#include <cassert>

int main() {
  assert(divisor_count(1) == 1);
  assert(divisor_count(12) == 6);  // 1, 2, 3, 4, 6, and 12.
  assert(divisor_count(999999937) == 2);
  assert(divisor_sum(1) == 1);
  assert(divisor_sum(12) == 28);
  assert(divisor_sum(6) == 12);  // 6 is perfect, since its aliquot sum is 6.

  // 220 and 284 are the smallest amicable pair.
  assert(divisor_sum(220) - 220 == 284);
  assert(divisor_sum(284) - 284 == 220);

  auto counts = divisor_count_table(100);
  auto sums = divisor_sum_table(100);
  assert(counts[0] == 0 && sums[0] == 0);
  for (int i = 1; i <= 100; i++) {
    assert(counts[i] == divisor_count(i));
    assert(sums[i] == divisor_sum(i));
  }
  return 0;
}
\end{lstlisting}
\subsection{Mobius Function and Inclusion-Exclusion}
\label{sec:6.3.12}
The Mobius function $\mu(n)$ is the multiplicative function defined as $1$ if $n = 1$, $0$ if $n$ is divisible by the square of a prime, and $(-1)^k$ if $n$ is a product of $k$ distinct primes. It is the coefficient that appears in Mobius inversion: if $g(n) = \sum_{d \mid n} f(d)$ for all $n$, then $f(n) = \sum_{d \mid n} \mu(d) \, g(n / d)$. This formalizes the principle of inclusion-exclusion in the divisor lattice, letting one recover a summand function from its divisor sums.

Inclusion-exclusion itself counts the size of a union by alternately adding and subtracting the sizes of intersections. A canonical number-theoretic instance is counting integers in $[1, n]$ that are coprime to a modulus $m$: subtract the multiples of each prime factor of $m$, add back those counted twice, and so on. The signed terms are precisely the values of $\mu$ over the squarefree divisors of $m$, so \inlinecode{count\_coprime(n, m)} equals $\sum_{d \mid m} \mu(d) \lfloor n / d \rfloor$.

\begin{compactitem}
  \item \inlinecode{mobius(n)} returns $\mu(n)$ for a single positive \inlinecode{n} by trial division.
  \item \inlinecode{mobius\_sieve(n)} returns a vector of length $n + 1$ where index $i$ holds $\mu(i)$, computed together with a linear sieve, where \inlinecode{n} must be nonnegative.
  \item \inlinecode{count\_coprime(n, m)} returns the number of integers in $[1, {\inlinecodemath{n}}]$ that are coprime to \inlinecode{m}, via inclusion-exclusion over the distinct prime factors of \inlinecode{m}, where \inlinecode{n} must be nonnegative and \inlinecode{m} must be positive.
\end{compactitem}

\Needspace*{3\baselineskip}
\textbf{Time Complexity}:
\begin{compactitem}
  \item $O(\sqrt{n})$ per call to \inlinecode{mobius()}.
  \item $O(n)$ per call to \inlinecode{mobius\_sieve()}.
  \item $O(\sqrt{m} + 2^{k} \cdot k)$ per call to \inlinecode{count\_coprime()}, where $k$ is the number of distinct prime factors of $m$ (at most 15 for $m$ within 64-bit range).
\end{compactitem}

\Needspace*{3\baselineskip}
\textbf{Space Complexity}:
\begin{compactitem}
  \item $O(n)$ auxiliary for \inlinecode{mobius\_sieve()}.
  \item $O(k)$ auxiliary for \inlinecode{count\_coprime()}.
  \item $O(1)$ auxiliary for \inlinecode{mobius()}.
\end{compactitem}

\begin{lstlisting}[style=cpp]
#include <cassert>
#include <cstdint>
#include <vector>

int mobius(int n) {
  assert(n > 0);
  if (n == 1) {
    return 1;
  }
  int res = 1;
  for (int p = 2; p <= n / p; p++) {
    if (n % p == 0) {
      n /= p;
      if (n % p == 0) {
        return 0;  // p^2 divides the original n.
      }
      res = -res;
    }
  }
  if (n > 1) {
    res = -res;  // One remaining prime factor larger than its square root.
  }
  return res;
}

std::vector<int> mobius_sieve(int n) {
  assert(n >= 0);
  std::vector<int> mu(n + 1), primes;
  std::vector<char> composite(n + 1);
  if (n >= 1) {
    mu[1] = 1;
  }
  for (int i = 2; i <= n; i++) {
    if (!composite[i]) {
      primes.push_back(i);
      mu[i] = -1;
    }
    for (int p : primes) {
      if (p > n / i) {
        break;
      }
      composite[i * p] = true;
      if (i % p == 0) {
        mu[i * p] = 0;  // p^2 divides i * p.
        break;
      }
      mu[i * p] = -mu[i];
    }
  }
  return mu;
}

int64_t count_coprime(int64_t n, int64_t m) {
  assert(n >= 0 && m > 0);
  std::vector<int64_t> primes;
  for (int64_t p = 2; p <= m / p; p++) {
    if (m % p == 0) {
      primes.push_back(p);
      while (m % p == 0) {
        m /= p;
      }
    }
  }
  if (m > 1) {
    primes.push_back(m);
  }
  int k = static_cast<int>(primes.size());
  int64_t total = 0;
  for (int mask = 0; mask < (1 << k); mask++) {
    int64_t product = 1;
    int bits = 0;
    for (int i = 0; i < k; i++) {
      if ((mask >> i) & 1) {
        product *= primes[i];
        bits++;
      }
    }
    total += (bits & 1 ? -1 : 1) * (n / product);
  }
  return total;
}
\end{lstlisting}
\Needspace*{4\baselineskip}
\textbf{Example Usage}\par
\begin{lstlisting}[style=cpp]
#include <cassert>
using namespace std;

int main() {
  assert(mobius(1) == 1);
  assert(mobius(2) == -1);
  assert(mobius(6) == 1);    // 2 * 3, two distinct primes.
  assert(mobius(12) == 0);   // Divisible by 2^2.
  assert(mobius(30) == -1);  // 2 * 3 * 5, three distinct primes.

  assert((mobius_sieve(12) == vector<int>{0, 1, -1, -1, 0, -1, 1, -1, 0, 0, 1, -1, 0}));

  assert(count_coprime(10, 6) == 3);  // 1, 5, 7
  assert(count_coprime(100, 1) == 100);
  return 0;
}
\end{lstlisting}

\section{\texorpdfstring{Arbitrary Precision Arithmetic}{6.4 Arbitrary Precision Arithmetic}}
\setcounter{section}{4}
\setcounter{subsection}{0}
\subsection{Big Integer (Simple)}
\label{sec:6.4.1}
Perform simple arithmetic operations on arbitrary precision big integers whose digits are internally represented as an \inlinecode{std::string} in little-endian order.

This version is intentionally small and decimal-oriented. Addition and subtraction are the usual schoolbook digit scans with carry or borrow. Multiplication keeps a shifted copy of the left operand for each digit of the right operand, adds that row once per digit value, then shifts by one decimal place. Division is long division: it scans the dividend from most significant digit to least, maintains a running remainder, and repeatedly subtracts the divisor to discover each quotient digit.

\begin{compactitem}
  \item \inlinecode{BigInt(n = 0)} constructs a big integer from an integer \inlinecode{n}.
  \item \inlinecode{BigInt(s)} constructs a big integer from a string \inlinecode{s}, which must strictly consist of a sequence of numeric digits, optionally preceded by a minus sign.
  \item \inlinecode{to\_string()} returns the string representation of the big integer.
  \item \inlinecode{comp(a, b)} returns $-1$, $0$, or $1$ depending on whether the big integers \inlinecode{a} and \inlinecode{b} compare less, equal, or greater, respectively.
  \item \inlinecode{add(a, b)} returns the sum of big integers \inlinecode{a} and \inlinecode{b}.
  \item \inlinecode{sub(a, b)} returns the difference of big integers \inlinecode{a} and \inlinecode{b}.
  \item \inlinecode{mul(a, b)} returns the product of big integers \inlinecode{a} and \inlinecode{b}.
  \item \inlinecode{div(a, b)} returns the quotient of big integers \inlinecode{a} and \inlinecode{b}.
\end{compactitem}

\Needspace*{3\baselineskip}
\textbf{Time Complexity}:
\begin{compactitem}
  \item $O(n)$ per call to the constructor, \inlinecode{to\_string()}, \inlinecode{comp()}, \inlinecode{add()}, and \inlinecode{sub()}, where $n$ is the total number of digits in the arguments and result for each operation.
  \item $O(n \cdot m)$ per call to \inlinecode{mul(a, b)} and \inlinecode{div(a, b)} where $n$ is the number of digits in \inlinecode{a} and $m$ is the number of digits in \inlinecode{b}.
\end{compactitem}

\Needspace*{3\baselineskip}
\textbf{Space Complexity}:
\begin{compactitem}
  \item $O(n)$ for storage of the big integer, where $n$ is the number of digits.
  \item $O(n)$ auxiliary for \inlinecode{to\_string()}, \inlinecode{add()}, \inlinecode{sub()}, \inlinecode{mul()}, and \inlinecode{div()}, where $n$ is the total number of digits in the arguments and result for each operation.
\end{compactitem}

\begin{lstlisting}[style=cpp]
#include <algorithm>
#include <cassert>
#include <cctype>
#include <cstddef>
#include <cstdint>
#include <stdexcept>
#include <string>

class BigInt {
  std::string digits;
  int sign;

  void normalize() {
    std::size_t pos = digits.find_last_not_of('0');
    if (pos != std::string::npos) {
      digits.erase(pos + 1);
    } else {
      digits = "0";  // An all-zero (or empty) magnitude collapses to a single "0".
    }
    if (digits.size() == 1 && digits[0] == '0') {
      sign = 1;
    }
  }

  static int comp(const std::string &a, const std::string &b, int asign, int bsign) {
    if (asign != bsign) {
      return asign < bsign ? -1 : 1;
    }
    if (a.size() != b.size()) {
      return a.size() < b.size() ? -asign : asign;
    }
    for (int i = static_cast<int>(a.size()) - 1; i >= 0; i--) {
      if (a[i] != b[i]) {
        return a[i] < b[i] ? -asign : asign;
      }
    }
    return 0;
  }

  static BigInt add(const std::string &a, const std::string &b, int asign, int bsign) {
    if (asign != bsign) {
      return (asign == 1) ? sub(a, b, asign, 1) : sub(b, a, bsign, 1);
    }
    BigInt res;
    res.sign = asign;
    res.digits.resize(std::max(a.size(), b.size()) + 1, '0');
    for (int i = 0, carry = 0; i < static_cast<int>(res.digits.size()); i++) {
      int d = carry;
      if (i < static_cast<int>(a.size())) {
        d += a[i] - '0';
      }
      if (i < static_cast<int>(b.size())) {
        d += b[i] - '0';
      }
      res.digits[i] = '0' + (d % 10);
      carry = d / 10;
    }
    res.normalize();
    return res;
  }

  static BigInt sub(const std::string &a, const std::string &b, int asign, int bsign) {
    if (asign == -1 || bsign == -1) {
      return add(a, b, asign, -bsign);
    }
    BigInt res;
    if (comp(a, b, asign, bsign) < 0) {
      res = sub(b, a, bsign, asign);
      res.sign = -1;
      return res;
    }
    res.digits.assign(a.size(), '0');
    for (int i = 0, borrow = 0; i < static_cast<int>(res.digits.size()); i++) {
      int d = (i < static_cast<int>(b.size()) ? a[i] - b[i] : a[i] - '0') - borrow;
      if (a[i] > '0') {
        borrow = 0;
      }
      if (d < 0) {
        d += 10;
        borrow = 1;
      }
      res.digits[i] = '0' + (d % 10);
    }
    res.normalize();
    return res;
  }

 public:
  BigInt(int64_t n = 0) {
    sign = (n < 0) ? -1 : 1;
    if (n == 0) {
      digits = "0";
      return;
    }
    for (; n != 0; n /= 10) {
      int digit = static_cast<int>(n % 10);
      digits += '0' + (digit < 0 ? -digit : digit);
    }
    normalize();
  }

  BigInt(const std::string &s) {
    if (s.empty() || (s[0] == '-' && s.size() == 1)) {
      throw std::runtime_error("Invalid string format to construct BigInt.");
    }
    digits.assign(s.rbegin(), s.rend());
    if (s[0] == '-') {
      sign = -1;
      digits.erase(digits.size() - 1);
    } else {
      sign = 1;
    }
    if (digits.find_first_not_of("0123456789") != std::string::npos) {
      throw std::runtime_error("Invalid string format to construct BigInt.");
    }
    normalize();
  }

  std::string to_string() const {
    return (sign < 0 ? "-" : "") + std::string(digits.rbegin(), digits.rend());
  }

  friend int comp(const BigInt &a, const BigInt &b) {
    return comp(a.digits, b.digits, a.sign, b.sign);
  }

  friend BigInt add(const BigInt &a, const BigInt &b) {
    return add(a.digits, b.digits, a.sign, b.sign);
  }

  friend BigInt sub(const BigInt &a, const BigInt &b) {
    return sub(a.digits, b.digits, a.sign, b.sign);
  }

  friend BigInt mul(const BigInt &a, const BigInt &b) {
    BigInt res, row(a);
    for (char dc : b.digits) {
      for (int j = 0; j < (dc - '0'); j++) {
        res = add(res.digits, row.digits, res.sign, row.sign);
      }
      if (row.digits.size() > 1 || row.digits[0] != '0') {
        row.digits.insert(0, "0");
      }
    }
    res.sign = a.sign * b.sign;
    res.normalize();
    return res;
  }

  friend BigInt div(const BigInt &a, const BigInt &b) {
    assert(b.digits.size() != 1 || b.digits[0] != '0');
    BigInt res, row;
    res.digits.assign(a.digits.size(), '0');
    row.digits.clear();  // Running remainder; kept normalized and nonnegative below.
    for (int i = static_cast<int>(a.digits.size()) - 1; i >= 0; i--) {
      row.digits.insert(row.digits.begin(), a.digits[i]);
      row.normalize();  // Strip the spurious high-order zero left by the previous remainder.
      // Subtract while the remainder is >= b. A strict ">" undercounts the quotient digit when the
      // remainder equals b, and leaves the remainder too large, overflowing the next digit past 9.
      while (comp(row.digits, b.digits, 1, 1) >= 0) {
        res.digits[i]++;
        row = sub(row.digits, b.digits, 1, 1);
      }
    }
    res.sign = a.sign * b.sign;
    res.normalize();
    return res;
  }
};
\end{lstlisting}
\Needspace*{4\baselineskip}
\textbf{Example Usage}\par
\begin{lstlisting}[style=cpp]
int main() {
  BigInt a("-9899819294989142124"), b("12398124981294214");
  assert(add(a, b).to_string() == "-9887421170007847910");
  assert(sub(a, b).to_string() == "-9912217419970436338");
  assert(mul(a, b).to_string() == "-122739196911503356525379735104870536");
  assert(div(a, b).to_string() == "-798");
  assert(comp(a, b) == -1 && comp(a, a) == 0 && comp(b, a) == 1);
  assert(BigInt(INT64_MIN).to_string() == "-9223372036854775808");
  return 0;
}
\end{lstlisting}
\subsection{Big Integer}
\label{sec:6.4.2}
Perform operations on arbitrary precision big integers internally represented as a vector of base-\inlinecode{BASE} digits in little-endian order. \inlinecode{BASE} defaults to $10^9$ and must be a power of 10 no larger than $10^9$. Multiplication selects schoolbook, Karatsuba, or FFT from the operand size; \inlinecode{MUL\_CUTOFF} is the padded transform length where FFT takes over and normally needs no adjustment. Division likewise splits into a quadratic loop for small divisors and a recursive one past \inlinecode{DIV\_CUTOFF} limbs, so it inherits whichever multiplication the size selects. Typical arithmetic operations involving mixed numeric primitives and strings are supported through implicit construction and hidden friend operators, as long as at least one operand is a \inlinecode{BigInt} at any given level of evaluation.

Addition, subtraction, and comparisons are performed using the standard linear algorithms. Multiplication first converts limbs from base \inlinecode{BASE} to the narrower base \inlinecode{MUL\_BASE} so coefficient products fit safely, then multiplies with schoolbook, Karatsuba, or (if beyond \inlinecode{MUL\_CUTOFF}) complex FFT convolution according to the operand size, converting the result back. Division and modulo are computed together by normalized long division: scale the operands so the divisor's leading limb is large, estimate each quotient limb from the top one or two limbs of the running remainder, correct the estimate by adding the divisor back if necessary, then unscale the remainder. Beyond \inlinecode{DIV\_CUTOFF} limbs, that loop becomes the base case of Burnikel-Ziegler recursion, which halves both operands and recovers each half of the quotient from a division of the same shape, paying one multiplication per level rather than one estimate per limb.

\begin{compactitem}
  \item \inlinecode{BigInt()} constructs zero, and \inlinecode{BigInt(n)} constructs a big integer from an integer \inlinecode{n}.
  \item \inlinecode{BigInt(s)} constructs a big integer from a C string or an \inlinecode{std::string} \inlinecode{s}.
  \item \inlinecode{operator=} is defined to copy from another big integer or to assign from a 64-bit integer primitive.
  \item \inlinecode{size()} returns the number of digits in the base-10 representation.
  \item Operators \inlinecode{\textgreater{}\textgreater{}} and \inlinecode{\textless{}\textless{}} are defined to support stream-based input and output.
  \item \inlinecode{to\_string()}, \inlinecode{to\_llong()}, \inlinecode{to\_double()}, and \inlinecode{to\_ldouble()} return the big integer converted to an \inlinecode{std::string}, \inlinecode{int64\_t}, \inlinecode{double}, and \inlinecode{long double} respectively. For the latter three data types, overflow behavior is based on that of inputting from \inlinecode{std::istream}. Equivalent conversions are also available through explicit casts to \inlinecode{int}, \inlinecode{long long}, \inlinecode{double}, and \inlinecode{long double}.
  \item \inlinecode{to\_arithmetic\textless{}T\textgreater{}()} returns the value converted to arithmetic type \inlinecode{T} by streaming the base-10 representation into it, and is what the three conversions above are defined in terms of.
  \item \inlinecode{abs()} returns the absolute value.
  \item \inlinecode{comp(v)} returns $-1$, $0$, or $1$ depending on whether the big integer compares less than, equal to, or greater than \inlinecode{v}, respectively.
  \item Operators \inlinecode{\textless{}}, \inlinecode{\textgreater{}}, \inlinecode{\textless{}=}, \inlinecode{\textgreater{}=}, \inlinecode{==}, \inlinecode{!=}, \inlinecode{+}, \inlinecode{-}, \inlinecode{*}, \inlinecode{/}, \inlinecode{\%}, \inlinecode{++}, \inlinecode{--}, \inlinecode{+=}, \inlinecode{-=}, \inlinecode{*=}, \inlinecode{/=}, and \inlinecode{\%=} are defined analogous to those on integer primitives.
  \item \inlinecode{div(v)} returns a pair consisting of the quotient and remainder after dividing by \inlinecode{v}.
  \item \inlinecode{pow(n)} returns the big integer raised to the nonnegative power \inlinecode{n} using binary exponentiation.
  \item \inlinecode{mul\_pow10(n)} returns the big integer multiplied by $10^{{\inlinecodemath{n}}}$ for nonnegative \inlinecode{n}, which costs a limb shift rather than a multiplication.
  \item \inlinecode{sqrt()} returns the integral part of the square root using a digit-by-digit algorithm in the internal base.
  \item \inlinecode{nth\_root(n)} returns the integral part of the \inlinecode{n}-th root using Newton's method. Negative inputs require odd \inlinecode{n}.
  \item \inlinecode{rand(d)} returns a random, positive big integer with \inlinecode{d} digits.
\end{compactitem}

\Needspace*{3\baselineskip}
\textbf{Time Complexity}:
\begin{compactitem}
  \item $O(1)$ per call to \inlinecode{size()}.
  \item $O(d)$ per call to the constructors, \inlinecode{to\_string()}, \inlinecode{to\_llong()}, \inlinecode{to\_double()}, \inlinecode{to\_ldouble()}, \inlinecode{to\_arithmetic()}, \inlinecode{abs()}, \inlinecode{comp()}, \inlinecode{rand()}, and all comparison and arithmetic operators except multiplication, division, and modulo, where $d$ is the total number of digits in the arguments and result for each operation.
  \item $O(d \cdot \log\left(d\right) \cdot \log\left(\log\left(d\right)\right))$ per call to multiplication operations once operands reach \inlinecode{MUL\_CUTOFF}, and $O(d^{1.585})$ below it.
  \item $O((d/m) \cdot M(m) \cdot \log m)$ per call to division and modulo operations once the divisor reaches \inlinecode{DIV\_CUTOFF}, and $O(d \cdot m)$ below it, where $d$ and $m$ are the number of digits in the dividend and divisor respectively, and $M(m)$ is the cost of one multiplication of two $m$-digit big integers. Both forms consume the dividend in $d/m$ blocks of the divisor's size; only the per-block cost differs.
  \item $O(M(d) \log n)$ per call to \inlinecode{pow()}, where $n$ is the exponent and $d$ is the maximum digit length reached during exponentiation.
  \item $O(d + n)$ per call to \inlinecode{mul\_pow10()}, where $n$ is the power of ten.
  \item $O(d^{2})$ per call to \inlinecode{sqrt()}.
  \item $O(M(d) \cdot \log\left(d\right) \cdot \log\left(d \cdot n\right))$ per call to \inlinecode{nth\_root()}, where $n$ is the root. Newton's iteration starts within a factor of ten of the answer and doubles the correct digits each step, so it converges in $O(\log d)$ steps, each costing one exponentiation and one division.
\end{compactitem}

\Needspace*{3\baselineskip}
\textbf{Space Complexity}:
\begin{compactitem}
  \item $O(d)$ for storage of the big integer.
  \item $O(d)$ auxiliary for negation, addition, subtraction, multiplication, division, \inlinecode{abs()}, \inlinecode{sqrt()}, \inlinecode{pow()}, and \inlinecode{nth\_root()}.
  \item $O(1)$ auxiliary for all other operations.
\end{compactitem}

\begin{lstlisting}[style=cpp]
#include <algorithm>
#include <cassert>
#include <chrono>
#include <cmath>
#include <complex>
#include <cstdint>
#include <cstring>
#include <iomanip>
#include <istream>
#include <iterator>
#include <ostream>
#include <random>
#include <sstream>
#include <string>
#include <tuple>
#include <utility>
#include <vector>

class BigInt {
  static const int BASE = 1000000000, BASE_DIGITS = 9;  // BASE must equal 10^BASE_DIGITS.
  // Multiplication rebases to half as many decimal digits per limb, so that coefficient products
  // stay exact both in int64_t and within the FFT's double precision.
  static const int MUL_BASE = 10000, MUL_BASE_DIGITS = 4;  // Keep consistent with BASE_DIGITS.
  static const int MUL_CUTOFF = 256;  // Padded MUL_BASE length past which FFT beats Karatsuba.
  static const int DIV_CUTOFF = 64;   // Divisor limb count past which recursion wins.
  static_assert(
      BASE >= 100 && BASE <= 1000000000,
      "BASE must be between 10^2 and 10^9; below that MUL_BASE_DIGITS would round down to zero"
  );
  static_assert(
      MUL_BASE_DIGITS == BASE_DIGITS / 2 && MUL_BASE <= 10000,
      "MUL_BASE must hold half of BASE's digits and stay exact under the FFT"
  );

  using vint = std::vector<int>;
  using vint64 = std::vector<int64_t>;
  using vcd = std::vector<std::complex<double>>;

  vint digits;
  int sign;

  void normalize() {
    while (!digits.empty() && digits.back() == 0) {
      digits.pop_back();
    }
    if (digits.empty()) {
      sign = 1;
    }
  }

  void read(int n, const char *s) {
    sign = 1;
    digits.clear();
    int pos = 0;
    while (pos < n && (s[pos] == '-' || s[pos] == '+')) {
      if (s[pos] == '-') {
        sign = -sign;
      }
      pos++;
    }
    for (int i = n - 1; i >= pos; i -= BASE_DIGITS) {
      int x = 0;
      for (int j = std::max(pos, i - BASE_DIGITS + 1); j <= i; j++) {
        x = x * 10 + s[j] - '0';
      }
      digits.push_back(x);
    }
    normalize();
  }

  static int comp(const vint &a, const vint &b, int asign, int bsign) {
    if (asign != bsign) {
      return asign < bsign ? -1 : 1;
    }
    if (a.size() != b.size()) {
      return a.size() < b.size() ? -asign : asign;
    }
    for (int i = static_cast<int>(a.size()) - 1; i >= 0; i--) {
      if (a[i] != b[i]) {
        return a[i] < b[i] ? -asign : asign;
      }
    }
    return 0;
  }

  static BigInt add(const vint &a, const vint &b, int asign, int bsign) {
    if (asign != bsign) {
      return (asign == 1) ? sub(a, b, asign, 1) : sub(b, a, bsign, 1);
    }
    BigInt res;
    res.digits = a;
    res.sign = asign;
    int carry = 0, size = static_cast<int>(std::max(a.size(), b.size()));
    for (int i = 0; i < size || carry; i++) {
      if (i == static_cast<int>(res.digits.size())) {
        res.digits.push_back(0);
      }
      res.digits[i] += carry + (i < static_cast<int>(b.size()) ? b[i] : 0);
      carry = (res.digits[i] >= BASE) ? 1 : 0;
      if (carry) {
        res.digits[i] -= BASE;
      }
    }
    return res;
  }

  static BigInt sub(const vint &a, const vint &b, int asign, int bsign) {
    if (asign == -1 || bsign == -1) {
      return add(a, b, asign, -bsign);
    }
    BigInt res;
    if (comp(a, b, asign, bsign) < 0) {
      res = sub(b, a, bsign, asign);
      res.sign = -1;
      return res;
    }
    res.digits = a;
    res.sign = asign;
    for (int i = 0, borrow = 0; i < static_cast<int>(a.size()) || borrow; i++) {
      res.digits[i] -= borrow + (i < static_cast<int>(b.size()) ? b[i] : 0);
      borrow = res.digits[i] < 0;
      if (borrow) {
        res.digits[i] += BASE;
      }
    }
    res.normalize();
    return res;
  }

  static vint convert_base(const vint &digits, int l1, int l2) {
    vint64 p(std::max(l1, l2) + 1);
    p[0] = 1;
    for (int i = 1; i < static_cast<int>(p.size()); i++) {
      p[i] = p[i - 1] * 10;
    }
    vint res;
    int64_t curr = 0;
    for (int i = 0, curr_digits = 0; i < static_cast<int>(digits.size()); i++) {
      curr += digits[i] * p[curr_digits];
      curr_digits += l1;
      while (curr_digits >= l2) {
        res.push_back(static_cast<int>((curr % p[l2])));
        curr /= p[l2];
        curr_digits -= l2;
      }
    }
    res.push_back(static_cast<int>(curr));
    while (!res.empty() && res.back() == 0) {
      res.pop_back();
    }
    return res;
  }

  template<typename It>
  static vint64 karatsuba(It alo, It ahi, It blo, It bhi) {
    int n = std::distance(alo, ahi), k = n / 2;
    vint64 res(n * 2);
    if (n <= 32) {
      for (int i = 0; i < n; i++) {
        for (int j = 0; j < n; j++) {
          res[i + j] += alo[i] * blo[j];
        }
      }
      return res;
    }
    auto a1b1 = karatsuba(alo, alo + k, blo, blo + k);
    auto a2b2 = karatsuba(alo + k, ahi, blo + k, bhi);
    vint64 a2(alo + k, ahi), b2(blo + k, bhi);
    for (int i = 0; i < k; i++) {
      a2[i] += alo[i];
      b2[i] += blo[i];
    }
    auto r = karatsuba(a2.begin(), a2.end(), b2.begin(), b2.end());
    for (int i = 0; i < static_cast<int>(a1b1.size()); i++) {
      r[i] -= a1b1[i];
      res[i] += a1b1[i];
    }
    for (int i = 0; i < static_cast<int>(a2b2.size()); i++) {
      r[i] -= a2b2[i];
      res[i + n] += a2b2[i];
    }
    for (int i = 0; i < static_cast<int>(r.size()); i++) {
      res[i + k] += r[i];
    }
    return res;
  }

  template<typename It>
  static vcd fft(It lo, It hi, bool invert = false) {
    int n = std::distance(lo, hi), k = 0, high1 = -1;
    while ((1 << k) < n) {
      k++;
    }
    std::vector<int> rev(n);
    for (int i = 1; i < n; i++) {
      if (!(i & (i - 1))) {
        high1++;
      }
      rev[i] = rev[i ^ (1 << high1)];
      rev[i] |= (1 << (k - high1 - 1));
    }
    vcd roots(n), res(n);
    for (int i = 0; i < n; i++) {
      double alpha = 2 * 3.14159265358979323846 * i / n;
      roots[i] = std::complex<double>(std::cos(alpha), std::sin(alpha));
      res[i] = *(lo + rev[i]);
    }
    for (int len = 1; len < n; len <<= 1) {
      vcd tmp(n);
      int rstep = static_cast<int>(roots.size()) / (len << 1);
      for (int pdest = 0; pdest < n; pdest += len) {
        int p = pdest;
        for (int i = 0; i < len; i++) {
          std::complex<double> c = roots[i * rstep] * res[p + len];
          tmp[pdest] = res[p] + c;
          tmp[pdest + len] = res[p] - c;
          pdest++;
          p++;
        }
      }
      res.swap(tmp);
    }
    if (invert) {
      for (int i = 0; i < static_cast<int>(res.size()); i++) {
        res[i] /= n;
      }
      std::reverse(res.begin() + 1, res.end());
    }
    return res;
  }

  // Schoolbook long division of nonnegative values, quadratic in the limb counts.
  std::pair<BigInt, BigInt> div_schoolbook(const BigInt &v) const {
    int norm = BASE / (v.digits.back() + 1);
    BigInt an = *this * norm, bn = v * norm, q, r;
    q.digits.resize(an.digits.size());
    for (int i = static_cast<int>(an.digits.size()) - 1; i >= 0; i--) {
      r *= BASE;
      r += an.digits[i];
      int s1 = (r.digits.size() <= bn.digits.size()) ? 0 : r.digits[bn.digits.size()];
      int s2 = (r.digits.size() <= bn.digits.size() - 1) ? 0 : r.digits[bn.digits.size() - 1];
      int d = (static_cast<int64_t>(s1) * BASE + s2) / bn.digits.back();
      for (r -= bn * d; r < 0; r += bn) {
        d--;
      }
      q.digits[i] = d;
    }
    q.normalize();
    r.normalize();
    return {q, r / norm};
  }

  // Burnikel-Ziegler recursive division of nonnegative values. Balanced halves replace the
  // quadratic inner loop, paying one multiplication per level instead of one estimate per limb, so
  // division inherits whichever multiplication the operand size selects.
  std::pair<BigInt, BigInt> div_recursive(const BigInt &v) const {
    auto shifted = [](const BigInt &x, int k) {  // Multiplies by BASE^k.
      BigInt res(x);
      if (!res.digits.empty()) {
        res.digits.insert(res.digits.begin(), k, 0);
      }
      return res;
    };
    auto slice = [](const BigInt &x, int lo, int hi) {  // Limbs [lo, hi), or 0 if empty range.
      BigInt res;
      lo = std::max(lo, 0);
      hi = std::min(hi, static_cast<int>(x.digits.size()));
      if (lo < hi) {
        res.digits.assign(x.digits.begin() + lo, x.digits.begin() + hi);
      }
      res.normalize();
      return res;
    };
    // Divides at most 2n limbs by exactly n limbs whose leading limb is at least BASE / 2, where
    // the quotient is known to be below BASE^n. Each half of the quotient comes from one
    // three-by-two block division, which itself divides two half-size blocks by one.
    auto rec = [&](auto &&rec, const BigInt &a, const BigInt &b,
                   int n) -> std::pair<BigInt, BigInt> {
      if (n % 2 != 0 || n < DIV_CUTOFF) {
        return a.div_schoolbook(b);
      }
      int k = n / 2;
      BigInt hi = slice(b, k, 2 * k), lo = slice(b, 0, k);
      auto block = [&](const BigInt &x) {  // At most 3k limbs by b, with quotient below BASE^k.
        BigInt q, r;
        if (slice(x, 2 * k, 3 * k).comp(hi) < 0) {
          std::tie(q, r) = rec(rec, slice(x, k, 3 * k), hi, k);
        } else {
          // The estimate saturates: the leading limbs already reach the divisor, so take the
          // largest quotient the block can hold and recover the matching remainder.
          q = shifted(BigInt(1), k) - 1;
          r = slice(x, k, 3 * k) - hi * q;
        }
        r = shifted(r, k) + slice(x, 0, k) - q * lo;
        while (r < 0) {  // At most twice, since the divisor's leading limb is at least BASE / 2.
          r += b;
          q -= 1;
        }
        return std::make_pair(q, r);
      };
      auto top = block(slice(a, k, 4 * k));
      auto bottom = block(shifted(top.second, k) + slice(a, 0, k));
      return {shifted(top.first, k) + bottom.first, bottom.second};
    };
    // Scale so the divisor's leading limb exceeds BASE / 2, then pad it to a multiple of a power
    // of two so that every recursive split stays even until the blocks fall under DIV_CUTOFF, and
    // consume the dividend one block at a time.
    int scale = BASE / (v.digits.back() + 1);
    BigInt a = *this * scale, b = v * scale;
    int limbs = static_cast<int>(b.digits.size()), block = 1;
    while (block * DIV_CUTOFF < limbs) {
      block *= 2;
    }
    int n = (limbs + block - 1) / block * block, sigma = n - limbs;
    a = shifted(a, sigma);
    b = shifted(b, sigma);
    BigInt q, r;
    for (int i = (static_cast<int>(a.digits.size()) + n - 1) / n - 1; i >= 0; i--) {
      auto step = rec(rec, shifted(r, n) + slice(a, i * n, (i + 1) * n), b, n);
      q = shifted(q, n) + step.first;
      r = step.second;
    }
    // Both scalings scale the remainder and leave the quotient alone.
    return {q, slice(r, sigma, static_cast<int>(r.digits.size())) / scale};
  }

 public:
  BigInt() : sign(1) {}
  BigInt(int v) { *this = static_cast<int64_t>(v); }
  BigInt(int64_t v) { *this = v; }
  BigInt(const char *s) { read(strlen(s), s); }
  BigInt(const std::string &s) { read(s.size(), s.c_str()); }

  BigInt &operator=(int64_t v) {
    sign = v < 0 ? -1 : 1;
    digits.clear();
    for (; v != 0; v /= BASE) {
      int limb = static_cast<int>(v % BASE);
      digits.push_back(limb < 0 ? -limb : limb);
    }
    return *this;
  }

  int size() const {
    return digits.empty() ? 1
                          : static_cast<int>(std::to_string(digits.back()).size()) +
                                BASE_DIGITS * (static_cast<int>(digits.size()) - 1);
  }

  friend std::istream &operator>>(std::istream &in, BigInt &v) {
    std::string s;
    in >> s;
    v.read(s.size(), s.c_str());
    return in;
  }

  friend std::ostream &operator<<(std::ostream &out, const BigInt &v) {
    return out << v.to_string();
  }

  std::string to_string() const {
    std::ostringstream oss;
    if (sign == -1) {
      oss << '-';
    }
    oss << (digits.empty() ? 0 : digits.back());
    for (int i = static_cast<int>(digits.size()) - 2; i >= 0; i--) {
      oss << std::setw(BASE_DIGITS) << std::setfill('0') << digits[i];
    }
    return oss.str();
  }

  template<typename T>
  T to_arithmetic() const {
    std::stringstream ss(to_string());
    T res;
    ss >> res;
    return res;
  }

  int64_t to_llong() const { return to_arithmetic<int64_t>(); }
  double to_double() const { return to_arithmetic<double>(); }
  long double to_ldouble() const { return to_arithmetic<long double>(); }

  explicit operator int() const { return static_cast<int>(to_llong()); }
  explicit operator long long() const { return static_cast<long long>(to_llong()); }
  explicit operator double() const { return to_double(); }
  explicit operator long double() const { return to_ldouble(); }

  int comp(const BigInt &v) const { return comp(digits, v.digits, sign, v.sign); }

  // The comparison and binary arithmetic operators are hidden friends, so a raw integer operand on
  // either side converts through the implicit constructor.
  friend bool operator<(const BigInt &a, const BigInt &b) { return a.comp(b) < 0; }
  friend bool operator>(const BigInt &a, const BigInt &b) { return a.comp(b) > 0; }
  friend bool operator<=(const BigInt &a, const BigInt &b) { return a.comp(b) <= 0; }
  friend bool operator>=(const BigInt &a, const BigInt &b) { return a.comp(b) >= 0; }
  friend bool operator==(const BigInt &a, const BigInt &b) { return a.comp(b) == 0; }
  friend bool operator!=(const BigInt &a, const BigInt &b) { return a.comp(b) != 0; }

  BigInt abs() const {
    BigInt res(*this);
    res.sign = 1;
    return res;
  }

  BigInt operator-() const {
    BigInt res(*this);
    if (!digits.empty()) res.sign = -sign;
    return res;
  }

  friend BigInt operator+(const BigInt &a, const BigInt &b) {
    return add(a.digits, b.digits, a.sign, b.sign);
  }

  friend BigInt operator-(const BigInt &a, const BigInt &b) {
    return sub(a.digits, b.digits, a.sign, b.sign);
  }

  BigInt &operator*=(int v) {
    int64_t factor = v;
    if (factor < 0) {
      sign = -sign;
      factor = -factor;
    }
    int64_t carry = 0;
    for (int i = 0; i < static_cast<int>(digits.size()) || carry; i++) {
      if (i == static_cast<int>(digits.size())) {
        digits.push_back(0);
      }
      int64_t curr = digits[i] * factor + carry;
      carry = curr / BASE;
      digits[i] = static_cast<int>((curr % BASE));
    }
    normalize();
    return *this;
  }

  BigInt operator*(int v) const {
    BigInt res(*this);
    res *= v;
    return res;
  }

  friend BigInt operator*(const BigInt &u, const BigInt &v) {
    vint a = convert_base(u.digits, BASE_DIGITS, MUL_BASE_DIGITS);
    vint b = convert_base(v.digits, BASE_DIGITS, MUL_BASE_DIGITS);
    int n = 1;
    while (n < 2 * static_cast<int>(std::max(a.size(), b.size()))) {
      n <<= 1;
    }
    a.resize(n, 0);
    b.resize(n, 0);
    vint64 c;
    if (n < MUL_CUTOFF) {
      c = karatsuba(a.begin(), a.end(), b.begin(), b.end());
    } else {
      auto at = fft(a.begin(), a.end()), bt = fft(b.begin(), b.end());
      for (int i = 0; i < n; i++) {
        at[i] *= bt[i];
      }
      at = fft(at.begin(), at.end(), true);
      c.resize(n);
      for (int i = 0; i < n; i++) {
        c[i] = at[i].real() + 0.5;
      }
    }
    BigInt res;
    res.sign = u.sign * v.sign;
    for (int i = 0, carry = 0; i < static_cast<int>(c.size()); i++) {
      int64_t d = c[i] + carry;
      res.digits.push_back(d % MUL_BASE);
      carry = d / MUL_BASE;
    }
    res.digits = convert_base(res.digits, MUL_BASE_DIGITS, BASE_DIGITS);
    res.normalize();
    return res;
  }

  BigInt &operator/=(int v) {
    assert(v != 0);
    int64_t divisor = v;
    if (divisor < 0) {
      sign = -sign;
      divisor = -divisor;
    }
    int64_t rem = 0;
    for (int i = static_cast<int>(digits.size()) - 1; i >= 0; i--) {
      int64_t curr = digits[i] + rem * static_cast<int64_t>(BASE);
      digits[i] = static_cast<int>((curr / divisor));
      rem = curr % divisor;
    }
    normalize();
    return *this;
  }

  BigInt operator/(int v) const {
    BigInt res(*this);
    res /= v;
    return res;
  }

  int operator%(int v) const {
    assert(v != 0);
    int64_t divisor = v;
    if (divisor < 0) {
      divisor = -divisor;
    }
    int64_t m = 0;
    for (int i = static_cast<int>(digits.size()) - 1; i >= 0; i--) {
      m = (digits[i] + m * static_cast<int64_t>(BASE)) % divisor;
    }
    return static_cast<int>(m * sign);
  }

  std::pair<BigInt, BigInt> div(const BigInt &v) const {
    assert(v != 0);
    if (comp(digits, v.digits, 1, 1) < 0) {
      return {BigInt(0), *this};
    }
    auto res = static_cast<int>(v.digits.size()) < DIV_CUTOFF ? abs().div_schoolbook(v.abs())
                                                              : abs().div_recursive(v.abs());
    res.first.sign = sign * v.sign;
    res.second.sign = sign;
    res.first.normalize();
    res.second.normalize();
    return res;
  }

  friend BigInt operator/(const BigInt &a, const BigInt &b) { return a.div(b).first; }
  friend BigInt operator%(const BigInt &a, const BigInt &b) { return a.div(b).second; }
  BigInt operator++(int){ BigInt t(*this); operator++(); return t; }
  BigInt operator--(int){ BigInt t(*this); operator--(); return t; }
  BigInt &operator++() { *this = *this + BigInt(1); return *this; }
  BigInt &operator--() { *this = *this - BigInt(1); return *this; }
  BigInt &operator+=(const BigInt &v) { *this = *this + v; return *this; }
  BigInt &operator-=(const BigInt &v) { *this = *this - v; return *this; }
  BigInt &operator*=(const BigInt &v) { *this = *this * v; return *this; }
  BigInt &operator/=(const BigInt &v) { *this = *this / v; return *this; }
  BigInt &operator%=(const BigInt &v) { *this = *this % v; return *this; }

  BigInt pow(int n) const {
    assert(n >= 0);
    if (n == 0) {
      return BigInt(1);
    }
    if (*this == 0) {
      return BigInt(0);
    }
    BigInt x(*this), res(1);
    for (; n != 0; n >>= 1) {
      if (n & 1) {
        res *= x;
      }
      x *= x;
    }
    return res;
  }

  // Multiplying by a power of ten is a limb shift, since BASE is a power of ten by construction.
  BigInt mul_pow10(int n) const {
    assert(n >= 0);
    static const int POW10[] = {1, 10, 100, 1000, 10000, 100000, 1000000, 10000000, 100000000};
    BigInt res(*this);
    res.digits.insert(res.digits.begin(), n / BASE_DIGITS, 0);
    res *= POW10[n % BASE_DIGITS];  // Also normalizes, which is what clears a shifted zero.
    return res;
  }

  BigInt sqrt() const {
    assert(sign != -1);
    BigInt v(*this);
    while (v.digits.empty() || v.digits.size() % 2 == 1) {
      v.digits.push_back(0);
    }
    int n = static_cast<int>(v.digits.size());
    int ldig =
        static_cast<int>(::sqrt(static_cast<double>(v.digits[n - 1]) * BASE + v.digits[n - 2]));
    int norm = BASE / (ldig + 1);
    v *= norm;
    v *= norm;
    while (v.digits.empty() || v.digits.size() % 2 == 1) {
      v.digits.push_back(0);
    }
    BigInt r(static_cast<int64_t>(v.digits[n - 1]) * BASE + v.digits[n - 2]);
    int q = ldig =
        static_cast<int>(::sqrt(static_cast<double>(v.digits[n - 1]) * BASE + v.digits[n - 2]));
    BigInt res;
    for (int j = n / 2 - 1; j >= 0; j--) {
      for (;; q--) {
        BigInt r1 =
            (r - (res * 2 * BASE + q) * q) * BASE * BASE +
            (j > 0 ? static_cast<int64_t>(v.digits[2 * j - 1]) * BASE + v.digits[2 * j - 2] : 0);
        if (r1 >= 0) {
          r = r1;
          break;
        }
      }
      res = res * BASE + q;
      if (j > 0) {
        int sz1 = static_cast<int>(res.digits.size());
        int sz2 = static_cast<int>(r.digits.size());
        int d1 = (sz1 + 2 < sz2) ? r.digits[sz1 + 2] : 0;
        int d2 = (sz1 + 1 < sz2) ? r.digits[sz1 + 1] : 0;
        int d3 = (sz1 < sz2) ? r.digits[sz1] : 0;
        q = (static_cast<int64_t>(d1) * BASE * BASE + static_cast<int64_t>(d2) * BASE + d3) /
            (ldig * 2);
      }
    }
    res.normalize();
    return res / norm;
  }

  BigInt nth_root(int n) const {
    assert(n > 0 && (sign != -1 || n % 2 == 1));
    if (*this == 0) {
      return BigInt(0);
    }
    // Newton's iteration, started just above the answer and decreasing to it.
    BigInt magnitude = abs();
    BigInt x = BigInt(10).pow(static_cast<int>(std::ceil(static_cast<double>(size()) / n))), y;
    while (x.comp(y = (x * (n - 1) + magnitude / x.pow(n - 1)) / n) > 0) {
      x = y;
    }
    return sign == -1 ? -x : x;
  }

  static BigInt rand(int d) {
    if (d == 0) {
      return BigInt(0);
    }
    static std::mt19937 rng(std::chrono::steady_clock::now().time_since_epoch().count());
    std::uniform_int_distribution<int> first_digit(1, 9), digit(0, 9);
    std::string s(1, static_cast<char>('0' + first_digit(rng)));
    for (int i = 1; i < d; i++) {
      s += static_cast<char>('0' + digit(rng));
    }
    return BigInt(s);
  }

  friend int comp(const BigInt &a, const BigInt &b) { return a.comp(b); }
  friend BigInt abs(const BigInt &v) { return v.abs(); }
  friend BigInt pow(const BigInt &v, int n) { return v.pow(n); }
  friend BigInt sqrt(const BigInt &v) { return v.sqrt(); }
  friend BigInt nth_root(const BigInt &v, int n) { return v.nth_root(n); }
};
\end{lstlisting}
\Needspace*{4\baselineskip}
\textbf{Example Usage}\par
\begin{lstlisting}[style=cpp]
#include <climits>

int main() {
  BigInt a("-9899819294989142124"), b("12398124981294214");
  assert(a + b == "-9887421170007847910");
  assert(a - b == "-9912217419970436338");
  assert(a * b == "-122739196911503356525379735104870536");
  assert(a / b == "-798");

  // Raw integers and strings work on either side of comparisons and arithmetic.
  assert("12398124981294214" == b);
  assert(5 + BigInt(7) == 12);
  assert(-100 < b && 5 % BigInt(3) == 2);

  assert(BigInt(20).pow(12345).size() == 16062);
  assert(BigInt("9812985918924981892491829").nth_root(4) == 1769906);
  assert(BigInt(-8).nth_root(3) == -2);
  for (int i = -100; i <= 100; i++) {
    if (i >= 0) {
      assert(BigInt(i).sqrt() == static_cast<int>(sqrt(i)));
    }
    for (int j = -100; j <= 100; j++) {
      assert(BigInt(i) + BigInt(j) == i + j);
      assert(BigInt(i) - BigInt(j) == i - j);
      assert(BigInt(i) * BigInt(j) == i * j);
      if (j != 0) {
        assert(BigInt(i) / BigInt(j) == i / j);
      }
      if (0 < i && i <= 10 && 0 < j && j <= 10) {
        assert(BigInt(i).nth_root(j) == static_cast<int64_t>((pow(i, 1.0 / j) + 1E-5)));
        int64_t p = 1;
        for (int k = 0; k < j; k++) {
          p *= i;
        }
        assert(BigInt(i).pow(j) == p);
      }
    }
  }
  std::mt19937 rng(1234567);  // Fixed seed for reproducibility.
  std::uniform_int_distribution<int> len_dist(1, 100);
  for (int i = 0; i < 20; i++) {
    int d = len_dist(rng);
    BigInt val(BigInt::rand(d)), root(val.sqrt()), lower(root * root), upper(root + 1);
    upper *= upper;
    assert(lower <= val && val < upper);
    std::uniform_int_distribution<int> divisor_len_dist(1, d);
    BigInt divisor(BigInt::rand(divisor_len_dist(rng)) + 1), quotient(val / divisor);
    lower = quotient * divisor;
    upper = divisor * (quotient + 1);
    assert(val >= lower && val < upper);
  }
  BigInt x(-6);
  assert(x.to_string() == "-6");
  assert(x.to_llong() == -6LL);
  assert(x.to_double() == -6.0);
  assert(x.to_ldouble() == -6.0);
  assert(static_cast<int>(x) == -6);
  assert(static_cast<long long>(x) == -6LL);
  assert(static_cast<double>(x) == -6.0);
  assert(BigInt(INT64_MIN).to_string() == "-9223372036854775808");
  assert(BigInt(INT64_MIN).to_llong() == INT64_MIN);
  assert(-BigInt(0) == 0);
  assert(BigInt(3) * INT_MIN == "-6442450944");
  assert(BigInt("6442450944") / INT_MIN == -3);
  assert(BigInt("6442450945") % INT_MIN == 1);
  return 0;
}
\end{lstlisting}
\subsection{Big Decimal}
\label{sec:6.4.3}
Performs exact decimal arithmetic using the full \inlinecode{BigInt} implementation from \secref{6.4.2} as its coefficient type, together with a nonnegative scale. Unlike binary floating point, every parsed decimal is represented exactly: for example, \inlinecode{BigDecimal("1.20")} stores coefficient $120$ with scale $2$. Addition and subtraction preserve the larger operand scale, multiplication adds the scales, and formatting preserves trailing zeros unless \inlinecode{normalized()} is called. Only finite, well-formed decimal strings are supported; \inlinecode{NaN} and infinities are not represented.

Division and square root generally have nonterminating decimal expansions, so both operations require an explicit result scale. The rounding mode defaults to \inlinecode{HalfAway}, which rounds ties away from zero. The other modes are \inlinecode{TowardZero}, \inlinecode{AwayFromZero}, \inlinecode{Floor}, \inlinecode{Ceil}, and \inlinecode{HalfEven}.

\begin{compactitem}
  \item \inlinecode{BigDecimal(n = 0)} constructs a decimal from integer \inlinecode{n}.
  \item \inlinecode{BigDecimal(s)} parses decimal string \inlinecode{s}, optionally written in scientific notation, and preserves trailing zeros in its mantissa.
  \item \inlinecode{decimal\_places()}, \inlinecode{precision()}, and \inlinecode{sign()} return the scale, the number of coefficient digits, and the sign as $-1$, $0$, or $1$, respectively.
  \item \inlinecode{unscaled\_value()} returns the underlying \inlinecode{BigInt} coefficient.
  \item \inlinecode{abs()} and \inlinecode{normalized()} return the absolute value and a value with trailing fractional zeros removed.
  \item \inlinecode{move\_point(places)} moves the decimal point right by \inlinecode{places} positions, or left if \inlinecode{places} is negative.
  \item \inlinecode{rescale(new\_scale, mode = Round::HalfAway)} returns the value with exactly \inlinecode{new\_scale} decimal places, rounding if the scale is reduced.
  \item \inlinecode{trunc()}, \inlinecode{floor()}, and \inlinecode{ceil()} return the corresponding integer as a \inlinecode{BigInt}.
  \item \inlinecode{divide(v, result\_scale, mode = Round::HalfAway)} divides by nonzero \inlinecode{v} and rounds to exactly \inlinecode{result\_scale} decimal places.
  \item \inlinecode{sqrt(result\_scale, mode = Round::HalfAway)} returns the nonnegative square root rounded to exactly \inlinecode{result\_scale} decimal places.
  \item \inlinecode{pow(n)} returns the value raised to nonnegative integer power \inlinecode{n}.
  \item \inlinecode{to\_string()} returns ordinary decimal notation with the stored number of decimal places.
  \item Stream operators and comparisons are defined, as are \inlinecode{+}, \inlinecode{-}, \inlinecode{*}, \inlinecode{\%}, and their compound assignments.
\end{compactitem}

Overflow warning: The stored scale and every decimal-shift magnitude must fit in \inlinecode{int}.

\Needspace*{3\baselineskip}
\textbf{Time Complexity}:
\begin{compactitem}
  \item $O(1)$ per call to \inlinecode{decimal\_places()}, \inlinecode{precision()}, \inlinecode{sign()}, and \inlinecode{unscaled\_value()}.
  \item $O(d)$ per call to parsing, formatting, normalization, addition, subtraction, comparisons, and decimal-point movement, where $d$ is the number of relevant decimal digits.
  \item $O(M(d) \log n)$ per call to \inlinecode{pow(n)}, where $M(d)$ is the multiplication cost for the largest intermediate coefficient.
  \item $O(d^{2})$ per call to square root, matching the \inlinecode{BigInt} operation used here.
  \item Division, remainder, and rescaling to fewer places inherit the \inlinecode{BigInt} division cost of \secref{6.4.2}: $O((d/m) \cdot M(m) \cdot \log m)$ once the divisor reaches its \inlinecode{DIV\_CUTOFF} and $O(d \cdot m)$ below it, where $d$ and $m$ are the number of digits in the dividend and divisor.
\end{compactitem}

\textbf{Space Complexity}: $O(d)$ for stored digits, returned values, and arithmetic results.

\begin{lstlisting}[style=cpp]
#define main bigint_example
#include "6.4.2_Big_Integer.cpp"
#undef main

#include <climits>

class BigDecimal {
 public:
  enum class Round { TowardZero, AwayFromZero, Floor, Ceil, HalfAway, HalfEven };

 private:
  BigInt coeff;
  int scale;

  static int checked_scale(int64_t scale) {
    assert(0 <= scale && scale <= INT_MAX);
    return static_cast<int>(scale);
  }

  BigDecimal(BigInt coeff, int scale) : coeff(std::move(coeff)), scale(scale) {}

  BigInt scaled(int new_scale) const { return coeff.mul_pow10(new_scale - scale); }
  static BigInt pow10(int n) { return BigInt(10).pow(n); }

  static bool round_up(bool negative, int half_cmp, const BigInt &q, Round mode) {
    return mode == Round::AwayFromZero || (mode == Round::Floor && negative) ||
           (mode == Round::Ceil && !negative) || (mode == Round::HalfAway && half_cmp >= 0) ||
           (mode == Round::HalfEven && (half_cmp > 0 || (half_cmp == 0 && q % 2 != 0)));
  }

  static BigInt rounded_quotient(const BigInt &num, const BigInt &den, Round mode) {
    assert(den > 0);
    bool negative = num < 0;
    auto [q, r] = num.abs().div(den);
    bool increment = r != 0 && round_up(negative, (2 * r).comp(den), q, mode);
    return negative ? -(q + increment) : q + increment;
  }

 public:
  BigDecimal(int64_t n = 0) : coeff(n), scale(0) {}

  explicit BigDecimal(const std::string &s) : coeff(0), scale(0) {
    assert(s.size() <= INT_MAX);
    size_t e = s.find_first_of("eE");
    std::string mantissa = s.substr(0, e);
    int exponent = e == std::string::npos ? 0 : std::stoi(s.substr(e + 1));
    int pos = 0;
    bool negative = false, point = false;
    if (pos < static_cast<int>(mantissa.size()) && (mantissa[pos] == '+' || mantissa[pos] == '-')) {
      negative = mantissa[pos++] == '-';
    }
    std::string digits;
    for (; pos < static_cast<int>(mantissa.size()); pos++) {
      if (mantissa[pos] == '.') {
        assert(!point);
        point = true;
      } else {
        assert('0' <= mantissa[pos] && mantissa[pos] <= '9');
        digits += mantissa[pos];
        scale += point;
      }
    }
    assert(!digits.empty());
    coeff = negative ? -BigInt(digits) : BigInt(digits);
    int64_t adjusted_scale = static_cast<int64_t>(scale) - exponent;
    if (adjusted_scale < 0) {
      coeff = coeff.mul_pow10(checked_scale(-adjusted_scale));
      scale = 0;
    } else {
      scale = checked_scale(adjusted_scale);
    }
  }

  int decimal_places() const { return scale; }
  int precision() const { return coeff.size(); }
  int sign() const { return (coeff > 0) - (coeff < 0); }
  const BigInt &unscaled_value() const { return coeff; }
  BigDecimal abs() const { return BigDecimal(coeff.abs(), scale); }

  BigDecimal normalized() const {
    BigInt c = coeff;
    int s = scale;
    while (s > 0 && c % 10 == 0) {
      c /= 10;
      s--;
    }
    return BigDecimal(c, s);
  }

  BigDecimal move_point(int places) const {
    int64_t new_scale = static_cast<int64_t>(scale) - places;
    return new_scale >= 0 ? BigDecimal(coeff, checked_scale(new_scale))
                          : BigDecimal(coeff.mul_pow10(checked_scale(-new_scale)), 0);
  }

  BigDecimal rescale(int new_scale, Round mode = Round::HalfAway) const {
    assert(new_scale >= 0);
    return new_scale >= scale
               ? BigDecimal(scaled(new_scale), new_scale)
               : BigDecimal(rounded_quotient(coeff, pow10(scale - new_scale), mode), new_scale);
  }

  BigInt trunc() const { return rescale(0, Round::TowardZero).coeff; }
  BigInt floor() const { return rescale(0, Round::Floor).coeff; }
  BigInt ceil() const { return rescale(0, Round::Ceil).coeff; }

  BigDecimal divide(const BigDecimal &v, int result_scale, Round mode = Round::HalfAway) const {
    assert(v.coeff != 0 && result_scale >= 0);
    BigInt num = coeff.abs(), den = v.coeff.abs();
    int64_t shift = static_cast<int64_t>(v.scale) + result_scale - scale;
    (shift >= 0 ? num : den) =
        (shift >= 0 ? num : den).mul_pow10(checked_scale(shift < 0 ? -shift : shift));
    if ((coeff < 0) != (v.coeff < 0)) {
      num = -num;
    }
    return BigDecimal(rounded_quotient(num, den, mode), result_scale);
  }

  BigDecimal sqrt(int result_scale, Round mode = Round::HalfAway) const {
    assert(coeff >= 0 && result_scale >= 0);
    BigInt num = coeff, den = 1;
    int64_t shift = 2LL * result_scale - scale;
    (shift >= 0 ? num : den) =
        (shift >= 0 ? num : den).mul_pow10(checked_scale(shift < 0 ? -shift : shift));
    BigInt q = (num / den).sqrt();
    bool exact = q * q * den == num;
    int half_cmp = (4 * num).comp(den * (4 * q * q + 4 * q + 1));
    return BigDecimal(q + (!exact && round_up(false, half_cmp, q, mode)), result_scale);
  }

  std::string to_string() const {
    std::string s = coeff.abs().to_string();
    if (static_cast<int>(s.size()) <= scale) {
      s.insert(0, scale + 1 - s.size(), '0');
    }
    if (scale > 0) {
      s.insert(s.size() - scale, 1, '.');
    }
    return coeff < 0 ? "-" + s : s;
  }

  friend std::istream &operator>>(std::istream &in, BigDecimal &v) {
    std::string s;
    if (in >> s) {
      v = BigDecimal(s);
    }
    return in;
  }

  friend std::ostream &operator<<(std::ostream &out, const BigDecimal &v) {
    return out << v.to_string();
  }

  friend bool operator<(const BigDecimal &a, const BigDecimal &b) {
    int s = std::max(a.scale, b.scale);
    return a.scaled(s) < b.scaled(s);
  }

  friend bool operator==(const BigDecimal &a, const BigDecimal &b) {
    int s = std::max(a.scale, b.scale);
    return a.scaled(s) == b.scaled(s);
  }

  friend bool operator>(const BigDecimal &a, const BigDecimal &b) { return b < a; }
  friend bool operator<=(const BigDecimal &a, const BigDecimal &b) { return !(b < a); }
  friend bool operator>=(const BigDecimal &a, const BigDecimal &b) { return !(a < b); }
  friend bool operator!=(const BigDecimal &a, const BigDecimal &b) { return !(a == b); }

  friend BigDecimal operator+(const BigDecimal &a, const BigDecimal &b) {
    int s = std::max(a.scale, b.scale);
    return BigDecimal(a.scaled(s) + b.scaled(s), s);
  }

  friend BigDecimal operator-(const BigDecimal &a, const BigDecimal &b) { return a + -b; }

  friend BigDecimal operator*(const BigDecimal &a, const BigDecimal &b) {
    return BigDecimal(a.coeff * b.coeff, checked_scale(static_cast<int64_t>(a.scale) + b.scale));
  }

  friend BigDecimal operator%(const BigDecimal &a, const BigDecimal &b) {
    assert(b.coeff != 0);
    int s = std::max(a.scale, b.scale);
    return BigDecimal(a.scaled(s) % b.scaled(s), s);
  }

  BigDecimal operator-() const { return BigDecimal(-coeff, scale); }

  BigDecimal pow(int n) const {
    assert(n >= 0);
    return BigDecimal(coeff.pow(n), checked_scale(static_cast<int64_t>(scale) * n));
  }

  BigDecimal &operator+=(const BigDecimal &v) { return *this = *this + v; }
  BigDecimal &operator-=(const BigDecimal &v) { return *this = *this - v; }
  BigDecimal &operator*=(const BigDecimal &v) { return *this = *this * v; }
  BigDecimal &operator%=(const BigDecimal &v) { return *this = *this % v; }
};
\end{lstlisting}
\Needspace*{4\baselineskip}
\textbf{Example Usage}\par
\begin{lstlisting}[style=cpp]
using namespace std;

int main() {
  using Round = BigDecimal::Round;

  assert(BigDecimal().to_string() == "0");
  assert(BigDecimal(-12).to_string() == "-12");
  assert(BigDecimal("+.50").to_string() == "0.50");
  assert(BigDecimal("1.20e3").to_string() == "1200");
  assert(BigDecimal("1.20e-2").to_string() == "0.0120");

  BigDecimal a("1.20"), b("3.4");
  assert(a.decimal_places() == 2 && a.precision() == 3);
  assert(a.unscaled_value() == 120);
  assert(a.sign() == 1 && BigDecimal(0).sign() == 0 && BigDecimal(-1).sign() == -1);
  assert(BigDecimal("-1.20").abs().to_string() == "1.20");
  assert(a.normalized().to_string() == "1.2");
  assert(a == BigDecimal("1.2") && a != b);
  assert(a < b && b > a && a <= BigDecimal("1.200") && b >= a);

  assert((a + b).to_string() == "4.60");
  assert((b - a).to_string() == "2.20");
  assert((a * b).to_string() == "4.080");
  assert((-a).to_string() == "-1.20");
  assert((BigDecimal("5.50") % 2).to_string() == "1.50");

  BigDecimal x = a;
  x += b;
  assert(x.to_string() == "4.60");
  x -= b;
  assert(x.to_string() == "1.20");
  x *= b;
  assert(x.to_string() == "4.080");
  x %= 1.0;
  assert(x.to_string() == "0.080");

  assert(a.move_point(2).to_string() == "120");
  assert(a.move_point(-1).to_string() == "0.120");
  assert(a.rescale(4).to_string() == "1.2000");
  assert(BigDecimal("1.29").rescale(1, Round::TowardZero).to_string() == "1.2");
  assert(BigDecimal("1.21").rescale(1, Round::AwayFromZero).to_string() == "1.3");
  assert(BigDecimal("-1.21").rescale(1, Round::Floor).to_string() == "-1.3");
  assert(BigDecimal("-1.21").rescale(1, Round::Ceil).to_string() == "-1.2");
  assert(BigDecimal("1.25").rescale(1, Round::HalfAway).to_string() == "1.3");
  assert(BigDecimal("2.685").rescale(2, Round::HalfEven).to_string() == "2.68");

  assert(BigDecimal(1).divide(8, 4).to_string() == "0.1250");
  assert(BigDecimal(2).divide(3, 4).to_string() == "0.6667");
  assert(BigDecimal(-1).divide(8, 2).to_string() == "-0.13");
  assert(BigDecimal(1).divide(8, 2, Round::HalfEven).to_string() == "0.12");

  assert(BigDecimal("-1.2").trunc() == -1);
  assert(BigDecimal("-1.2").floor() == -2);
  assert(BigDecimal("-1.2").ceil() == -1);
  assert(a.pow(3).to_string() == "1.728000");

  assert(BigDecimal(2).sqrt(5).to_string() == "1.41421");
  assert(BigDecimal("2.25").sqrt(2).to_string() == "1.50");
  assert(BigDecimal("2.25").sqrt(0, Round::HalfAway).to_string() == "2");
  assert(BigDecimal("6.25").sqrt(0, Round::HalfEven).to_string() == "2");
  assert(BigDecimal(2).sqrt(0, Round::Ceil).to_string() == "2");
  for (int i = 0; i <= 100; i++) {
    assert(BigDecimal(i * i).sqrt(0) == BigDecimal(i));
  }

  stringstream in("-12.340"), out;
  BigDecimal streamed;
  in >> streamed;
  out << streamed;
  assert(out.str() == "-12.340");

  assert((BigDecimal("12345678901234567890.12") * 10).to_string() == "123456789012345678901.20");
  return 0;
}
\end{lstlisting}
\subsection{Rational Numbers}
\label{sec:6.4.4}
Perform operations on rational numbers internally represented as two integers: a numerator and a denominator. The template integer type must support streamed input/output, comparisons, and arithmetic operations.

\begin{compactitem}
  \item \inlinecode{Rational\textless{}Int\textgreater{}(n)} constructs a rational number with numerator \inlinecode{n} and denominator $1$.
  \item \inlinecode{Rational\textless{}Int\textgreater{}(n, d)} constructs a rational number with numerator \inlinecode{n} and denominator \inlinecode{d}.
  \item \inlinecode{operator\textgreater{}\textgreater{}} inputs a rational number using the next integer from the stream as the numerator and 1 as the denominator.
  \item \inlinecode{operator\textless{}\textless{}} outputs the rational number as a string consisting of possibly a minus sign followed by the numerator, followed by a slash, followed by the denominator.
  \item \inlinecode{to\_string()}, \inlinecode{to\_llong()}, \inlinecode{to\_double()}, and \inlinecode{to\_ldouble()} return the rational converted to an \inlinecode{std::string}, \inlinecode{int64\_t}, \inlinecode{double}, and \inlinecode{long double} respectively with potential truncation or approximation for the primitive types. Equivalent conversions are also available through explicit casts to \inlinecode{int}, \inlinecode{long long}, \inlinecode{double}, and \inlinecode{long double}.
  \item \inlinecode{to\_arithmetic\textless{}T\textgreater{}()} is the general form behind the three numeric conversions above, returning the rational as any arithmetic type \inlinecode{T} readable from a stream. Integral \inlinecode{T} truncates toward zero.
  \item \inlinecode{abs()}, \inlinecode{floor()}, and \inlinecode{ceil()} return the absolute value, floor, and ceiling.
  \item Operators \inlinecode{\textless{}}, \inlinecode{\textgreater{}}, \inlinecode{\textless{}=}, \inlinecode{\textgreater{}=}, \inlinecode{==}, \inlinecode{!=}, \inlinecode{+}, \inlinecode{-}, \inlinecode{*}, \inlinecode{/}, \inlinecode{\%}, \inlinecode{++}, \inlinecode{--}, \inlinecode{+=}, \inlinecode{-=}, \inlinecode{*=}, \inlinecode{/=}, and \inlinecode{\%=} are defined analogous to those on numerical primitives. The comparisons and binary arithmetic are hidden friends, so a raw integer operand works on either side.
\end{compactitem}

Overflow warning: Internal operations do not check for overflow. Comparisons and arithmetic cross-multiply numerators and denominators, so instantiate with a wider integer type (such as \inlinecode{\_\_int128}, or even \inlinecode{BigInt}) if the values may grow large.

\Needspace*{3\baselineskip}
\textbf{Time Complexity}:
\begin{compactitem}
  \item $O(\log s)$ operations on \inlinecode{Int} per call to the constructor \inlinecode{Rational(n, d)}, where $s = |n| + |d|$, spent reducing to lowest terms.
  \item $O(1)$ operations on \inlinecode{Int} per call to everything else.
\end{compactitem}

\Needspace*{3\baselineskip}
\textbf{Space Complexity}:
\begin{compactitem}
  \item Two \inlinecode{Int} values for storage of the rational number.
  \item $O(1)$ auxiliary \inlinecode{Int} values for all operations.
\end{compactitem}

\begin{lstlisting}[style=cpp]
#include <cassert>
#include <cstdint>
#include <istream>
#include <ostream>
#include <sstream>
#include <string>

template<typename Int>
class Rational {
  Int num, den;

 public:
  Rational() : num(0), den(1) {}
  Rational(const Int &n) : num(n), den(1) {}

  Rational(const Int &n, const Int &d) : num(n), den(d) {
    assert(den != 0);
    if (den < 0) {
      num = -num;
      den = -den;
    }
    Int a(num < 0 ? -num : num), b(den), tmp;
    while (a != 0 && b != 0) {
      tmp = a % b;
      a = b;
      b = tmp;
    }
    Int gcd = (b == 0) ? a : b;
    num /= gcd;
    den /= gcd;
  }

  friend std::istream &operator>>(std::istream &in, Rational &r) {
    in >> r.num;
    r.den = 1;
    return in;
  }

  friend std::ostream &operator<<(std::ostream &out, const Rational &r) {
    out << r.num << "/" << r.den;
    return out;
  }

  std::string to_string() const {
    std::stringstream ss;
    ss << num << "/" << den;
    return ss.str();
  }

  // Round-trips through a stringstream so that any Int supporting streamed I/O (e.g. a big-integer
  // type) converts, even without a direct cast to the target type. Integral T truncates.
  template<typename T>
  T to_arithmetic() const {
    std::stringstream ss;
    ss << num << " " << den;
    T n, d;
    ss >> n >> d;
    return n / d;
  }

  int64_t to_llong() const { return to_arithmetic<int64_t>(); }
  double to_double() const { return to_arithmetic<double>(); }
  long double to_ldouble() const { return to_arithmetic<long double>(); }

  explicit operator int() const { return static_cast<int>(to_llong()); }
  explicit operator long long() const { return static_cast<long long>(to_llong()); }
  explicit operator double() const { return to_double(); }
  explicit operator long double() const { return to_ldouble(); }

  // The comparison and binary arithmetic operators are hidden friends, so a raw integer operand on
  // either side converts through the implicit constructor.
  friend bool operator<(const Rational &a, const Rational &b) {
    return a.num * b.den < b.num * a.den;
  }

  friend bool operator==(const Rational &a, const Rational &b) {
    return a.num == b.num && a.den == b.den;
  }

  friend bool operator>(const Rational &a, const Rational &b) { return b < a; }
  friend bool operator<=(const Rational &a, const Rational &b) { return !(b < a); }
  friend bool operator>=(const Rational &a, const Rational &b) { return !(a < b); }
  friend bool operator!=(const Rational &a, const Rational &b) { return !(a == b); }

  Rational abs() const { return Rational(num < 0 ? -num : num, den); }
  friend Rational abs(const Rational &r) { return r.abs(); }
  Int floor() const { return num < 0 ? -((-num + den - 1) / den) : num / den; }
  Int ceil() const { return num < 0 ? -(-num / den) : (num + den - 1) / den; }

  friend Rational operator+(const Rational &a, const Rational &b) {
    return Rational(a.num * b.den + b.num * a.den, a.den * b.den);
  }

  friend Rational operator-(const Rational &a, const Rational &b) {
    return Rational(a.num * b.den - b.num * a.den, a.den * b.den);
  }

  friend Rational operator*(const Rational &a, const Rational &b) {
    return Rational(a.num * b.num, a.den * b.den);
  }

  friend Rational operator/(const Rational &a, const Rational &b) {
    return Rational(a.num * b.den, a.den * b.num);
  }

  friend Rational operator%(const Rational &a, const Rational &b) {
    return a - b * Rational(a.num * b.den / (b.num * a.den), 1);
  }

  Rational operator-() const { return Rational(-num, den); }
  Rational operator++(int) { Rational t(*this); operator++(); return t; }
  Rational operator--(int) { Rational t(*this); operator--(); return t; }
  Rational &operator++() { *this = *this + 1; return *this; }
  Rational &operator--() { *this = *this - 1; return *this; }
  Rational &operator+=(const Rational &r) { *this = *this + r; return *this; }
  Rational &operator-=(const Rational &r) { *this = *this - r; return *this; }
  Rational &operator*=(const Rational &r) { *this = *this * r; return *this; }
  Rational &operator/=(const Rational &r) { *this = *this / r; return *this; }
  Rational &operator%=(const Rational &r) { *this = *this % r; return *this; }
};
\end{lstlisting}
\Needspace*{4\baselineskip}
\textbf{Example Usage}\par
\begin{lstlisting}[style=cpp]
#include <cmath>
using namespace std;

bool EQ(double a, double b) {
  return fabs(a - b) < 1e-9;
}

int main() {
  using Rational = Rational<int64_t>;
  assert(Rational(-21, 1) % 2 == -1);
  Rational r(Rational(-53, 10) % Rational(-17, 10));
  assert(EQ(r.to_ldouble(), fmod(-5.3, -1.7)));
  assert(r.to_string() == "-1/5");
  assert(static_cast<int>(Rational(7, 2)) == 3);
  assert(static_cast<long long>(Rational(7, 2)) == 3LL);
  assert(EQ(static_cast<double>(Rational(1, 2)), 0.5));

  // Raw integers work on either side of comparisons and arithmetic.
  assert(2 + Rational(1, 2) == Rational(5, 2));
  assert(1 < Rational(3, 2) && Rational(3, 2) < 2);
  assert(2 % Rational(3, 4) == Rational(1, 2));

  // Construction reduces the fraction and moves any sign onto the numerator.
  assert(Rational(2, 4).to_string() == "1/2");
  assert(Rational(1, -2).to_string() == "-1/2");
  assert(Rational(5).to_string() == "5/1");

  // Arithmetic between rationals reduces the result too.
  assert(Rational(1, 2) + Rational(1, 3) == Rational(5, 6));
  assert(Rational(1, 2) - Rational(1, 3) == Rational(1, 6));
  assert(Rational(2, 3) * Rational(3, 4) == Rational(1, 2));
  assert(Rational(1, 2) / Rational(3, 4) == Rational(2, 3));

  // Rounding takes a different branch on each sign, and is exact on whole numbers.
  assert(Rational(7, 2).floor() == 3 && Rational(7, 2).ceil() == 4);
  assert(Rational(-7, 2).floor() == -4 && Rational(-7, 2).ceil() == -3);
  assert(Rational(4, 2).floor() == 2 && Rational(4, 2).ceil() == 2);
  assert(Rational(-7, 2).abs() == Rational(7, 2) && abs(Rational(-7, 2)) == Rational(7, 2));

  // The named conversions and the general form behind them.
  assert(Rational(7, 2).to_llong() == 3 && EQ(Rational(7, 2).to_double(), 3.5));
  assert(Rational(7, 2).to_arithmetic<int>() == 3);

  // Stream input reads one integer, and output prints the reduced fraction.
  Rational streamed;
  istringstream in("7");
  in >> streamed;
  assert(streamed == 7);
  ostringstream out;
  out << Rational(-3, 6);
  assert(out.str() == "-1/2");
  return 0;
}
\end{lstlisting}

\section{\texorpdfstring{Linear Algebra}{6.5 Linear Algebra}}
\setcounter{section}{5}
\setcounter{subsection}{0}
\subsection{Matrix Operations}
\label{sec:6.5.1}
Basic matrix operations defined on a two-dimensional vector of numeric values. Matrix arithmetic also supports modular value types such as \inlinecode{Modular\textless{}MOD\textgreater{}} from \secref{6.3.3}. The element type must be constructible from $0$ and $1$ and support the arithmetic operators used by the requested operation.

\begin{compactitem}
  \item \inlinecode{make\_matrix\textless{}T\textgreater{}(m, n, v)} constructs and returns an $m$ by $n$ matrix with 0-based indices (row indices $[0, {\inlinecodemath{m}})$ and column indices $[0, {\inlinecodemath{n}})$), where every value is initialized to \inlinecode{v}.
  \item \inlinecode{identity\_matrix\textless{}T\textgreater{}(n)} returns the \inlinecode{n} by \inlinecode{n} identity matrix, that is, a matrix where each \inlinecode{a[i][j]} equals $1$ if \inlinecode{i} = \inlinecode{j}, or $0$ otherwise.
  \item \inlinecode{rows(a)} returns the number of rows in matrix \inlinecode{a}.
  \item \inlinecode{cols(a)} returns the number of columns in matrix \inlinecode{a}.
  \item \inlinecode{a[i][j]} may be used to access or modify the specified entry of \inlinecode{a}.
  \item Operators \inlinecode{\textless{}}, \inlinecode{\textgreater{}}, \inlinecode{\textless{}=}, \inlinecode{\textgreater{}=}, \inlinecode{==}, and \inlinecode{!=} define lexicographical comparison based on that of \inlinecode{std::vector}.
  \item Operators \inlinecode{+}, \inlinecode{-}, \inlinecode{*}, \inlinecode{/}, \inlinecode{+=}, \inlinecode{-=}, \inlinecode{*=}, and \inlinecode{/=} define scalar addition, subtraction, multiplication, and division involving a matrix and a numeric scalar value.
  \item \inlinecode{a * v} returns the matrix-vector product as a vector.
  \item Operators \inlinecode{*} and \inlinecode{*=} define matrix multiplication.
\end{compactitem}

Multiplication is the schoolbook triple loop, which is $O(m \cdot n \cdot p)$ and is what contest constraints assume. Strassen's algorithm reaches $O(n^{2.807})$ by splitting each matrix into quarters and combining them with seven recursive products instead of eight, and Toom-Cook style refinements go lower still, but the additions that replace the saved multiplication carry large constants and worsen numerical stability. The crossover sits in the hundreds of rows even in tuned libraries, so the plain loop below is the right default; reach for a blocked or cache-oblivious ordering before reaching for a subcubic algorithm.

Exponentiation uses iterative binary exponentiation: keep an accumulated result, square the base each round, and multiply it when the current exponent bit is set. The power sum uses the block identity that $\begin{bmatrix} A & A \\ 0 & I \end{bmatrix}^p$ has $A + A^2 + \dots + A^p$ in its upper-right block.

\begin{compactitem}
  \item Operators \inlinecode{\textasciicircum{}} and \inlinecode{\textasciicircum{}=} define iterative binary matrix exponentiation of a square matrix \inlinecode{a} by a \inlinecode{uint64\_t} power \inlinecode{p}.
  \item \inlinecode{power\_sum(a, p)} returns the power sum of a square matrix \inlinecode{a} up to a \inlinecode{uint64\_t} power \inlinecode{p}, that is, $a + a^2 + \ldots + a^p$.
\end{compactitem}

\Needspace*{3\baselineskip}
\textbf{Time Complexity}:
\begin{compactitem}
  \item $O(m \cdot n)$ per construction, output, comparison, or scalar-arithmetic operation on $m$ by $n$ matrices.
  \item $O(1)$ per call to \inlinecode{rows()} and \inlinecode{cols()}.
  \item $O(m \cdot n)$ per matrix-matrix addition or subtraction of $m$ by $n$ matrices.
  \item $O(n^{3} \log p)$ per exponentiation of an $n$ by $n$ matrix to power $p$.
  \item $O(n^{3} \log p)$ per call to \inlinecode{power\_sum()} for an $n$ by $n$ matrix and power $p$.
  \item $O(m \cdot n)$ per multiplication of an $m$ by $n$ matrix by a vector of length $n$.
  \item $O(m \cdot n \cdot k)$ per multiplication of an $m$ by $n$ matrix by an $n$ by $k$ matrix.
\end{compactitem}

\Needspace*{3\baselineskip}
\textbf{Space Complexity}:
\begin{compactitem}
  \item $O(1)$ auxiliary for \inlinecode{rows()}, \inlinecode{cols()}, \inlinecode{a[i][j]} access, comparison and scalar compound operators.
  \item $O(n^{2})$ auxiliary for exponentiation of an $n$ by $n$ matrix to power $p$.
  \item $O(n^{2})$ auxiliary for \inlinecode{power\_sum()} of an $n$ by $n$ matrix up to power $p$.
  \item $O(m)$ for the vector returned by matrix-vector multiplication.
  \item $O(m \cdot n)$ auxiliary for all non-in-place operations returning an $m$ by $n$ matrix.
\end{compactitem}

\begin{lstlisting}[style=cpp]
#include <cassert>
#include <cstdint>
#include <iomanip>
#include <ostream>
#include <vector>

template<typename T>
using Matrix = std::vector<std::vector<T>>;

template<typename T = int>
Matrix<T> make_matrix(int m, int n, const T &v = T{}) {
  assert(m >= 0 && n >= 0);
  return Matrix<T>(m, std::vector<T>(n, v));
}

template<typename T = int>
Matrix<T> identity_matrix(int n) {
  Matrix<T> res = make_matrix<T>(n, n);
  for (int i = 0; i < n; i++) {
    res[i][i] = 1;
  }
  return res;
}

template<typename T>
int rows(const Matrix<T> &a) {
  return static_cast<int>(a.size());
}

template<typename T>
int cols(const Matrix<T> &a) {
  return a.empty() ? 0 : static_cast<int>(a[0].size());
}

template<typename T>
std::ostream &operator<<(std::ostream &out, const Matrix<T> &a) {
  auto flags = out.flags();
  auto precision = out.precision();
  out << std::fixed << std::setprecision(5);
  for (int i = 0; i < rows(a); i++) {
    for (int j = 0; j < cols(a); j++) {
      out << std::setw(10) << a[i][j];
    }
    out << std::endl;
  }
  out.flags(flags);
  out.precision(precision);
  return out;
}

template<typename T, typename U>
Matrix<T> &operator+=(Matrix<T> &a, const U &v) {
  for (int i = 0; i < rows(a); i++) {
    for (int j = 0; j < cols(a); j++) {
      a[i][j] += v;
    }
  }
  return a;
}

template<typename T, typename U>
Matrix<T> &operator-=(Matrix<T> &a, const U &v) {
  for (int i = 0; i < rows(a); i++) {
    for (int j = 0; j < cols(a); j++) {
      a[i][j] -= v;
    }
  }
  return a;
}

template<typename T, typename U>
Matrix<T> &operator*=(Matrix<T> &a, const U &v) {
  for (int i = 0; i < rows(a); i++) {
    for (int j = 0; j < cols(a); j++) {
      a[i][j] *= v;
    }
  }
  return a;
}

template<typename T, typename U>
Matrix<T> &operator/=(Matrix<T> &a, const U &v) {
  for (int i = 0; i < rows(a); i++) {
    for (int j = 0; j < cols(a); j++) {
      a[i][j] /= v;
    }
  }
  return a;
}

template<typename T>
Matrix<T> &operator+=(Matrix<T> &a, const Matrix<T> &b) {
  assert(rows(a) == rows(b) && cols(a) == cols(b));
  for (int i = 0; i < rows(a); i++) {
    for (int j = 0; j < cols(a); j++) {
      a[i][j] += b[i][j];
    }
  }
  return a;
}

template<typename T>
Matrix<T> &operator-=(Matrix<T> &a, const Matrix<T> &b) {
  assert(rows(a) == rows(b) && cols(a) == cols(b));
  for (int i = 0; i < rows(a); i++) {
    for (int j = 0; j < cols(a); j++) {
      a[i][j] -= b[i][j];
    }
  }
  return a;
}

template<typename T>
Matrix<T> operator+(const Matrix<T> &a, const Matrix<T> &b) {
  Matrix<T> c(a);
  return c += b;
}

template<typename T>
Matrix<T> operator-(const Matrix<T> &a, const Matrix<T> &b) {
  Matrix<T> c(a);
  return c -= b;
}

template<typename T>
std::vector<T> operator*(const Matrix<T> &a, const std::vector<T> &v) {
  assert(cols(a) == static_cast<int>(v.size()));
  std::vector<T> res(rows(a));
  for (int i = 0; i < rows(a); i++) {
    for (int j = 0; j < cols(a); j++) {
      res[i] += a[i][j] * v[j];
    }
  }
  return res;
}

template<typename T>
Matrix<T> operator*(const Matrix<T> &a, const Matrix<T> &b) {
  assert(cols(a) == rows(b));
  Matrix<T> res = make_matrix<T>(rows(a), cols(b));
  // The k loop sits outside j so the inner loop walks res[i] and b[k] contiguously instead of
  // striding down a column of b. Each res[i][j] still accumulates in increasing k, so results are
  // unchanged, including for floating-point types.
  for (int i = 0; i < rows(a); i++) {
    for (int k = 0; k < rows(b); k++) {
      for (int j = 0; j < cols(b); j++) {
        res[i][j] += a[i][k] * b[k][j];
      }
    }
  }
  return res;
}

template<typename T>
Matrix<T> &operator*=(Matrix<T> &a, const Matrix<T> &b) { return a = a * b; }

template<typename T, typename U>
Matrix<T> operator+(const Matrix<T> &a, const U &v) { Matrix<T> m(a); return m += v; }

template<typename T, typename U>
Matrix<T> operator-(const Matrix<T> &a, const U &v) { Matrix<T> m(a); return m -= v; }

template<typename T, typename U>
Matrix<T> operator*(const Matrix<T> &a, const U &v) { Matrix<T> m(a); return m *= v; }

template<typename T, typename U>
Matrix<T> operator/(const Matrix<T> &a, const U &v) { Matrix<T> m(a); return m /= v; }

template<typename T, typename U>
Matrix<T> operator+(const U &v, const Matrix<T> &a) { return a + v; }

template<typename T, typename U>
Matrix<T> operator*(const U &v, const Matrix<T> &a) { return a * v; }

template<typename T, typename U>
Matrix<T> operator-(const U &v, const Matrix<T> &a) {
  Matrix<T> m = make_matrix<T>(rows(a), cols(a), v);
  return m -= a;
}

template<typename T>
Matrix<T> operator^(Matrix<T> a, uint64_t p) {
  assert(rows(a) == cols(a));
  Matrix<T> res = identity_matrix<T>(rows(a));
  while (p > 0) {
    if (p & 1) {
      res *= a;
    }
    p >>= 1;
    if (p > 0) {
      a *= a;
    }
  }
  return res;
}

template<typename T>
Matrix<T> &operator^=(Matrix<T> &a, uint64_t p) {
  return a = a ^ p;
}

template<typename T>
Matrix<T> power_sum(const Matrix<T> &a, uint64_t p) {
  assert(rows(a) == cols(a));
  int n = rows(a);
  if (p == 0) {
    return make_matrix<T>(n, n);
  }
  // The upper-right block of [[A, A], [0, I]]^p is A + A^2 + ... + A^p.
  Matrix<T> block = make_matrix<T>(2 * n, 2 * n);
  for (int i = 0; i < n; i++) {
    block[i + n][i + n] = 1;
    for (int j = 0; j < n; j++) {
      block[i][j] = a[i][j];
      block[i][j + n] = a[i][j];
    }
  }
  block = block ^ p;
  Matrix<T> res = make_matrix<T>(n, n);
  for (int i = 0; i < n; i++) {
    for (int j = 0; j < n; j++) {
      res[i][j] = block[i][j + n];
    }
  }
  return res;
}
\end{lstlisting}
\Needspace*{4\baselineskip}
\textbf{Example Usage}\par
\begin{lstlisting}[style=cpp]
using namespace std;

int main() {
  using matrix = Matrix<int>;
  matrix m = make_matrix(5, 5, 10) + 10;
  vector<int> v{1, 2, 3, 4, 5};
  assert((m * v == vector<int>{300, 300, 300, 300, 300}));

  m[0][0] += 5;
  assert(m[0][0] == 25 && m[1][1] == 20);
  assert((m ^ 0) == identity_matrix(5));
  assert((identity_matrix<int>(3) ^ 5) == identity_matrix<int>(3));
  assert(power_sum(m, 3) == m + m * m + (m ^ 3));

  Matrix<double> d = make_matrix<double>(2, 2, 0.5);
  assert(rows(d) == 2 && cols(d) == 2);
  assert((d + 0.25)[0][0] == 0.75);
  return 0;
}
\end{lstlisting}
\subsection{Row Reduction}
\label{sec:6.5.2}
Converts a matrix to reduced row echelon form using Gaussian elimination to solve a system of linear equations and compute rank. Each round finds a row with a nonzero entry in the current leading column, normalizes that row, and subtracts multiples of it from every other row to clear the column. Floating-point matrices use the largest-magnitude candidate for partial pivoting, though elimination can still be prone to rounding error; use LU decomposition for greater numerical stability. Exact field types such as \inlinecode{Modular\textless{}MOD\textgreater{}} instead use any nonzero pivot and perform exact arithmetic. Such types must support construction from $0$ and $1$, equality, addition, subtraction, multiplication, and division.

\begin{compactitem}
  \item \inlinecode{row\_reduce(a)} assigns the matrix \inlinecode{a} to its reduced row echelon form, returning a reference to the modified argument itself.
  \item \inlinecode{matrix\_rank(a)} returns the rank of matrix \inlinecode{a}, i.e. the number of nonzero rows after row reduction.
  \item \inlinecode{solve\_system(a, b, \&x)} solves the system of linear equations $ax = b$ given an $m$ by $n$ matrix \inlinecode{a} over the real numbers or an exact field, and a length $m$ vector \inlinecode{b}, returning $0$ if there is one solution, $-1$ if there are zero solutions, or $-2$ if there are infinite solutions. If there is exactly one solution, then the output vector \inlinecode{x} is populated with the solution of length $n$.
\end{compactitem}

\Needspace*{3\baselineskip}
\textbf{Time Complexity}:
\begin{compactitem}
  \item $O(m \cdot n \cdot \min\left(m, n\right))$ per call to \inlinecode{row\_reduce()} and \inlinecode{matrix\_rank()}, where $m$ and $n$ are the numbers of rows and columns of \inlinecode{a}, respectively.
  \item $O(m \cdot n \cdot \min\left(m, n\right))$ per call to \inlinecode{solve\_system()}.
\end{compactitem}

\Needspace*{3\baselineskip}
\textbf{Space Complexity}:
\begin{compactitem}
  \item $O(1)$ auxiliary for \inlinecode{row\_reduce()}.
  \item $O(m \cdot n)$ auxiliary for \inlinecode{matrix\_rank()} and \inlinecode{solve\_system()}.
\end{compactitem}

\begin{lstlisting}[style=cpp]
#include <algorithm>
#include <cmath>
#include <type_traits>
#include <utility>
#include <vector>

const double EPS = 1e-9;

template<typename T>
bool rref_is_zero(const T &v) {
  if constexpr (std::is_floating_point_v<T>) {
    return std::fabs(v) < EPS;
  }
  return v == T{0};
}

template<typename Matrix>
Matrix &row_reduce(Matrix &a) {
  if (a.empty()) {
    return a;
  }
  using T = std::decay_t<decltype(a[0][0])>;
  int rows = static_cast<int>(a.size()), cols = static_cast<int>(a[0].size());
  for (int r = 0, lead = 0; r < rows && lead < cols; lead++) {
    int pivot = r;
    if constexpr (std::is_floating_point_v<T>) {
      for (int i = r + 1; i < rows; i++) {
        if (std::fabs(a[i][lead]) > std::fabs(a[pivot][lead])) {
          pivot = i;
        }
      }
    } else {
      while (pivot < rows && rref_is_zero(a[pivot][lead])) {
        pivot++;
      }
    }
    if (pivot == rows || rref_is_zero(a[pivot][lead])) {
      continue;
    }
    std::swap(a[pivot], a[r]);
    auto lv = a[r][lead];
    for (int j = 0; j < cols; j++) {
      a[r][j] /= lv;
    }
    for (int i = 0; i < rows; i++) {
      if (i != r) {
        lv = a[i][lead];
        for (int j = 0; j < cols; j++) {
          a[i][j] -= lv * a[r][j];
        }
      }
    }
    std::fill(a[r].begin(), a[r].begin() + lead, 0);
    a[r][lead] = 1;
    r++;
  }
  return a;
}

template<typename Matrix>
int matrix_rank(Matrix a) {
  row_reduce(a);
  int rank = 0;
  for (int i = 0; i < static_cast<int>(a.size()); i++) {
    for (int j = 0; j < static_cast<int>(a[i].size()); j++) {
      if (!rref_is_zero(a[i][j])) {
        rank++;
        break;
      }
    }
  }
  return rank;
}

template<typename Matrix, typename T>
int solve_system(const Matrix &a, const std::vector<T> &b, std::vector<T> *x) {
  if (x == nullptr || a.empty() || a.size() != b.size()) {
    return -1;
  }
  int rows = static_cast<int>(a.size()), cols = static_cast<int>(a[0].size());
  Matrix m(a);
  for (int i = 0; i < rows; i++) {
    m[i].push_back(b[i]);
  }
  row_reduce(m);
  int rank = 0;
  for (int i = 0; i < rows; i++) {
    int lead = -1;
    for (int j = 0; j < cols && lead < 0; j++) {
      if (!rref_is_zero(m[i][j])) {
        lead = j;
      }
    }
    if (lead < 0) {
      // A zero coefficient row with a nonzero constant means 0 = nonzero (no solution).
      if (!rref_is_zero(m[i][cols])) {
        return -1;
      }
    } else {
      rank++;
    }
  }
  // A unique solution requires a pivot in every column; fewer means free variables remain. This
  // also catches trailing free columns, which a per-row "lead > i" test alone would miss.
  if (rank < cols) {
    return -2;
  }
  x->resize(cols);
  for (int i = 0; i < cols; i++) {
    (*x)[i] = m[i][cols];
  }
  return 0;
}
\end{lstlisting}
\Needspace*{4\baselineskip}
\textbf{Example Usage}\par
\begin{lstlisting}[style=cpp]
#include <cassert>
using namespace std;

int main() {
  vector<vector<double>> a{{-1, 2, 5}, {1, 0, -6}, {-4, 2, 2}};
  vector<double> b{3, 1, -2};
  vector<double> x;
  assert(solve_system(a, b, &x) == 0);
  assert(matrix_rank(a) == 3);
  assert(matrix_rank(vector<vector<double>>{{1, 2, 3}, {2, 4, 6}, {0, 1, 1}}) == 2);
  assert(solve_system(vector<vector<double>>{{0, 0}}, vector<double>{1}, &x) == -1);
  assert(solve_system(vector<vector<double>>{{1, 0}}, vector<double>{1}, &x) == -2);
  for (int i = 0; i < static_cast<int>(a.size()); i++) {
    double sum = 0;
    for (int j = 0; j < static_cast<int>(a[i].size()); j++) {
      sum += a[i][j] * x[j];
    }
    assert(fabs(sum - b[i]) < EPS);
  }
  return 0;
}
\end{lstlisting}
\subsection{Determinant and Inverse}
\label{sec:6.5.3}
Computes the determinant and inverse of a square matrix using Gaussian elimination. The inverse of a matrix $a$ is another matrix $b$ such that $ab$ equals the identity matrix. The inverse of $a$ exists if and only if the determinant of $a$ is nonzero. In this case, $a$ is called invertible or non-singular. The determinant falls out of elimination as the product of the pivots, negated once per row swap. The inverse is found by appending the identity matrix and row-reducing the combined matrix: the operations that turn $a$ into the identity turn the identity into the inverse.

Floating-point matrices use the largest-magnitude pivot in each column, although Gaussian elimination can still suffer from rounding error; use LU decomposition for greater numerical stability. Exact field types such as \inlinecode{Modular\textless{}MOD\textgreater{}} use any nonzero pivot and perform exact arithmetic. These types must support construction from $0$ and $1$, equality, addition, subtraction, multiplication, and division. For an integer matrix whose determinant is wanted exactly, the Bareiss algorithm runs a fraction-free elimination: every intermediate value is itself a determinant of a submatrix, so the divisions are always exact and the arithmetic stays in integers.

\begin{compactitem}
  \item \inlinecode{det\_naive(a)} returns the determinant of an $n$ by $n$ matrix \inlinecode{a}, using the classic divide-and-conquer algorithm by Laplace expansions. It is division-free and computes in the matrix's element type, so an integer matrix yields an exact integer determinant; but at $O(n!)$ it is only practical for tiny $n$, so prefer \inlinecode{det\_bareiss} for exact integer determinants.
  \item \inlinecode{det(a, eps = 1e-10)} returns the determinant of an $n$ by $n$ floating-point or exact-field matrix \inlinecode{a} using Gaussian elimination. Floating-point pivots within \inlinecode{eps} of zero are treated as singular; exact fields use equality with zero.
  \item \inlinecode{det\_bareiss(a)} returns the exact determinant of an integer matrix \inlinecode{a} using fraction-free elimination, with no rounding error.
  \item \inlinecode{inverse(a, eps = 1e-10)} returns the inverse of square matrix \inlinecode{a} with floating-point or exact-field elements, or \inlinecode{std::nullopt} if the matrix is singular. Floating-point pivots within \inlinecode{eps} of zero are treated as singular.
\end{compactitem}

Overflow warning: In \inlinecode{det\_bareiss()}, stored entries are minors of \inlinecode{a}, but the products in each update can be larger. They use 128-bit intermediates when available; otherwise those products must fit in \inlinecode{int64\_t}.

\Needspace*{3\baselineskip}
\textbf{Time Complexity}:
\begin{compactitem}
  \item $O(n!)$ per call to \inlinecode{det\_naive()}, where $n$ is the dimension of the matrix.
  \item $O(n^{3})$ per call to \inlinecode{det()}, \inlinecode{det\_bareiss()}, and \inlinecode{inverse()} where $n$ is the dimension of the matrix.
\end{compactitem}

\Needspace*{3\baselineskip}
\textbf{Space Complexity}:
\begin{compactitem}
  \item $O(n)$ auxiliary stack space and $O(n^{3})$ auxiliary heap space for \inlinecode{det\_naive()}, where $n$ is the dimension of the matrix.
  \item $O(n^{2})$ auxiliary heap space for \inlinecode{det()}, \inlinecode{det\_bareiss()}, and \inlinecode{inverse()}.
\end{compactitem}

\begin{lstlisting}[style=cpp]
#include <cassert>
#include <cmath>
#include <cstdint>
#include <optional>
#include <type_traits>
#include <utility>
#include <vector>

template<typename SquareMatrix>
auto det_naive(const SquareMatrix &a) {
  int n = static_cast<int>(a.size());
  assert(n == 0 || static_cast<int>(a[0].size()) == n);
  using T = std::decay_t<decltype(a[0][0])>;
  if (n == 0) {
    return T{1};
  }
  if (n == 1) {
    return a[0][0];
  }
  if (n == 2) {
    return a[0][0] * a[1][1] - a[0][1] * a[1][0];
  }
  T res = 0;
  SquareMatrix temp(n - 1, typename SquareMatrix::value_type(n - 1));
  for (int p = 0; p < n; p++) {
    int h = 0, k = 0;
    for (int i = 1; i < n; i++) {
      for (int j = 0; j < n; j++) {
        if (j == p) {
          continue;
        }
        temp[h][k++] = a[i][j];
        if (k == n - 1) {
          h++;
          k = 0;
        }
      }
    }
    res += (p % 2 == 0 ? 1 : -1) * a[0][p] * det_naive(temp);
  }
  return res;
}

template<typename SquareMatrix>
int elimination_pivot(const SquareMatrix &a, int row, int col, double eps) {
  using T = std::decay_t<decltype(a[0][0])>;
  int n = static_cast<int>(a.size());
  if constexpr (std::is_floating_point_v<T>) {
    int pivot = row;
    for (int r = row + 1; r < n; r++) {
      if (std::fabs(a[r][col]) > std::fabs(a[pivot][col])) {
        pivot = r;
      }
    }
    return std::fabs(a[pivot][col]) <= eps ? -1 : pivot;
  }
  for (int r = row; r < n; r++) {
    if (a[r][col] != T{0}) {
      return r;
    }
  }
  return -1;
}

template<typename SquareMatrix>
auto det(const SquareMatrix &a, double eps = 1e-10) {
  int n = static_cast<int>(a.size());
  assert(n == 0 || static_cast<int>(a[0].size()) == n);
  using T = std::decay_t<decltype(a[0][0])>;
  SquareMatrix b(a);
  T res = 1;
  for (int i = 0; i < n; i++) {
    int p = elimination_pivot(b, i, i, eps);
    if (p == -1) {
      return T{0};
    }
    if (p != i) {
      std::swap(b[p], b[i]);
      res = T{0} - res;
    }
    res *= b[i][i];
    for (int j = i + 1; j < n; j++) {
      T z = b[j][i] / b[i][i];
      for (int k = i; k < n; k++) {
        b[j][k] -= z * b[i][k];
      }
    }
  }
  return res;
}

int64_t det_bareiss(std::vector<std::vector<int64_t>> a) {
  int n = static_cast<int>(a.size());
  assert(n == 0 || static_cast<int>(a[0].size()) == n);
  int64_t prev = 1, sign = 1;
  for (int k = 0; k < n; k++) {
    if (a[k][k] == 0) {  // Swap in a nonzero pivot; each swap flips the sign.
      int p = -1;
      for (int r = k + 1; r < n; r++) {
        if (a[r][k] != 0) {
          p = r;
          break;
        }
      }
      if (p == -1) {
        return 0;
      }
      std::swap(a[k], a[p]);
      sign = -sign;
    }
    for (int i = k + 1; i < n; i++) {
      for (int j = k + 1; j < n; j++) {
#if defined(__SIZEOF_INT128__)
        __extension__ typedef __int128 int128_t;
        int128_t numerator =
            static_cast<int128_t>(a[i][j]) * a[k][k] - static_cast<int128_t>(a[i][k]) * a[k][j];
        a[i][j] = static_cast<int64_t>(numerator / prev);
#else
        a[i][j] = (a[i][j] * a[k][k] - a[i][k] * a[k][j]) / prev;  // Overflow warning.
#endif
      }
    }
    prev = a[k][k];
  }
  return n == 0 ? 1 : sign * a[n - 1][n - 1];
}

template<typename SquareMatrix>
std::optional<SquareMatrix> inverse(SquareMatrix a, double eps = 1e-10) {
  int n = static_cast<int>(a.size());
  assert(n == 0 || static_cast<int>(a[0].size()) == n);
  using T = std::decay_t<decltype(a[0][0])>;
  for (int i = 0; i < n; i++) {
    a[i].resize(2 * n);
    for (int j = n; j < n * 2; j++) {
      a[i][j] = (i == j - n ? 1 : 0);
    }
  }
  for (int i = 0; i < n; i++) {
    int p = elimination_pivot(a, i, i, eps);
    if (p == -1) {
      return std::nullopt;
    }
    std::swap(a[p], a[i]);
    T pivot = a[i][i];
    for (int j = i; j < n * 2; j++) {
      a[i][j] /= pivot;
    }
    for (int j = 0; j < n; j++) {
      if (i != j) {
        T factor = a[j][i];
        for (int k = 0; k < n * 2; k++) {
          a[j][k] -= factor * a[i][k];
        }
      }
    }
  }
  for (int i = 0; i < n; i++) {
    a[i].erase(a[i].begin(), a[i].begin() + n);
  }
  return a;
}
\end{lstlisting}
\Needspace*{4\baselineskip}
\textbf{Example Usage}\par
\begin{lstlisting}[style=cpp]
#include <cassert>
using namespace std;

int main() {
  vector<vector<double>> a{{6, 1, 1}, {4, -2, 5}, {2, 8, 7}};
  int n = static_cast<int>(a.size());
  vector<vector<double>> res(n, vector<double>(n));
  assert(fabs(det(a) - det_naive(a)) < 1e-10);

  // Bareiss gives the determinant of an integer matrix exactly, with no rounding.
  vector<vector<int64_t>> ai{{6, 1, 1}, {4, -2, 5}, {2, 8, 7}};
  assert(det_naive(ai) == -306);  // Division-free, so also exact in the int64_t element type.
  assert(det_naive(vector<vector<int64_t>>{}) == 1);
  assert(det_bareiss(ai) == -306);
  assert(det_bareiss({{2, 0, 0}, {0, 3, 0}, {0, 0, 5}}) == 30);
  assert(det_bareiss({{1, 2}, {2, 4}}) == 0);  // Singular.
  auto inv = inverse(a);
  assert(inv);
  assert(!inverse(vector<vector<double>>{{1, 2}, {2, 4}}));
  for (int i = 0; i < n; i++) {
    for (int j = 0; j < n; j++) {
      for (int k = 0; k < n; k++) {
        res[i][j] += a[i][k] * (*inv)[k][j];
      }
    }
  }
  for (int i = 0; i < n; i++) {
    for (int j = 0; j < n; j++) {
      assert(fabs(res[i][j] - (i == j ? 1 : 0)) < 1e-10);
    }
  }
  return 0;
}
\end{lstlisting}
\subsection{LU Decomposition}
\label{sec:6.5.4}
The LU decomposition of a matrix $A$ with row-partial pivoting is a factorization of $A$ (after some rows are possibly permuted by a permutation matrix $p$) as a product of a lower triangular matrix $L$ and an upper triangular matrix $U$. This factorization can be used to tackle many common problems in linear algebra such as solving systems of linear equations and computing determinants. An improvement on basic row reduction, LU decomposition by row-partial pivoting keeps the relative magnitude of matrix values small, thus reducing the relative error due to rounding in computed solutions. These routines require floating-point matrix elements; use \secref{6.5.3} for exact-field or integer determinant calculations.

\begin{compactitem}
  \item \inlinecode{lu\_decompose(a, \&perm, eps = 1e-10)} assigns floating-point matrix \inlinecode{a} to merged LU decomposition matrix \inlinecode{lu}; it returns $0$ for even row-swap parity, $1$ for odd parity, or $-1$ for a degenerate matrix (i.e. singular for square matrices). The merged matrix stores \inlinecode{l[i][j]} in \inlinecode{lu[i][j]} when \inlinecode{i \textgreater{} j}, and \inlinecode{u[i][j]} otherwise. The diagonal entries \inlinecode{l[i][i]} are always $1$, so they are not explicitly stored. Access general entries of the lower and upper triangular matrices via \inlinecode{getl(lu, i, j)} and \inlinecode{getu(lu, i, j)}. Optional output vector \inlinecode{perm} represents $p$: \inlinecode{perm[i]} is the only column containing $1$ in row \inlinecode{i}. Left-multiplying \inlinecode{a} by this permutation matrix gives the product of the separate lower and upper triangular matrices, not the merged storage matrix \inlinecode{lu} itself.
  \item \inlinecode{solve\_system(a, b, \&x, eps = 1e-10)} solves the system of linear equations $Ax = b$ given an $m$ by $n$ floating-point matrix \inlinecode{a} and a length $m$ vector \inlinecode{b}, returning $0$ if there is one solution or $-1$ if there are zero or infinite solutions. If there is exactly one solution, then the output vector \inlinecode{x} is populated with the solution of length $n$; otherwise, \inlinecode{x} is unchanged.
  \item \inlinecode{det(a)} returns the determinant of an $n$ by $n$ floating-point matrix \inlinecode{a} using LU decomposition.
  \item \inlinecode{inverse(a)} returns the inverse of the $n$ by $n$ floating-point matrix \inlinecode{a}, or \inlinecode{std::nullopt} if \inlinecode{a} is singular.
\end{compactitem}

\Needspace*{3\baselineskip}
\textbf{Time Complexity}:
\begin{compactitem}
  \item $O(m \cdot n \cdot \min\left(m, n\right))$ per call to \inlinecode{lu\_decompose(a)}, where $m$ and $n$ are the numbers of rows and columns of \inlinecode{a}, respectively.
  \item $O(m \cdot n^{2})$ per call to \inlinecode{solve\_system()}, which requires $m \geq n$.
  \item $O(n^{3})$ per call to \inlinecode{det(a)} and \inlinecode{inverse(a)}, where $n$ is the length of square matrix \inlinecode{a}.
\end{compactitem}

\Needspace*{3\baselineskip}
\textbf{Space Complexity}:
\begin{compactitem}
  \item $O(1)$ auxiliary for \inlinecode{lu\_decompose()}.
  \item $O(n^{2})$ auxiliary for \inlinecode{det()} and \inlinecode{inverse()}.
  \item $O(m \cdot n)$ auxiliary for \inlinecode{solve\_system()}.
\end{compactitem}

\begin{lstlisting}[style=cpp]
#include <algorithm>
#include <cassert>
#include <cmath>
#include <numeric>
#include <optional>
#include <type_traits>
#include <vector>

template<typename Matrix>
int lu_decompose(Matrix &a, std::vector<int> *perm = nullptr, double eps = 1e-10) {
  int rows = static_cast<int>(a.size());
  int cols = a.empty() ? 0 : static_cast<int>(a[0].size());
  int parity = 0;
  if (perm != nullptr) {
    perm->resize(rows);
    std::iota(perm->begin(), perm->end(), 0);
  }
  for (int i = 0; i < rows && i < cols; i++) {
    int pi = i;
    for (int k = i + 1; k < rows; k++) {
      if (std::fabs(a[k][i]) > std::fabs(a[pi][i])) {
        pi = k;
      }
    }
    if (std::fabs(a[pi][i]) <= eps) {
      return -1;
    }
    if (pi != i) {
      if (perm != nullptr) {
        std::iter_swap(perm->begin() + i, perm->begin() + pi);
      }
      std::iter_swap(a.begin() + i, a.begin() + pi);
      parity = 1 - parity;
    }
    for (int j = i + 1; j < rows; j++) {
      a[j][i] /= a[i][i];
      for (int k = i + 1; k < cols; k++) {
        a[j][k] -= a[j][i] * a[i][k];
      }
    }
  }
  return parity;
}

template<typename Matrix>
auto getl(const Matrix &lu, int i, int j) {
  using T = std::decay_t<decltype(lu[0][0])>;
  return i > j ? lu[i][j] : T(i == j);
}

template<typename Matrix>
auto getu(const Matrix &lu, int i, int j) {
  using T = std::decay_t<decltype(lu[0][0])>;
  return i <= j ? lu[i][j] : T{0};
}

template<typename Matrix, typename T>
int solve_system(const Matrix &a, const std::vector<T> &b, std::vector<T> *x, double eps = 1e-10) {
  int rows = static_cast<int>(a.size());
  int cols = a.empty() ? 0 : static_cast<int>(a[0].size());
  if (x == nullptr || a.empty() || a.size() != b.size() || rows < cols) {
    return -1;
  }
  Matrix lu(a);
  std::vector<int> perm;
  int status = lu_decompose(lu, &perm, eps);
  if (status < 0) {
    return status;
  }
  std::vector<T> solution(cols);
  for (int i = 0; i < cols; i++) {
    solution[i] = b[perm[i]];
    for (int k = 0; k < i; k++) {
      solution[i] -= getl(lu, i, k) * solution[k];
    }
  }
  for (int i = cols - 1; i >= 0; i--) {
    for (int k = i + 1; k < cols; k++) {
      solution[i] -= getu(lu, i, k) * solution[k];
    }
    solution[i] /= getu(lu, i, i);
  }
  for (int i = 0; i < rows; i++) {
    T val = 0;
    for (int j = 0; j < cols; j++) {
      val += a[i][j] * solution[j];
    }
    // Mixed absolute/relative tolerance: dividing by b[i] alone would skip the check for negative
    // b[i] (ratio goes negative) and divide by zero when b[i] == 0.
    if (std::fabs(val - b[i]) > eps * (1.0 + std::fabs(b[i]))) {
      return -1;
    }
  }
  x->swap(solution);
  return 0;
}

template<typename Matrix>
auto det(const Matrix &a) {
  int n = static_cast<int>(a.size());
  assert(n == 0 || static_cast<int>(a[0].size()) == n);
  using T = std::decay_t<decltype(a[0][0])>;
  Matrix lu(a);
  int status = lu_decompose(lu);
  if (status < 0) {
    return T{0};
  }
  T res = 1;
  for (int i = 0; i < n; i++) {
    res *= lu[i][i];
  }
  return status == 0 ? res : -res;
}

template<typename SquareMatrix>
std::optional<SquareMatrix> inverse(SquareMatrix a) {
  int n = static_cast<int>(a.size());
  assert(n == 0 || static_cast<int>(a[0].size()) == n);
  std::vector<int> perm;
  int status = lu_decompose(a, &perm);
  if (status < 0) {
    return std::nullopt;
  }
  SquareMatrix ia(n, typename SquareMatrix::value_type(n, 0));
  for (int j = 0; j < n; j++) {
    for (int i = 0; i < n; i++) {
      if (perm[i] == j) {
        ia[i][j] = 1.0;
      } else {
        ia[i][j] = 0.0;
      }
      for (int k = 0; k < i; k++) {
        ia[i][j] -= getl(a, i, k) * ia[k][j];
      }
    }
    for (int i = n - 1; i >= 0; i--) {
      for (int k = i + 1; k < n; k++) {
        ia[i][j] -= getu(a, i, k) * ia[k][j];
      }
      ia[i][j] /= getu(a, i, i);
    }
  }
  return ia;
}
\end{lstlisting}
\Needspace*{4\baselineskip}
\textbf{Example Usage}\par
\begin{lstlisting}[style=cpp]
#include <cassert>
using namespace std;

bool EQ(double a, double b) {
  return fabs(a - b) < 1e-9;
}

int main() {
  {  // Solve a system.
    vector<vector<double>> a{{-1, 2, 5}, {1, 0, -6}, {-4, 2, 2}};
    vector<double> b{3, 1, -2};
    vector<double> x;
    assert(solve_system(a, b, &x) == 0);
    for (int i = 0; i < static_cast<int>(a.size()); i++) {
      double sum = 0;
      for (int j = 0; j < static_cast<int>(a[0].size()); j++) {
        sum += a[i][j] * x[j];
      }
      assert(EQ(sum, b[i]));
    }
  }
  {  // Find the determinant.
    vector<vector<double>> a{{1, 3, 5}, {2, 4, 7}, {1, 1, 0}};
    assert(EQ(det(a), 4));
  }
  {  // Find the inverse.
    vector<vector<double>> a{{6, 1, 1}, {4, -2, 5}, {2, 8, 7}};
    auto inv = inverse(a);
    int n = static_cast<int>(a.size());
    vector<vector<double>> res(n, vector<double>(n));
    assert(inv);
    for (int i = 0; i < n; i++) {
      for (int j = 0; j < n; j++) {
        for (int k = 0; k < n; k++) {
          res[i][j] += a[i][k] * (*inv)[k][j];
        }
      }
    }
    for (int i = 0; i < n; i++) {
      for (int j = 0; j < n; j++) {
        assert(EQ(res[i][j], i == j ? 1 : 0));
      }
    }
  }
  {  // Singular matrices have no inverse.
    vector<vector<double>> singular{{1, 2}, {2, 4}};
    assert(!inverse(singular));
    // Failed operations leave their output arguments unchanged.
    vector<vector<double>> inconsistent{{1}, {1}};
    vector<double> b{1, 2}, x{99, 100};
    assert(solve_system(inconsistent, b, &x) == -1);
    assert((x == vector<double>{99, 100}));
  }
  return 0;
}
\end{lstlisting}
\subsection{QR Decomposition and Least Squares}
\label{sec:6.5.5}
The QR decomposition factors an $m$ by $n$ matrix $A$ with $m \geq n$ into $A = QR$, where $Q$ has orthonormal columns and $R$ is upper triangular. Orthonormal columns are what make the factorization useful: $Q$ preserves lengths, so multiplying by $Q^T$ rewrites a problem in a new basis without distorting distances, and a triangular system is solved by back substitution.

This implementation uses Householder reflections. A Householder reflection is the mirror across the hyperplane orthogonal to a chosen vector $v$, written $I - 2vv^T/(v^Tv)$, and $v$ is chosen so that the reflection maps the part of a column below the diagonal onto the axis, zeroing it in one step. Applying one reflection per column leaves $R$, and accumulating the reflections gives $Q$. Reflecting is stable because it never scales anything: the alternative of orthogonalizing column by column with Gram-Schmidt loses orthogonality as the columns approach linear dependence, whereas reflections keep it to within rounding error.

The overdetermined system $Ax = b$ usually has no solution, so least squares asks for the $x$ whose residual is shortest. Because $Q$ preserves lengths, $\|Ax - b\|$ equals $\|Rx - Q^Tb\|$, and the entries of $Q^Tb$ below row $n$ cannot be changed by any $x$; minimizing therefore means solving the square triangular system formed by the first $n$ rows, which back substitution does directly.

\begin{compactitem}
  \item \inlinecode{qr\_decompose(a, \&q, \&r)} factors the $m$ by $n$ matrix \inlinecode{a} with $m \geq n$ into an $m$ by $m$ orthogonal matrix \inlinecode{q} and an $m$ by $n$ upper triangular matrix \inlinecode{r}, whose product is \inlinecode{a}.
  \item \inlinecode{least\_squares(a, b, eps = 1e-10)} returns the vector $x$ minimizing the Euclidean norm of ${\inlinecodemath{a}}x - {\inlinecodemath{b}}$, or \inlinecode{std::nullopt} if the columns of \inlinecode{a} are linearly dependent to within \inlinecode{eps}. The matrix \inlinecode{a} must have at least as many rows as columns, and \inlinecode{b} must have one entry per row.
\end{compactitem}

Fitting any linear model reduces to this: one row per data point, one column per basis function. Solving the normal equations $A^TAx = A^Tb$ by the row reduction of section \secref{6.5.2} gives the same answer but squares the condition number, losing roughly twice as many digits.

\Needspace*{3\baselineskip}
\textbf{Time Complexity}:
\begin{compactitem}
  \item $O(m^{2} \cdot n)$ per call to \inlinecode{qr\_decompose()}, where $m$ and $n$ are the numbers of rows and columns.
  \item $O(m \cdot n^{2})$ per call to \inlinecode{least\_squares()}, since it never forms \inlinecode{q} explicitly.
\end{compactitem}

\Needspace*{3\baselineskip}
\textbf{Space Complexity}:
\begin{compactitem}
  \item $O(m^{2} + m \cdot n)$ for the matrices returned by \inlinecode{qr\_decompose()}.
  \item $O(m \cdot n)$ auxiliary and $O(n)$ for the vector returned by \inlinecode{least\_squares()}.
\end{compactitem}

\begin{lstlisting}[style=cpp]
#include <cassert>
#include <cmath>
#include <optional>
#include <vector>

using Matrix = std::vector<std::vector<double>>;

// Reflects the trailing rows of the given column onto the axis, returning the reflection vector.
std::vector<double> householder_vector(const Matrix &a, int col) {
  int m = static_cast<int>(a.size());
  std::vector<double> v(m);
  double norm = 0;
  for (int i = col; i < m; i++) {
    v[i] = a[i][col];
    norm += v[i] * v[i];
  }
  norm = std::sqrt(norm);
  // Reflect away from a[col][col] so that the subtraction below never cancels.
  v[col] += a[col][col] >= 0 ? norm : -norm;
  return v;
}

void qr_decompose(const Matrix &a, Matrix *q, Matrix *r) {
  int m = static_cast<int>(a.size()), n = a.empty() ? 0 : static_cast<int>(a[0].size());
  assert(m >= n && q != nullptr && r != nullptr);
  *r = a;
  q->assign(m, std::vector<double>(m));
  for (int i = 0; i < m; i++) {
    (*q)[i][i] = 1;
  }
  for (int col = 0; col < n; col++) {
    std::vector<double> v = householder_vector(*r, col);
    double vv = 0;
    for (int i = col; i < m; i++) {
      vv += v[i] * v[i];
    }
    if (vv == 0) {
      continue;  // The column is already zero below the diagonal.
    }
    for (int j = 0; j < n; j++) {  // Apply the reflection to R from the left.
      double dot = 0;
      for (int i = col; i < m; i++) {
        dot += v[i] * (*r)[i][j];
      }
      for (int i = col; i < m; i++) {
        (*r)[i][j] -= 2 * dot / vv * v[i];
      }
    }
    for (int j = 0; j < m; j++) {  // Accumulate the reflection into Q from the right.
      double dot = 0;
      for (int i = col; i < m; i++) {
        dot += v[i] * (*q)[j][i];
      }
      for (int i = col; i < m; i++) {
        (*q)[j][i] -= 2 * dot / vv * v[i];
      }
    }
  }
}

std::optional<std::vector<double>> least_squares(
    const Matrix &a, const std::vector<double> &b, double eps = 1e-10
) {
  int m = static_cast<int>(a.size()), n = a.empty() ? 0 : static_cast<int>(a[0].size());
  assert(m >= n && static_cast<int>(b.size()) == m);
  Matrix r = a;
  std::vector<double> rhs = b;
  for (int col = 0; col < n; col++) {
    std::vector<double> v = householder_vector(r, col);
    double vv = 0;
    for (int i = col; i < m; i++) {
      vv += v[i] * v[i];
    }
    if (vv == 0) {
      continue;
    }
    for (int j = col; j < n; j++) {
      double dot = 0;
      for (int i = col; i < m; i++) {
        dot += v[i] * r[i][j];
      }
      for (int i = col; i < m; i++) {
        r[i][j] -= 2 * dot / vv * v[i];
      }
    }
    double dot = 0;  // The same reflection applied to the right-hand side.
    for (int i = col; i < m; i++) {
      dot += v[i] * rhs[i];
    }
    for (int i = col; i < m; i++) {
      rhs[i] -= 2 * dot / vv * v[i];
    }
  }
  std::vector<double> x(n);
  for (int i = n - 1; i >= 0; i--) {  // Back substitution over the triangular top block.
    if (std::fabs(r[i][i]) <= eps) {
      return std::nullopt;
    }
    double sum = rhs[i];
    for (int j = i + 1; j < n; j++) {
      sum -= r[i][j] * x[j];
    }
    x[i] = sum / r[i][i];
  }
  return x;
}
\end{lstlisting}
\Needspace*{4\baselineskip}
\textbf{Example Usage}\par
\begin{lstlisting}[style=cpp]
#include <cmath>
using namespace std;

bool EQ(double a, double b) {
  return fabs(a - b) < 1e-9;
}

int main() {
  Matrix a{{12, -51, 4}, {6, 167, -68}, {-4, 24, -41}};
  Matrix q, r;
  qr_decompose(a, &q, &r);
  for (int i = 0; i < 3; i++) {
    for (int j = 0; j < 3; j++) {
      double product = 0, orthogonal = 0;
      for (int k = 0; k < 3; k++) {
        product += q[i][k] * r[k][j];
        orthogonal += q[k][i] * q[k][j];
      }
      assert(EQ(product, a[i][j]));                // The factors multiply back to a.
      assert(EQ(orthogonal, i == j ? 1.0 : 0.0));  // The columns of q are orthonormal.
      assert(i <= j || EQ(r[i][j], 0.0));          // r is upper triangular.
    }
  }

  // An exactly solvable square system.
  auto exact = least_squares(Matrix{{2, 1}, {1, 3}}, vector<double>{5, 10});
  assert(exact.has_value() && EQ((*exact)[0], 1.0) && EQ((*exact)[1], 3.0));

  // Fit y = c0 + c1*x to four collinear points, then to points with one outlier.
  Matrix design{{1, 0}, {1, 1}, {1, 2}, {1, 3}};
  auto line = least_squares(design, vector<double>{1, 3, 5, 7});
  assert(line.has_value() && EQ((*line)[0], 1.0) && EQ((*line)[1], 2.0));
  auto noisy = least_squares(design, vector<double>{1, 3, 5, 8});
  assert(noisy.has_value() && EQ((*noisy)[0], 0.8) && EQ((*noisy)[1], 2.3));

  // Dependent columns have no unique least-squares solution.
  assert(!least_squares(Matrix{{1, 2}, {2, 4}, {3, 6}}, vector<double>{1, 2, 3}).has_value());
  return 0;
}
\end{lstlisting}
\subsection{Characteristic Polynomial}
\label{sec:6.5.6}
Computes the characteristic polynomial of a square matrix over a prime field: the coefficients of $\det(xI - A)$, a monic degree-$n$ polynomial whose roots are the eigenvalues of $A$. Its constant term recovers the determinant (up to sign), its second-highest coefficient is the negated trace, and evaluating it at a scalar gives $\det(xI - A)$ for that $x$. It is the standard tool behind eigenvalue counting, matrix-power recurrences via Cayley-Hamilton, and problems that reduce to a determinant in an unknown.

A direct expansion would cost a determinant per coefficient. Instead the matrix is first reduced to upper Hessenberg form, where every entry more than one place below the diagonal is zero. This reduction uses similarity transforms (a row operation paired with the inverse column operation), which leave the characteristic polynomial unchanged. The characteristic polynomial of a Hessenberg matrix then follows from a short recurrence over its leading principal submatrices, adding one row and column at a time, for an overall cost of $O(n^{3})$. The arithmetic is exact modulo a prime, since the Hessenberg reduction divides by pivots and is not numerically stable over the reals.

\begin{compactitem}
  \item \inlinecode{characteristic\_polynomial(a)} returns a vector $c$ of length $n + 1$ with $\det(xI - {\inlinecodemath{a}}) = c_0 + c_1 x + ... + c_n x^n$ modulo \inlinecode{MOD}, where $c_n = 1$. The matrix \inlinecode{a} must be square; its entries are reduced modulo \inlinecode{MOD} internally.
\end{compactitem}

\textbf{Time Complexity}: $O(n^{3})$ per call, where $n$ is the dimension of the matrix.

\textbf{Space Complexity}: $O(n^{2})$ auxiliary.

\begin{lstlisting}[style=cpp]
#include <cassert>
#include <cstdint>
#include <utility>
#include <vector>

const int64_t MOD = 998244353;

int64_t powmod(int64_t x, int64_t n, int64_t m) {
  int64_t res = 1 % m;
  for (x %= m; n > 0; n >>= 1) {
    if (n & 1) {
      res = res * x % m;
    }
    x = x * x % m;
  }
  return res;
}

void to_hessenberg(std::vector<std::vector<int64_t>> &m) {
  int n = static_cast<int>(m.size());
  for (int j = 0; j + 2 < n; j++) {
    int pivot = -1;
    for (int i = j + 1; i < n; i++) {
      if (m[i][j] != 0) {
        pivot = i;
        break;
      }
    }
    if (pivot == -1) {
      continue;
    }
    if (pivot != j + 1) {  // Move the pivot to the subdiagonal via a similarity swap.
      std::swap(m[pivot], m[j + 1]);
      for (int r = 0; r < n; r++) {
        std::swap(m[r][pivot], m[r][j + 1]);
      }
    }
    int64_t pinv = powmod(m[j + 1][j], MOD - 2, MOD);
    for (int k = j + 2; k < n; k++) {
      int64_t u = m[k][j] * pinv % MOD;
      if (u == 0) {
        continue;
      }
      for (int c = 0; c < n; c++) {  // Row k -= u * row (j + 1).
        m[k][c] = (m[k][c] - u * m[j + 1][c] % MOD + MOD) % MOD;
      }
      for (int r = 0; r < n; r++) {  // Inverse column operation keeps it a similarity.
        m[r][j + 1] = (m[r][j + 1] + u * m[r][k]) % MOD;
      }
    }
  }
}

std::vector<int64_t> characteristic_polynomial(std::vector<std::vector<int64_t>> a) {
  int n = static_cast<int>(a.size());
  assert(n == 0 || static_cast<int>(a[0].size()) == n);
  for (auto &row : a) {
    for (int64_t &x : row) {
      x = (x % MOD + MOD) % MOD;
    }
  }
  to_hessenberg(a);
  std::vector<std::vector<int64_t>> p(n + 1);
  p[0] = {1};
  for (int k = 0; k < n; k++) {
    p[k + 1].assign(k + 2, 0);
    for (int i = 0; i < static_cast<int>(p[k].size()); i++) {
      p[k + 1][i + 1] = (p[k + 1][i + 1] + p[k][i]) % MOD;                // x * p[k].
      p[k + 1][i] = (p[k + 1][i] - a[k][k] * p[k][i] % MOD + MOD) % MOD;  // -a[k][k] * p[k].
    }
    int64_t t = 1;
    for (int i = 0; i < k; i++) {
      t = t * a[k - i][k - i - 1] % MOD;  // Running product of subdiagonal entries.
      int64_t coeff = t * a[k - i - 1][k] % MOD;
      for (int j = 0; j < static_cast<int>(p[k - i - 1].size()); j++) {
        p[k + 1][j] = (p[k + 1][j] - coeff * p[k - i - 1][j] % MOD + MOD) % MOD;
      }
    }
  }
  return p[n];
}
\end{lstlisting}
\Needspace*{4\baselineskip}
\textbf{Example Usage}\par
\begin{lstlisting}[style=cpp]
#include <cassert>
using namespace std;

int main() {
  // [[1, 2], [3, 4]]: det(xI - A) = x^2 - 5x - 2.
  vector<vector<int64_t>> a{{1, 2}, {3, 4}};
  assert((characteristic_polynomial(a) == vector<int64_t>{MOD - 2, MOD - 5, 1}));  // -2, -5, 1.

  // det(xI - B) = (x - 3)(x^2 - 7x + 6) = x^3 - 10x^2 + 27x - 18.
  vector<vector<int64_t>> b{{2, 0, 1}, {0, 3, 0}, {4, 1, 5}};
  assert((characteristic_polynomial(b) == vector<int64_t>{MOD - 18, 27, MOD - 10, 1}));
  return 0;
}
\end{lstlisting}
\subsection{Eigenvalues (Jacobi)}
\label{sec:6.5.7}
An eigenvector of a square matrix is a direction the matrix does not rotate, only rescales, and the scale factor is its eigenvalue. For a real symmetric matrix, the spectral theorem guarantees that a full set of real eigenvalues exists and that eigenvectors for distinct eigenvalues are orthogonal, so the matrix is diagonalized by an orthogonal change of basis.

The Jacobi eigenvalue algorithm builds that basis one rotation at a time. It repeatedly picks the largest off-diagonal entry and applies a rotation in the plane of that entry's row and column, choosing the angle that sets the entry to zero. A rotation is a similarity transform, so it preserves the eigenvalues, and it also preserves the sum of squares of all entries; since zeroing an entry removes its square from the off-diagonal total and moves it onto the diagonal, every rotation strictly reduces the off-diagonal weight. Earlier zeros are partially refilled by later rotations, but the total decays geometrically, and the matrix converges to a diagonal one whose entries are the eigenvalues. Accumulating the rotations gives the eigenvectors.

\begin{compactitem}
  \item \inlinecode{jacobi\_eigen(a, eps = 1e-12, iterations = 100)} returns the pair (\inlinecode{values}, \inlinecode{vectors}) for a real symmetric $n$ by $n$ matrix \inlinecode{a}. The eigenvalues in \inlinecode{values} are sorted in decreasing order, and \inlinecode{vectors[i]} is a unit eigenvector for \inlinecode{values[i]}, so the eigenvectors are rows rather than columns. Each iteration rotates away the currently largest off-diagonal entry, stopping early once every such entry is smaller than \inlinecode{eps}; if the \inlinecode{iterations} limit is reached first, the current approximation is returned. \inlinecode{iterations} must be nonnegative. The matrix must be symmetric, which is not checked; an unsymmetric one silently produces nonsense.
\end{compactitem}

A general matrix instead needs the roots of the characteristic polynomial of section \secref{6.5.6}, found including complex ones by the Ehrlich-Aberth iteration of section \secref{5.4.5}, though accuracy degrades with degree since polynomial roots are far more sensitive to coefficients than eigenvalues are to the matrix. When only the dominant eigenvalue is wanted, power iteration finds it with no decomposition at all, by repeatedly multiplying a random vector and renormalizing.

\textbf{Time Complexity}: $O(I \cdot n^{2})$ per call for at most $I$ iterations, where $n$ is the matrix dimension. Each iteration searches $O(n^{2})$ entries for the largest off-diagonal one and applies one $O(n)$ rotation.

\textbf{Space Complexity}: $O(n^{2})$ auxiliary, and $O(n^{2})$ for the returned eigenvectors.

\begin{lstlisting}[style=cpp]
#include <algorithm>
#include <cassert>
#include <cmath>
#include <numeric>
#include <utility>
#include <vector>

using Matrix = std::vector<std::vector<double>>;

std::pair<std::vector<double>, Matrix> jacobi_eigen(
    Matrix a, double eps = 1e-12, int iterations = 100
) {
  int n = static_cast<int>(a.size());
  assert((n == 0 || static_cast<int>(a[0].size()) == n) && iterations >= 0);
  Matrix vectors(n, std::vector<double>(n));
  for (int i = 0; i < n; i++) {
    vectors[i][i] = 1;
  }
  while (iterations--) {
    int p = 0, q = 1;
    double largest = 0;
    for (int i = 0; i < n; i++) {
      for (int j = i + 1; j < n; j++) {
        if (std::fabs(a[i][j]) > largest) {
          largest = std::fabs(a[i][j]);
          p = i;
          q = j;
        }
      }
    }
    if (n < 2 || largest <= eps) {
      break;
    }
    // Choose the rotation angle that zeroes a[p][q], using the numerically safe form of
    // tan(theta) rather than computing the angle itself.
    double theta = (a[q][q] - a[p][p]) / (2 * a[p][q]);
    double t = (theta >= 0 ? 1 : -1) / (std::fabs(theta) + std::sqrt(theta * theta + 1));
    double c = 1 / std::sqrt(t * t + 1), s = t * c;
    for (int k = 0; k < n; k++) {
      double akp = a[k][p], akq = a[k][q];
      a[k][p] = c * akp - s * akq;
      a[k][q] = s * akp + c * akq;
    }
    for (int k = 0; k < n; k++) {
      double apk = a[p][k], aqk = a[q][k];
      a[p][k] = c * apk - s * aqk;
      a[q][k] = s * apk + c * aqk;
      double vpk = vectors[p][k], vqk = vectors[q][k];
      vectors[p][k] = c * vpk - s * vqk;
      vectors[q][k] = s * vpk + c * vqk;
    }
  }
  std::vector<int> order(n);
  std::iota(order.begin(), order.end(), 0);
  std::sort(order.begin(), order.end(), [&](int x, int y) { return a[x][x] > a[y][y]; });
  std::vector<double> values(n);
  Matrix sorted(n);
  for (int i = 0; i < n; i++) {
    values[i] = a[order[i]][order[i]];
    sorted[i] = vectors[order[i]];
  }
  return {values, sorted};
}
\end{lstlisting}
\Needspace*{4\baselineskip}
\textbf{Example Usage}\par
\begin{lstlisting}[style=cpp]
#include <cassert>
#include <cmath>
using namespace std;

bool EQ(double a, double b) {
  return fabs(a - b) < 1e-9;
}

int main() {
  // A symmetric matrix with eigenvalues 3, 2, and 1.
  Matrix a{{2, 0, -1}, {0, 2, 0}, {-1, 0, 2}};
  auto [values, vectors] = jacobi_eigen(a);
  assert(EQ(values[0], 3) && EQ(values[1], 2) && EQ(values[2], 1));

  for (int i = 0; i < 3; i++) {
    double norm = 0;
    for (int k = 0; k < 3; k++) {
      norm += vectors[i][k] * vectors[i][k];
    }
    assert(EQ(norm, 1));  // Eigenvectors come back normalized.
    for (int r = 0; r < 3; r++) {
      double image = 0;
      for (int k = 0; k < 3; k++) {
        image += a[r][k] * vectors[i][k];
      }
      assert(EQ(image, values[i] * vectors[i][r]));  // The defining equation a*v = lambda*v.
    }
  }

  // A diagonal matrix is already solved, and a 1-by-1 matrix is its own eigenvalue.
  auto diagonal = jacobi_eigen(Matrix{{5, 0}, {0, -4}});
  assert(EQ(diagonal.first[0], 5) && EQ(diagonal.first[1], -4));
  assert(EQ(jacobi_eigen(Matrix{{7}}).first[0], 7));
  return 0;
}
\end{lstlisting}
\subsection{Sparse Matrix}
\label{sec:6.5.8}
Provides a sparse matrix type together with arithmetic and elimination routines. A sparse matrix stores only nonzero entries, which is useful when the matrix is large but most values are zero.

The matrix keeps both row maps and column maps in sync: \inlinecode{row(i)} can iterate all nonzeros in row \inlinecode{i}, while \inlinecode{col(j)} can iterate all nonzeros in column \inlinecode{j}. That bidirectional storage is especially useful for sparse elimination or graph-like matrix operations, where both row updates and column pivot lookups are common.

The class treats \inlinecode{T\{\}} as additive zero and requires \inlinecode{value == T\{\}} to be a valid zero test.

\begin{compactitem}
  \item \inlinecode{SparseMatrix\textless{}T\textgreater{}(rows, cols)} constructs a \inlinecode{rows} by \inlinecode{cols} sparse matrix.
  \item \inlinecode{num\_rows()} returns the number of rows.
  \item \inlinecode{num\_cols()} returns the number of columns.
  \item \inlinecode{nonzeros()} returns the number of stored nonzero entries.
  \item \inlinecode{get(i, j)} returns the value at row \inlinecode{i}, column \inlinecode{j}, or $0$ if the entry is not stored.
  \item \inlinecode{set(i, j, value)} assigns an entry. Assigning zero erases it from both maps.
  \item \inlinecode{add(i, j, delta)} adds \inlinecode{delta} to an entry. If the result becomes zero, the entry is erased.
  \item \inlinecode{row(i)} returns a map from column index to value for the nonzero entries in row \inlinecode{i}.
  \item \inlinecode{col(j)} returns a map from row index to value for the nonzero entries in column \inlinecode{j}.
  \item \inlinecode{swap\_rows(i, j)} swaps two rows while keeping the column maps synchronized.
  \item \inlinecode{transpose()} transposes the matrix in place.
  \item \inlinecode{a * x} returns the matrix-vector product with vector \inlinecode{x}.
  \item \inlinecode{a + b} and \inlinecode{a - b} return the sum and difference of two matrices with identical dimensions.
  \item \inlinecode{a * b} returns the matrix product, requiring \inlinecode{a.num\_cols()} to equal \inlinecode{b.num\_rows()}. Only pairs of stored entries are ever multiplied.
  \item \inlinecode{a * k} returns \inlinecode{a} scaled by the value \inlinecode{k}. Scaling by zero yields an empty matrix.
  \item Operators \inlinecode{+=}, \inlinecode{-=}, and \inlinecode{*=} assign the corresponding result back into the left operand. In all of these operators, entries that cancel to zero are not stored.
\end{compactitem}

\Needspace*{3\baselineskip}
\textbf{Time Complexity}:
\begin{compactitem}
  \item $O(\log d)$ per call to \inlinecode{get()}, \inlinecode{set()}, and \inlinecode{add()}, where $d$ is the number of nonzeros in the touched row or column.
  \item $O(1)$ per call to \inlinecode{transpose()} and $O(z)$ per matrix-vector product, where $z$ is the number of stored nonzero entries.
  \item $O((r_i + r_j) \cdot \log d)$ per call to \inlinecode{swap\_rows(i, j)}, where $r_i$ and $r_j$ are the sizes of the two rows.
  \item $O((z_a + z_b) \cdot \log d)$ per call to \inlinecode{operator+} and \inlinecode{operator-}, where $z_a$ and $z_b$ are the operand nonzero counts.
  \item $O(z \cdot \log d)$ per call to scalar \inlinecode{operator*}, where $z$ is the number of stored nonzero entries.
  \item $O(f \cdot \log d)$ per call to \inlinecode{operator*}, where $f$ is the number of scalar products accumulated, that is, the sum of \inlinecode{b.row(k).size()} over every stored entry \inlinecode{a[i][k]}.
  \item $O(1)$ per call to \inlinecode{num\_rows()}, \inlinecode{num\_cols()}, \inlinecode{nonzeros()}, \inlinecode{row()}, and \inlinecode{col()}.
\end{compactitem}

\Needspace*{3\baselineskip}
\textbf{Space Complexity}:
\begin{compactitem}
  \item $O(z)$ storage.
  \item $O(1)$ auxiliary for \inlinecode{get()}, \inlinecode{set()}, \inlinecode{add()}, \inlinecode{row()}, \inlinecode{col()}, and \inlinecode{transpose()}.
  \item $O(r_i + r_j)$ auxiliary for \inlinecode{swap\_rows(i, j)}.
  \item $O(m)$ for the vector returned by matrix-vector multiplication, where $m$ is the number of rows.
  \item $O(c)$ auxiliary for \inlinecode{operator*}, where $c$ is the largest number of nonzeros in one result row.
\end{compactitem}

\begin{lstlisting}[style=cpp]
#include <cassert>
#include <map>
#include <unordered_set>
#include <utility>
#include <vector>

template<typename T>
class SparseMatrix {
  int rows = 0, cols = 0, stored = 0;
  std::vector<std::map<int, T>> rvals, cvals;

  bool is_zero(const T &value) const { return value == T{}; }

 public:
  SparseMatrix(int rows, int cols) : rows(rows), cols(cols) {
    assert(rows >= 0 && cols >= 0);
    rvals.resize(rows);
    cvals.resize(cols);
  }

  int num_rows() const { return rows; }
  int num_cols() const { return cols; }
  int nonzeros() const { return stored; }

  T get(int i, int j) const {
    auto it = rvals[i].find(j);
    return it == rvals[i].end() ? T{} : it->second;
  }

  void set(int i, int j, const T &value) {
    auto it = rvals[i].find(j);
    if (is_zero(value)) {
      if (it != rvals[i].end()) {
        rvals[i].erase(it);
        cvals[j].erase(i);
        stored--;
      }
      return;
    }
    if (it == rvals[i].end()) {
      stored++;
    }
    rvals[i][j] = value;
    cvals[j][i] = value;
  }

  void add(int i, int j, const T &delta) {
    if (is_zero(delta)) {
      return;
    }
    set(i, j, get(i, j) + delta);
  }

  const std::map<int, T> &row(int i) const { return rvals[i]; }
  const std::map<int, T> &col(int j) const { return cvals[j]; }

  void swap_rows(int i, int j) {
    if (i == j) {
      return;
    }
    std::unordered_set<int> touched;
    for (const auto &[col, value] : rvals[i]) {
      touched.insert(col);
    }
    for (const auto &[col, value] : rvals[j]) {
      touched.insert(col);
    }
    for (int col : touched) {
      T x = get(i, col), y = get(j, col);
      set(i, col, y);
      set(j, col, x);
    }
  }

  void transpose() {
    std::swap(rows, cols);
    std::swap(rvals, cvals);
  }

  std::vector<T> operator*(const std::vector<T> &x) const {
    assert(static_cast<int>(x.size()) == cols);
    std::vector<T> res(rows);
    for (int i = 0; i < rows; i++) {
      for (const auto &[j, value] : rvals[i]) {
        res[i] += value * x[j];
      }
    }
    return res;
  }

  SparseMatrix operator+(const SparseMatrix &b) const {
    assert(rows == b.rows && cols == b.cols);
    SparseMatrix res(rows, cols);
    for (int i = 0; i < rows; i++) {
      for (const auto &[j, value] : rvals[i]) {
        res.set(i, j, value);
      }
      for (const auto &[j, value] : b.rvals[i]) {
        res.add(i, j, value);  // Cancellation to zero erases the entry.
      }
    }
    return res;
  }

  SparseMatrix operator-(const SparseMatrix &b) const {
    assert(rows == b.rows && cols == b.cols);
    SparseMatrix res(rows, cols);
    for (int i = 0; i < rows; i++) {
      for (const auto &[j, value] : rvals[i]) {
        res.set(i, j, value);
      }
      for (const auto &[j, value] : b.rvals[i]) {
        res.add(i, j, T{} - value);
      }
    }
    return res;
  }

  // Scaling by zero cancels every entry, so the result stores nothing at all.
  SparseMatrix operator*(const T &k) const {
    SparseMatrix res(rows, cols);
    if (is_zero(k)) {
      return res;
    }
    for (int i = 0; i < rows; i++) {
      for (const auto &[j, value] : rvals[i]) {
        res.set(i, j, value * k);
      }
    }
    return res;
  }

  // Row-by-row product: each stored a[i][k] scatters row k of b into an accumulator for row i, so
  // only nonzero-by-nonzero pairs are ever multiplied.
  SparseMatrix operator*(const SparseMatrix &b) const {
    assert(cols == b.rows);
    SparseMatrix res(rows, b.cols);
    for (int i = 0; i < rows; i++) {
      std::map<int, T> acc;
      for (const auto &[k, aik] : rvals[i]) {
        for (const auto &[j, bkj] : b.rvals[k]) {
          acc[j] = acc[j] + aik * bkj;
        }
      }
      for (const auto &[j, value] : acc) {
        res.set(i, j, value);
      }
    }
    return res;
  }

  SparseMatrix &operator+=(const SparseMatrix &b) { return *this = *this + b; }
  SparseMatrix &operator-=(const SparseMatrix &b) { return *this = *this - b; }
  SparseMatrix &operator*=(const SparseMatrix &b) { return *this = *this * b; }
  SparseMatrix &operator*=(const T &k) { return *this = *this * k; }
};
\end{lstlisting}
The elimination routines require field-like arithmetic because they divide by pivots, so use types such as \inlinecode{double}, rational numbers, or modular integers. For floating-point matrices, replace \inlinecode{is\_zero()} with an epsilon comparison if small roundoff values should be erased. At each pivot column, sparse elimination chooses the available row with the fewest stored entries to reduce fill-in, although an unfriendly matrix can still become dense.

\begin{compactitem}
  \item \inlinecode{row\_reduce(a, limit)} converts columns $[0, {\inlinecodemath{limit}})$ of \inlinecode{a} to sparse row echelon form, returning the rank found in those columns.
  \item \inlinecode{det(a)} returns the determinant of a square sparse matrix.
  \item \inlinecode{sparse\_rank(a)} returns the rank of a sparse matrix.
  \item \inlinecode{solve\_system(a, b, \&x)} solves the system $ax = b$, returning $0$ for one solution, $-1$ for no solution, or $-2$ for infinitely many solutions. When one solution exists, \inlinecode{x} is populated.
\end{compactitem}

\Needspace*{3\baselineskip}
\textbf{Time Complexity}:
\begin{compactitem}
  \item $O(f \cdot \log d)$ for sparse elimination, where $f$ is the number of entry updates performed after fill-in and $d$ is a touched row or column size. In the worst case this is still cubic.
  \item $O(f \cdot \log d)$ per call to \inlinecode{det()}, \inlinecode{sparse\_rank()}, and \inlinecode{solve\_system()}.
\end{compactitem}

\textbf{Space Complexity}: $O(z + m)$ auxiliary for elimination, in addition to the copied matrix used by \inlinecode{det()} and \inlinecode{sparse\_rank()} or the augmented matrix used by \inlinecode{solve\_system()}.

\begin{lstlisting}[style=cpp]
template<typename T>
int row_reduce(SparseMatrix<T> &a, int limit) {
  assert(0 <= limit && limit <= a.num_cols());
  int r = 0;
  for (int c = 0; c < limit && r < a.num_rows(); c++) {
    int pivot = -1;
    for (const auto &[i, value] : a.col(c)) {
      if (i >= r && (pivot == -1 || a.row(i).size() < a.row(pivot).size())) {
        pivot = i;
      }
    }
    if (pivot == -1) {
      continue;
    }
    a.swap_rows(r, pivot);
    T inv_pivot = T{1} / a.get(r, c);
    std::vector<std::pair<int, T>> pivot_row(a.row(r).begin(), a.row(r).end());
    std::vector<int> touched_rows;
    for (const auto &[i, value] : a.col(c)) {
      if (i > r) {
        touched_rows.push_back(i);
      }
    }
    for (int i : touched_rows) {
      T coeff = -a.get(i, c) * inv_pivot;
      for (const auto &[j, value] : pivot_row) {
        if (j != c) {
          a.add(i, j, coeff * value);
        }
      }
      a.set(i, c, T{});
    }
    r++;
  }
  return r;
}

template<typename T>
T det(SparseMatrix<T> a) {
  assert(a.num_rows() == a.num_cols());
  int swaps = 0, r = 0, n = a.num_rows();
  T det = 1;
  for (int c = 0; c < n && r < n; c++) {
    int pivot = -1;
    for (const auto &[i, value] : a.col(c)) {
      if (i >= r && (pivot == -1 || a.row(i).size() < a.row(pivot).size())) {
        pivot = i;
      }
    }
    if (pivot == -1) {
      return T{};
    }
    if (pivot != r) {
      a.swap_rows(r, pivot);
      swaps++;
    }
    T pivot_value = a.get(r, c);
    det *= pivot_value;
    T inv_pivot = T{1} / pivot_value;
    std::vector<std::pair<int, T>> pivot_row(a.row(r).begin(), a.row(r).end());
    std::vector<int> touched_rows;
    for (const auto &[i, value] : a.col(c)) {
      if (i > r) {
        touched_rows.push_back(i);
      }
    }
    for (int i : touched_rows) {
      T coeff = -a.get(i, c) * inv_pivot;
      for (const auto &[j, value] : pivot_row) {
        if (j != c) {
          a.add(i, j, coeff * value);
        }
      }
      a.set(i, c, T{});
    }
    r++;
  }
  return (r == n ? (swaps % 2 == 0 ? det : -det) : T{});
}

template<typename T>
int sparse_rank(SparseMatrix<T> a) {
  return row_reduce(a, a.num_cols());
}

template<typename T>
int solve_system(const SparseMatrix<T> &a, const std::vector<T> &b, std::vector<T> *x) {
  if (x == nullptr || a.num_rows() != static_cast<int>(b.size())) {
    return -1;
  }
  int rows = a.num_rows(), cols = a.num_cols();
  SparseMatrix<T> aug(rows, cols + 1);
  for (int i = 0; i < rows; i++) {
    for (const auto &[j, value] : a.row(i)) {
      aug.set(i, j, value);
    }
    aug.set(i, cols, b[i]);
  }
  row_reduce(aug, cols);
  std::vector<int> pivot_col;
  for (int i = 0; i < rows; i++) {
    int lead = -1;
    for (const auto &[j, value] : aug.row(i)) {
      if (j < cols) {
        lead = j;
        break;
      }
    }
    if (lead == -1) {
      if (aug.get(i, cols) != T{}) {
        return -1;
      }
    } else {
      pivot_col.push_back(lead);
    }
  }
  if (static_cast<int>(pivot_col.size()) < cols) {
    return -2;
  }
  x->assign(cols, T{});
  for (int row = static_cast<int>(pivot_col.size()) - 1; row >= 0; row--) {
    int c = pivot_col[row];
    T rhs = aug.get(row, cols);
    for (const auto &[j, value] : aug.row(row)) {
      if (j > c && j < cols) {
        rhs -= value * (*x)[j];
      }
    }
    (*x)[c] = rhs / aug.get(row, c);
  }
  return 0;
}
\end{lstlisting}
\Needspace*{4\baselineskip}
\textbf{Example Usage}\par
\begin{lstlisting}[style=cpp]
#include <cassert>
#include <cstdint>
using namespace std;

int main() {
  SparseMatrix<int64_t> a(3, 4);
  a.set(0, 1, 5);
  a.set(2, 3, 7);
  a.add(0, 1, -2);
  a.add(1, 2, 4);

  assert(a.num_rows() == 3 && a.num_cols() == 4);
  assert(a.nonzeros() == 3);
  assert(a.get(0, 1) == 3);
  assert(a.get(0, 0) == 0);
  assert(a.row(0).begin()->first == 1);
  assert(a.col(2).begin()->first == 1);

  vector<int64_t> x{10, 20, 30, 40};
  assert((a * x == vector<int64_t>{60, 120, 280}));

  a.add(0, 1, -3);  // Entries that become zero are removed from both row and column maps.
  assert(a.get(0, 1) == 0);
  assert(a.nonzeros() == 2);
  assert(a.row(0).empty());
  assert(a.col(1).empty());

  a.transpose();
  assert(a.num_rows() == 4 && a.num_cols() == 3);
  assert(a.get(3, 2) == 7);
  assert(a.get(2, 1) == 4);

  //  p = [1 2]   q = [0 3]   p + q = [1 5]   p * q = [4 3]
  //      [0 1]       [2 0]           [2 1]           [2 0]
  SparseMatrix<int64_t> p(2, 2), q(2, 2);
  p.set(0, 0, 1);
  p.set(0, 1, 2);
  p.set(1, 1, 1);
  q.set(0, 1, 3);
  q.set(1, 0, 2);
  SparseMatrix<int64_t> sum = p + q;
  assert(sum.get(0, 0) == 1 && sum.get(0, 1) == 5);
  assert(sum.get(1, 0) == 2 && sum.get(1, 1) == 1);
  SparseMatrix<int64_t> product = p * q;
  assert(product.get(0, 0) == 4 && product.get(0, 1) == 3);
  assert(product.get(1, 0) == 2 && product.get(1, 1) == 0);
  assert(product.nonzeros() == 3);  // The zero entry is never stored.

  // Cancelling entries are dropped rather than stored as explicit zeros.
  SparseMatrix<int64_t> negated(2, 2);
  negated.set(0, 0, -1);
  negated.set(0, 1, -2);
  negated.set(1, 1, -1);
  assert((p + negated).nonzeros() == 0);
  assert((p - p).nonzeros() == 0);

  SparseMatrix<int64_t> difference = p - q;
  assert(difference.get(0, 0) == 1 && difference.get(0, 1) == -1);
  assert(difference.get(1, 0) == -2 && difference.get(1, 1) == 1);

  SparseMatrix<int64_t> scaled = p * 3;
  assert(scaled.get(0, 0) == 3 && scaled.get(0, 1) == 6 && scaled.get(1, 1) == 3);
  assert((p * 0LL).nonzeros() == 0);

  SparseMatrix<int64_t> acc = p;
  acc += q;
  assert(acc.get(0, 1) == 5 && acc.get(1, 0) == 2);
  acc -= q;
  acc *= 2;
  assert(acc.get(0, 0) == 2 && acc.get(0, 1) == 4 && acc.get(1, 1) == 2);

  SparseMatrix<double> d(3, 3);
  d.set(0, 0, 2);
  d.set(0, 1, 1);
  d.set(1, 1, 3);
  d.set(2, 0, 1);
  d.set(2, 2, 4);
  double determinant = det(d);
  assert(determinant > 24 - 1e-9 && determinant < 24 + 1e-9);
  assert(sparse_rank(d) == 3);

  SparseMatrix<double> sys(3, 3);
  sys.set(0, 0, -1);
  sys.set(0, 1, 2);
  sys.set(0, 2, 5);
  sys.set(1, 0, 1);
  sys.set(1, 2, -6);
  sys.set(2, 0, -4);
  sys.set(2, 1, 2);
  sys.set(2, 2, 2);
  vector<double> rhs{3, 1, -2}, sol;
  assert(solve_system(sys, rhs, &sol) == 0);
  for (int i = 0; i < sys.num_rows(); i++) {
    double sum = 0;
    for (const auto &[j, value] : sys.row(i)) {
      sum += value * sol[j];
    }
    assert(sum > rhs[i] - 1e-9 && sum < rhs[i] + 1e-9);
  }
  return 0;
}
\end{lstlisting}
\subsection{Matrix-Tree Theorem}
\label{sec:6.5.9}
Kirchhoff's matrix-tree theorem counts the spanning trees of an undirected graph. Form the Laplacian matrix $L = D - A$, where $D$ is the diagonal matrix of degrees and $A$ is the adjacency matrix (with multiplicities for parallel edges). Every cofactor of $L$ is equal, and that common value is the number of spanning trees. Equivalently, delete any one row and the matching column from $L$ and take the determinant of what remains.

The determinant is evaluated modulo a prime by Gaussian elimination, so the count is returned modulo \inlinecode{MOD}. Spanning-tree counts grow astronomically, so a modular count is the usual contest form; run the routine under a few different primes and combine with the Chinese remainder theorem to recover a large exact value. The same construction handles weighted graphs (sum edge weights into the Laplacian to get the weighted tree count) and, with the directed Laplacian and a fixed deleted row, counts rooted spanning arborescences.

\begin{compactitem}
  \item \inlinecode{count\_spanning\_trees(n, edges)} returns the number of spanning trees of the undirected graph on nodes $[0, {\inlinecodemath{n}})$ with the given edge list, modulo \inlinecode{MOD}. Parallel edges are honored; self-loops do not affect the count. The result is $0$ when the graph is disconnected, and $1$ when \inlinecode{n} is $1$.
\end{compactitem}

\textbf{Time Complexity}: $O(n^{3} + m)$ per call, where $n$ is the number of nodes and $m$ is the number of edges.

\textbf{Space Complexity}: $O(n^{2})$ auxiliary.

\begin{lstlisting}[style=cpp]
#include <algorithm>
#include <cassert>
#include <cstdint>
#include <utility>
#include <vector>

const int64_t MOD = 998244353;

int64_t powmod(int64_t x, int64_t n, int64_t m) {
  int64_t res = 1 % m;
  for (x %= m; n > 0; n >>= 1) {
    if (n & 1) {
      res = res * x % m;
    }
    x = x * x % m;
  }
  return res;
}

int64_t count_spanning_trees(int n, const std::vector<std::pair<int, int>> &edges) {
  assert(n >= 0);
  assert(std::all_of(edges.begin(), edges.end(), [n](auto edge) {
    auto [u, v] = edge;
    return 0 <= u && u < n && 0 <= v && v < n;
  }));
  if (n <= 1) {
    return n == 1 ? 1 : 0;
  }
  // Build the Laplacian, then drop the last row and column to form an (n-1) by (n-1) cofactor.
  int m = n - 1;
  std::vector<std::vector<int64_t>> lap(m, std::vector<int64_t>(m));
  for (const auto &e : edges) {
    int u = e.first, v = e.second;
    if (u == v) {
      continue;  // Self-loops contribute nothing to the Laplacian.
    }
    if (u < m) {
      lap[u][u] = (lap[u][u] + 1) % MOD;
    }
    if (v < m) {
      lap[v][v] = (lap[v][v] + 1) % MOD;
    }
    if (u < m && v < m) {
      lap[u][v] = (lap[u][v] - 1 + MOD) % MOD;
      lap[v][u] = (lap[v][u] - 1 + MOD) % MOD;
    }
  }
  // Determinant modulo MOD by Gaussian elimination.
  int64_t det = 1;
  for (int col = 0; col < m; col++) {
    int pivot = -1;
    for (int r = col; r < m; r++) {
      if (lap[r][col] != 0) {
        pivot = r;
        break;
      }
    }
    if (pivot == -1) {
      return 0;  // Singular cofactor: the graph is disconnected.
    }
    if (pivot != col) {
      std::swap(lap[pivot], lap[col]);
      det = (MOD - det) % MOD;
    }
    det = det * lap[col][col] % MOD;
    int64_t pinv = powmod(lap[col][col], MOD - 2, MOD);
    for (int r = col + 1; r < m; r++) {
      int64_t factor = lap[r][col] * pinv % MOD;
      if (factor == 0) {
        continue;
      }
      for (int k = col; k < m; k++) {
        lap[r][k] = (lap[r][k] - factor * lap[col][k]) % MOD;
        if (lap[r][k] < 0) {
          lap[r][k] += MOD;
        }
      }
    }
  }
  return det;
}
\end{lstlisting}
\Needspace*{4\baselineskip}
\textbf{Example Usage}\par
\begin{lstlisting}[style=cpp]
#include <cassert>
using namespace std;

int main() {
  // A path 0-1-2 and a single triangle each have a known number of spanning trees.
  assert(count_spanning_trees(3, {{0, 1}, {1, 2}}) == 1);          // A tree has exactly one.
  assert(count_spanning_trees(3, {{0, 1}, {1, 2}, {2, 0}}) == 3);  // Triangle: drop any one edge.

  // Complete graphs: Cayley's formula gives n^(n-2) spanning trees.
  auto complete = [](int n) {
    vector<pair<int, int>> e;
    for (int i = 0; i < n; i++) {
      for (int j = i + 1; j < n; j++) {
        e.emplace_back(i, j);
      }
    }
    return count_spanning_trees(n, e);
  };
  assert(complete(4) == 16);   // 4^2.
  assert(complete(5) == 125);  // 5^3.

  // Disconnected graphs have no spanning tree.
  assert(count_spanning_trees(4, {{0, 1}, {2, 3}}) == 0);
  return 0;
}
\end{lstlisting}

\section{\texorpdfstring{Polynomials and Linear Recurrences}{6.6 Polynomials and Linear Recurrences}}
\setcounter{section}{6}
\setcounter{subsection}{0}
\subsection{Number Theoretic Transform}
\label{sec:6.6.1}
The number theoretic transform (NTT) is the discrete Fourier transform carried out in modular arithmetic instead of over the complex numbers. It multiplies two polynomials (equivalently, computes the convolution of two sequences) modulo a prime in $O(n \log n)$ time, with no floating point error. The modulus must be of the form $c \cdot 2^k + 1$ so that the ring of integers modulo it contains a $2^k$-th root of unity, which limits transform sizes to powers of two up to $2^k$. The code below uses the common prime $998244353 = 119 \cdot 2^{23} + 1$ with primitive root $3$, supporting transforms of any power-of-two length up to $2^{23}$.

The transform replaces a coefficient vector with its evaluations at the powers of a root of unity; pointwise multiplying two such evaluation vectors and transforming back gives their convolution. This is the same divide-and-conquer butterfly as the fast Fourier transform, with the root of unity and all arithmetic taken modulo the prime.

\begin{compactitem}
  \item \inlinecode{ntt(a, invert = false)} transforms the vector \inlinecode{a} in place, whose length must be a power of two. The forward transform uses \inlinecode{invert = false}; the inverse uses \inlinecode{invert = true}.
  \item \inlinecode{convolve(a, b)} returns the convolution of \inlinecode{a} and \inlinecode{b} modulo the prime, that is, the coefficients of the product of the two polynomials whose coefficients are the inputs. The result has length ${\inlinecodemath{a.size()}} + {\inlinecodemath{b.size()}} - 1$, or is empty if either input is empty. Input values are assumed to lie in $[0, {\inlinecodemath{MOD}})$. Small inputs use direct multiplication to avoid NTT overhead.
\end{compactitem}

\Needspace*{3\baselineskip}
\textbf{Time Complexity}:
\begin{compactitem}
  \item $O(n \log n)$ per call to \inlinecode{ntt()}, where $n$ is the length of the vector.
  \item $O(|a||b|)$ per call to \inlinecode{convolve()} on small inputs and $O(n \log n)$ otherwise, where $n$ is the padded transform length.
\end{compactitem}

\Needspace*{3\baselineskip}
\textbf{Space Complexity}:
\begin{compactitem}
  \item $O(1)$ auxiliary for \inlinecode{ntt()}, transforming in place.
  \item $O(n)$ auxiliary for \inlinecode{convolve()}.
\end{compactitem}

\begin{lstlisting}[style=cpp]
#include <algorithm>
#include <cassert>
#include <cstdint>
#include <utility>
#include <vector>

const int64_t MOD = 998244353;
const int64_t ROOT = 3;
const int MAX_POWER_OF_TWO = 23;
const int NAIVE_CUTOFF = 150;

int64_t powmod(int64_t x, int64_t n, int64_t m) {
  int64_t res = 1 % m;
  for (x %= m; n > 0; n >>= 1) {
    if (n & 1) {
      res = res * x % m;
    }
    x = x * x % m;
  }
  return res;
}

void ntt(std::vector<int64_t> &a, bool invert = false) {
  int n = static_cast<int>(a.size());
  assert(n > 0 && (n & (n - 1)) == 0 && n <= (1 << MAX_POWER_OF_TWO));
  for (int i = 1, j = 0; i < n; i++) {
    int bit = n >> 1;
    for (; j & bit; bit >>= 1) {
      j ^= bit;
    }
    j ^= bit;
    if (i < j) {
      std::swap(a[i], a[j]);
    }
  }
  for (int len = 2; len <= n; len <<= 1) {
    int64_t root = powmod(ROOT, (MOD - 1) / len, MOD);
    if (invert) {
      root = powmod(root, MOD - 2, MOD);
    }
    for (int i = 0; i < n; i += len) {
      int64_t w = 1;
      for (int k = 0; k < len / 2; k++) {
        int64_t u = a[i + k], v = a[i + k + len / 2] * w % MOD;
        a[i + k] = (u + v) % MOD;
        a[i + k + len / 2] = (u - v + MOD) % MOD;
        w = w * root % MOD;
      }
    }
  }
  if (invert) {
    int64_t n_inv = powmod(n, MOD - 2, MOD);
    for (int64_t &x : a) {
      x = x * n_inv % MOD;
    }
  }
}

std::vector<int64_t> convolve(std::vector<int64_t> a, std::vector<int64_t> b) {
  if (a.empty() || b.empty()) {
    return {};
  }
  int result_size = static_cast<int>(a.size() + b.size()) - 1;
  if (std::min(a.size(), b.size()) < NAIVE_CUTOFF) {
    std::vector<int64_t> res(result_size);
    for (int i = 0; i < static_cast<int>(a.size()); i++) {
      for (int j = 0; j < static_cast<int>(b.size()); j++) {
        res[i + j] = (res[i + j] + a[i] * b[j]) % MOD;
      }
    }
    return res;
  }
  int n = 1;
  while (n < result_size) {
    n <<= 1;
  }
  assert(n <= (1 << MAX_POWER_OF_TWO));
  a.resize(n);
  b.resize(n);
  ntt(a);
  ntt(b);
  for (int i = 0; i < n; i++) {
    a[i] = a[i] * b[i] % MOD;
  }
  ntt(a, true);
  a.resize(result_size);
  return a;
}
\end{lstlisting}
\Needspace*{4\baselineskip}
\textbf{Example Usage}\par
\begin{lstlisting}[style=cpp]
#include <cassert>
using namespace std;

int main() {
  // (1 + 2x + 3x^2)(4 + 5x + 6x^2) = 4 + 13x + 28x^2 + 27x^3 + 18x^4.
  vector<int64_t> a{1, 2, 3}, b{4, 5, 6};
  assert((convolve(a, b) == vector<int64_t>{4, 13, 28, 27, 18}));
  assert(convolve({}, b).empty());
  assert((convolve(vector<int64_t>{5}, vector<int64_t>{7}) == vector<int64_t>{35}));
  vector<int64_t> ones(160, 1);
  vector<int64_t> d = convolve(ones, ones);
  assert(d[0] == 1 && d[159] == 160 && d[318] == 1);
  return 0;
}
\end{lstlisting}
\subsection{Fast Fourier Transform}
\label{sec:6.6.2}
The fast Fourier transform (FFT) multiplies two polynomials, or equivalently convolves two sequences, in $O(n \log n)$ time by evaluating them at the complex roots of unity. Like the number theoretic transform of the previous section, it replaces a coefficient vector with its evaluations at the powers of a root of unity, multiplies two such evaluation vectors pointwise, and transforms back. The difference is the arithmetic: the FFT works over the complex numbers with $e^{2\pi i / n}$ as the root of unity, so it convolves arbitrary real coefficients rather than residues modulo a prime.

The trade-off is precision. Complex arithmetic in double precision accumulates rounding error, so the result is recovered by rounding each output to the nearest integer; this is reliable as long as the true coefficients stay well within the roughly 15 significant decimal digits a \inlinecode{double} can hold. When exact results modulo a prime are required, or coefficients are large, prefer the number theoretic transform; use the FFT for real-valued convolution or big-integer multiplication.

\begin{compactitem}
  \item \inlinecode{fft(a, invert = false)} transforms the complex vector \inlinecode{a} in place, whose length must be a power of two. The forward transform uses \inlinecode{invert = false}; the inverse uses \inlinecode{invert = true}.
  \item \inlinecode{convolve(a, b)} returns the convolution of two integer sequences \inlinecode{a} and \inlinecode{b}, that is, the coefficients of the product of the two polynomials whose coefficients are the inputs. The result has length ${\inlinecodemath{a.size()}} + {\inlinecodemath{b.size()}} - 1$, or is empty if either input is empty, with each entry rounded to the nearest integer. Small inputs use direct multiplication to avoid FFT overhead.
\end{compactitem}

Overflow warning: The exact coefficients and the direct branch's intermediate products and sums must fit in \inlinecode{int64\_t}.

\Needspace*{3\baselineskip}
\textbf{Time Complexity}:
\begin{compactitem}
  \item $O(n \log n)$ per call to \inlinecode{fft()}, where $n$ is the length of the vector.
  \item $O(|a||b|)$ per call to \inlinecode{convolve()} on small inputs and $O(n \log n)$ otherwise, where $n$ is the padded transform length.
\end{compactitem}

\Needspace*{3\baselineskip}
\textbf{Space Complexity}:
\begin{compactitem}
  \item $O(1)$ auxiliary for \inlinecode{fft()}, transforming in place.
  \item $O(n)$ auxiliary for \inlinecode{convolve()}.
\end{compactitem}

\begin{lstlisting}[style=cpp]
#include <algorithm>
#include <cassert>
#include <cmath>
#include <complex>
#include <cstdint>
#include <utility>
#include <vector>

using cdbl = std::complex<double>;
const double PI = std::acos(-1.0);
const int NAIVE_CUTOFF = 150;

void fft(std::vector<cdbl> &a, bool invert = false) {
  int n = static_cast<int>(a.size());
  assert(n > 0 && (n & (n - 1)) == 0);
  for (int i = 1, j = 0; i < n; i++) {
    int bit = n >> 1;
    for (; j & bit; bit >>= 1) {
      j ^= bit;
    }
    j ^= bit;
    if (i < j) {
      std::swap(a[i], a[j]);
    }
  }
  for (int len = 2; len <= n; len <<= 1) {
    double angle = 2 * PI / len * (invert ? -1 : 1);
    cdbl root(std::cos(angle), std::sin(angle));
    for (int i = 0; i < n; i += len) {
      cdbl w(1);
      for (int k = 0; k < len / 2; k++) {
        cdbl u = a[i + k], v = a[i + k + len / 2] * w;
        a[i + k] = u + v;
        a[i + k + len / 2] = u - v;
        w *= root;
      }
    }
  }
  if (invert) {
    for (cdbl &x : a) {
      x /= n;
    }
  }
}

std::vector<int64_t> convolve(const std::vector<int64_t> &a, const std::vector<int64_t> &b) {
  if (a.empty() || b.empty()) {
    return {};
  }
  int result_size = static_cast<int>(a.size() + b.size() - 1);
  if (std::min(a.size(), b.size()) < NAIVE_CUTOFF) {
    std::vector<int64_t> res(result_size);
    for (int i = 0; i < static_cast<int>(a.size()); i++) {
      for (int j = 0; j < static_cast<int>(b.size()); j++) {
        res[i + j] += a[i] * b[j];  // Overflow warning.
      }
    }
    return res;
  }
  int n = 1;
  while (n < result_size) {
    n <<= 1;
  }
  std::vector<cdbl> fa(a.begin(), a.end()), fb(b.begin(), b.end());
  fa.resize(n);
  fb.resize(n);
  fft(fa);
  fft(fb);
  for (int i = 0; i < n; i++) {
    fa[i] *= fb[i];
  }
  fft(fa, true);
  std::vector<int64_t> res(result_size);
  for (int i = 0; i < result_size; i++) {
    res[i] = std::llround(fa[i].real());
  }
  return res;
}
\end{lstlisting}
\Needspace*{4\baselineskip}
\textbf{Example Usage}\par
\begin{lstlisting}[style=cpp]
using namespace std;

int main() {
  // (1 + 2x + 3x^2)(4 + 5x + 6x^2) = 4 + 13x + 28x^2 + 27x^3 + 18x^4.
  vector<int64_t> a{1, 2, 3}, b{4, 5, 6};
  assert((convolve(a, b) == vector<int64_t>{4, 13, 28, 27, 18}));
  assert(convolve({}, b).empty());
  assert((convolve(vector<int64_t>{5}, vector<int64_t>{7}) == vector<int64_t>{35}));
  vector<int64_t> ones(160, 1);
  vector<int64_t> d = convolve(ones, ones);
  assert(d[0] == 1 && d[159] == 160 && d[318] == 1);
  return 0;
}
\end{lstlisting}
\subsection{Polynomial Operations}
\label{sec:6.6.3}
Polynomial operations treat a vector \inlinecode{a} as coefficients of $a_0 + a_1 x + \dots + a_{n-1} x^{n-1}$ over a field. For formal power series operations, vectors represent coefficients modulo $x^n$; no notion of analytic convergence is involved. This section implements the core algebra needed for formal power series and polynomial division modulo the prime $998244353$. Multiplication uses the number theoretic transform, so inverse and division are fast enough for large polynomials while still keeping the API close to the underlying coefficient vectors.

The code aliases one modular field element as \inlinecode{Coeff} and one coefficient vector as \inlinecode{Poly}. Input coefficients must lie in $[0, {\inlinecodemath{MOD}})$.

\begin{compactitem}
  \item \inlinecode{eval(a, x)} returns the value of \inlinecode{a} at \inlinecode{x} modulo \inlinecode{MOD}; \inlinecode{x} is reduced modulo \inlinecode{MOD}.
  \item \inlinecode{add(a, b)} returns \inlinecode{a + b}.
  \item \inlinecode{subtract(a, b)} returns \inlinecode{a - b}.
  \item \inlinecode{multiply(a, b)} returns \inlinecode{a * b}.
  \item \inlinecode{derivative(a)} returns the formal derivative of \inlinecode{a}.
  \item \inlinecode{antiderivative(a)} returns the formal antiderivative of \inlinecode{a} with constant term zero.
\end{compactitem}

For formal power series, the inverse modulo $x^n$ is a series $b$ satisfying $a \cdot b \equiv 1 \pmod{x^n}$. Series division uses this inverse, and logarithm follows from $(\log a)' = a'/a$. Exponential and square root use Newton iteration, doubling the number of correct coefficients each round. Power factors $a = x^t c b$, where $b(0) = 1$, and applies $a^k = x^{tk} c^k \exp(k \log b)$. To extend the normalized square-root implementation, factor out the lowest power of $x$, normalize its constant term, then multiply back its modular square root (see \secref{6.3.5}); an odd lowest nonzero power has no square root.

\begin{compactitem}
  \item \inlinecode{inverse(a, n)} returns the first \inlinecode{n} coefficients of \inlinecode{1 / a}, requiring \inlinecode{a[0]} to be nonzero.
  \item \inlinecode{series\_divide(a, b, n)} returns the first \inlinecode{n} coefficients of the formal power series \inlinecode{a / b}, requiring \inlinecode{b[0]} to be nonzero.
  \item \inlinecode{log(a, n)} returns the first \inlinecode{n} coefficients of $\log(a)$, requiring \inlinecode{a[0]} to be $1$.
  \item \inlinecode{exp(a, n)} returns the first \inlinecode{n} coefficients of $\exp(a)$, requiring \inlinecode{a[0]} to be $0$.
  \item \inlinecode{sqrt(a, n)} returns the first \inlinecode{n} coefficients of the square root of \inlinecode{a} whose constant term is $1$, requiring \inlinecode{a[0]} to be $1$.
  \item \inlinecode{power(a, k, n)} returns the first \inlinecode{n} coefficients of $a^k$ for a nonnegative \inlinecode{uint64\_t} exponent \inlinecode{k}; \inlinecode{k = 0} returns the constant series $1$.
\end{compactitem}

Euclidean polynomial division uses the reversal trick: reverse both polynomials, compute a truncated series inverse of the reversed divisor, multiply, truncate, and reverse back.

\begin{compactitem}
  \item \inlinecode{div(a, b)} returns the polynomial quotient of \inlinecode{a / b}, requiring \inlinecode{b} to be nonzero.
  \item \inlinecode{mod(a, b)} returns the polynomial remainder of \inlinecode{a / b}, requiring \inlinecode{b} to be nonzero.
\end{compactitem}

\Needspace*{3\baselineskip}
\textbf{Time Complexity}:
\begin{compactitem}
  \item $O(n)$ per call to \inlinecode{eval()}, \inlinecode{add()}, \inlinecode{subtract()}, \inlinecode{derivative()}, and \inlinecode{antiderivative()}.
  \item $O(|a||b|)$ per call to \inlinecode{multiply()} on small inputs and $O(n \log n)$ otherwise, where $n$ is the padded transform length.
  \item $O(n \log n)$ per call to \inlinecode{inverse()} and \inlinecode{series\_divide()}, where $n$ is the requested length.
  \item $O(n \log n)$ per call to \inlinecode{log()}, \inlinecode{exp()}, and \inlinecode{sqrt()}, where $n$ is the requested length. The exponential carries the largest constant, since every Newton round computes a logarithm.
  \item $O(n \log n + \log k)$ per call to \inlinecode{power()}, where $n$ is the requested length.
  \item $O(n \log n)$ per call to \inlinecode{div()} and \inlinecode{mod()}, where $n$ is the padded multiplication length.
\end{compactitem}

\textbf{Space Complexity}: $O(1)$ auxiliary per call to \inlinecode{eval()} and $O(n)$ auxiliary for all other operations.

\begin{lstlisting}[style=cpp]
#include <algorithm>
#include <cassert>
#include <cstdint>
#include <utility>
#include <vector>

using Coeff = int64_t;
using Poly = std::vector<Coeff>;

const Coeff MOD = 998244353;
const Coeff ROOT = 3;
const int MAX_POWER_OF_TWO = 23;
const int NAIVE_CUTOFF = 150;

Coeff powmod(Coeff x, uint64_t n, Coeff m = MOD) {
  Coeff res = 1 % m;
  for (x %= m; n > 0; n >>= 1) {
    if (n & 1) {
      res = res * x % m;
    }
    x = x * x % m;
  }
  return res;
}

void ntt(Poly &a, bool invert) {
  int n = static_cast<int>(a.size());
  assert(n > 0 && (n & (n - 1)) == 0 && n <= (1 << MAX_POWER_OF_TWO));
  for (int i = 1, j = 0; i < n; i++) {
    int bit = n >> 1;
    for (; j & bit; bit >>= 1) {
      j ^= bit;
    }
    j ^= bit;
    if (i < j) {
      std::swap(a[i], a[j]);
    }
  }
  for (int len = 2; len <= n; len <<= 1) {
    Coeff root = powmod(ROOT, (MOD - 1) / len);
    if (invert) {
      root = powmod(root, MOD - 2);
    }
    for (int i = 0; i < n; i += len) {
      Coeff w = 1;
      for (int k = 0; k < len / 2; k++) {
        Coeff u = a[i + k], v = a[i + k + len / 2] * w % MOD;
        a[i + k] = (u + v) % MOD;
        a[i + k + len / 2] = (u - v + MOD) % MOD;
        w = w * root % MOD;
      }
    }
  }
  if (invert) {
    Coeff n_inv = powmod(n, MOD - 2);
    for (Coeff &x : a) {
      x = x * n_inv % MOD;
    }
  }
}

void trim(Poly &a) {
  while (!a.empty() && a.back() == 0) {
    a.pop_back();
  }
}

Coeff eval(const Poly &a, Coeff x) {
  x = (x % MOD + MOD) % MOD;
  Coeff res = 0;
  for (auto it = a.rbegin(); it != a.rend(); ++it) {
    res = (res * x + *it) % MOD;
  }
  return res;
}

Poly add(const Poly &a, const Poly &b) {
  Poly res(std::max(a.size(), b.size()));
  for (int i = 0; i < static_cast<int>(res.size()); i++) {
    Coeff x = i < static_cast<int>(a.size()) ? a[i] : 0;
    Coeff y = i < static_cast<int>(b.size()) ? b[i] : 0;
    res[i] = (x + y) % MOD;
  }
  trim(res);
  return res;
}

Poly subtract(const Poly &a, const Poly &b) {
  Poly res(std::max(a.size(), b.size()));
  for (int i = 0; i < static_cast<int>(res.size()); i++) {
    Coeff x = i < static_cast<int>(a.size()) ? a[i] : 0;
    Coeff y = i < static_cast<int>(b.size()) ? b[i] : 0;
    res[i] = (x - y + MOD) % MOD;
  }
  trim(res);
  return res;
}

Poly multiply(Poly a, Poly b) {
  if (a.empty() || b.empty()) {
    return {};
  }
  int result_size = static_cast<int>(a.size() + b.size() - 1);
  if (std::min(a.size(), b.size()) < NAIVE_CUTOFF) {
    Poly res(result_size);
    for (int i = 0; i < static_cast<int>(a.size()); i++) {
      for (int j = 0; j < static_cast<int>(b.size()); j++) {
        res[i + j] = (res[i + j] + a[i] * b[j]) % MOD;
      }
    }
    trim(res);
    return res;
  }
  int n = 1;
  while (n < result_size) {
    n <<= 1;
  }
  assert(n <= (1 << MAX_POWER_OF_TWO));
  a.resize(n);
  b.resize(n);
  ntt(a, false);
  ntt(b, false);
  for (int i = 0; i < n; i++) {
    a[i] = a[i] * b[i] % MOD;
  }
  ntt(a, true);
  a.resize(result_size);
  trim(a);
  return a;
}

Poly derivative(const Poly &a) {
  if (a.size() <= 1) {
    return {};
  }
  Poly res(a.size() - 1);
  for (int i = 1; i < static_cast<int>(a.size()); i++) {
    res[i - 1] = a[i] * i % MOD;
  }
  trim(res);
  return res;
}

Poly antiderivative(const Poly &a) {
  Poly res(a.size() + 1);
  for (int i = 0; i < static_cast<int>(a.size()); i++) {
    res[i + 1] = a[i] * powmod(i + 1, MOD - 2) % MOD;
  }
  return res;
}

Poly inverse(const Poly &a, int n) {
  assert(n >= 0 && !a.empty() && a[0] != 0);
  if (n == 0) {
    return {};
  }
  Poly res{powmod(a[0], MOD - 2)};
  while (static_cast<int>(res.size()) < n) {
    int len = static_cast<int>(res.size()) << 1;
    Poly cut(a.begin(), a.begin() + std::min(static_cast<int>(a.size()), len));
    Poly prod = multiply(multiply(res, res), cut);
    res.resize(len);
    for (int i = len / 2; i < std::min(len, static_cast<int>(prod.size())); i++) {
      res[i] = (MOD - prod[i]) % MOD;
    }
  }
  res.resize(n);
  return res;
}

Poly series_divide(const Poly &a, const Poly &b, int n) {
  assert(n >= 0 && !b.empty() && b[0] != 0);
  if (n == 0) {
    return {};
  }
  Poly cut(a.begin(), a.begin() + std::min(static_cast<int>(a.size()), n));
  Poly res = multiply(cut, inverse(b, n));
  res.resize(n);
  return res;
}

Poly log(const Poly &a, int n) {
  assert(n >= 0 && !a.empty() && a[0] == 1);
  if (n == 0) {
    return {};
  }
  // Since (log a)' = a'/a, integrating the truncated quotient recovers log a.
  Poly cut(a.begin(), a.begin() + std::min(static_cast<int>(a.size()), n));
  Poly res = multiply(derivative(cut), inverse(cut, n));
  res.resize(n - 1);
  res = antiderivative(res);
  res.resize(n);
  return res;
}

Poly exp(const Poly &a, int n) {
  assert(n >= 0 && (a.empty() || a[0] == 0));
  if (n == 0) {
    return {};
  }
  // Newton iteration on log(f) - a = 0 gives f <- f*(1 - log(f) + a), doubling the number of
  // correct coefficients each round.
  Poly res{1};
  while (static_cast<int>(res.size()) < n) {
    int len = static_cast<int>(res.size()) << 1;
    Poly cut(a.begin(), a.begin() + std::min(static_cast<int>(a.size()), len));
    Poly correction = subtract(cut, log(res, len));
    correction.resize(std::max(1, static_cast<int>(correction.size())));
    correction[0] = (correction[0] + 1) % MOD;
    res = multiply(res, correction);
    res.resize(len);
  }
  res.resize(n);
  return res;
}

Poly sqrt(const Poly &a, int n) {
  assert(n >= 0 && !a.empty() && a[0] == 1);
  if (n == 0) {
    return {};
  }
  // Newton iteration on f^2 - a = 0 gives f <- (f + a/f) / 2.
  const Coeff INV2 = powmod(2, MOD - 2);
  Poly res{1};
  while (static_cast<int>(res.size()) < n) {
    int len = static_cast<int>(res.size()) << 1;
    Poly cut(a.begin(), a.begin() + std::min(static_cast<int>(a.size()), len));
    res = add(res, multiply(cut, inverse(res, len)));
    res.resize(len);
    for (Coeff &c : res) {
      c = c * INV2 % MOD;
    }
  }
  res.resize(n);
  return res;
}

Poly power(const Poly &a, uint64_t k, int n) {
  assert(n >= 0);
  if (n == 0) {
    return {};
  }
  Poly res(n);
  if (k == 0) {
    res[0] = 1;
    return res;
  }
  int limit = std::min(static_cast<int>(a.size()), n);
  int first = 0;
  while (first < limit && a[first] == 0) {
    first++;
  }
  if (first == limit || (first > 0 && k > static_cast<uint64_t>((n - 1) / first))) {
    return res;
  }
  int shift = static_cast<int>(static_cast<uint64_t>(first) * k);
  int len = n - shift;
  int terms = std::min(static_cast<int>(a.size()) - first, len);
  Poly normalized(a.begin() + first, a.begin() + first + terms);
  Coeff leading = normalized[0];
  Coeff leading_inv = powmod(leading, MOD - 2);
  for (Coeff &c : normalized) {
    c = c * leading_inv % MOD;
  }
  Poly exponent = log(normalized, len);
  Coeff k_mod = static_cast<Coeff>(k % static_cast<uint64_t>(MOD));
  for (Coeff &c : exponent) {
    c = c * k_mod % MOD;
  }
  Poly body = exp(exponent, len);
  Coeff leading_power = powmod(leading, k);
  for (int i = 0; i < len; i++) {
    res[i + shift] = body[i] * leading_power % MOD;
  }
  return res;
}

Poly div(Poly a, Poly b) {
  trim(a);
  trim(b);
  assert(!b.empty());
  if (a.size() < b.size()) {
    return {};
  }
  int quotient_size = static_cast<int>(a.size() - b.size() + 1);
  std::reverse(a.begin(), a.end());
  std::reverse(b.begin(), b.end());
  b.resize(quotient_size);
  Poly q = multiply(a, inverse(b, quotient_size));
  q.resize(quotient_size);
  std::reverse(q.begin(), q.end());
  trim(q);
  return q;
}

Poly mod(const Poly &a, const Poly &b) {
  Poly q = div(a, b);
  Poly res = subtract(a, multiply(q, b));
  if (res.size() >= b.size()) {
    res.resize(b.size() - 1);
  }
  trim(res);
  return res;
}
\end{lstlisting}
\Needspace*{4\baselineskip}
\textbf{Example Usage}\par
\begin{lstlisting}[style=cpp]
#include <cassert>

int main() {
  // 5 + 7x + 11x^2 evaluates to 63 at x = 2 and 9 at x = -1.
  Poly p{5, 7, 11};
  assert(eval(p, 2) == 63 && eval(p, -1) == 9);

  // (1 + 2x + 3x^2)(4 + 5x + 6x^2) = 4 + 13x + 28x^2 + 27x^3 + 18x^4.
  Poly a{1, 2, 3}, b{4, 5, 6};
  assert((add(a, b) == Poly{5, 7, 9}));
  assert((subtract(b, a) == Poly{3, 3, 3}));
  assert((multiply(a, b) == Poly{4, 13, 28, 27, 18}));

  // d/dx (5 + 7x + 11x^2) = 7 + 22x; integrating back chooses constant term 0.
  assert((derivative(p) == Poly{7, 22}));
  assert((antiderivative(derivative(p)) == Poly{0, 7, 11}));

  // (1 - x)^-1 = 1 + x + x^2 + x^3 + ... modulo x^6.
  assert((inverse({1, MOD - 1}, 6) == Poly{1, 1, 1, 1, 1, 1}));

  // (1 + x) / (1 - x) = 1 + 2x + 2x^2 + ... modulo x^5.
  assert((series_divide({1, 1}, {1, MOD - 1}, 5) == Poly{1, 2, 2, 2, 2}));

  // (x^3 - 1) / (x - 1) = x^2 + x + 1, remainder 0.
  Poly dividend{MOD - 1, 0, 0, 1};
  Poly divisor{MOD - 1, 1};
  assert((div(dividend, divisor) == Poly{1, 1, 1}));
  assert(mod(dividend, divisor).empty());

  // exp(x) = 1 + x + x^2/2 + x^3/6 + ..., checked by clearing each denominator.
  Poly exp_x = exp({0, 1}, 4);
  assert(exp_x[0] == 1 && exp_x[1] == 1);
  assert(exp_x[2] * 2 % MOD == 1 && exp_x[3] * 6 % MOD == 1);

  // log(1 + x) = x - x^2/2 + x^3/3 - ...
  Poly log_1x = log({1, 1}, 4);
  assert(log_1x[0] == 0 && log_1x[1] == 1);
  assert((log_1x[2] * 2 + 1) % MOD == 0 && log_1x[3] * 3 % MOD == 1);

  // Truncated to the same length, exp and log invert each other.
  Poly series{1, 5, 9, 2, 7};
  Poly logged = log(series, 5);
  assert(logged[0] == 0);
  assert(exp(logged, 5) == series);

  // A series square root squares back to the original, modulo x^5.
  Poly root = sqrt(series, 5);
  Poly root_squared = multiply(root, root);
  root_squared.resize(5);
  assert(root_squared == series);

  // (2x + x^2)^3 = 8x^3 + 12x^4 + 6x^5 + x^6.
  assert((power({0, 2, 1}, 3, 7) == Poly{0, 0, 0, 8, 12, 6, 1}));
  assert((power({}, 0, 4) == Poly{1, 0, 0, 0}));

  Poly ones(160, 1);
  Poly square = multiply(ones, ones);
  assert(square[0] == 1 && square[159] == 160 && square[318] == 1);
  return 0;
}
\end{lstlisting}
\subsection{Polynomial Interpolation}
\label{sec:6.6.4}
Recovers the unique polynomial of degree less than $n$ that passes through $n$ given points, working modulo a prime. The $x$ coordinates must be distinct; the $y$ coordinates may be anything. Two operations are provided: reconstructing the full coefficient vector, and evaluating the interpolating polynomial at a single point without ever forming the coefficients.

Coefficient recovery uses Newton's divided differences. It first replaces the sample values by their successive divided differences in place, then expands the Newton basis $1$, $(x - x_0)$, $(x - x_0)(x - x_1)$, $\ldots$ into ordinary coefficients by multiplying through one factor at a time. Point evaluation instead applies the Lagrange formula directly, summing each sample weighted by the product of $(t - x_j) / (x_k - x_j)$ over the other nodes. Coefficient recovery performs $O(n^{2})$ modular inverses, while point evaluation performs $n$.

\begin{compactitem}
  \item \inlinecode{interpolate(x, y)} returns the coefficients $a_0, ..., a_{n-1}$ (constant term first) of the polynomial $P$ with $P({\inlinecodemath{x[i]}}) = {\inlinecodemath{y[i]}}$, modulo \inlinecode{MOD}. The entries of \inlinecode{x} must be distinct modulo \inlinecode{MOD}.
  \item \inlinecode{interpolate\_at(x, y, t)} returns $P({\inlinecodemath{t}})$ modulo \inlinecode{MOD} for that same polynomial, without building its coefficients.
\end{compactitem}

\textbf{Time Complexity}: $O(n^{2} \log p)$ per call to \inlinecode{interpolate()} and $O(n^{2} + n \log p)$ per call to \inlinecode{interpolate\_at()}, where $n$ is the number of points and $p$ is \inlinecode{MOD}.

\textbf{Space Complexity}: $O(n)$ auxiliary.

\begin{lstlisting}[style=cpp]
#include <cassert>
#include <cstdint>
#include <vector>

const int64_t MOD = 998244353;

int64_t powmod(int64_t x, int64_t n, int64_t m) {
  int64_t res = 1 % m;
  for (x %= m; n > 0; n >>= 1) {
    if (n & 1) {
      res = res * x % m;
    }
    x = x * x % m;
  }
  return res;
}

int64_t modnorm(int64_t a) { return (a %= MOD) < 0 ? a + MOD : a; }
int64_t addmod(int64_t a, int64_t b) { return modnorm(modnorm(a) + modnorm(b)); }
int64_t submod(int64_t a, int64_t b) { return modnorm(modnorm(a) - modnorm(b)); }
int64_t inv(int64_t a) { return powmod(modnorm(a), MOD - 2, MOD); }

void validate_points(const std::vector<int64_t> &x) {
  for (int i = 0; i < static_cast<int>(x.size()); i++) {
    for (int j = 0; j < i; j++) {
      assert(modnorm(x[i]) != modnorm(x[j]));
    }
  }
}

std::vector<int64_t> interpolate(std::vector<int64_t> x, std::vector<int64_t> y) {
  assert(x.size() == y.size());
  validate_points(x);
  int n = static_cast<int>(x.size());
  for (int i = 0; i < n; i++) {
    x[i] = modnorm(x[i]);
    y[i] = modnorm(y[i]);
  }
  // Replace y by its divided differences in place.
  for (int k = 0; k < n; k++) {
    for (int i = k + 1; i < n; i++) {
      y[i] = submod(y[i], y[k]) * inv(submod(x[i], x[k])) % MOD;
    }
  }
  // Expand the Newton basis into ordinary coefficients.
  std::vector<int64_t> res(n), temp(n);
  temp[0] = 1;
  for (int k = 0; k < n; k++) {
    int64_t last = 0;
    for (int i = 0; i < n; i++) {
      res[i] = addmod(res[i], y[k] * temp[i] % MOD);
      int64_t old = temp[i];
      temp[i] = last;
      last = old;
      temp[i] = submod(temp[i], last * x[k] % MOD);
    }
  }
  return res;
}

int64_t interpolate_at(const std::vector<int64_t> &x, const std::vector<int64_t> &y, int64_t t) {
  assert(x.size() == y.size());
  validate_points(x);
  int n = static_cast<int>(x.size());
  t = modnorm(t);
  int64_t res = 0;
  for (int k = 0; k < n; k++) {
    int64_t num = 1, den = 1;
    for (int j = 0; j < n; j++) {
      if (j != k) {
        num = num * submod(t, x[j]) % MOD;
        den = den * submod(x[k], x[j]) % MOD;
      }
    }
    res = addmod(res, modnorm(y[k]) * num % MOD * inv(den) % MOD);
  }
  return res;
}
\end{lstlisting}
\Needspace*{4\baselineskip}
\textbf{Example Usage}\par
\begin{lstlisting}[style=cpp]
#include <cassert>
using namespace std;

int main() {
  // P(x) = x^2 + 2x + 3 sampled at x = 0, 1, 2 gives y = 3, 6, 11.
  vector<int64_t> x{0, 1, 2}, y{3, 6, 11};
  assert((interpolate(x, y) == vector<int64_t>{3, 2, 1}));
  assert(interpolate_at(x, y, 5) == 38);  // 25 + 10 + 3.

  // Non-consecutive nodes for P(x) = 2x^3 - x + 4.
  auto poly = [](int64_t t) { return modnorm(2 * t * t * t - t + 4); };
  vector<int64_t> xs{2, 5, 7, 10}, ys;
  for (int64_t v : xs) {
    ys.push_back(poly(v));
  }
  assert((interpolate(xs, ys) == vector<int64_t>{4, MOD - 1, 0, 2}));  // 4 - x + 2x^3.
  assert(interpolate_at(xs, ys, 13) == poly(13));
  return 0;
}
\end{lstlisting}
\subsection{Linear Recurrence (Berlekamp-Massey)}
\label{sec:6.6.5}
The Berlekamp-Massey algorithm finds the shortest linear recurrence that generates a given sequence over a field. Given the first terms of a sequence taken modulo a prime, it returns coefficients $\Lambda_0, \dots, \Lambda_{L-1}$ of the minimal recurrence $s_i = \Lambda_0 \, s_{i-1} + \Lambda_1 \, s_{i-2} + \dots + \Lambda_{L-1} \, s_{i-L} \pmod{p}$, holding for every $i \geq L$. It scans the sequence once, maintaining a current candidate recurrence and correcting it whenever a term fails to match the prediction, using the last position where a correction was needed. The order $L$ never decreases, and $2L$ terms suffice to determine a recurrence of order $L$. This is the standard tool for guessing a linear recurrence from sampled values, after which the sequence can be extended or its $n$-th term found by fast methods.

An order-$L$ recurrence generally needs at least $2L$ terms to be determined, so the result is most trustworthy when \inlinecode{s} is comfortably longer than twice the returned length. If the returned $L$ is close to $n / 2$, the sequence is under-determined and the recurrence may be a spurious overfit.

The code operates modulo the prime $998244353$; change \inlinecode{MOD} to use a different prime field.

\begin{compactitem}
  \item \inlinecode{berlekamp\_massey(s)} returns the coefficients of the shortest recurrence that the sequence \inlinecode{s} satisfies, as described above. The returned length is the order $L$ of the recurrence. An empty vector is returned when \inlinecode{s} is entirely zero (no recurrence is needed). Entries are reduced modulo \inlinecode{MOD} internally.
\end{compactitem}

\textbf{Time Complexity}: $O(n^{2} + L \log p)$ per call, where $n$ is the length of \inlinecode{s}, $L$ is the order of the recurrence found, and $p$ is the modulus.

\textbf{Space Complexity}: $O(n)$ auxiliary.

\begin{lstlisting}[style=cpp]
#include <cstdint>
#include <vector>

const int64_t MOD = 998244353;

int64_t powmod(int64_t x, int64_t n, int64_t m) {
  int64_t res = 1 % m;
  for (x %= m; n > 0; n >>= 1) {
    if (n & 1) {
      res = res * x % m;
    }
    x = x * x % m;
  }
  return res;
}

std::vector<int64_t> berlekamp_massey(const std::vector<int64_t> &s) {
  int n = static_cast<int>(s.size()), len = 0, m = 0;
  if (n == 0) {
    return {};
  }
  std::vector<int64_t> cur(n), last(n), prev;
  cur[0] = last[0] = 1;
  int64_t inv_last_delta = 1;
  for (int i = 0; i < n; i++) {
    m++;
    int64_t delta = (s[i] % MOD + MOD) % MOD;
    for (int j = 1; j <= len; j++) {
      int64_t value = (s[i - j] % MOD + MOD) % MOD;
      delta = (delta + cur[j] * value) % MOD;
    }
    if (delta == 0) {
      continue;
    }
    prev = cur;
    int64_t coeff = delta * inv_last_delta % MOD;
    for (int j = m; j < n; j++) {
      cur[j] = (cur[j] - coeff * last[j - m] % MOD + MOD) % MOD;
    }
    if (2 * len > i) {
      continue;
    }
    len = i + 1 - len;
    last = prev;
    inv_last_delta = powmod(delta, MOD - 2, MOD);  // Refresh only when the length increases.
    m = 0;
  }
  cur.resize(len + 1);
  cur.erase(cur.begin());
  for (int64_t &x : cur) {
    x = (MOD - x) % MOD;
  }
  return cur;
}
\end{lstlisting}
\Needspace*{4\baselineskip}
\textbf{Example Usage}\par
\begin{lstlisting}[style=cpp]
#include <cassert>
using namespace std;

int main() {
  assert(berlekamp_massey({}).empty());
  // Fibonacci: s[i] = s[i-1] + s[i-2].
  vector<int64_t> fib{1, 1, 2, 3, 5, 8, 13, 21};
  assert((berlekamp_massey(fib) == vector<int64_t>{1, 1}));

  // Geometric: s[i] = 2 * s[i-1].
  vector<int64_t> geo{1, 2, 4, 8, 16, 32};
  assert((berlekamp_massey(geo) == vector<int64_t>{2}));
  assert((berlekamp_massey(vector<int64_t>{1, -1, 1, -1, 1, -1}) == vector<int64_t>{MOD - 1}));

  // Tribonacci: s[i] = s[i-1] + s[i-2] + s[i-3].
  vector<int64_t> trib{0, 1, 1, 2, 4, 7, 13, 24, 44};
  assert((berlekamp_massey(trib) == vector<int64_t>{1, 1, 1}));
  return 0;
}
\end{lstlisting}
\subsection{Linear Recurrence (Kitamasa)}
\label{sec:6.6.6}
Computes the $n$-th term of a linear recurrence modulo a prime in logarithmic time. Given a recurrence of order $L$, $s_i = c_0 s_{i-1} + c_1 s_{i-2} + \ldots + c_{L-1} s_{i-L}$, together with the first $L$ terms, the naive approach unrolls $n$ steps. Kitamasa's method instead jumps directly to index $n$ by working with the characteristic polynomial. The coefficient layout matches the output of Berlekamp-Massey, so the two compose directly: guess the recurrence from sampled values, then jump to any index.

The key identity is that the $n$-th term is a fixed linear combination of the first $L$ terms, with weights given by $x^n$ reduced modulo the characteristic polynomial $f(x) = x^L - c_0 x^{L-1} - ... - c_{L-1}$. The reduction $x^L = c_0 x^{L-1} + ... + c_{L-1}$ lets any product of two degree-$<L$ polynomials be folded back to degree $<L$, so $x^n \mod f$ is obtained by binary exponentiation, and the $n$-th term is the dot product of its coefficients with the initial terms.

Each doubling step is one polynomial multiplication followed by one remainder, both taken from section \secref{6.6.3}, where the multiplication is a number theoretic transform and the remainder is a Newton-iterated division. Squaring and reducing by hand instead costs $O(L^{2})$ per step, which is simpler to write but asymptotically worse; \inlinecode{MOD} is inherited from that section and must stay transform-friendly.

\begin{compactitem}
  \item \inlinecode{kth\_term(rec, init, n)} returns $s_n$ modulo \inlinecode{MOD}, where \inlinecode{rec} holds the coefficients $c_0, ..., c_{L-1}$ and \inlinecode{init} holds at least $L$ initial terms $s_0, ..., s_{L-1}$. Indexing is 0-based and \inlinecode{n} $\geq 0$. Values are reduced modulo \inlinecode{MOD} internally.
\end{compactitem}

\textbf{Time Complexity}: $O(L \log L \log n)$ per call.

\textbf{Space Complexity}: $O(L)$ auxiliary.

\begin{lstlisting}[style=cpp]
#include <cassert>
#include <cstdint>
#include <vector>

#define main polynomial_example
#include "6.6.3_Polynomial_Operations.cpp"
#undef main

int64_t kth_term(const std::vector<int64_t> &rec, const std::vector<int64_t> &init, int64_t n) {
  assert(n >= 0 && init.size() >= rec.size());
  int L = static_cast<int>(rec.size());
  if (n < static_cast<int64_t>(init.size())) {
    return ((init[n] % MOD) + MOD) % MOD;
  }
  if (L == 0) {
    return 0;
  }
  // The characteristic polynomial f(x) = x^L - c_0 x^(L-1) - ... - c_(L-1), indexed by degree.
  Poly f(L + 1);
  f[L] = 1;
  for (int j = 0; j < L; j++) {
    f[L - 1 - j] = ((-rec[j]) % MOD + MOD) % MOD;
  }
  Poly result{1}, base = mod(Poly{0, 1}, f);  // The polynomials x^0 and x, each modulo f.
  for (int64_t e = n; e > 0; e >>= 1) {
    if (e & 1) {
      result = mod(multiply(result, base), f);
    }
    base = mod(multiply(base, base), f);
  }
  int64_t ans = 0;
  for (int j = 0; j < L && j < static_cast<int>(result.size()); j++) {
    ans = (ans + result[j] * (((init[j] % MOD) + MOD) % MOD)) % MOD;
  }
  return ans;
}
\end{lstlisting}
\Needspace*{4\baselineskip}
\textbf{Example Usage}\par
\begin{lstlisting}[style=cpp]
using namespace std;

int main() {
  // Fibonacci s_i = s_{i-1} + s_{i-2} with s_0 = s_1 = 1: 1, 1, 2, 3, 5, 8, ...
  vector<int64_t> fib_rec{1, 1}, fib_init{1, 1};
  assert(kth_term(fib_rec, fib_init, 0) == 1);
  assert(kth_term(fib_rec, fib_init, 9) == 55);
  assert(kth_term(fib_rec, fib_init, 10) == 89);

  // Tribonacci s_i = s_{i-1} + s_{i-2} + s_{i-3} with s_0,s_1,s_2 = 0,1,1.
  vector<int64_t> trib_rec{1, 1, 1}, trib_init{0, 1, 1};
  assert(kth_term(trib_rec, trib_init, 8) == 44);

  // Geometric s_i = 3 s_{i-1}, s_0 = 1: matches 3^n modulo MOD.
  vector<int64_t> geo_rec{3}, geo_init{1};
  int64_t pw = 1;
  for (int i = 0; i < 50; i++) {
    assert(kth_term(geo_rec, geo_init, i) == pw);
    pw = pw * 3 % MOD;
  }
  assert(kth_term({-1}, {1}, 5) == MOD - 1);
  return 0;
}
\end{lstlisting}

\section{\texorpdfstring{Probability and Statistics}{6.7 Probability and Statistics}}
\setcounter{section}{7}
\setcounter{subsection}{0}
\subsection{Probability Distributions}
\label{sec:6.7.1}
The following functions evaluate the standard probability distributions. For a discrete distribution, the probability mass function gives the probability of an exact outcome and the cumulative distribution function gives the probability of that outcome or any smaller one. For a continuous distribution, the density integrates to those probabilities rather than being one, so only the cumulative function is a probability. Counts must be nonnegative, and every probability argument must lie in $[0, 1]$.

Every mass function here is evaluated in logarithms and exponentiated once at the end. Computing $\binom{n}{k} p^k (1 - p)^{n-k}$ directly overflows the binomial coefficient and underflows the powers long before the product itself becomes small, whereas the logarithms stay near the magnitude of the answer. \inlinecode{std::lgamma} supplies $\ln n!$ as $\ln \Gamma(n + 1)$ without forming $n!$, so the factorials never exist, and outcomes are ordered so that no cancellation occurs.

Each cumulative function sums its mass function term by term, which keeps the code short and the error small but costs one term per outcome. For a single tail probability at a very large count the normal approximation with a continuity correction is the usual substitute, while \inlinecode{quantile()} inverts any of these by the bisection of section \secref{5.1.1}.

\begin{compactitem}
  \item \inlinecode{log\_factorial(n)} returns $\ln(n!)$, and \inlinecode{log\_choose(n, k)} returns $\ln \binom{n}{k}$, which is $-\infty$ when \inlinecode{k} lies outside $[0, {\inlinecodemath{n}}]$.
  \item \inlinecode{binomial\_pmf(n, k, p)} and \inlinecode{binomial\_cdf(n, k, p)} evaluate the number of successes in \inlinecode{n} independent trials that each succeed with probability \inlinecode{p}.
  \item \inlinecode{geometric\_pmf(k, p)} and \inlinecode{geometric\_cdf(k, p)} evaluate the number of trials up to and including the first success, so \inlinecode{k} is at least $1$ and \inlinecode{p} must be positive.
  \item \inlinecode{negative\_binomial\_pmf(k, r, p)} and \inlinecode{negative\_binomial\_cdf(k, r, p)} evaluate the number of trials up to and including the \inlinecode{r}-th success, so \inlinecode{k} is at least \inlinecode{r}, which is at least $1$.
  \item \inlinecode{poisson\_pmf(k, lambda)} and \inlinecode{poisson\_cdf(k, lambda)} evaluate the number of events in one unit of time when events arrive independently at nonnegative mean rate \inlinecode{lambda}.
  \item \inlinecode{hypergeometric\_pmf(n, k, draws, hits)} and \inlinecode{hypergeometric\_cdf(n, k, draws, hits)} evaluate the number of marked items among \inlinecode{draws} items taken without replacement from a population of \inlinecode{n} items of which \inlinecode{k} are marked.
  \item \inlinecode{normal\_pdf(x, mu = 0, sigma = 1)} and \inlinecode{normal\_cdf(x, mu = 0, sigma = 1)} evaluate the normal distribution of mean \inlinecode{mu} and positive standard deviation \inlinecode{sigma}, defaulting to the standard normal distribution.
  \item \inlinecode{exponential\_pdf(x, lambda)} and \inlinecode{exponential\_cdf(x, lambda)} evaluate the waiting time until the next event when events arrive at positive mean rate \inlinecode{lambda}.
  \item \inlinecode{quantile(cdf, p, lo, hi, iterations = 100)} returns the smallest $x$ in [\inlinecode{lo}, \inlinecode{hi}] whose cumulative probability \inlinecode{cdf(x)} reaches \inlinecode{p}, for any nondecreasing \inlinecode{cdf}. The bounds must satisfy \inlinecode{lo} $\leq$ \inlinecode{hi}, \inlinecode{cdf(hi)} must reach \inlinecode{p}, and \inlinecode{iterations} must be positive.
\end{compactitem}

Sampling is left to \inlinecode{\textless{}random\textgreater{}}, whose \inlinecode{std::binomial\_distribution}, \inlinecode{std::poisson\_distribution}, \inlinecode{std::normal\_distribution}, and their siblings draw from exactly these families. What the standard library does not provide, and this section supplies, is the ability to evaluate and invert them.

Overflow warning: \inlinecode{log\_choose()} is exact only while \inlinecode{std::lgamma} resolves neighboring integers, which holds well past the range where a \inlinecode{double} probability retains any precision.

\Needspace*{3\baselineskip}
\textbf{Time Complexity}:
\begin{compactitem}
  \item $O(1)$ per call to every mass and density function, and to \inlinecode{log\_factorial()} and \inlinecode{log\_choose()}.
  \item $O(k)$ per call to \inlinecode{binomial\_cdf()}, \inlinecode{poisson\_cdf()}, and \inlinecode{negative\_binomial\_cdf()}, where $k$ is the outcome whose tail is summed, and $O(h)$ per call to \inlinecode{hypergeometric\_cdf()}, where $h$ is its \inlinecode{hits} argument.
  \item $O(1)$ per call to \inlinecode{geometric\_cdf()}, \inlinecode{normal\_cdf()}, and \inlinecode{exponential\_cdf()}, which have closed forms.
  \item $O(n)$ calls to \inlinecode{cdf()} per call to \inlinecode{quantile()}, where $n$ is the number of iterations.
\end{compactitem}

\textbf{Space Complexity}: $O(1)$ auxiliary for all operations.

\begin{lstlisting}[style=cpp]
#include <algorithm>
#include <cassert>
#include <cmath>
#include <limits>

double log_factorial(int n) {
  assert(n >= 0);
  return std::lgamma(n + 1.0);
}

double log_choose(int n, int k) {
  assert(n >= 0);
  if (k < 0 || k > n) {
    return -std::numeric_limits<double>::infinity();
  }
  return log_factorial(n) - log_factorial(k) - log_factorial(n - k);
}

double binomial_pmf(int n, int k, double p) {
  assert(n >= 0 && 0 <= p && p <= 1);
  if (k < 0 || k > n) {
    return 0;
  }
  if (p <= 0) {
    return k == 0 ? 1 : 0;
  }
  if (p >= 1) {
    return k == n ? 1 : 0;
  }
  return std::exp(log_choose(n, k) + k * std::log(p) + (n - k) * std::log1p(-p));
}

double binomial_cdf(int n, int k, double p) {
  assert(n >= 0 && 0 <= p && p <= 1);
  double res = 0;
  int last = std::min(k, n);
  for (int i = 0; i <= last; i++) {
    res += binomial_pmf(n, i, p);
    if (i == last) {
      break;
    }
  }
  return std::min(res, 1.0);
}

double geometric_pmf(int k, double p) {
  assert(0 < p && p <= 1);
  return k < 1 ? 0 : p * std::pow(1 - p, k - 1);
}

double geometric_cdf(int k, double p) {
  assert(0 < p && p <= 1);
  return k < 1 ? 0 : -std::expm1(k * std::log1p(-p));
}

double negative_binomial_pmf(int k, int r, double p) {
  assert(r >= 1 && 0 < p && p <= 1);
  if (k < r) {
    return 0;
  }
  if (p >= 1) {
    return k == r ? 1 : 0;
  }
  return std::exp(log_choose(k - 1, r - 1) + r * std::log(p) + (k - r) * std::log1p(-p));
}

double negative_binomial_cdf(int k, int r, double p) {
  assert(r >= 1 && 0 < p && p <= 1);
  if (k < r) {
    return 0;
  }
  double res = 0;
  for (int i = r;; i++) {
    res += negative_binomial_pmf(i, r, p);
    if (i == k) {
      break;
    }
  }
  return std::min(res, 1.0);
}

double poisson_pmf(int k, double lambda) {
  assert(lambda >= 0);
  if (k < 0) {
    return 0;
  }
  if (lambda <= 0) {
    return k == 0 ? 1 : 0;
  }
  return std::exp(-lambda + k * std::log(lambda) - log_factorial(k));
}

double poisson_cdf(int k, double lambda) {
  assert(lambda >= 0);
  if (k < 0) {
    return 0;
  }
  double res = 0;
  for (int i = 0;; i++) {
    res += poisson_pmf(i, lambda);
    if (i == k) {
      break;
    }
  }
  return std::min(res, 1.0);
}

double hypergeometric_pmf(int n, int k, int draws, int hits) {
  assert(n >= 0 && 0 <= k && k <= n && 0 <= draws && draws <= n);
  if (hits < 0 || hits > draws || hits > k || draws - hits > n - k) {
    return 0;
  }
  return std::exp(log_choose(k, hits) + log_choose(n - k, draws - hits) - log_choose(n, draws));
}

double hypergeometric_cdf(int n, int k, int draws, int hits) {
  assert(n >= 0 && 0 <= k && k <= n && 0 <= draws && draws <= n);
  if (hits < 0) {
    return 0;
  }
  double res = 0;
  int last = std::min({hits, k, draws});
  for (int i = 0; i <= last; i++) {
    res += hypergeometric_pmf(n, k, draws, i);
    if (i == last) {
      break;
    }
  }
  return std::min(res, 1.0);
}

double normal_pdf(double x, double mu = 0, double sigma = 1) {
  assert(sigma > 0);
  static const double INV_SQRT_2PI = 0.3989422804014327;
  double z = (x - mu) / sigma;
  return INV_SQRT_2PI / sigma * std::exp(-0.5 * z * z);
}

double normal_cdf(double x, double mu = 0, double sigma = 1) {
  assert(sigma > 0);
  static const double INV_SQRT_2 = 0.7071067811865476;
  return 0.5 * std::erfc(-(x - mu) * INV_SQRT_2 / sigma);
}

double exponential_pdf(double x, double lambda) {
  assert(lambda > 0);
  return x < 0 ? 0 : lambda * std::exp(-lambda * x);
}

double exponential_cdf(double x, double lambda) {
  assert(lambda > 0);
  return x < 0 ? 0 : -std::expm1(-lambda * x);
}

template<typename Cdf>
double quantile(Cdf cdf, double p, double lo, double hi, int iterations = 100) {
  assert(0 <= p && p <= 1 && lo <= hi && iterations > 0);
  if (cdf(lo) >= p) {
    return lo;
  }
  assert(cdf(hi) >= p);
  for (int i = 0; i < iterations; i++) {
    double mid = lo + (hi - lo) / 2;
    if (cdf(mid) >= p) {
      hi = mid;
    } else {
      lo = mid;
    }
  }
  return hi;
}
\end{lstlisting}
\Needspace*{4\baselineskip}
\textbf{Example Usage}\par
\begin{lstlisting}[style=cpp]
#include <cassert>
using namespace std;

bool EQ(double a, double b) {
  return fabs(a - b) < 1e-9;
}

int main() {
  // Three fair coin flips give two heads with probability 3/8.
  assert(EQ(binomial_pmf(3, 2, 0.5), 0.375));
  assert(EQ(binomial_cdf(3, 3, 0.5), 1));
  // Large counts stay finite where the binomial coefficient alone would not.
  assert(EQ(binomial_cdf(1000, 1000, 0.5), 1));
  assert(binomial_pmf(1000, 500, 0.5) > 0.02);

  // A fair die shows its first six on the third roll with probability 25/216.
  assert(EQ(geometric_pmf(3, 1.0 / 6), 25.0 / 216));
  assert(EQ(geometric_cdf(1, 0.25), 0.25));
  assert(EQ(negative_binomial_pmf(3, 2, 0.5), 0.25));

  assert(EQ(poisson_pmf(0, 2), exp(-2.0)));
  assert(EQ(poisson_cdf(1, 1), 2 * exp(-1.0)));

  // Two marked items among five drawn from ten, of which four are marked.
  assert(EQ(hypergeometric_pmf(10, 4, 5, 2), 10.0 / 21));
  assert(EQ(hypergeometric_cdf(10, 4, 5, 4), 1));

  assert(EQ(normal_cdf(0), 0.5));
  assert(fabs(normal_cdf(1.959963984540054) - 0.975) < 1e-9);
  assert(EQ(exponential_cdf(0, 3), 0));

  // Inverting a cumulative function recovers the value that produced it.
  assert(
      fabs(quantile([](double x) { return normal_cdf(x); }, 0.975, -10, 10) - 1.959963984540054) <
      1e-9
  );
  return 0;
}
\end{lstlisting}
\subsection{Expected Value DP}
\label{sec:6.7.2}
A random process that moves between states accumulates an expected cost that satisfies one linear equation per state. If leaving state $u$ costs $c_u$ and sends the process to state $v$ with probability $p_{uv}$, then linearity of expectation gives $E[u] = c_u + \sum_v p_{uv} E[v]$, where the expected cost of a state is the immediate cost plus the average of the expected costs of its successors. Counting steps is the special case $c_u = 1$ for every non-absorbing state, and an absorbing state is one with no outgoing transitions, so that $E = c$ there.

When the transition graph is acyclic, the equations can be evaluated instead of solved: process the states in reverse topological order and every successor on the right-hand side is already known. This is ordinary backward induction, and it is the reason expected-value problems are usually solved from the goal backwards rather than from the start forwards.

A self-loop is the one cycle this method still handles. If the process stays in $u$ with probability $q$, then $E[u] = c_u + q E[u] + \sum_{v \neq u} p_{uv} E[v]$, and isolating $E[u]$ gives $E[u] = (c_u + \sum_{v \neq u} p_{uv} E[v]) / (1 - q)$. Equivalently, the process retries until it leaves, so the expected number of attempts is the geometric mean $1/(1 - q)$. Self-loops appear whenever a move can be rejected and repeated, such as a die roll that would overshoot the goal or a draw that repeats a coupon already held.

\begin{compactitem}
  \item \inlinecode{expected\_cost(trans, cost)} returns a vector \inlinecode{e} such that \inlinecode{e[u]} is the expected total cost accumulated from state \inlinecode{u} until the process reaches a state with no outgoing transitions. States are the indices of \inlinecode{trans}, each entry \inlinecode{trans[u]} lists the outgoing transitions of \inlinecode{u} as pairs of (\inlinecode{next\_state}, \inlinecode{probability}), and \inlinecode{cost[u]} is the cost charged on leaving \inlinecode{u}. Transition destinations must be valid state indices, and their nonnegative probabilities must sum to at most $1$, with any missing probability treated as ending the process.
\end{compactitem}

Aside from self-loops the transition graph must be acyclic, since a genuine cycle leaves the equations mutually dependent and they must then be solved as a linear system; two states that reach each other call for the absorbing Markov chains of section \secref{6.7.3} instead.

\textbf{Time Complexity}: $O(n + m)$ per call, where $n$ is the number of states and $m$ is the total number of transitions.

\Needspace*{3\baselineskip}
\textbf{Space Complexity}:
\begin{compactitem}
  \item $O(n + m)$ auxiliary.
  \item $O(n)$ for the returned vector.
\end{compactitem}

\begin{lstlisting}[style=cpp]
#include <cassert>
#include <utility>
#include <vector>

std::vector<double> expected_cost(
    const std::vector<std::vector<std::pair<int, double>>> &trans, const std::vector<double> &cost
) {
  int n = static_cast<int>(trans.size());
  assert(static_cast<int>(cost.size()) == n);
  std::vector<int> indeg(n);
  for (int u = 0; u < n; u++) {
    double prob = 0;
    for (auto [v, p] : trans[u]) {
      assert(v >= 0 && v < n && p >= 0);
      prob += p;
      if (v != u) {  // Self-loops are resolved algebraically, so they do not order the states.
        indeg[v]++;
      }
    }
    assert(prob <= 1 + 1e-12);
    (void)prob;
  }
  std::vector<int> order, ready;
  order.reserve(n);
  for (int u = 0; u < n; u++) {
    if (indeg[u] == 0) {
      ready.push_back(u);
    }
  }
  while (!ready.empty()) {
    int u = ready.back();
    ready.pop_back();
    order.push_back(u);
    for (auto [v, p] : trans[u]) {
      if (v != u && --indeg[v] == 0) {
        ready.push_back(v);
      }
    }
  }
  assert(static_cast<int>(order.size()) == n);  // Cycles other than self-loops are unsupported.
  std::vector<double> res(n);
  for (int i = n - 1; i >= 0; i--) {
    int u = order[i];
    double stay = 0, total = cost[u];
    for (auto [v, p] : trans[u]) {
      if (v == u) {
        stay += p;
      } else {
        total += p * res[v];
      }
    }
    assert(stay < 1);  // A state that never leaves itself has infinite expected cost.
    res[u] = total / (1 - stay);
  }
  return res;
}
\end{lstlisting}
\Needspace*{4\baselineskip}
\textbf{Example Usage}\par
\begin{lstlisting}[style=cpp]
#include <cmath>
using namespace std;

bool EQ(double a, double b) {
  return fabs(a - b) < 1e-9;
}

int main() {
  // Roll a fair die to advance along squares [0, n], where a roll that would pass the last square
  // is rerolled. State n is absorbing, and every roll costs one step.
  const int n = 2;
  vector<vector<pair<int, double>>> board(n + 1);
  for (int i = 0; i < n; i++) {
    for (int face = 1; face <= 6; face++) {
      board[i].emplace_back(i + face <= n ? i + face : i, 1.0 / 6);
    }
  }
  vector<double> steps(n + 1, 1);
  steps[n] = 0;
  auto rolls = expected_cost(board, steps);
  assert(EQ(rolls[n], 0));
  assert(EQ(rolls[1], 6.0));  // Only one of six faces advances from square 1.
  assert(EQ(rolls[0], 6.0));
  // The same equation before the self-loop is eliminated: four of six faces reroll square 0.
  assert(EQ(rolls[0], 1 + rolls[1] / 6 + rolls[2] / 6 + 4.0 / 6 * rolls[0]));

  // The coupon collector, where state k is the number of distinct coupons already held. A draw
  // repeats a held coupon with probability k/m, which is a self-loop.
  const int m = 8;
  vector<vector<pair<int, double>>> coupons(m + 1);
  for (int k = 0; k < m; k++) {
    coupons[k].emplace_back(k, static_cast<double>(k) / m);
    coupons[k].emplace_back(k + 1, static_cast<double>(m - k) / m);
  }
  vector<double> draws(m + 1, 1);
  draws[m] = 0;
  auto collected = expected_cost(coupons, draws);
  double harmonic = 0;
  for (int i = 1; i <= m; i++) {
    harmonic += 1.0 / i;
  }
  assert(EQ(collected[0], m * harmonic));  // The classic m*H_m coupon collector bound.
  return 0;
}
\end{lstlisting}
\subsection{Absorbing Markov Chains}
\label{sec:6.7.3}
An absorbing Markov chain splits its states into absorbing states, which are never left once entered, and transient states, which are left forever with probability $1$. Ordering the transient states first splits the transition matrix into four blocks: $Q$ holds the probabilities between transient states, $R$ holds the probabilities from transient states into absorbing ones, and the two remaining blocks are a zero matrix and an identity matrix, since an absorbing state only ever moves to itself.

The expected number of visits to transient state $j$ starting from transient state $i$ is the entry $(i, j)$ of $I + Q + Q^2 + \dots$, since the entry $(i, j)$ of $Q^k$ is the probability of being at $j$ after exactly $k$ steps. Because absorption is certain, the powers of $Q$ tend to zero and the series converges to the fundamental matrix $N = (I - Q)^{-1}$. Every other quantity follows from $N$: summing a row counts the expected visits to all transient states, which is the expected number of steps before absorption, and multiplying by $R$ takes one final step into an absorbing state.

\begin{compactitem}
  \item \inlinecode{fundamental\_matrix(q)} returns $N = (I - {\inlinecodemath{q}})^{-1}$ for the transient-to-transient matrix \inlinecode{q}, as a matrix whose entry in row $i$ and column $j$ is the expected number of visits to transient state $j$ when starting at transient state $i$, counting the starting visit itself.
  \item \inlinecode{absorption\_steps(q)} returns a vector $t$ where $t_i$ is the expected number of steps taken from transient state $i$ until absorption.
  \item \inlinecode{absorption\_probabilities(q, r)} returns a matrix $B$ where $B_{i,k}$ is the probability that a chain starting at transient state $i$ is eventually absorbed at absorbing state $k$, given the transient-to-absorbing matrix \inlinecode{r}.
\end{compactitem}

Every transient state must reach an absorbing one, which is exactly what makes $I - Q$ invertible; a state that cannot is not transient but part of a closed group to be modeled as one absorbing state, so the zero pivot is asserted rather than reported. Unlike the backward induction of section \secref{6.7.2}, cycles are no obstacle here, since the system is solved rather than evaluated.

\Needspace*{3\baselineskip}
\textbf{Time Complexity}:
\begin{compactitem}
  \item $O(n^{3})$ per call to \inlinecode{fundamental\_matrix()} and \inlinecode{absorption\_steps()}, where $n$ is the number of transient states.
  \item $O(n^{3} + n^{2} \cdot k)$ per call to \inlinecode{absorption\_probabilities()}, where $k$ is the number of absorbing states.
\end{compactitem}

\Needspace*{3\baselineskip}
\textbf{Space Complexity}:
\begin{compactitem}
  \item $O(n^{2})$ auxiliary for \inlinecode{fundamental\_matrix()} and \inlinecode{absorption\_steps()}, and $O(n^{2} + n \cdot k)$ for \inlinecode{absorption\_probabilities()}.
  \item $O(n^{2})$ for the returned matrix from \inlinecode{fundamental\_matrix()}, $O(n)$ for the returned vector from \inlinecode{absorption\_steps()}, and $O(n \cdot k)$ for the returned matrix from \inlinecode{absorption\_probabilities()}.
\end{compactitem}

\begin{lstlisting}[style=cpp]
#include <cassert>
#include <cmath>
#include <utility>
#include <vector>

const double EPS = 1e-9;

std::vector<std::vector<double>> fundamental_matrix(const std::vector<std::vector<double>> &q) {
  int n = static_cast<int>(q.size());
  assert(n == 0 || static_cast<int>(q[0].size()) == n);
  // Reduce the augmented matrix [I - Q | I] to [I | N] by Gauss-Jordan elimination.
  std::vector<std::vector<double>> a(n, std::vector<double>(2 * n));
  for (int i = 0; i < n; i++) {
    for (int j = 0; j < n; j++) {
      a[i][j] = (i == j ? 1.0 : 0.0) - q[i][j];
    }
    a[i][n + i] = 1;
  }
  for (int col = 0; col < n; col++) {
    int pivot = col;
    for (int i = col + 1; i < n; i++) {
      if (fabs(a[i][col]) > fabs(a[pivot][col])) {
        pivot = i;
      }
    }
    assert(fabs(a[pivot][col]) > EPS);  // Some transient state never reaches an absorbing state.
    std::swap(a[col], a[pivot]);
    double scale = a[col][col];
    for (int j = col; j < 2 * n; j++) {
      a[col][j] /= scale;
    }
    for (int i = 0; i < n; i++) {
      if (i != col && a[i][col] != 0) {
        double factor = a[i][col];
        for (int j = col; j < 2 * n; j++) {
          a[i][j] -= factor * a[col][j];
        }
      }
    }
  }
  std::vector<std::vector<double>> res(n, std::vector<double>(n));
  for (int i = 0; i < n; i++) {
    for (int j = 0; j < n; j++) {
      res[i][j] = a[i][n + j];
    }
  }
  return res;
}

std::vector<double> absorption_steps(const std::vector<std::vector<double>> &q) {
  auto visits = fundamental_matrix(q);
  int n = static_cast<int>(visits.size());
  std::vector<double> res(n);
  for (int i = 0; i < n; i++) {
    for (int j = 0; j < n; j++) {
      res[i] += visits[i][j];
    }
  }
  return res;
}

std::vector<std::vector<double>> absorption_probabilities(
    const std::vector<std::vector<double>> &q, const std::vector<std::vector<double>> &r
) {
  assert(q.size() == r.size());
  auto visits = fundamental_matrix(q);
  int n = static_cast<int>(visits.size()), k = n == 0 ? 0 : static_cast<int>(r[0].size());
  std::vector<std::vector<double>> res(n, std::vector<double>(k));
  for (int i = 0; i < n; i++) {
    for (int j = 0; j < n; j++) {
      for (int c = 0; c < k; c++) {
        res[i][c] += visits[i][j] * r[j][c];
      }
    }
  }
  return res;
}
\end{lstlisting}
\Needspace*{4\baselineskip}
\textbf{Example Usage}\par
\begin{lstlisting}[style=cpp]
#include <algorithm>
#include <cmath>
using namespace std;

bool EQ(double a, double b) {
  return fabs(a - b) < 1e-9;
}

int main() {
  // A gambler with i of 4 dollars bets a dollar on a fair coin until going broke or reaching 4.
  // Transient states 0, 1, and 2 stand for holding 1, 2, and 3 dollars, and the two absorbing
  // states are ruin and victory. States 0 and 1 can each reach the other, so the transition graph
  // is cyclic and backward induction alone does not apply.
  vector<vector<double>> q{{0, 0.5, 0}, {0.5, 0, 0.5}, {0, 0.5, 0}};
  vector<vector<double>> r{{0.5, 0}, {0, 0}, {0, 0.5}};

  auto visits = fundamental_matrix(q);
  assert(EQ(visits[0][0], 1.5) && EQ(visits[0][1], 1.0) && EQ(visits[0][2], 0.5));
  assert(EQ(visits[1][0], 1.0) && EQ(visits[1][1], 2.0) && EQ(visits[1][2], 1.0));
  assert(EQ(visits[2][0], 0.5) && EQ(visits[2][1], 1.0) && EQ(visits[2][2], 1.5));

  // Starting with i dollars, the game lasts i*(4 - i) rounds on average.
  auto steps = absorption_steps(q);
  assert(EQ(steps[0], 3.0) && EQ(steps[1], 4.0) && EQ(steps[2], 3.0));

  // A fair game is won with probability i/4.
  auto odds = absorption_probabilities(q, r);
  assert(EQ(odds[0][1], 0.25) && EQ(odds[1][1], 0.5) && EQ(odds[2][1], 0.75));
  assert(all_of(odds.begin(), odds.end(), [](const auto &row) {
    return EQ(row[0] + row[1], 1.0);  // Absorption is certain.
  }));
  return 0;
}
\end{lstlisting}
\subsection{Stationary Distribution}
\label{sec:6.7.4}
A Markov chain with no absorbing states never settles on a single state, but the fraction of time it spends in each state does settle. That limit is the stationary distribution: the row vector $\pi$ with $\pi P = \pi$ and $\sum_i \pi_i = 1$, which is the one distribution left unchanged by a step of the chain. It exists and is unique whenever the chain is irreducible, meaning every state can reach every other, and it is the limit of the state distribution from any starting point once the chain is also aperiodic.

The defining equations are one short of determining $\pi$. Each column of $\pi P = \pi$ says that the probability flowing into a state equals the probability flowing out, and those $n$ equations sum to the trivial identity $1 = 1$, leaving the system rank $n - 1$ and any scalar multiple of a solution also a solution. Substituting the normalization $\sum_i \pi_i = 1$ for one of them pins down the scale and makes the system uniquely solvable by ordinary elimination.

\begin{compactitem}
  \item \inlinecode{stationary\_distribution(p, eps = 1e-10)} returns the stationary distribution of the $n$ by $n$ row-stochastic transition matrix \inlinecode{p}, where \inlinecode{p[i][j]} is the probability of stepping from state \inlinecode{i} to state \inlinecode{j}. It returns \inlinecode{std::nullopt} when the system is singular to within \inlinecode{eps}, which happens when the chain has more than one closed recurrent class and therefore no unique stationary distribution. A reducible chain may still have a unique distribution when all states eventually enter the same closed class.
\end{compactitem}

Two quantities follow at once: the expected number of steps to return to state $i$ is $1/\pi_i$, and on an undirected graph walked by a uniformly random incident edge, $\pi_v$ is proportional to the degree of $v$, which is why such a walk is uniform in the limit on a regular graph. Where the absorbing chains of section \secref{6.7.3} ask what happens before the process stops, this asks what happens when it never does.

\textbf{Time Complexity}: $O(n^{3})$ per call, where $n$ is the number of states.

\textbf{Space Complexity}: $O(n^{2})$ auxiliary and $O(n)$ for the returned distribution.

\begin{lstlisting}[style=cpp]
#include <cassert>
#include <cmath>
#include <optional>
#include <utility>
#include <vector>

std::optional<std::vector<double>> stationary_distribution(
    const std::vector<std::vector<double>> &p, double eps = 1e-10
) {
  int n = static_cast<int>(p.size());
  assert(n == 0 || static_cast<int>(p[0].size()) == n);
  if (n == 0) {
    return std::vector<double>{};
  }
  // Build the transpose of P - I, then replace its last row with the normalization equation.
  std::vector<std::vector<double>> a(n, std::vector<double>(n + 1));
  for (int i = 0; i < n; i++) {
    for (int j = 0; j < n; j++) {
      a[i][j] = p[j][i] - (i == j ? 1.0 : 0.0);
    }
  }
  for (int j = 0; j < n; j++) {
    a[n - 1][j] = 1;
  }
  a[n - 1][n] = 1;
  for (int col = 0; col < n; col++) {
    int pivot = col;
    for (int i = col + 1; i < n; i++) {
      if (fabs(a[i][col]) > fabs(a[pivot][col])) {
        pivot = i;
      }
    }
    if (fabs(a[pivot][col]) <= eps) {
      return std::nullopt;  // Multiple closed classes make the stationary distribution nonunique.
    }
    std::swap(a[col], a[pivot]);
    double scale = a[col][col];  // Saved, since the loop below overwrites it on its first step.
    for (int j = col; j <= n; j++) {
      a[col][j] /= scale;
    }
    for (int i = 0; i < n; i++) {
      if (i != col && a[i][col] != 0) {
        double factor = a[i][col];
        for (int j = col; j <= n; j++) {
          a[i][j] -= factor * a[col][j];
        }
      }
    }
  }
  std::vector<double> res(n);
  for (int i = 0; i < n; i++) {
    res[i] = a[i][n];
  }
  return res;
}
\end{lstlisting}
\Needspace*{4\baselineskip}
\textbf{Example Usage}\par
\begin{lstlisting}[style=cpp]
#include <algorithm>
#include <cassert>
#include <cmath>
using namespace std;

bool EQ(double a, double b) {
  return fabs(a - b) < 1e-9;
}

int main() {
  // A two-state chain that leaves state 0 with probability 1/4 and state 1 with probability 1/2,
  // so it rests in state 0 twice as often.
  auto two = stationary_distribution({{0.75, 0.25}, {0.5, 0.5}});
  assert(two.has_value());
  assert(EQ((*two)[0], 2.0 / 3) && EQ((*two)[1], 1.0 / 3));

  // A random walk on a path of three vertices, where the middle vertex has twice the degree and so
  // twice the stationary probability. Its expected return time is the reciprocal, 2.
  auto path = stationary_distribution({{0, 1, 0}, {0.5, 0, 0.5}, {0, 1, 0}});
  assert(path.has_value());
  assert(EQ((*path)[0], 0.25) && EQ((*path)[1], 0.5) && EQ((*path)[2], 0.25));
  assert(EQ(1 / (*path)[1], 2.0));

  // A doubly stochastic chain is uniform in the limit no matter the transition probabilities.
  auto cycle = stationary_distribution({{0, 0.3, 0.7}, {0.7, 0, 0.3}, {0.3, 0.7, 0}});
  assert(cycle.has_value());
  assert(all_of(cycle->begin(), cycle->end(), [](double p) { return EQ(p, 1.0 / 3); }));

  // Two states that never reach each other admit no unique distribution.
  assert(!stationary_distribution({{1, 0}, {0, 1}}).has_value());

  // Reducibility alone is not an obstruction: state 0 is transient and state 1 is the sole class.
  auto reducible = stationary_distribution({{0, 1}, {0, 1}});
  assert(reducible.has_value() && EQ((*reducible)[0], 0) && EQ((*reducible)[1], 1));
  return 0;
}
\end{lstlisting}
\subsection{Statistical Learning}
\label{sec:6.7.5}
The following routines fit a model to observed data and use it to describe or predict unobserved data. Each takes its input as a \inlinecode{Points} matrix whose rows are samples and whose columns are features, so a row is one observation of the same measurements taken in the same order. Supervised routines additionally take a \inlinecode{labels} vector, parallel to the rows, holding each sample's class as an integer in $[0, {\inlinecodemath{classes}})$.

Features measured on different scales distort every distance here, so call \inlinecode{standardize()} before \inlinecode{kmeans()}, \inlinecode{knn\_classify()}, or \inlinecode{logistic\_regression()} and apply the returned pair to later queries. Naive Bayes is exempt, since it fits a separate variance per feature. None of these choose \inlinecode{k} or \inlinecode{l2} for you: hold out part of the data and keep the value predicting it best.

\begin{lstlisting}[style=cpp]
#include <algorithm>
#include <cassert>
#include <cmath>
#include <cstddef>
#include <limits>
#include <numeric>
#include <random>
#include <utility>
#include <vector>

using Points = std::vector<std::vector<double>>;

double sqdist(const std::vector<double> &a, const std::vector<double> &b) {
  assert(a.size() == b.size());
  double res = 0;
  for (size_t i = 0; i < a.size(); i++) {
    res += (a[i] - b[i]) * (a[i] - b[i]);
  }
  return res;
}
\end{lstlisting}
Standardizing subtracts each column's mean and divides by its standard deviation, which is what makes one Euclidean distance meaningful across features recorded in different units. The population standard deviation is used, dividing by $n$ rather than $n - 1$, since these columns are the data being fit and not a sample standing in for a larger population. A column that never varies would give a zero divisor, so it reports $1$ instead and standardizes to all zeros, contributing nothing to any distance, which is exactly what a constant feature deserves.

\begin{compactitem}
  \item \inlinecode{column\_mean(points)} and \inlinecode{column\_stddev(points, mean)} return the per-column mean and population standard deviation, the latter reporting $1$ for a column that never varies.
  \item \inlinecode{standardize(points)} rescales every column in place to zero mean and unit variance and returns the (\inlinecode{mean}, \inlinecode{stddev}) it used, while \inlinecode{standardize(x, mean, stddev)} applies a returned pair to one later sample.
\end{compactitem}

\textbf{Time Complexity}: $O(n \cdot d)$ per call over $n$ rows of $d$ columns, and $O(d)$ for one sample.

\textbf{Space Complexity}: $O(d)$ auxiliary.

\begin{lstlisting}[style=cpp]
std::vector<double> column_mean(const Points &points) {
  assert(!points.empty());
  int n = static_cast<int>(points.size()), d = static_cast<int>(points[0].size());
  std::vector<double> res(d);
  for (int i = 0; i < n; i++) {
    for (int j = 0; j < d; j++) {
      res[j] += points[i][j];
    }
  }
  for (int j = 0; j < d; j++) {
    res[j] /= n;
  }
  return res;
}

std::vector<double> column_stddev(const Points &points, const std::vector<double> &mean) {
  assert(!points.empty() && mean.size() == points[0].size());
  int n = static_cast<int>(points.size()), d = static_cast<int>(points[0].size());
  std::vector<double> res(d);
  for (int i = 0; i < n; i++) {
    for (int j = 0; j < d; j++) {
      res[j] += (points[i][j] - mean[j]) * (points[i][j] - mean[j]);
    }
  }
  for (int j = 0; j < d; j++) {
    res[j] = std::sqrt(res[j] / n);
    if (res[j] < 1e-12) {
      res[j] = 1;
    }
  }
  return res;
}

void standardize(
    std::vector<double> &x, const std::vector<double> &mean, const std::vector<double> &stddev
) {
  assert(x.size() == mean.size() && x.size() == stddev.size());
  assert(std::all_of(stddev.begin(), stddev.end(), [](double value) { return value > 0; }));
  for (size_t j = 0; j < x.size(); j++) {
    x[j] = (x[j] - mean[j]) / stddev[j];
  }
}

std::pair<std::vector<double>, std::vector<double>> standardize(Points &points) {
  std::vector<double> mean = column_mean(points), stddev = column_stddev(points, mean);
  for (std::vector<double> &row : points) {
    standardize(row, mean, stddev);
  }
  return std::make_pair(mean, stddev);
}
\end{lstlisting}
$k$-means partitions the rows into $k$ clusters so that each row lies nearer to its own cluster's mean than to any other. Lloyd's algorithm alternates the two halves of that condition: assign every row to the nearest centroid, then move every centroid to the mean of the rows assigned to it. Each half can only lower the total squared distance, so the process converges, but only to a local optimum that the starting centroids decide.

Choosing those starting centroids is called seeding, and $k$-means++ is the standard rule: take the first uniformly at random, then draw each remaining one with probability proportional to its squared distance from the nearest already chosen. That spreads them out instead of letting several land in one dense region, bounding the expected cost at $O(\log k)$ times optimal, which uniform choice cannot promise. The loop stops once an assignment pass changes nothing, and a cluster that loses all its rows keeps its previous centroid.

\begin{compactitem}
  \item \inlinecode{kmeans(points, k, rng, iterations = 100)} returns (\inlinecode{centroids}, \inlinecode{assignment}), partitioning the rows into \inlinecode{k} clusters, where \inlinecode{assignment[i]} is the cluster of row \inlinecode{i} and \inlinecode{centroids[c]} is the mean of cluster \inlinecode{c}. The number of rows must be at least \inlinecode{k}; both \inlinecode{k} and \inlinecode{iterations} must be positive.
\end{compactitem}

\textbf{Time Complexity}: $O(i \cdot n \cdot k \cdot d)$ per call for \inlinecode{iterations} $i$, $n$ rows, and $d$ columns.

\textbf{Space Complexity}: $O(n + k \cdot d)$ auxiliary.

\begin{lstlisting}[style=cpp]
std::pair<Points, std::vector<int>> kmeans(
    const Points &points, int k, std::mt19937 &rng, int iterations = 100
) {
  int n = static_cast<int>(points.size());
  assert(k > 0 && k <= n && iterations > 0);
  Points centroids;
  centroids.push_back(points[std::uniform_int_distribution<int>(0, n - 1)(rng)]);
  std::vector<double> best(n);
  for (int i = 0; i < n; i++) {
    best[i] = sqdist(points[i], centroids[0]);
  }
  while (static_cast<int>(centroids.size()) < k) {
    std::discrete_distribution<int> pick(best.begin(), best.end());
    centroids.push_back(points[pick(rng)]);
    for (int i = 0; i < n; i++) {
      best[i] = std::min(best[i], sqdist(points[i], centroids.back()));
    }
  }
  std::vector<int> assignment(n, -1);
  for (int it = 0; it < iterations; it++) {
    bool changed = false;
    for (int i = 0; i < n; i++) {
      int nearest = 0;
      double nearest_dist = sqdist(points[i], centroids[0]);
      for (int c = 1; c < k; c++) {
        double dist = sqdist(points[i], centroids[c]);
        if (dist < nearest_dist) {
          nearest = c;
          nearest_dist = dist;
        }
      }
      if (assignment[i] != nearest) {
        assignment[i] = nearest;
        changed = true;
      }
    }
    if (!changed) {
      break;
    }
    Points sums(k, std::vector<double>(points[0].size()));
    std::vector<int> counts(k);
    for (int i = 0; i < n; i++) {
      for (size_t j = 0; j < points[i].size(); j++) {
        sums[assignment[i]][j] += points[i][j];
      }
      counts[assignment[i]]++;
    }
    for (int c = 0; c < k; c++) {
      if (counts[c] > 0) {
        for (size_t j = 0; j < sums[c].size(); j++) {
          centroids[c][j] = sums[c][j] / counts[c];
        }
      }
    }
  }
  return std::make_pair(centroids, assignment);
}
\end{lstlisting}
$k$-nearest neighbors classifies a query by letting the $k$ nearest training rows vote, which fits no model at all: the training data is the model, and all the work happens at query time. Small $k$ follows the data closely and is swayed by noise, while large $k$ smooths across genuine boundaries, so $k$ trades variance against bias. Averaging the neighbors' values rather than counting their votes turns the same routine into regression.

Selecting the $k$ nearest with \inlinecode{std::nth\_element} costs linear time rather than the $O(n \log n)$ of a full sort, since the neighbors' relative order is irrelevant to the vote. Scanning every row makes each query linear in the data, which the k-d tree of section \secref{2.6.5} improves to logarithmic on low-dimensional data; above roughly ten features the two converge and the scan below is preferable.

\begin{compactitem}
  \item \inlinecode{knn\_classify(points, labels, query, k)} returns the majority label among the \inlinecode{k} rows nearest to \inlinecode{query}, breaking a tie toward the smaller label. \inlinecode{k} must lie in $[1, n]$ for \inlinecode{n} rows.
\end{compactitem}

\textbf{Time Complexity}: $O(n \cdot d)$ per call, for $n$ rows of $d$ columns.

\textbf{Space Complexity}: $O(n)$ auxiliary.

\begin{lstlisting}[style=cpp]
int knn_classify(
    const Points &points, const std::vector<int> &labels, const std::vector<double> &query, int k
) {
  int n = static_cast<int>(points.size());
  assert(n > 0 && static_cast<int>(labels.size()) == n && query.size() == points[0].size());
  assert(k >= 1 && k <= n);
  assert(std::all_of(labels.begin(), labels.end(), [](int label) { return label >= 0; }));
  std::vector<int> order(n);
  std::iota(order.begin(), order.end(), 0);
  std::vector<double> dist(n);
  for (int i = 0; i < n; i++) {
    dist[i] = sqdist(points[i], query);
  }
  std::nth_element(order.begin(), order.begin() + (k - 1), order.end(), [&](int a, int b) {
    return dist[a] < dist[b];
  });
  std::vector<int> votes(*std::max_element(labels.begin(), labels.end()) + 1);
  for (int i = 0; i < k; i++) {
    votes[labels[order[i]]]++;
  }
  return static_cast<int>(std::max_element(votes.begin(), votes.end()) - votes.begin());
}
\end{lstlisting}
Naive Bayes picks the class maximizing $P(c) \prod_j P(x_j \mid c)$, which is Bayes' rule with the denominator dropped because it does not depend on the class. The independence it assumes among features given the class is almost never true, and yet the classifier is accurate anyway: the products it computes are poorly calibrated as probabilities, but the class that maximizes them is usually still the right one, since the errors tend to move every class the same way.

The Gaussian variant below models each feature within each class by a normal distribution, so fitting is one pass accumulating a mean and variance per (class, feature) pair. Scores accumulate as logarithms, both because a product of hundreds of densities underflows and because the logarithm of a normal density is a plain quadratic. A feature that is constant within a class would give a zero variance and an infinite score, so variances are floored at a small fraction of the overall spread.

\begin{compactitem}
  \item \inlinecode{NaiveBayes(points, labels, classes)} fits a Gaussian naive Bayes classifier, \inlinecode{predict(x)} returns its most probable class for a sample \inlinecode{x}, and \inlinecode{log\_likelihood(x, c)} returns the unnormalized log probability that \inlinecode{x} belongs to class \inlinecode{c}. A class carried by no row is given zero prior, so it scores $-\infty$ and is never predicted.
\end{compactitem}

\Needspace*{3\baselineskip}
\textbf{Time Complexity}:
\begin{compactitem}
  \item $O(n \cdot d)$ per construction.
  \item $O(c \cdot d)$ per call to \inlinecode{predict()} and $O(d)$ per call to \inlinecode{log\_likelihood()}, for $n$ rows of $d$ columns and $c$ classes.
\end{compactitem}

\textbf{Space Complexity}: $O(c \cdot d)$ for storage.

\begin{lstlisting}[style=cpp]
class NaiveBayes {
  std::vector<double> log_prior;
  Points mean, var;

 public:
  NaiveBayes(const Points &points, const std::vector<int> &labels, int classes) {
    assert(!points.empty() && labels.size() == points.size() && classes > 0);
    assert(std::all_of(labels.begin(), labels.end(), [&](int label) {
      return 0 <= label && label < classes;
    }));
    int n = static_cast<int>(points.size()), d = static_cast<int>(points[0].size());
    std::vector<int> counts(classes);
    mean.assign(classes, std::vector<double>(d));
    var.assign(classes, std::vector<double>(d));
    log_prior.assign(classes, 0);
    for (int i = 0; i < n; i++) {
      counts[labels[i]]++;
      for (int j = 0; j < d; j++) {
        mean[labels[i]][j] += points[i][j];
      }
    }
    for (int c = 0; c < classes; c++) {
      if (counts[c] == 0) {  // A class with no rows has zero prior, so it is never predicted.
        log_prior[c] = -std::numeric_limits<double>::infinity();
        continue;
      }
      log_prior[c] = std::log(static_cast<double>(counts[c]) / n);
      for (int j = 0; j < d; j++) {
        mean[c][j] /= counts[c];
      }
    }
    double spread = 0;
    for (int i = 0; i < n; i++) {
      for (int j = 0; j < d; j++) {
        double diff = points[i][j] - mean[labels[i]][j];
        var[labels[i]][j] += diff * diff;
        spread += diff * diff;
      }
    }
    double floor_var = 1e-9 * spread / n + 1e-12;
    for (int c = 0; c < classes; c++) {
      for (int j = 0; j < d; j++) {
        var[c][j] = counts[c] == 0 ? floor_var : std::max(var[c][j] / counts[c], floor_var);
      }
    }
  }

  double log_likelihood(const std::vector<double> &x, int c) const {
    assert(c >= 0 && c < static_cast<int>(log_prior.size()) && x.size() == mean[c].size());
    static const double LOG_2PI = 1.8378770664093453;
    double res = log_prior[c];
    for (size_t j = 0; j < x.size(); j++) {
      double diff = x[j] - mean[c][j];
      res -= 0.5 * (LOG_2PI + std::log(var[c][j]) + diff * diff / var[c][j]);
    }
    return res;
  }

  int predict(const std::vector<double> &x) const {
    int best = 0;
    double best_score = log_likelihood(x, 0);
    for (size_t c = 1; c < log_prior.size(); c++) {
      double score = log_likelihood(x, static_cast<int>(c));
      if (score > best_score) {
        best = static_cast<int>(c);
        best_score = score;
      }
    }
    return best;
  }
};
\end{lstlisting}
Logistic regression models the probability of class $1$ as $\sigma(w \cdot x)$ for the logistic function $\sigma(z) = 1/(1 + e^{-z})$, which maps any real score into $(0, 1)$. Fitting minimizes the average negative log-likelihood, a convex function whose gradient is $X^{\top}(\sigma(Xw) - y)/n$: each row pushes the weights by its own features, scaled by how wrong the current prediction is on it. Convexity makes every local minimum global, so plain gradient descent suffices.

Unlike the least squares of section \secref{6.5.5} there is no closed form, and on perfectly separable data the weights diverge as every predicted probability is driven to $0$ or $1$. The \inlinecode{l2} penalty adds $\lambda\|w\|^2$ to keep them finite, excluding the intercept so that shifting all labels does not distort it. Section \secref{5.2.4} supplies the adaptive steps worth substituting when convergence is slow.

\begin{compactitem}
  \item \inlinecode{logistic\_predict(w, x)} returns the probability that sample \inlinecode{x} belongs to class $1$ under a weight vector \inlinecode{w} whose index $0$ is the intercept.
  \item \inlinecode{logistic\_regression(points, labels, ...)} returns such a weight vector, fit by gradient descent over labels that are $0$ or $1$. Optional parameters default to \inlinecode{rate = 0.1}, \inlinecode{iterations = 1000}, and \inlinecode{l2 = 0}.
\end{compactitem}

\textbf{Time Complexity}: $O(i \cdot n \cdot d)$ per call for \inlinecode{iterations} $i$, $n$ rows, and $d$ columns.

\textbf{Space Complexity}: $O(d)$ auxiliary.

\begin{lstlisting}[style=cpp]
double sigmoid(double z) {
  return z >= 0 ? 1 / (1 + std::exp(-z)) : std::exp(z) / (1 + std::exp(z));
}

double logistic_predict(const std::vector<double> &w, const std::vector<double> &x) {
  assert(w.size() == x.size() + 1);
  double z = w[0];
  for (size_t j = 0; j < x.size(); j++) {
    z += w[j + 1] * x[j];
  }
  return sigmoid(z);
}

std::vector<double> logistic_regression(
    const Points &points, const std::vector<int> &labels, double rate = 0.1, int iterations = 1000,
    double l2 = 0
) {
  assert(!points.empty() && labels.size() == points.size());
  assert(rate > 0 && iterations > 0 && l2 >= 0);
  assert(std::all_of(labels.begin(), labels.end(), [](int label) {
    return label == 0 || label == 1;
  }));
  int n = static_cast<int>(points.size()), d = static_cast<int>(points[0].size());
  std::vector<double> w(d + 1), grad(d + 1);
  for (int it = 0; it < iterations; it++) {
    std::fill(grad.begin(), grad.end(), 0.0);
    for (int i = 0; i < n; i++) {
      double error = logistic_predict(w, points[i]) - labels[i];
      grad[0] += error;
      for (int j = 0; j < d; j++) {
        grad[j + 1] += error * points[i][j];
      }
    }
    w[0] -= rate * grad[0] / n;
    for (int j = 1; j <= d; j++) {
      w[j] -= rate * (grad[j] / n + l2 * w[j]);
    }
  }
  return w;
}
\end{lstlisting}
\Needspace*{4\baselineskip}
\textbf{Example Usage}\par
\begin{lstlisting}[style=cpp]
#include <cassert>
using namespace std;

int main() {
  // Two well-separated blobs: the first three rows lie near (0, 0) and the last three near
  // (10, 10).
  Points points{{0, 0}, {1, 0}, {0, 1}, {10, 10}, {11, 10}, {10, 11}};
  vector<int> labels{0, 0, 0, 1, 1, 1};

  // Standardizing puts both columns on one scale; the same pair transforms later queries.
  Points scaled(points);
  pair<vector<double>, vector<double>> shift = standardize(scaled);
  vector<double> query{9, 9};
  standardize(query, shift.first, shift.second);
  assert(knn_classify(scaled, labels, query, 3) == 1);

  mt19937 rng(1234567);
  pair<Points, vector<int>> clustered = kmeans(points, 2, rng);
  assert(clustered.first.size() == 2);
  assert(clustered.second[0] == clustered.second[1] && clustered.second[1] == clustered.second[2]);
  assert(clustered.second[3] == clustered.second[4] && clustered.second[4] == clustered.second[5]);
  assert(clustered.second[0] != clustered.second[3]);

  assert(knn_classify(points, labels, {0.5, 0.5}, 3) == 0);
  assert(knn_classify(points, labels, {9, 9}, 3) == 1);

  NaiveBayes nb(points, labels, 2);
  assert(nb.predict({0.5, 0.5}) == 0);
  assert(nb.predict({9, 9}) == 1);
  // Declaring a third class that no row carries leaves it with zero prior, never predicted.
  NaiveBayes sparse(points, labels, 3);
  assert(isinf(sparse.log_likelihood({0.5, 0.5}, 2)));
  assert(sparse.predict({0.5, 0.5}) == 0 && sparse.predict({9, 9}) == 1);

  vector<double> w = logistic_regression(points, labels, 0.5, 2000);
  assert(logistic_predict(w, {0.5, 0.5}) < 0.5);
  assert(logistic_predict(w, {9, 9}) > 0.5);
  return 0;
}
\end{lstlisting}

\section{\texorpdfstring{Game Theory}{6.8 Game Theory}}
\setcounter{section}{8}
\setcounter{subsection}{0}
\subsection{Nim}
\label{sec:6.8.1}
Solves normal-play nim and related impartial games where each position is the disjoint sum of independent piles. In normal play, the player who makes the last move wins. A nim position is losing exactly when the XOR of all pile sizes is $0$.

This is the base case of the Sprague-Grundy theorem: every finite impartial game position is equivalent to a Nim pile of size equal to its Grundy number, and sums of games combine by XOR.

\begin{compactitem}
  \item \inlinecode{nim\_sum(piles)} returns the XOR of all pile sizes.
  \item \inlinecode{first\_player\_wins(piles)} returns whether the player to move has a winning strategy.
  \item \inlinecode{winning\_move(piles)} returns a move (\inlinecode{pile}, \inlinecode{new\_size}) that moves to XOR $0$, or $(-1, -1)$ if the position is losing.
\end{compactitem}

\textbf{Time Complexity}: $O(n)$ per call, where $n$ is the number of piles.

\textbf{Space Complexity}: $O(1)$ auxiliary.

\begin{lstlisting}[style=cpp]
#include <utility>
#include <vector>

int nim_sum(const std::vector<int> &piles) {
  int x = 0;
  for (int p : piles) {
    x ^= p;
  }
  return x;
}

bool first_player_wins(const std::vector<int> &piles) {
  return nim_sum(piles) != 0;
}

std::pair<int, int> winning_move(const std::vector<int> &piles) {
  int x = nim_sum(piles);
  if (x == 0) {
    return {-1, -1};
  }
  for (int i = 0; i < static_cast<int>(piles.size()); i++) {
    if (int target = piles[i] ^ x; target < piles[i]) {
      return {i, target};
    }
  }
  return {-1, -1};
}
\end{lstlisting}
\Needspace*{4\baselineskip}
\textbf{Example Usage}\par
\begin{lstlisting}[style=cpp]
#include <cassert>
using namespace std;

int main() {
  vector<int> piles{3, 4, 5};
  // 3 ^ 4 ^ 5 = 2, so the position is winning.
  assert(nim_sum(piles) == 2);
  assert(first_player_wins(piles));
  auto [pile_idx, new_size] = winning_move(piles);
  piles[pile_idx] = new_size;
  // A winning move always leaves XOR 0 for the opponent.
  assert(nim_sum(piles) == 0);

  vector<int> losing{1, 2, 3};
  assert(!first_player_wins(losing));
  assert((winning_move(losing) == pair<int, int>{-1, -1}));
  return 0;
}
\end{lstlisting}
\subsection{Nim Product}
\label{sec:6.8.2}
Nimbers are the values of impartial games, added by XOR. Under the additional operation of nim multiplication they form a field: the nonnegative integers below $2^{2^k}$ are closed under both operations, with XOR as addition and nim product as multiplication. Writing $\otimes$ for the nim product and $\oplus$ for XOR, nim multiplication is commutative, associative, and distributive over $\oplus$, and is defined recursively by $a \otimes b = \operatorname{mex}\{(a' \otimes b) \oplus (a \otimes b') \oplus (a' \otimes b') : 0 \le a' < a, \, 0 \le b' < b\}$. Its main use is multiplying the Grundy values of independent ``coin-turning'' games whose move sets combine as a product, such as two-dimensional turning games built from one-dimensional ones.

Rather than evaluate the recursive definition directly, this implementation precomputes the nim product of every pair of single bits $2^i \otimes 2^j$ for $i, j < 64$. Each such product follows from the Fermat-like rule $2^{2^k} \otimes 2^{2^k} = 2^{2^k} \oplus 2^{2^k - 1}$ and the fact that distinct Fermat powers multiply like ordinary powers of two. A full product then XORs together the bit-pair products selected by the set bits of the operands. Because every 64-bit value lies in the field of size $2^{64}$, the operation is total on \inlinecode{uint64\_t} and the nonzero values have multiplicative inverses.

\begin{compactitem}
  \item \inlinecode{nim\_product(x, y)} returns the nim product of \inlinecode{x} and \inlinecode{y}.
  \item \inlinecode{nim\_pow(b, e)} returns \inlinecode{b} raised to the \inlinecode{e}-th power under nim multiplication.
  \item \inlinecode{nim\_inverse(b)} returns the multiplicative inverse of a nonzero \inlinecode{b}, so that \inlinecode{nim\_product(b, nim\_inverse(b)) == 1}.
\end{compactitem}

\Needspace*{3\baselineskip}
\textbf{Time Complexity}:
\begin{compactitem}
  \item $O(1)$ table construction on first use (a fixed 64 by 64 table).
  \item $O(64^{2})$ per call to \inlinecode{nim\_product()}; $O(64^{2} \log e)$ per call to \inlinecode{nim\_pow()} and \inlinecode{nim\_inverse()}.
\end{compactitem}

\textbf{Space Complexity}: $O(1)$ for a fixed 64 by 64 table of 64-bit products.

\begin{lstlisting}[style=cpp]
#include <cassert>
#include <cstdint>

uint64_t bit_prod[64][64];

// Build nim products for pairs of single-bit values. Each entry depends only on entries with a
// smaller first index, so a single ascending pass fills the table.
const bool nim_product_init = [] {
  for (int i = 0; i < 64; i++) {
    for (int j = 0; j < 64; j++) {
      if ((i & j) == 0) {
        bit_prod[i][j] = 1ULL << (i | j);  // Distinct Fermat powers multiply like 2^(i|j).
      } else {
        int a = (i & j) & -(i & j);  // Lowest Fermat power shared by both exponents.
        bit_prod[i][j] = bit_prod[i ^ a][j] ^ bit_prod[(i ^ a) | (a - 1)][(j ^ a) | (i & (a - 1))];
      }
    }
  }
  return true;
}();

uint64_t nim_product(uint64_t x, uint64_t y) {
  uint64_t res = 0;
  for (int i = 0; i < 64 && (x >> i) != 0; i++) {
    if ((x >> i) & 1) {
      for (int j = 0; j < 64 && (y >> j) != 0; j++) {
        if ((y >> j) & 1) {
          res ^= bit_prod[i][j];
        }
      }
    }
  }
  return res;
}

uint64_t nim_pow(uint64_t b, uint64_t e) {
  uint64_t res = 1;
  for (; e > 0; e >>= 1) {
    if (e & 1) {
      res = nim_product(res, b);
    }
    b = nim_product(b, b);
  }
  return res;
}

uint64_t nim_inverse(uint64_t b) {
  assert(b != 0);
  // The multiplicative group of the size-2^64 field has order 2^64 - 1, so b^(2^64 - 2) = b^{-1}.
  return nim_pow(b, static_cast<uint64_t>(-2));
}
\end{lstlisting}
\Needspace*{4\baselineskip}
\textbf{Example Usage}\par
\begin{lstlisting}[style=cpp]
#include <cassert>

int main() {
  // Smallest nontrivial products: {0,1,2,3} form the field GF(4) under XOR and nim product.
  assert(nim_product(2, 2) == 3);
  assert(nim_product(2, 3) == 1);
  assert(nim_product(3, 3) == 2);
  assert(nim_product(1, 12345) == 12345);  // 1 is the multiplicative identity.
  assert(nim_product(0, 12345) == 0);

  // Field axioms on a range of values: commutativity, distributivity over XOR, and inverses.
  for (uint64_t a = 0; a < 64; a++) {
    for (uint64_t b = 0; b < 64; b++) {
      assert(nim_product(a, b) == nim_product(b, a));
      for (uint64_t c = 0; c < 8; c++) {
        assert(nim_product(a, b ^ c) == (nim_product(a, b) ^ nim_product(a, c)));
      }
    }
    if (a != 0) {
      assert(nim_product(a, nim_inverse(a)) == 1);
    }
  }
  assert(nim_product(0x1234567890abcdefULL, nim_inverse(0x1234567890abcdefULL)) == 1);
  return 0;
}
\end{lstlisting}
\subsection{Sprague-Grundy Theorem}
\label{sec:6.8.3}
Computes Grundy numbers for finite impartial games. The Grundy number of a game position is the minimum excluded nonnegative integer (MEX) of the Grundy numbers reachable in one move. A position is losing exactly when its Grundy number is $0$.

For a sum of independent impartial games, the combined Grundy number is the XOR of the component Grundy numbers. This is the Sprague-Grundy theorem.

\begin{compactitem}
  \item \inlinecode{subtraction\_game\_grundy(max\_stones, moves)} computes Grundy numbers for a game where a move subtracts one positive value in \inlinecode{moves} from the pile. It returns a vector of length \inlinecode{max\_stones + 1}, whose element \inlinecode{grundy[stones]} is the Grundy number for a pile of that size. Nonpositive moves are ignored.
  \item \inlinecode{sum\_grundy(grundies)} returns the XOR of component Grundy numbers.
\end{compactitem}

\Needspace*{3\baselineskip}
\textbf{Time Complexity}:
\begin{compactitem}
  \item $O(n \cdot m)$ per call to \inlinecode{subtraction\_game\_grundy(max\_stones, moves)}, where $n$ is \inlinecode{max\_stones} and $m$ is the number of allowed moves.
  \item $O(k)$ per call to \inlinecode{sum\_grundy()}, where $k$ is the input length.
\end{compactitem}

\Needspace*{3\baselineskip}
\textbf{Space Complexity}:
\begin{compactitem}
  \item $O(n)$ output and $O(m)$ auxiliary for \inlinecode{subtraction\_game\_grundy(max\_stones, moves)}.
  \item $O(1)$ auxiliary for \inlinecode{sum\_grundy(grundies)}.
\end{compactitem}

\begin{lstlisting}[style=cpp]
#include <cassert>
#include <vector>

int mex(const std::vector<int> &values) {
  std::vector<char> seen(values.size() + 1);
  for (int v : values) {
    if (0 <= v && v < static_cast<int>(seen.size())) {
      seen[v] = true;
    }
  }
  for (int i = 0; i < static_cast<int>(seen.size()); i++) {
    if (!seen[i]) {
      return i;
    }
  }
  return static_cast<int>(seen.size());
}

std::vector<int> subtraction_game_grundy(int max_stones, const std::vector<int> &moves) {
  assert(max_stones >= 0);
  std::vector<int> grundy(max_stones + 1);
  for (int stones = 1; stones <= max_stones; stones++) {
    std::vector<int> reachable;
    for (int m : moves) {
      if (0 < m && m <= stones) {
        reachable.push_back(grundy[stones - m]);
      }
    }
    grundy[stones] = mex(reachable);
  }
  return grundy;
}

int sum_grundy(const std::vector<int> &grundies) {
  int x = 0;
  for (int g : grundies) {
    x ^= g;
  }
  return x;
}
\end{lstlisting}
\Needspace*{4\baselineskip}
\textbf{Example Usage}\par
\begin{lstlisting}[style=cpp]
#include <cassert>
using namespace std;

int main() {
  vector<int> moves{1, 3, 4};
  auto g = subtraction_game_grundy(10, moves);
  // With moves 1, 3, 4, pile size 2 is losing and pile size 4 has mex {0, 1} = 2.
  assert((g == vector<int>{0, 1, 0, 1, 2, 3, 2, 0, 1, 0, 1}));

  vector<int> heaps{g[5], g[7], g[10]};
  // Independent games combine by xor of their Grundy numbers.
  assert(sum_grundy(heaps) == (g[5] ^ g[7] ^ g[10]));
  return 0;
}
\end{lstlisting}
\subsection{Grundy Numbers on DAG}
\label{sec:6.8.4}
Computes Grundy numbers for impartial games whose positions form a directed acyclic graph. Each node is a game position, and each outgoing edge is a legal move. Terminal nodes have Grundy number $0$; all other nodes take the MEX of their successors' Grundy numbers.

The implementation uses memoized DFS. It assumes the graph is acyclic; if cycles are present, the usual finite impartial-game Grundy recurrence is not directly valid without additional analysis.

\begin{compactitem}
  \item \inlinecode{grundy\_on\_dag(g)} returns the Grundy number of every node in graph \inlinecode{g}, where \inlinecode{g[u]} contains all positions reachable from position \inlinecode{u} in one move and every endpoint is a valid node index.
\end{compactitem}

\textbf{Time Complexity}: $O(n + m)$ per call, where $n$ is the number of nodes and $m$ is the number of edges.

\textbf{Space Complexity}: $O(d)$ auxiliary stack space and $O(n + m)$ auxiliary heap space, where $d$ is DFS depth.

\begin{lstlisting}[style=cpp]
#include <algorithm>
#include <cassert>
#include <vector>

std::vector<int> grundy_on_dag(const std::vector<std::vector<int>> &g) {
  assert(std::all_of(g.begin(), g.end(), [&](const auto &edges) {
    return std::all_of(edges.begin(), edges.end(), [&](int v) {
      return 0 <= v && v < static_cast<int>(g.size());
    });
  }));
  std::vector<int> memo(g.size(), -1);
  auto dfs = [&](auto &&dfs, int u) {
    if (memo[u] != -1) {
      return;
    }
    std::vector<char> seen(g[u].size() + 1);
    for (int v : g[u]) {
      dfs(dfs, v);
      int x = memo[v];
      if (x < static_cast<int>(seen.size())) {
        seen[x] = true;
      }
    }
    int res = 0;
    while (res < static_cast<int>(seen.size()) && seen[res]) {
      res++;
    }
    memo[u] = res;
  };
  for (int u = 0; u < static_cast<int>(g.size()); u++) {
    dfs(dfs, u);
  }
  return memo;
}
\end{lstlisting}
\Needspace*{4\baselineskip}
\textbf{Example Usage}\par
\begin{lstlisting}[style=cpp]
#include <cassert>
using namespace std;

int main() {
  vector<vector<int>> g{{1, 2}, {3}, {3, 4}, {}, {}};
  // Terminal nodes have Grundy 0, their predecessors have mex {0} = 1, and node 0 has mex {1} = 0.
  assert((grundy_on_dag(g) == vector<int>{0, 1, 1, 0, 0}));
  return 0;
}
\end{lstlisting}
\subsection{Minimax and Alpha-Beta Pruning}
\label{sec:6.8.5}
Searches a finite two-player, zero-sum, perfect-information game tree using minimax and alpha-beta pruning. Minimax assumes both players play optimally: the maximizing player chooses the child with largest value, and the minimizing player chooses the child with smallest value.

Alpha-beta pruning computes the same value as minimax, but avoids exploring branches that cannot affect the final answer. It is most effective when good moves are searched first.

The example game is ``take $1$ or $2$ stones'': starting with \inlinecode{stones} stones, players alternate taking $1$ or $2$ stones, and the player who takes the last stone wins. The evaluation returns $1$ when the maximizing player wins and $-1$ when the minimizing player wins.

\begin{compactitem}
  \item \inlinecode{minimax(stones, maximizing)} returns the exact game-tree value, where \inlinecode{maximizing} indicates whose turn it is.
  \item \inlinecode{alpha\_beta(stones, maximizing)} returns the same value using alpha-beta pruning.
  \item \inlinecode{best\_take(stones)} returns an optimal first move using alpha-beta pruning, either $1$ or $2$, for positive \inlinecode{stones}.
\end{compactitem}

\textbf{Time Complexity}: $O(b^{d})$ per call in the worst case, where $b$ is branching factor and $d$ is search depth. Alpha-beta can achieve $O(b^{d/2})$ with optimal move ordering.

\textbf{Space Complexity}: $O(d)$ auxiliary recursion stack space.

\begin{lstlisting}[style=cpp]
#include <algorithm>

int minimax(int stones, bool maximizing) {
  if (stones == 0) {
    return maximizing ? -1 : 1;
  }
  int best = maximizing ? -2 : 2;
  for (int take = 1; take <= 2 && take <= stones; take++) {
    int child = minimax(stones - take, !maximizing);
    best = maximizing ? std::max(best, child) : std::min(best, child);
  }
  return best;
}

int alpha_beta(int stones, bool maximizing, int alpha = -2, int beta = 2) {
  if (stones == 0) {
    return maximizing ? -1 : 1;
  }
  int value = maximizing ? -2 : 2;
  for (int take = 1; take <= 2 && take <= stones; take++) {
    int child = alpha_beta(stones - take, !maximizing, alpha, beta);
    if (maximizing) {
      value = std::max(value, child);
      alpha = std::max(alpha, value);
    } else {
      value = std::min(value, child);
      beta = std::min(beta, value);
    }
    if (alpha >= beta) {
      break;
    }
  }
  return value;
}

int best_take(int stones) {
  int best_move = 1, best_value = -2;
  for (int take = 1; take <= 2 && take <= stones; take++) {
    int value = alpha_beta(stones - take, false);
    if (value > best_value) {
      best_value = value;
      best_move = take;
    }
  }
  return best_move;
}
\end{lstlisting}
\Needspace*{4\baselineskip}
\textbf{Example Usage}\par
\begin{lstlisting}[style=cpp]
#include <cassert>

int main() {
  assert(minimax(1, true) == 1);
  assert(minimax(3, true) == -1);
  for (int stones = 1; stones <= 10; stones++) {
    assert(minimax(stones, true) == alpha_beta(stones, true));
  }
  assert(best_take(4) == 1);  // Move to losing state with 3 stones.
  assert(best_take(5) == 2);  // Move to losing state with 3 stones.
  return 0;
}
\end{lstlisting}
